\documentclass{amsart}
\usepackage{amssymb}
\newtheorem{theorem}{Theorem}[section]
\newtheorem{proposition}[theorem]{Proposition}
\newtheorem{lemma}[theorem]{Lemma}
\newtheorem{corollary}[theorem]{Corollary}
\theoremstyle{definition}
\newtheorem{definition}[theorem]{Definition}
\theoremstyle{remark}
\newtheorem*{remark}{Remark}
\numberwithin{equation}{section}
\makeatletter
\def\subsection{\@startsection{subsection}{2}\z@{.5\linespacing\@plus.7\linespacing}{-.5em}{\normalfont\bfseries\boldmath}}
\makeatother
\newcommand{\RunningHead}{Unimodality of parking functions by maximum displacement}
\uppercasenonmath\RunningHead
\begin{document}
\title{Unimodality of the number of parking functions with a given maximum displacement}
\begin{abstract}
We prove that, for every $n$, the numbers of parking functions of length $n$ in which the maximum displacement of a car from its preferred spot is $0,1,\dots,n-1$ form a unimodal sequence. This confirms a conjecture of Celano, Elder, Hadaway, Harris, Martin, Priestley and Udell. The proof follows, for each value of the maximum displacement, the sign of the difference of two neighbouring terms as $n$ grows. It is computer-assisted: explicit estimates are used when the displacement is large, and exact computations with integers, carried out by computer, when the displacement or the length is small.
\end{abstract}
\maketitle
\sbox0{\normalfont\scriptsize\RunningHead}\ifdim\wd0>.85\textwidth\markboth{}{}\else\markboth{\RunningHead}{\RunningHead}\fi

\section{Introduction}\label{sec:intro}
Parking functions were introduced by Konheim and Weiss \cite{konheimweiss1966}, who studied a scheme for the storage of records that is now called hashing with linear probing, and who applied their results, at the end of their paper, to a problem of parking cars in a street. Parking functions have since become standard objects of combinatorics: they are equinumerous with the labelled trees on one more vertex, and they occur in the study of noncrossing partitions, of hyperplane arrangements and of the diagonal harmonics; see the survey of Yan \cite{yan2015}. A parking function describes a way in which $n$ cars can park one after another in a row of $n$ spots when each car has a preferred spot and takes the first empty spot at or after it. This is the rule of hashing with linear probing, in which items are stored one after another in a table whose cells lie in a row, each item in the first empty cell at or after its hash address: the cars are the items, the spots are the cells, and the number of spots by which a car is pushed beyond the spot it prefers, its displacement, is the distance between the address of an item and its cell.

The displacements measure the cost of storing and of finding an item, and they have been studied in both settings. Their sum, the total displacement, is the cost of building the table. Knuth \cite{knuth1998} and Flajolet, Poblete and Viola \cite{flajolet1998} find its moments, the latter also its limit laws. Pappe \cite[Theorem~3.1]{pappe2024} shows that the generating polynomial of the sum of the preferences of the cars, taken over the parking functions of length $n$, has log-concave coefficients, and his note remarks that this statistic is essentially equivalent to the displacement. For the largest of the displacements, the maximum displacement, Petersson \cite[Theorem~2.1]{petersson2009} determines the limit distribution in the setting of hashing, and formulas for the number of parking functions in which no displacement exceeds a given bound are given in \cite{aguilarfraga2024} and \cite{celano2025}; we recall these results in Section~\ref{subsec:intro-related}. The theorem of Pappe is a statement about the shape of the distribution of the sum of the preferences. The present paper is about the shape of the distribution of the maximum displacement: we ask how the numbers of parking functions of a given length whose maximum displacement is $0,1,2,\dots$ compare with one another.

To fix the notation, let $n\ge 1$ be an integer. The cars are numbered $1,\dots,n$ in the order in which they arrive, and the spots are numbered $1,\dots,n$ along the row and are all empty at the start. Car~$i$ prefers a spot $a_i\in\{1,\dots,n\}$ and parks in the first empty spot among $a_i,a_i+1,\dots,n$; if these spots are all occupied, car~$i$ cannot park. The sequence $\alpha=(a_1,\dots,a_n)$ is called a \emph{preference sequence} of length $n$, and it is called a \emph{parking function} if every car parks; we write $\mathrm{PF}_n$ for the set of all parking functions of length $n$. For $\alpha\in\mathrm{PF}_n$ the \emph{displacement} of car~$i$ is $d_i(\alpha)=\operatorname{pos}_i(\alpha)-a_i\ge 0$, where $\operatorname{pos}_i(\alpha)$ is the spot in which car~$i$ parks, and the \emph{maximum displacement} of $\alpha$ is
\[
M(\alpha)=\max_{1\le i\le n}d_i(\alpha)\in\{0,1,\dots,n-1\}.
\]
Parking functions with $M(\alpha)\le\ell$ are called \emph{$\ell$-interval parking functions} \cite[Definition~2.7]{celano2025}, and it is their number that the formulas mentioned above give. They range from the permutations of $1,\dots,n$, for $\ell=0$, where every car parks in the spot it prefers, to all parking functions of length $n$, for $\ell=n-1$. Here and below we cite \cite{celano2025} by the numbers of its second version; the numbers that we cite are the same in the first version except in one place, which is noted in Section~\ref{subsec:intro-related}.

Sorting the parking functions of length $n$ by their maximum displacement gives the numbers
\[
b_{n,\ell}=\#\{\alpha\in\mathrm{PF}_n:\ M(\alpha)=\ell\},\qquad 0\le\ell\le n-1.
\]
For $\ell\ge 1$ this is the number of $\ell$-interval parking functions of length $n$ that are not $(\ell-1)$-interval. The sum of these numbers is the number of all parking functions of length $n$, which is $(n+1)^{n-1}$, as follows from \cite[Lemma~2 and (6.1)]{konheimweiss1966}; a short proof, due to Pollak, is given in \cite[Section~2]{foatariordan1974}. The two ends of the sequence have closed forms: $b_{n,0}=n!$, the number of permutations, and $b_{n,n-1}=n^{n-2}$ for $n\ge 2$ \cite[Example~2.8 and Section~6]{celano2025}. For $n\le 7$ the numbers are
\begin{equation}\label{eq:intro-table}
\begin{array}{c|l}
n & b_{n,0},\ b_{n,1},\ \dots,\ b_{n,n-1}\\ \hline
1 & 1\\
2 & 2,\ 1\\
3 & 6,\ 7,\ 3\\
4 & 24,\ 51,\ 34,\ 16\\
5 & 120,\ 421,\ 377,\ 253,\ 125\\
6 & 720,\ 3963,\ 4594,\ 3688,\ 2546,\ 1296\\
7 & 5040,\ 42253,\ 62145,\ 57398,\ 46142,\ 32359,\ 16807
\end{array}
\end{equation}

The sum and the two ends of the sequence $(b_{n,0},\dots,b_{n,n-1})$ are thus known in closed form, and the question is how the terms in between are ordered. In \eqref{eq:intro-table} each row increases up to some entry and decreases from there on. The largest entry does not stand at a fixed place, however: it is the second entry of the rows $n=3,4,5$ and the third entry of the rows $n=6,7$. Celano, Elder, Hadaway, Harris, Martin, Priestley and Udell, who tabulate $b_{n,\ell}$ for $n\le 9$ \cite[Table~1]{celano2025}, conjecture, among the problems with which their paper closes, that for every $n\ge 1$ the sequence $(b_{n,0},\dots,b_{n,n-1})$ is unimodal \cite[Conjecture~6.5]{celano2025}. Here a finite sequence $(u_0,\dots,u_m)$ of real numbers is called \emph{unimodal} if there is some $\ell^*\in\{0,\dots,m\}$ such that $u_{\ell-1}\le u_\ell$ holds for $1\le\ell\le\ell^*$ and $u_\ell\ge u_{\ell+1}$ holds for $\ell^*\le\ell\le m-1$; equality is allowed. Unimodality and log-concavity are known for many sequences that arise in algebra, combinatorics and geometry, and they have been proved by methods of many kinds, which Stanley surveys in \cite{stanley1989}. The conjecture asks for the maximum of the displacements what the theorem of Pappe quoted above answers for the sum of the preferences; that theorem gives log-concavity, which is the stronger property, since a log-concave sequence of positive numbers is unimodal.

In the sentence before their conjecture, the authors of \cite{celano2025} state three inequalities at the two ends of a row, for all $n\ge 1$ and as consequences of their enumerative results, namely $b_{n,0}\le b_{n,1}$, $b_{n,1}\le b_{n,2}$ and $b_{n,n-2}\ge b_{n,n-1}$, and they add that experimentally the pattern seems to continue \cite[Section~6]{celano2025}. As \eqref{eq:intro-table} shows, the first of these inequalities fails for $n=2$ and the second for $n=3,4,5$; both hold in the later rows of the table. For the terms between the two ends of a row, the support for the conjecture is thus the computed rows. In this paper we prove the conjecture.

\begin{theorem}\label{thm:main}
For every integer $n\ge 1$ the sequence $(b_{n,0},b_{n,1},\dots,b_{n,n-1})$ is unimodal.
\end{theorem}

Theorem~\ref{thm:main} establishes for every length the shape that is seen in the tabulated rows. Its proof gives more on the left of a row: for $n\ge 17$ the row is nondecreasing on its left half, i.e.\ $b_{n,\ell-1}\le b_{n,\ell}$ whenever $1\le\ell\le n/2$ (Proposition~\ref{prop:leftpart} below). In small rows this can fail, as $b_{5,1}=421>b_{5,2}=377$ shows. Where the maximum of a row lies is determined neither by the theorem nor by its proof; we return to this question at the end of Section~\ref{subsec:intro-outline}.

Our proof is analytic. It constructs no correspondence between the parking functions that are counted by two neighbouring terms, and unimodality is not obtained from a stronger property such as log-concavity: we work directly with the signs of the differences of neighbouring terms. The combinatorics of parking functions enters through a single known relation, the recursion \eqref{eq:intro-lastcar} of Aguilar-Fraga et al.\ \cite[Theorem~3.7]{aguilarfraga2024}, which is recalled below. It expresses the number of parking functions of length $n+1$ in which no displacement exceeds a bound $k$ through the same numbers for smaller lengths and the same bound. Such a recursion relates different lengths for one bound, whereas the theorem compares different values of the maximum displacement for one length. We therefore read the table \eqref{eq:intro-table} by its columns and follow, for a fixed $\ell$, the sign of $b_{n,\ell}-b_{n,\ell-1}$ as $n$ grows. The two exceptions to the inequalities of \cite{celano2025} that were noted above are instances of the pattern that the proof follows: in a fixed column the difference is negative for the first lengths and positive for the later ones. The theorem follows when we know that the difference, once it is positive, does not become negative again, and that it does not become positive in one column before it has done so in the column to the left. This is what we prove, for the rows and the columns that the theorem needs.

Several of the tools that we apply to the solutions of this recursion are classical, among them the binomial identity of Abel \cite{abel1826}, the bounds of Robbins for the factorial \cite{robbins1955} and the theorem of de~Bruijn and Erd\H{o}s on renewal sequences with a log-convex kernel \cite{debruijn1953}; they are cited where they are used, and the last of them is proved there in the form that we need. The steps on which the argument turns are developed in the paper. They include a linear equation that replaces the recursion, which is quadratic, in a range of rows of each column; a comparison with renewal sequences, by which a column is shown to remain positive; a property that takes the place of convexity where the recursion does not preserve convexity; and a polynomial in a real variable through which two neighbouring columns are compared. Section~\ref{subsec:intro-outline} describes them.

Part of the proof is carried out by computer. The estimates are made with explicit constants, and they hold for all columns from an explicit one on, namely from column $82$ or from column $265$, depending on the statement that is proved. For the columns before these, and for the rows $n\le 16$, the proof relies on exact computations with integers. No parking functions are enumerated in them: the numbers that are needed are obtained from recursions, up to row $300$, and the integers that occur have up to about $14700$ decimal digits. In Section~\ref{subsec:intro-outline}, after the description of the estimates, we say what is computed, why finitely many computed numbers suffice for statements about all rows of a column, and which statements of the paper depend on the computations. The computations themselves are the subject of Section~\ref{sec:tables}, which contains the data that are needed to repeat them.

\subsection{Cumulative counts and related results}\label{subsec:intro-related}

We now recall the formulas for the number of $\ell$-interval parking functions and the limit theorem for the maximum displacement that were mentioned at the beginning, and we fix the notation for these numbers, the cumulative counts, which is used throughout the paper. For integers $n\ge 0$ and $k\ge 0$ let
\[
A_{n,k}=\#\{\alpha\in\mathrm{PF}_n:\ M(\alpha)\le k\}
\]
be the number of parking functions of length $n$ in which every car parks at most $k$ spots beyond the spot it prefers. Here and throughout the paper we admit $n=0$: the empty sequence is the one parking function of length~$0$, and its maximum displacement is~$0$. We also use the definition of $b_{n,\ell}$ for all integers $n\ge 0$ and $\ell\ge 0$. Thus $b_{0,0}=1$, $b_{0,\ell}=0$ for $\ell\ge 1$, and $b_{n,\ell}=0$ for $\ell\ge n\ge 1$, and we have
\[
A_{n,k}=\sum_{\ell=0}^{k}b_{n,\ell},\qquad b_{n,0}=A_{n,0},\qquad b_{n,\ell}=A_{n,\ell}-A_{n,\ell-1}\quad(\ell\ge 1).
\]
In particular $A_{0,k}=1$, while $A_{n,k}$ is nondecreasing in $k$ and equals $(n+1)^{n-1}$ for $n\le k+1$, where the bound on the displacements is no restriction.

One description of $A_{n,k}$ is a sum over the permutations $\sigma=(\sigma_1,\dots,\sigma_n)$ of $1,\dots,n$ \cite[Section~4.1, (6)--(8)]{celano2025}; in the first version of \cite{celano2025} these are the formulas (2)--(4) of Section~3. If $\lambda_i$ denotes the length of the longest run of consecutive spots that ends at spot~$i$, spot~$i$ included, and on which the values of $\sigma$ are at most $\sigma_i$, then there are $\prod_{i=1}^{n}\min(k+1,\lambda_i)$ parking functions with maximum displacement at most $k$ in which car~$\sigma_i$ parks in spot~$i$ for every~$i$, and $A_{n,k}$ is the sum of these products over all $\sigma$. Another description is a recursion in $n$, obtained by fixing the spot of the car that arrives last \cite[Theorem~3.7]{aguilarfraga2024}. In our notation it reads
\begin{equation}\label{eq:intro-lastcar}
A_{n+1,k}=\sum_{j=0}^{n}\min(j+1,k+1)\binom{n}{j}A_{j,k}\,A_{n-j,k}\qquad(n\ge 0,\ k\ge 0),
\end{equation}
with $A_{0,k}=1$. Below, \eqref{eq:intro-lastcar} is called the \emph{last-car recursion}; it is restated as Proposition~\ref{prop:lastcar}. Our proof starts from it: this recursion, the value $A_{0,k}=1$ and the monotonicity of $A_{n,k}$ in $k$ are all that the proof takes from the definition of a parking function.

The limit theorem is the theorem of Petersson on hashing with linear probing \cite[Theorem~2.1]{petersson2009}, in which the hash addresses of the items are independent and uniformly distributed over the cells. For a table with $m$ cells and $m-1$ items, conditioned on its last cell remaining empty, it says that the maximum displacement $d_{\max}$ of the items satisfies
\[
\Pr\Bigl(\frac{m-d_{\max}}{m^{2/3}}\le z\Bigr)\longrightarrow 1-\psi(z^{3/2})\quad(m\to\infty),\qquad 0<z<\infty,
\]
where $\psi$ is a function that the theorem defines by an explicit power series; by \cite[Corollary~2.4]{jansonpetersson2009} this series is $\psi(s)=\mathbb{E}\,e^{-s\mathcal{A}}$, where $\mathcal{A}=\int_0^1\max_{0\le s\le t}B_s\,dt$ for a standard Brownian motion $B$. Let $m=n+1$. Since an item that reaches an empty cell occupies it, the last cell remains empty exactly when the hash addresses of the $n$ items form a parking function of length $n$; the displacements of the items are then those of the cars, and all parking functions of length $n$ are equally likely. In our terms the theorem therefore says that all but an arbitrarily small fraction of the sum $(n+1)^{n-1}$ of the sequence $(b_{n,0},\dots,b_{n,n-1})$ lies at a distance of order $n^{2/3}$ from its end $\ell=n-1$. It is a statement of convergence in distribution, and it compares no two neighbouring terms of the sequence.

Among related results on special classes of parking functions, Selig and Zhu show that for nondecreasing parking functions the maximum displacement is the height of the corresponding {\L}ukasiewicz path \cite[Theorem~2.1]{seligzhu2024}, and Ager-Hart, Beerbower, Harris, Jakovleski and McClinton count, for each $\ell$, the $\ell$-interval Fubini rankings, the Fubini rankings in which at most $\ell+1$ competitors tie at a rank \cite[Corollary~4.1]{agerhart2026}; as they recall, Fubini rankings are parking functions.

\subsection{Outline of the proof}\label{subsec:intro-outline}

We describe the proof of Theorem~\ref{thm:main} in the order of the argument. The terms that we introduce on the way are used throughout the paper; the numerical constants of the proof are left, with few exceptions, to the statements in the body. We arrange the numbers $b_{n,k}$ in a table as in \eqref{eq:intro-table}, with rows indexed by the length $n$ and columns indexed by the displacement $k$. For integers $k\ge 1$ and $n\ge 0$ the \emph{step} of column $k$ in row $n$ is
\[
d_{n,k}=b_{n,k}-b_{n,k-1},
\]
and we say that column $k$ is \emph{positive}, \emph{nonpositive} or \emph{negative} in row $n$ according to the sign of this step. Row $n$ has a \emph{strict descent} at $i$ if $b_{n,i}<b_{n,i-1}$, and a \emph{strict ascent} at $j$ if $b_{n,j-1}<b_{n,j}$. A finite sequence is unimodal as soon as no strict descent is followed, further to the right, by a strict ascent. So the theorem follows once we show that $d_{n,i}<0<d_{n,j}$ is impossible for $1\le i<j\le n-1$ (Proposition~\ref{prop:descentascent}), and everything else in the paper is about the signs of the $d_{n,k}$.

These signs compare the entries of one row, i.e.\ counts for different displacements and one length, whereas the last-car recursion \eqref{eq:intro-lastcar}, as was said after Theorem~\ref{thm:main}, relates different lengths for one bound on the displacements; we have no recursion for the $b_{n,\ell}$ themselves. For $k\ge 2$ the step is a second difference in the bound, $d_{n,k}=A_{n,k}-2A_{n,k-1}+A_{n,k-2}$, of three sequences in $n$ each of which satisfies its own recursion. The proof therefore fixes a column $k$, follows the sign of $d_{n,k}$ as $n$ grows, and returns to the rows only at the end. A column $k\ge 2$ starts at row $n=k$, where $b_{k,k}=0$ and the step $d_{k,k}=-b_{k,k-1}=-k^{k-2}$ is negative. The picture behind the proof is that the column stays nonpositive for a while, turns positive at some row, and does not turn back. The first row $n>k$ with $d_{n,k}>0$ is called the \emph{crossing row} of column $k$, and the rows $n$ with $k+1\le n\le 2k-1$ form its \emph{chain range}.

Theorem~\ref{thm:main} rests on three propositions about columns, which make precise as much of this picture as the theorem needs. The first is the inequality on the left of a row that was mentioned after the theorem: in a row $n\ge 17$ no column $k$ with $2k\le n$ is negative, so that in these rows a column is not negative beyond its chain range.

\begin{proposition}\label{prop:leftpart}
Let $n$ and $k$ be integers with $n\ge 17$, $k\ge 1$ and $2k\le n$. Then $b_{n,k-1}\le b_{n,k}$.
\end{proposition}

The second proposition says that a column which is positive in one row of its chain range is positive in the next row of this range as well, so that there it changes its sign at most once, from nonpositive to positive.

\begin{proposition}\label{prop:chainrange}
Let $k\ge 2$ and $n$ be integers with $k+1\le n\le 2k-2$. If $b_{n,k-1}<b_{n,k}$, then $b_{n+1,k-1}<b_{n+1,k}$.
\end{proposition}

Applied in consecutive rows, Proposition~\ref{prop:chainrange} shows that a column which is positive in a row of its chain range is positive in every later row of this range: if $k\ge 2$, $m$ and $n$ are integers with $k+1\le m\le n\le 2k-1$ and $b_{m,k-1}<b_{m,k}$, then $b_{n,k-1}<b_{n,k}$.

The third says that as long as column $k$ has not turned positive, column $k+1$ has not turned positive either.

\begin{proposition}\label{prop:nextcolumn}
Let $k\ge 9$ and $n\ge k+2$ be integers. If $b_{j,k}\le b_{j,k-1}$ for every integer $j$ with $k+1\le j\le n$, then $b_{n,k+1}\le b_{n,k}$.
\end{proposition}

Propositions~\ref{prop:leftpart}, \ref{prop:chainrange} and~\ref{prop:nextcolumn} are proved in Sections~\ref{sec:left}, \ref{sec:far} and~\ref{sec:next}. The theorem follows from them in a few lines, which are given in Section~\ref{sec:proof}: in a row $n\ge 17$, Proposition~\ref{prop:leftpart} excludes a strict descent at $i\le n-1$ unless the row lies in the chain range of column $i$, and in that case Propositions~\ref{prop:chainrange} and~\ref{prop:nextcolumn} carry the nonpositive step from column $i$ to every later column of the row, so that no strict ascent follows. The rows $n\le 16$ are computed (Lemma~\ref{lem:sixteenrows}); they have to be treated apart, since the inequality of Proposition~\ref{prop:leftpart} fails in some of them and Proposition~\ref{prop:nextcolumn} starts at column~$9$. Each of the three propositions, in turn, is proved by estimates for the large columns and by exact computation, carried out by computer, for the small ones: the estimates apply to the columns $k\ge 82$ in Propositions~\ref{prop:leftpart} and~\ref{prop:chainrange}, and to the columns $k\ge 265$ in Proposition~\ref{prop:nextcolumn}. We describe the estimates first and the computations after them.

The estimates work with the counts divided by $n!$. The sequence $n\mapsto b_{n,\ell}$ of the entries of column $\ell$ of the table is called \emph{level}~$\ell$, and we write $\beta^{(\ell)}_n=b_{n,\ell}/n!$; for a column $k\ge 2$ we write $\delta_n=d_{n,k}/n!$. The last-car recursion is quadratic, but in the chain range it can be replaced by a linear relation, the column equation:
\[
\delta_n=\frac{k+1}{n}\,\delta_{n-1}+\frac{2}{n}\,\beta^{(k-1)}_{n-1}+\sum_{m=1}^{n-k}\varepsilon_m\,\delta_{n-m}\qquad(k+1\le n\le 2k-1),
\]
where $\varepsilon_1=1$ and $\varepsilon_m=(m-1)^{m-1}/m!$ for $m\ge 2$ (Proposition~\ref{prop:columnequation}). The step in row $n$ is fed by a nonnegative source, which comes from level $k-1$, and by the earlier steps of the same column, with positive weights. The negative start is fed back into the following rows, the source works against it, and the crossing row is the first row in which the source prevails. Summed over the first rows of the chain range, together with the corresponding equation for level $k-1$, the column equation shows that a large column $k$ cannot remain nonpositive in all the rows $k+1,\dots,k+r$ unless $r$ is at most a certain multiple of $k^{2/3}$ (Proposition~\ref{prop:strip}). Beyond the chain range there is no exact linear equation, and the proof works with inequalities that come from the generating series of the counts (Section~\ref{sec:series}).

The difficulty of the proof lies at three places, one in the proof of each of the three propositions. The first, in Proposition~\ref{prop:chainrange}, is that a column has to stay positive after its crossing row. The column equation keeps feeding the rows before the crossing, in which the step is nonpositive and of unknown size, into every later row, so that the positivity of the source does not decide the sign. On the \emph{near rows}, the rows of the chain range up to $k+0.47k$, we compare the column equation at row $n$ with the column equation at the crossing row, where the outcome was positive, and this controls what the rows before the crossing contribute (Proposition~\ref{prop:nearrows}). The comparison requires that the numbers $(k+m)\,\varepsilon_m\,\gamma_k^{-m}$, where $\gamma_k=(1+1/k)^{k-2}$, do not increase with $m$, and the estimate by which we show this reaches $m\le 0.47k$; this is the reason for the cut. On the rows of the chain range beyond these, the \emph{far rows}, we use the column equation at row $n$ alone, together with a run of positive rows before row $n$, and bound the step from below by a combination of values of level $k-1$ (Proposition~\ref{prop:farrows}). Up to row $k+0.6k$ the coefficients of this combination are nonnegative, and the claim reduces to a comparison of explicit numbers. Beyond row $k+0.6k$ some of the coefficients may be negative, and we bound the terms that carry them by a comparison with a renewal sequence. The comparison is proved in Section~\ref{sec:renewal} for arbitrary sequences, without reference to parking functions. Let $f=(f_i)_{i\ge 1}$ be a positive log-convex sequence, the kernel, and let $G$ be its renewal sequence, given by $G_0=1$ and $G_n=\sum_{i=1}^{n}f_i\,G_{n-i}$ for $n\ge 1$. If a sequence $v$ satisfies $\sum_{i=1}^{s}f_i\,v_{s-i}\le v_s$ for all $s$ up to some $\rho$, then for every $M\ge 1$ the sum $\sum_{m=M}^{\rho}\theta_m\,f_m\,v_{\rho-m}$, taken with nonnegative nondecreasing weights $\theta_m$, is largest in proportion to $v_\rho$ when $v=G$ (Proposition~\ref{prop:extremecolumn}). This bound rests on the classical theorem that the renewal sequence of a log-convex kernel is itself log-convex, which goes back to de~Bruijn and Erd\H{o}s \cite[Theorem~1]{debruijn1953}.

The second place is in Proposition~\ref{prop:leftpart} for a large column, in the rows beyond its chain range, where only inequalities are available. The rows before the crossing are still present in every later row, and we outweigh them through a single weighted sum of the steps of the column, which Proposition~\ref{prop:debt} bounds from below. This bound needs a convexity of the cumulative counts, suitably normalised, along a diagonal of their table, that is, when the length and the bound on the displacements are raised together. One would prove such a convexity by induction on the last-car recursion, but convexity does not pass through this induction: the recursion multiplies the terms by factors that increase with the bound on the displacements, and a factor of this kind does not preserve the convexity $2w_0\le w_-+w_+$ of three numbers $w_-$, $w_0$, $w_+$ that belong to three consecutive bounds. We therefore replace convexity by a family of inequalities, which we call property~K (Definition~\ref{def:propertyk}):
\[
2\,e^{\mu/k}\,w_0\le e^{\mu}\,w_++e^{-\mu}\,w_-\qquad\text{for every real number }\mu\ge 0,
\]
where $k$ is the index of the column. Convexity is the case $\mu=0$. Like convexity, property~K passes to sums and to nonnegative multiples of triples; unlike convexity, it also survives the multiplication by the factors that the recursion produces. Such a factor turns the inequality for one value of $\mu$ into the inequality for a larger value, so that the family as a whole passes through the induction although its single members do not (Lemma~\ref{lem:kratio}). One part of the recursion is not covered by this and needs an estimate of its own (Proposition~\ref{prop:kdirect}).

The third place is in Proposition~\ref{prop:nextcolumn} for a large column $k$, where two ranges have to meet. At a row $n=k+r$ of the chain range the steps of columns $k$ and $k+1$, divided by $n!$, are each written as a second difference of one polynomial, in which a real variable takes the place of the bound on the displacements, minus a sum over the compositions of an integer, in which a part of size $p$ carries the weight $p^{p-1}/p!$ and each cut between two parts carries a factor (Proposition~\ref{prop:curvatureboundary}). When the cuts carry no factor, the number of parts of a composition of $m$ that is drawn in proportion to its weight follows the law of the birthday problem, that is, the law of the number of draws from $m$ equally likely values before the first repetition (Proposition~\ref{prop:birthday}), and the two expressions are compared through estimates for this law. In this way we show that a nonpositive step of column $k$ forces a negative step of column $k+1$, and we can carry this through as long as $r$ is at most a certain multiple of $k^{2/3}$ (Proposition~\ref{prop:criterionpositive}). The hypothesis of Proposition~\ref{prop:nextcolumn}, however, says only that column $k$ has not turned positive up to row $k+r$. That this confines $r$ to the range just named is Proposition~\ref{prop:strip}, in which the multiple is the same. From Proposition~\ref{prop:strip}, or from a computation for the columns that it does not reach, the proofs of Propositions~\ref{prop:chainrange} and~\ref{prop:leftpart} also take the early crossing row from which they start.

It remains to describe the part of the proof that is carried out by computer. For the columns below those at which the estimates start, that is, below column $82$ in Propositions~\ref{prop:leftpart} and~\ref{prop:chainrange} and below column $265$ in Proposition~\ref{prop:nextcolumn}, and for the rows $n\le 16$, we rely on three computations, which are the subject of Section~\ref{sec:tables}. They are made in exact arithmetic with integers. No parking function is enumerated: the cumulative counts $A_{n,k}$ are obtained from the last-car recursion \eqref{eq:intro-lastcar} or, in the chain range, from a linear recursion that follows from it, and each computation ends in finitely many comparisons between integers that are formed from these counts. The first computation gives all entries of the rows $n\le 16$ and shows that each of these rows is unimodal (Lemma~\ref{lem:sixteenrows}). The second gives two signs for each of the columns $9\le k\le 264$: column $k$ is positive in a certain row of its chain range, and column $k+1$ is negative in the row before it (Lemma~\ref{lem:signpairs}). It reaches row $300$, and its integers have up to about $740$ decimal digits. The third gives, for each of the columns $2\le k\le 81$, the three counts $A_{n,k}$, $A_{n,k-1}$ and $A_{n,k-2}$, of which the step $d_{n,k}$ is the second difference, in all rows $n$ below a certain row; this row depends on the column and is at most row $175$. With these counts it verifies Proposition~\ref{prop:chainrange} for the column, verifies the inequality of Proposition~\ref{prop:leftpart} in those of the computed rows to which the proposition applies, and verifies that finitely many sums which are formed from the counts are nonnegative (Lemma~\ref{lem:columntable}). The integers that occur in these sums have up to about $14700$ decimal digits.

Finitely many computed numbers suffice for a different reason in each case. The rows $n\le 16$ are finitely many, and Proposition~\ref{prop:chainrange} concerns only the rows of the chain range, which are finitely many for each column. Proposition~\ref{prop:leftpart}, by contrast, concerns all rows $n\ge 2k$ of column $k$. Here a propagation theorem (Proposition~\ref{prop:propagation}) shows that, when the sums of the third computation are nonnegative, the step of the column is nonnegative in every row beyond the computed ones. Proposition~\ref{prop:nextcolumn} concerns infinitely many rows of each column as well, and for it the two signs of the second computation suffice because of Proposition~\ref{prop:chainrange}. By the first sign, the hypothesis of Proposition~\ref{prop:nextcolumn}, that column $k$ is nonpositive in all the rows $k+1,\dots,n$, can hold only for the rows $n$ up to the row of the second sign. If column $k+1$ were positive in such a row $n$, then by Proposition~\ref{prop:chainrange} it would still be positive in the row of the second sign, where it is negative. For the columns $82\le k\le 264$ the first sign has a second use, which was mentioned above: it supplies the positive row early in the chain range that the proofs of Propositions~\ref{prop:chainrange} and~\ref{prop:leftpart} need, and that Proposition~\ref{prop:strip} provides from column $265$ on (Corollary~\ref{cor:earlypositive}).

The statements of the paper that are obtained by computer are thus Lemmas~\ref{lem:sixteenrows}, \ref{lem:signpairs} and~\ref{lem:columntable}. Propositions~\ref{prop:leftpart}, \ref{prop:chainrange} and~\ref{prop:nextcolumn} depend on them only for the columns $k\le 264$; for the columns $k\ge 265$ they are proved without them. Theorem~\ref{thm:main} depends on them through these propositions and through Lemma~\ref{lem:sixteenrows}. Elsewhere in the paper the numerical verifications are comparisons between explicit numbers, which arise from the estimates and are made where they are needed.

The proof leaves open where the maximum of a row lies. In the picture described above, the columns $k\ge 2$ that are positive in row $n$ are those whose crossing row is at most $n$, and the maximum of the row stands at the last of them. Our argument, however, uses only bounds for the crossing row of a column and does not determine it; and, as was said in Section~\ref{subsec:intro-related}, the theorem of Petersson shows where most of the sum of a row lies and not which term is the largest. Still, the statements that we prove on the way confine the maximum from one side. In a row $n\ge 335$ the first column $p$ at which the maximum is attained stands at most $(121/25)^{1/3}\,n^{2/3}$ columns before the last column of the row, and the row increases strictly from column $81$ to column $p$ and decreases strictly from column $p+1$ on (Proposition~\ref{prop:rowpeak} and the paragraph after its proof, which the proof of Theorem~\ref{thm:main} does not use). The order $n^{2/3}$ of this bound is that of the distance from the end of the row at which, by the theorem of Petersson, most of the sum of the row lies. The constant is that of the range of rows which the estimates of Propositions~\ref{prop:strip} and~\ref{prop:criterionpositive} reach; it has no meaning beyond these estimates, and we have no bound for the distance from below.

\subsection{Organisation of the paper}\label{subsec:intro-organisation}

Propositions~\ref{prop:leftpart}, \ref{prop:chainrange} and~\ref{prop:nextcolumn} are stated in Section~\ref{subsec:intro-outline}. Sections~\ref{sec:prelim}--\ref{sec:tables} are preparatory: they treat the explicit sequences that occur, the last-car recursion, the column equation, the inequalities that hold in all rows, a monotonicity of the normalised counts on which the proof of Proposition~\ref{prop:leftpart} rests, and the facts that are obtained by exact computation. Sections~\ref{sec:crossing} and~\ref{sec:strip} bound the crossing row. Proposition~\ref{prop:chainrange} is proved in Sections~\ref{sec:near}--\ref{sec:far}, Proposition~\ref{prop:leftpart} in Sections~\ref{sec:propk}--\ref{sec:left}, and Proposition~\ref{prop:nextcolumn} in Sections~\ref{sec:comps}--\ref{sec:next}. Section~\ref{sec:proof} proves Theorem~\ref{thm:main} and closes with a bound for the position of the maximum of a row (Proposition~\ref{prop:rowpeak}).

\section{Preliminaries}\label{sec:prelim}
Throughout the paper sequences are indexed by the integers $n\ge 0$ unless another range is stated, and we use the convention $0^0=1$, so that $n^n/n!$ has the value $1$ at $n=0$. By a \emph{power series} we mean a formal power series in one variable $X$ with real coefficients; for a power series $F$ we write $[X^n]F$ for the coefficient of $X^n$ in $F$ and $F'$ for the formal derivative of $F$.

The proofs work with three explicit sequences, $(a_j)$, $(c_n)$ and $(\varepsilon_m)$, and the rest of this section collects what is needed about them: the ratios of their consecutive terms, the identities that connect them, bounds for their size, and two sums.

\subsection{The numbers $a_j$, $c_n$ and $\varepsilon_m$}\label{subsec:prelim-numbers}

The \emph{Cayley coefficients} are the numbers
\[
a_0=0,\qquad a_j=\frac{j^{\,j-1}}{j!}\quad(j\ge 1);
\]
by Cayley's formula, $j^{j-1}$ is the number of rooted labelled trees on $j$ vertices. The preferences of a parking function are also written $a_i$, as in the Introduction; no confusion will arise. Shifted by one, the Cayley coefficients are the numbers
\[
c_n=a_{n+1}=\frac{(n+1)^{n-1}}{n!}\qquad(n\ge 0),
\]
so that $c_0=1$. Since there are $(n+1)^{n-1}$ parking functions of length $n$ \cite[Section~2]{foatariordan1974}, $c_n$ is the number of parking functions of length $n$ divided by $n!$. The third sequence is the \emph{free kernel}
\[
\varepsilon_0=0,\qquad \varepsilon_m=\frac{(m-1)^{m-1}}{m!}\quad(m\ge 1),
\]
in which $\varepsilon_1=0^0/1!=1$; it is the kernel of a renewal equation satisfied by the $c_n$ (Corollary~\ref{cor:crenewal}). The first values are $a_1=a_2=1$, $a_3=\frac32$, $a_4=\frac83$ and $\varepsilon_1=1$, $\varepsilon_2=\frac12$, $\varepsilon_3=\frac23$, $\varepsilon_4=\frac98$. We often use the definitions in the forms
\begin{equation}\label{eq:prelim-forms}
(n+1)\,c_n=\frac{(n+1)^{n}}{n!},\qquad (m+1)\,\varepsilon_{m+1}=\frac{m^{m}}{m!},
\end{equation}
which hold for all integers $n\ge 0$ and $m\ge 0$. In particular, $a_j>0$ for $j\ge 1$, $c_n>0$ for $n\ge 0$, and $\varepsilon_m>0$ for $m\ge 1$.

The ratios of consecutive terms of these sequences are expressed by $(1+1/i)^{i}$ and $(1+1/i)^{i-1}$, and both are nondecreasing in $i\ge 1$. For the first this is classical: for every integer $i\ge 1$,
\begin{equation}\label{eq:prelim-euler}
\Bigl(1+\frac{1}{i}\Bigr)^{i}\le\Bigl(1+\frac{1}{i+1}\Bigr)^{i+1},
\end{equation}
by Bernoulli's inequality. For the second it is the following lemma.

\begin{lemma}\label{lem:powermono}
For all integers $i$ and $j$ with $1\le i\le j$,
\[
\Bigl(1+\frac{1}{i}\Bigr)^{i-1}\le\Bigl(1+\frac{1}{j}\Bigr)^{j-1}.
\]
\end{lemma}

\begin{proof}
By induction on $j$ it suffices to treat $j=i+1$. By the inequality stated before the lemma and since $1+\frac{1}{i+1}\le 1+\frac1i$,
\[
\Bigl(1+\frac{1}{i}\Bigr)^{i-1}\Bigl(1+\frac{1}{i}\Bigr)
\le\Bigl(1+\frac{1}{i+1}\Bigr)^{i+1}
\le\Bigl(1+\frac{1}{i+1}\Bigr)^{i}\Bigl(1+\frac{1}{i}\Bigr),
\]
and dividing by $1+\frac1i$ gives the claim.
\end{proof}

Since $a_{j+1}=(j+1)^{j}/(j+1)!=(j+1)^{j-1}/j!$, the ratio of two consecutive Cayley coefficients is
\begin{equation}\label{eq:prelim-ratio}
a_{j+1}=\Bigl(1+\frac{1}{j}\Bigr)^{j-1}a_j\qquad(j\ge 1).
\end{equation}
By Lemma~\ref{lem:powermono} and \eqref{eq:prelim-ratio} the ratio $a_{j+1}/a_j$ is nondecreasing in $j\ge 1$, that is, the Cayley coefficients $a_j$ with $j\ge 1$ form a log-convex sequence. We use this property in the following cross-multiplied form, which for $s=1$ and $v=u+1$ reads $a_{u+1}^2\le a_u\,a_{u+2}$.

\begin{lemma}\label{lem:cayleylogconvex}
For all integers $u$, $v$ and $s$ with $1\le u\le v$ and $s\ge 0$,
\[
a_{u+s}\,a_v\le a_u\,a_{v+s}.
\]
\end{lemma}

\begin{proof}
We argue by induction on $s$; the case $s=0$ is trivial. If the inequality holds for $s$, then \eqref{eq:prelim-ratio}, applied with $j=u+s$ and with $j=v+s$, gives
\begin{align*}
a_{u+s+1}\,a_v&=\Bigl(1+\frac{1}{u+s}\Bigr)^{u+s-1}a_{u+s}\,a_v\\
&\le\Bigl(1+\frac{1}{v+s}\Bigr)^{v+s-1}a_u\,a_{v+s}=a_u\,a_{v+s+1},
\end{align*}
where the inequality is obtained by multiplying the induction hypothesis by the inequality of Lemma~\ref{lem:powermono} for $u+s\le v+s$.
\end{proof}

\subsection{The binomial identity of Abel}\label{subsec:prelim-abel}

The identities among the numbers $a_j$, $c_n$ and $\varepsilon_m$ that we use all come from one classical identity, the binomial identity of Abel \cite{abel1826}. It has the form of the binomial theorem for $(x+y)^n$, with $(y-k)^{n-k}$ in place of $y^{n-k}$ and with the following weights in place of the powers $x^k$: for a real number $x$ and an integer $k\ge 0$ let
\[
w_0(x)=1,\qquad w_k(x)=x\,(x+k)^{k-1}\quad(k\ge 1).
\]
Note that $c_k=w_k(1)/k!$ for every $k\ge 0$. We use the identity with $x=1$ in this section and with general $x$ in Section~\ref{sec:renewal}.

\begin{lemma}[{Abel \cite[p.~159]{abel1826}}]\label{lem:abel}
For every integer $n\ge 0$ and all real numbers $x$ and $y$,
\[
\sum_{k=0}^{n}\binom{n}{k}\,w_k(x)\,(y-k)^{n-k}=(x+y)^n.
\]
\end{lemma}

In the notation of \cite{abel1826} the identity reads $(x+\alpha)^n=\sum_{\mu=0}^{n}\binom{n}{\mu}\,\alpha\,(\alpha-\mu\beta)^{\mu-1}(x+\mu\beta)^{n-\mu}$ for arbitrary $x$, $\alpha$ and $\beta$ and every positive integer $n$, the term with $\mu=0$ being $x^n$; the lemma is the case $\beta=-1$, with our $x$ and $y$ in the places of $\alpha$ and $x$, and for $n=0$ it is trivial.

Abel's identity yields the following recursion for the $c_n$. It is the last-car recursion of the Introduction, divided by $n!$, in the case in which the bound on the displacements is no restriction.

\begin{lemma}\label{lem:crecursion}
For every integer $n\ge 0$,
\[
(n+1)\,c_{n+1}=\sum_{j=0}^{n}(j+1)\,c_j\,c_{n-j}.
\]
\end{lemma}

\begin{proof}
We put $j=n-k$ on the right-hand side. By \eqref{eq:prelim-forms} and since $c_k=w_k(1)/k!$, it becomes
\begin{align*}
\sum_{k=0}^{n}\frac{(n+1-k)^{n-k}}{(n-k)!}\cdot\frac{w_k(1)}{k!}
&=\frac{1}{n!}\sum_{k=0}^{n}\binom{n}{k}\,w_k(1)\,(n+1-k)^{n-k}\\
&=\frac{(n+2)^n}{n!},
\end{align*}
where the last step is Lemma~\ref{lem:abel} with $x=1$ and $y=n+1$. This is the left-hand side, since $c_{n+1}=(n+2)^{n}/(n+1)!$.
\end{proof}

Written with the Cayley coefficients themselves, the recursion is the following convolution identity, on which the comparison of the weights of compositions in Section~\ref{sec:comps} rests.

\begin{lemma}\label{lem:cayleyconvolution}
For every integer $t\ge 0$,
\[
\sum_{i=0}^{t-1} i\,a_i\,a_{t-i}=(t-1)\,a_t.
\]
\end{lemma}

\begin{proof}
Since $a_0=0$, the cases $t\le 1$ are trivial. Let $t=n+2$ with $n\ge 0$. The term with $i=0$ vanishes, and with $i=j+1$ the others add up to $\sum_{j=0}^{n}(j+1)\,c_j\,c_{n-j}$, which by Lemma~\ref{lem:crecursion} equals $(n+1)\,c_{n+1}=(t-1)\,a_t$.
\end{proof}

\subsection{The series $C$ and the free kernel}\label{subsec:prelim-series}

Let
\[
C=\sum_{n\ge 0}c_nX^n=1+X+\tfrac32X^2+\tfrac83X^3+\cdots .
\]
Its constant coefficient is $c_0=1$, so $C$ has an inverse $1/C$. Since $c_n=a_{n+1}$ and $a_0=0$, we have $[X^j](XC)=a_j$ for every $j\ge 0$. With these series the recursion of Lemma~\ref{lem:crecursion} becomes the differential equation
\begin{equation}\label{eq:prelim-ode}
C'=(XC)'\,C,
\end{equation}
as one sees by comparing the coefficients of $X^n$, the coefficient of $X^j$ in $(XC)'$ being $(j+1)\,a_{j+1}=(j+1)\,c_j$. Classically, $XC$ is the tree function $T$, the exponential generating series of the rooted labelled trees, which satisfies $T=Xe^{T}$, and $C=e^{T}$; the differential equation is the derivative form of the second relation.

In Section~\ref{sec:chain} power series are divided by $C$. Since $(1/C)'=-C'/C^2$, the differential equation \eqref{eq:prelim-ode} gives
\begin{equation}\label{eq:prelim-inverse}
\Bigl(\frac{1}{C}\Bigr)'=-\frac{1}{C}\,(XC)',
\end{equation}
and the coefficients of a quotient by $C$ are given by the following lemma.

\begin{lemma}\label{lem:quotient}
Let $F$ be a power series. Then, for every integer $n\ge 0$,
\[
[X^n]\Bigl(\frac{1}{C}\,F\Bigr)=[X^n]F-\sum_{m=1}^{n}\varepsilon_m\,[X^{n-m}]F.
\]
\end{lemma}

\begin{proof}
We determine the coefficients of $1/C$. By the product rule $(XC)'=C+XC'$ and by the differential equation in the form $C'/C=(XC)'$,
\[
\frac{1}{C}\,(XC)'=1+X\,(XC)'=\sum_{n\ge 0}\frac{n^n}{n!}\,X^n,
\]
since the coefficient of $X^{n-1}$ in $(XC)'$ is $n\,a_n=n^n/n!$ for $n\ge 1$. By \eqref{eq:prelim-inverse}, $(1/C)'$ is the negative of this series. Comparing the coefficients of $X^m$ and using \eqref{eq:prelim-forms}, we find that
\[
(m+1)\,[X^{m+1}]\frac{1}{C}=-\frac{m^m}{m!}=-(m+1)\,\varepsilon_{m+1}\qquad(m\ge 0).
\]
Hence $[X^m](1/C)=-\varepsilon_m$ for every $m\ge 1$, while $[X^0](1/C)=1$, and the assertion is the formula for the coefficients of the product $\frac1C\,F$.
\end{proof}

It follows that the $c_n$ satisfy a renewal equation whose kernel is the free kernel.

\begin{corollary}\label{cor:crenewal}
For every integer $n\ge 1$,
\[
\sum_{m=1}^{n}\varepsilon_m\,c_{n-m}=c_n.
\]
\end{corollary}

\begin{proof}
This is Lemma~\ref{lem:quotient} with $F=C$, since the coefficient of $X^n$ in $\frac1C\,C=1$ is $0$ for $n\ge 1$.
\end{proof}

It remains to record the log-convexity of the free kernel, on which the arguments of Sections~\ref{sec:near} and~\ref{sec:far} rest. For an integer $m\ge 1$ let
\[
g_1=\frac12,\qquad g_m=\frac{m}{m+1}\Bigl(1+\frac{1}{m-1}\Bigr)^{m-1}\quad(m\ge 2).
\]
A direct computation from the definition of $\varepsilon$ shows that $g_m$ is the ratio of two consecutive terms of the free kernel:
\[
\varepsilon_{m+1}=g_m\,\varepsilon_m\qquad(m\ge 1).
\]

\begin{lemma}\label{lem:freelogconvex}
For every integer $m\ge 1$,
\[
\varepsilon_{m+1}^2\le\varepsilon_m\,\varepsilon_{m+2}.
\]
\end{lemma}

\begin{proof}
Since $\varepsilon_{m+1}=g_m\,\varepsilon_m$ and the $\varepsilon_m$ with $m\ge 1$ are positive, the assertion is that $g_m\le g_{m+1}$. We have $g_1=\frac12\le\frac43=g_2$. For $m\ge 2$ neither of the two positive factors of $g_m$ decreases when $m$ is replaced by $m+1$, the second by the inequality \eqref{eq:prelim-euler} with $i=m-1$.
\end{proof}

\subsection{Stirling bounds}\label{subsec:prelim-stirling}

For an integer $n\ge 0$ the number $e^{-n}n^n/n!$ is the probability that a Poisson variable with mean $n$ takes the value $n$; it is called the \emph{Poisson mass} at $n$, and its value at $n=0$ is $1$. The size of our sequences is estimated through the Poisson mass, since by their definitions
\[
e^{-j}a_j=\frac{1}{j}\cdot e^{-j}\frac{j^j}{j!}\qquad\text{and}\qquad e^{-(m+1)}\varepsilon_{m+1}=\frac{e^{-1}}{m+1}\cdot e^{-m}\frac{m^m}{m!}
\]
for all integers $j\ge 1$ and $m\ge 0$. By Stirling's formula $n!\sim\sqrt{2\pi n}\,(n/e)^n$, the Poisson mass at $n$ is asymptotic to $c_\ast/\sqrt n$, where
\[
c_\ast=\frac{1}{\sqrt{2\pi}}=0.398942\ldots,\qquad\text{so that}\qquad 0.39894\le c_\ast\le 0.39895.
\]
We use Stirling's formula with the bounds of Robbins \cite[p.~26, (1) and (2)]{robbins1955}: for every integer $n\ge 1$,
\begin{equation}\label{eq:prelim-robbins}
\sqrt{2\pi n}\,\Bigl(\frac{n}{e}\Bigr)^{n}e^{1/(12n+1)}<n!<\sqrt{2\pi n}\,\Bigl(\frac{n}{e}\Bigr)^{n}e^{1/(12n)}.
\end{equation}
Since $e^{-n}n^n/n!=(n/e)^n/n!$ and $e^{-1/(12n)}\ge 1-\frac{1}{12n}$, these bounds enclose the Poisson mass at $n$ between $\bigl(1-\frac{1}{12n}\bigr)\,c_\ast/\sqrt n$ and $c_\ast/\sqrt n$:
\begin{equation}\label{eq:prelim-stirling}
\Bigl(1-\frac{1}{12n}\Bigr)\frac{1}{\sqrt{2\pi n}}\le e^{-n}\,\frac{n^n}{n!}\le\frac{1}{\sqrt{2\pi n}}\qquad(n\ge 1).
\end{equation}

\subsection{Two sums}\label{subsec:prelim-sums}

We finally need the values of two series: $\sum_{j}e^{-j}a_j=1$ and $\sum_{m}e^{-m}\varepsilon_m=1-e^{-1}$. The first is classical. The numbers $e^{-j}a_j=e^{-j}j^{j-1}/j!$ with $j\ge 1$ are the probabilities of the Borel distribution with parameter $1$, which is the law of the total progeny of a Galton--Watson process whose offspring law is the Poisson law with mean $1$, and this total progeny is finite with probability $1$ \cite[Section~2.2 and (10)]{aldouspitman1998}. Hence the series of nonnegative real numbers $\sum_{j\ge 0}e^{-j}a_j$ converges, and
\begin{equation}\label{eq:prelim-borel}
\sum_{j\ge 0}e^{-j}\,a_j=1.
\end{equation}
The second sum is the one that the later sections use, and we deduce it from the first.

Both sums are values at $X=e^{-1}$ of series that Section~\ref{subsec:prelim-series} treats formally. By \eqref{eq:prelim-ratio} the ratio $a_{j+1}/a_j$ tends to $e$, so that $e^{-1}$ is the radius of convergence of $XC$. The second value is therefore not obtained by substituting $e^{-1}$ for $X$ in an identity that holds inside the disc of convergence, and we argue with partial sums instead, by means of the following elementary fact about a renewal equation.

Let $(f_m)_{m\ge 0}$ and $(u_n)_{n\ge 0}$ be sequences of nonnegative real numbers such that $f_0=0$, $u_0=1$ and the renewal equation
\[
u_n=\sum_{m=0}^{n}f_m\,u_{n-m}\qquad(n\ge 1)
\]
holds, and write $F_M=\sum_{m=0}^{M}f_m$ and $U_M=\sum_{n=0}^{M}u_n$ for the partial sums. Then, for all integers $M\ge 0$ and $k\ge 0$,
\begin{equation}\label{eq:prelim-partial}
1-\frac{1}{U_M}\le F_M\qquad\text{and}\qquad F_M\,U_k\le U_{M+k}-1.
\end{equation}
To see this, we sum the renewal equation over $1\le n\le M$ and exchange the order of summation. Since $f_0=0$, this gives
\[
U_M-1=\sum_{m=0}^{M}f_m\,U_{M-m}\qquad(M\ge 0).
\]
The sequence $(U_M)$ is nondecreasing, and $U_0=1$. Bounding $U_{M-m}$ by $U_M$ gives the first of the two inequalities. For the second we use the last display with $M+k$ in place of $M$, discard the terms with $m>M$, and bound $U_{M+k-m}$ from below by $U_k$ for $m\le M$.

\begin{lemma}\label{lem:epsmass}
The series of nonnegative real numbers
\[
\sum_{m\ge 0}e^{-m}\,\varepsilon_m
\]
converges, and its sum is $1-e^{-1}$.
\end{lemma}

\begin{proof}
We apply the inequalities \eqref{eq:prelim-partial} to $f_m=e^{-m}\varepsilon_m$ and $u_n=e^{-n}c_n$, for which $f_0=0$ and $u_0=1$; here the renewal equation is the equation of Corollary~\ref{cor:crenewal} multiplied by $e^{-n}$. Since $c_n=a_{n+1}$ and $a_0=0$, the sum \eqref{eq:prelim-borel} says that $\sum_{n\ge 0}u_n=e$. Hence $U_k\le e$ for every $k$, and $U_k\to e$.

The second of the two inequalities that we apply now gives $F_M\,U_k\le U_{M+k}-1\le e-1$ for all $M$ and $k$, and letting $k\to\infty$ we get $F_M\le 1-e^{-1}$ for every $M\ge 0$. The first inequality gives $F_M\ge 1-1/U_M$, and the right-hand side tends to $1-e^{-1}$ as $M\to\infty$. Hence $F_M\to 1-e^{-1}$.
\end{proof}

\section{The last-car recursion}\label{sec:recursion}
In this section we bring the last-car recursion \eqref{eq:intro-lastcar}, the one relation that the proof of Theorem~\ref{thm:main} takes from the parking process, into the form in which it is used: a recursion for the counts $A_{n,k}$ divided by $n!$. The last subsection shows that the ratio of these normalised counts to the numbers $c_n$ of Section~\ref{sec:prelim} does not decrease along the diagonal (Proposition~\ref{prop:diagonalmonotone}).

\subsection{The recursion}\label{subsec:recursion-lastcar}

The last-car recursion expresses $A_{n+1,k}$ through the numbers $A_{j,k}$ with $j\le n$, for the same bound $k$ on the displacements.

\begin{proposition}[{\cite[Theorem~3.7]{aguilarfraga2024}}]\label{prop:lastcar}
For all integers $n\ge 0$ and $k\ge 0$,
\[
A_{n+1,k}=\sum_{j=0}^{n}\min(j+1,k+1)\binom{n}{j}A_{j,k}\,A_{n-j,k}.
\]
\end{proposition}

The recursion is obtained by sorting the parking functions of length $n+1$ by the spot $s=j+1$ in which car~$n+1$ parks. The $j$ earlier cars that park on the left of $s$ can be chosen in $\binom{n}{j}$ ways. Their preferences form a parking function of length $j$, and those of the other $n-j$ earlier cars, reduced by $s$, form a parking function of length $n-j$; each of these cars has the same displacement there as in the given parking function. Car~$n+1$ prefers one of the spots $s-t$ with $0\le t\le j$, and its displacement is $t$, so that $\min(j+1,k+1)$ of its possible preferences are compatible with the bound $k$.

\subsection{The step of a column and the normalised counts}\label{subsec:recursion-normalised}

The recursion concerns the cumulative counts $A_{n,k}$, whereas Theorem~\ref{thm:main} is a statement about the signs of the steps $d_{n,k}$. Since $b_{n,\ell}=A_{n,\ell}-A_{n,\ell-1}$ for $\ell\ge 1$, a step is a second difference of the cumulative counts:
\begin{equation}\label{eq:recursion-second}
b_{n,k}-b_{n,k-1}=A_{n,k}-2A_{n,k-1}+A_{n,k-2}\qquad(k\ge 2,\ n\ge 0).
\end{equation}

The analytic arguments of the paper work with the counts divided by $n!$. An integer $y\ge 1$ is called a \emph{threshold}; it stands for the condition that all displacements are smaller than $y$, i.e.\ $M(\alpha)\le y-1$. For a threshold $y$ and an integer $n\ge 0$ let
\[
z^{(y)}_n=\frac{A_{n,y-1}}{n!}
\]
be the number of parking functions of length $n$ all of whose displacements are smaller than $y$, divided by $n!$. We call $z^{(y)}_n$ the \emph{normalised count} of threshold $y$ at length $n$. Since $A_{0,k}=1$ and $A_{n,k}$ is nondecreasing in $k$,
\begin{equation}\label{eq:recursion-monotone}
z^{(y)}_0=1\qquad\text{and}\qquad 0\le z^{(y)}_n\le z^{(y')}_n\quad\text{for } 1\le y\le y' \text{ and } n\ge 0.
\end{equation}

In the same way, for integers $\ell\ge 0$ and $n\ge 0$ let
\[
\beta^{(\ell)}_n=\frac{b_{n,\ell}}{n!}\ge 0
\]
be the normalised count of level $\ell$, and for integers $k\ge 1$ and $n\ge 0$ let
\[
\delta^{(k)}_n=\beta^{(k)}_n-\beta^{(k-1)}_n=\frac{d_{n,k}}{n!}
\]
be the \emph{normalised step} of column $k$ in row $n$, which has the sign of $d_{n,k}$. When the column is fixed, we write $\delta_n$ for $\delta^{(k)}_n$. Since $b_{n,0}=A_{n,0}$ and $b_{n,\ell}=A_{n,\ell}-A_{n,\ell-1}$ for $\ell\ge 1$, we have
\begin{equation}\label{eq:recursion-levels}
\beta^{(0)}_n=z^{(1)}_n\qquad\text{and}\qquad\beta^{(\ell)}_n=z^{(\ell+1)}_n-z^{(\ell)}_n\quad\text{for }\ell\ge 1,
\end{equation}
and \eqref{eq:recursion-second}, divided by $n!$, reads
\[
\delta^{(k)}_n=z^{(k+1)}_n-2\,z^{(k)}_n+z^{(k-1)}_n\qquad(k\ge 2).
\]
Thus a column $k\ge 2$ involves the three thresholds $k+1$, $k$ and $k-1$ together.

Divided by $n!$, the recursion of Proposition~\ref{prop:lastcar} loses its binomial coefficient and becomes a recursion for the sequence $(z^{(y)}_n)_{n\ge 0}$, into which the threshold enters only through the factor $\min(j+1,y)$. Without a bound on the displacements this factor is $j+1$, and the recursion is that of the $c_n$ (cf.\ Lemma~\ref{lem:crecursion}). We shall compare solutions of the two recursions, and such a comparison uses of the factor only the range in which it agrees with $j+1$; we therefore state the recursion for a general weight.

\begin{definition}\label{def:weighted}
Let $\omega=(\omega(t))_{t\ge 1}$ and $W=(W_n)_{n\ge 0}$ be sequences of real numbers. The sequence $W$ is a \emph{solution of the weighted last-car recursion with weight $\omega$} if $W_0=1$ and, for every integer $n\ge 0$,
\[
(n+1)\,W_{n+1}=\sum_{j=0}^{n}\omega(j+1)\,W_j\,W_{n-j}.
\]
\end{definition}

The recursion determines $W_{n+1}$ from $W_0,\dots,W_n$ and $\omega(1),\dots,\omega(n+1)$, so that every weight has exactly one solution. By Lemma~\ref{lem:crecursion} and $c_0=1$, the sequence $(c_n)_{n\ge 0}$ is the solution with weight $\omega(t)=t$.

\begin{corollary}\label{cor:normalised}
Let $y\ge 1$ be an integer. Then the sequence $z^{(y)}=(z^{(y)}_n)_{n\ge 0}$ is a solution of the weighted last-car recursion with weight $\omega(t)=\min(t,y)$, i.e.\ $z^{(y)}_0=1$ and, for every integer $n\ge 0$,
\begin{equation}\label{eq:recursion-threshold}
(n+1)\,z^{(y)}_{n+1}=\sum_{j=0}^{n}\min(j+1,y)\,z^{(y)}_j\,z^{(y)}_{n-j}.
\end{equation}
\end{corollary}

\begin{proof}
Since $z^{(y)}_0=A_{0,y-1}=1$, this is Proposition~\ref{prop:lastcar} with $k=y-1$, divided by $n!$.
\end{proof}

The two factors $z^{(y)}_j$ and $z^{(y)}_{n-j}$ of a term of this recursion belong to the cars on the left, resp.\ on the right, of the last car. Their lengths $j$ and $n-j$ add up to $n$, one less than the length on the left-hand side.

Let $x$ be a real number. The \emph{tilted form} of a sequence $(s_n)_{n\ge 0}$ at the \emph{tilt} $x$ is the sequence $(s_nx^n)_{n\ge 0}$; in particular, the numbers $z^{(y)}_nx^n$ are the tilted counts of threshold $y$ at tilt $x$. For $x>0$ a tilt changes no sign, so that a comparison within one row is the same statement before and after tilting, and the tilted counts are nonnegative and nondecreasing in the threshold, as the $z^{(y)}_n$ themselves are. What a tilt changes is the comparison of different rows, since the ratio of two consecutive terms is multiplied by $x$; this is what the tilts of Sections~\ref{sec:critical} and~\ref{sec:crossing} are chosen for. The recursion of Corollary~\ref{cor:normalised} keeps its form under a tilt: multiplying \eqref{eq:recursion-threshold}, with $N$ in place of $n$, by $x^{N+1}$, we see that for every integer $y\ge 1$, every real number $x$ and every integer $N\ge 0$,
\begin{equation}\label{eq:recursion-tilted}
(N+1)\,z^{(y)}_{N+1}x^{N+1}=x\sum_{i=0}^{N}\min(i+1,y)\;z^{(y)}_i x^{i}\;z^{(y)}_{N-i}x^{N-i}.
\end{equation}

\subsection{The free range}\label{subsec:recursion-free}

The lengths $n\le y$ form the \emph{free range} of threshold $y$. Recall that for these lengths the condition that all displacements are smaller than $y$ is no restriction, so that $A_{n,y-1}$ is the number $(n+1)^{n-1}$ of all parking functions of length $n$ \cite[Section~2]{foatariordan1974}. In normalised form this is the following lemma, which also follows directly from the recursions of Corollary~\ref{cor:normalised} and Lemma~\ref{lem:crecursion}; we give this short proof.

\begin{lemma}\label{lem:freerange}
Let $y\ge 1$ and $n\ge 0$ be integers with $n\le y$. Then
\[
z^{(y)}_n=c_n.
\]
\end{lemma}

\begin{proof}
By Corollary~\ref{cor:normalised} the sequences $z^{(y)}$ and $(c_n)_{n\ge 0}$ are the solutions of the weighted last-car recursion with weights $\min(t,y)$ and $t$, which agree for $1\le t\le y$. Since the recursion determines $W_{n+1}$ from $W_0,\dots,W_n$ and $\omega(1),\dots,\omega(n+1)$, induction on $n$ gives $z^{(y)}_n=c_n$ for all $n\le y$.
\end{proof}

The lengths $y+1\le n\le 2y+1$ form the \emph{first layer} of threshold $y$. For $n$ in the first layer the lengths of the two factors of a term in the recursion \eqref{eq:recursion-threshold} for $z^{(y)}_n$ add up to $n-1\le 2y$, so that at most one of the two factors lies outside the free range; the other factor is one of the $c_m$.

In the free range the normalised counts of a level and the normalised steps of a column vanish: by the relations \eqref{eq:recursion-levels} and Lemma~\ref{lem:freerange},
\[
\beta^{(\ell)}_n=0\quad(\ell\ge 1,\ n\le\ell)\qquad\text{and hence}\qquad\delta^{(k)}_n=0\quad(k\ge 2,\ n<k).
\]

\subsection{Monotonicity along the diagonal}\label{subsec:recursion-diagonal}

Since $n!\,c_n$ is the number of parking functions of length $n$, the ratio $z^{(y)}_n/c_n$ is the probability that all displacements of a uniformly chosen parking function of length $n$ are smaller than $y$. One more car can push the maximum displacement up. The following proposition says that allowing one more unit of displacement at least compensates for this: the probability does not decrease when $n$ and $y$ are both raised by one. The diagonal lower bound of Section~\ref{sec:series} (Proposition~\ref{prop:diagonallower}) is derived from it. The proof of the proposition is an induction on $n$, in which the last-car recursion writes each of the probabilities as an average over the spot in which the last car parks. The averages for the two sides of the inequality differ in their weights and in their values, and the second has one term more, so that the proof compares each term of the first with two neighbouring terms of the second.

\begin{proposition}\label{prop:diagonalmonotone}
For all integers $y\ge 1$ and $n\ge 0$,
\[
\frac{z^{(y)}_n}{c_n}\le\frac{z^{(y+1)}_{n+1}}{c_{n+1}}.
\]
\end{proposition}

\begin{proof}
For integers $y\ge 1$ and $0\le j\le n$ put
\[
\begin{gathered}
F^{(y)}_n=\frac{z^{(y)}_n}{c_n},\qquad
q_n(j)=\frac{(j+1)\,c_j\,c_{n-j}}{(n+1)\,c_{n+1}},\\
v^{(y)}_n(j)=\frac{\min(j+1,y)}{j+1}\,F^{(y)}_j\,F^{(y)}_{n-j}.
\end{gathered}
\]
We prove that $F^{(y)}_n\le F^{(y+1)}_{n+1}$ by induction on $n$, for all $y\ge 1$ at once. The recursion of Corollary~\ref{cor:normalised} expresses $F^{(y)}_{n+1}$ as an average of the $v^{(y)}_n(j)$ with weights $q_n(j)$. Intuitively, $q_n(j)$ is the probability that the last car of a uniformly chosen parking function of length $n+1$ parks in spot $j+1$, and $v^{(y)}_n(j)$ is the probability, given this, that all displacements are smaller than $y$.

Since the $c_i$ are positive, the inequalities \eqref{eq:recursion-monotone} give
\[
0\le F^{(y)}_i\le F^{(y+1)}_i\qquad(y\ge 1,\ i\ge 0),
\]
so that the ratios $F^{(y)}_i$ are nonnegative and nondecreasing in the threshold.

We begin with the representation of $F^{(y)}_{n+1}$ as an average: for $y\ge 1$ and $n\ge 0$,
\[
F^{(y)}_{n+1}=\sum_{j=0}^{n}q_n(j)\,v^{(y)}_n(j)\qquad\text{and}\qquad\sum_{j=0}^{n}q_n(j)=1.
\]
Indeed, $q_n(j)\,v^{(y)}_n(j)=\min(j+1,y)\,z^{(y)}_j\,z^{(y)}_{n-j}/((n+1)\,c_{n+1})$, so that the first identity is the recursion \eqref{eq:recursion-threshold} divided by $(n+1)\,c_{n+1}$; the second is Lemma~\ref{lem:crecursion} divided by the same number.

Next we compare the weights of two consecutive averages. For $0\le j\le n$,
\[
q_{n+1}(j)\le q_n(j)\qquad\text{and}\qquad q_{n+1}(j+1)\le q_n(j).
\]
After cancelling $(j+1)\,c_j$ and clearing the denominators, the first inequality reads $(n+1)\,c_{n+1-j}\,c_{n+1}\le(n+2)\,c_{n-j}\,c_{n+2}$. Since $c_i=a_{i+1}$ for every $i\ge 0$, this follows from the inequality $a_{n-j+2}\,a_{n+2}\le a_{n-j+1}\,a_{n+3}$, which is Lemma~\ref{lem:cayleylogconvex} with $u=n-j+1$, $v=n+2$ and $s=1$. In the same way, with $c_{n-j}$ cancelled, the second inequality reads
\[
(j+2)\,c_{j+1}\cdot(n+1)\,c_{n+1}\le(j+1)\,c_j\cdot(n+2)\,c_{n+2}.
\]
By \eqref{eq:prelim-ratio}, we have $c_i=(1+\frac1i)^{i-1}c_{i-1}$ for $i\ge 1$, and hence $(i+1)\,c_i=i\,(1+\frac1i)^{i}\,c_{i-1}$. We insert the second formula for $i=j+1$ on the left of the displayed inequality and the first formula for $i=n+2$ on the right, and cancel $(j+1)\,c_j\,c_{n+1}$. This turns it into the inequality between the first and the last member of
\begin{align*}
(n+1)\Bigl(1+\frac{1}{j+1}\Bigr)^{j+1}
&\le(n+1)\Bigl(1+\frac{1}{n+1}\Bigr)^{n+1}\\
&=(n+2)\Bigl(1+\frac{1}{n+1}\Bigr)^{n}
\le(n+2)\Bigl(1+\frac{1}{n+2}\Bigr)^{n+1}.
\end{align*}
Here the first step holds because $j\le n$ and $(1+\frac1i)^i$ is nondecreasing in $i\ge 1$ (Section~\ref{sec:prelim}), and the last step is Lemma~\ref{lem:powermono}.

The two averages will be compared by means of a general inequality between two weighted sums, which we prove before the induction. Let $m\ge 0$ be an integer. Let $q(j)$ and $v(j)$, for $0\le j\le m$, and $q'(j)$ and $v'(j)$, for $0\le j\le m-1$, be real numbers such that
\[
\sum_{j=0}^{m}q(j)=\sum_{j=0}^{m-1}q'(j)
\]
and, for every $0\le j\le m-1$,
\[
q(j)\le q'(j),\qquad q(j+1)\le q'(j),\qquad v'(j)\le v(j),\qquad v'(j)\le v(j+1).
\]
We claim that then
\[
\sum_{j=0}^{m-1}q'(j)\,v'(j)\le\sum_{j=0}^{m}q(j)\,v(j).
\]
For the proof let $\sigma_t=\sum_{j<t}q(j)$ for $0\le t\le m+1$ and $\sigma'_t=\sum_{j<t}q'(j)$ for $0\le t\le m$. These partial sums interlace:
\[
\sigma_t\le\sigma'_t\le\sigma_{t+1}\qquad(0\le t\le m).
\]
The first inequality follows from $q(j)\le q'(j)$, term by term. For the second, the inequalities $q(j+1)\le q'(j)$ give
\[
\sigma'_m-\sigma'_t=\sum_{t\le j<m}q'(j)\ge\sum_{t\le j<m}q(j+1)=\sigma_{m+1}-\sigma_{t+1},
\]
and $\sigma'_m=\sigma_{m+1}$ by the hypothesis on the two total sums. We now show by induction on $t$ that, for $0\le t\le m$,
\[
\sum_{j<t}q'(j)\,v'(j)\le\sum_{j<t}q(j)\,v(j)+(\sigma'_t-\sigma_t)\,v(t).
\]
The case $t=0$ is trivial. For the step from $t$ to $t+1\le m$ we split $q'(t)=(\sigma_{t+1}-\sigma'_t)+(\sigma'_{t+1}-\sigma_{t+1})$ into two parts, which are nonnegative because the partial sums interlace, and we use $v'(t)\le v(t)$ on the first part and $v'(t)\le v(t+1)$ on the second:
\[
q'(t)\,v'(t)\le(\sigma_{t+1}-\sigma'_t)\,v(t)+(\sigma'_{t+1}-\sigma_{t+1})\,v(t+1).
\]
When this is added to the induction hypothesis, the terms with $v(t)$ combine to $(\sigma_{t+1}-\sigma_t)\,v(t)=q(t)\,v(t)$, and the result is the inequality for $t+1$. For $t=m$ the last term of the inequality so proved is $(\sigma_{m+1}-\sigma_m)\,v(m)=q(m)\,v(m)$, so that the case $t=m$ is the claim.

We now carry out the induction on $n$. For $n=0$ both sides of the inequality $F^{(y)}_n\le F^{(y+1)}_{n+1}$ are $1$, by Lemma~\ref{lem:freerange}. Let $n\ge 0$, and suppose that $F^{(y')}_i\le F^{(y'+1)}_{i+1}$ for all $i\le n$ and all $y'\ge 1$. We show that $F^{(y)}_{n+1}\le F^{(y+1)}_{n+2}$ for every $y\ge 1$. Recall that $F^{(y)}_{n+1}$ is the average of the values $v^{(y)}_n(j)$, $0\le j\le n$, with the weights $q_n(j)$, and that $F^{(y+1)}_{n+2}$ is the average of the values $v^{(y+1)}_{n+1}(j)$, $0\le j\le n+1$, with the weights $q_{n+1}(j)$. We apply the general inequality with $m=n+1$, with $q=q_{n+1}$ and $q'=q_n$, and with $v=v^{(y+1)}_{n+1}$ and $v'=v^{(y)}_n$. The two total sums are $1$, because the weights of each of the two averages add up to $1$, as we saw when we wrote $F^{(y)}_{n+1}$ as an average, and the two conditions on the weights are the inequalities $q_{n+1}(j)\le q_n(j)$ and $q_{n+1}(j+1)\le q_n(j)$ proved above. It remains to verify the two conditions on the values, for $0\le j\le n$. Each value is a product of three nonnegative factors, and we compare the factors one by one.

The first condition is $v^{(y)}_n(j)\le v^{(y+1)}_{n+1}(j)$. Here $\min(j+1,y)\le\min(j+1,y+1)$, we have $F^{(y)}_j\le F^{(y+1)}_j$ by the monotonicity in the threshold, and we have $F^{(y)}_{n-j}\le F^{(y+1)}_{n+1-j}$ by the induction hypothesis for $i=n-j$.

The second condition, $v^{(y)}_n(j)\le v^{(y+1)}_{n+1}(j+1)$, is verified similarly. With $t=\min(j+1,y)$ we have $\min(j+2,y+1)=t+1$, so that the first factors of the two values are $\frac{t}{j+1}$ and $\frac{t+1}{j+2}$, and $\frac{t}{j+1}\le\frac{t+1}{j+2}$ because $t\le j+1$. Moreover $F^{(y)}_j\le F^{(y+1)}_{j+1}$ by the induction hypothesis for $i=j$, and $F^{(y)}_{n-j}\le F^{(y+1)}_{n-j}$ by the monotonicity in the threshold.

Hence the general inequality gives $F^{(y)}_{n+1}\le F^{(y+1)}_{n+2}$, which completes the induction.
\end{proof}

\section{The chain range and the column equation}\label{sec:chain}
The last-car recursion gives each normalised count as a sum of products of two earlier ones. Up to a certain length, however, it can be replaced by a linear relation, and this relation, combined for the three thresholds $k+1$, $k$ and $k-1$, is the column equation: an exact linear recursion for the normalised step $\delta^{(k)}_n$ of a column $k\ge 2$ on its chain range $k+1\le n\le 2k-1$ (Proposition~\ref{prop:columnequation}). The crossing criterion and the proof that a positive column stays positive in its chain range (Proposition~\ref{prop:chainrange}) start from this equation.

\subsection{The single-chain formula}\label{subsec:chain-formula}

Recall that, by Lemma~\ref{lem:crecursion} and $c_0=1$, the sequence $(c_n)$ is the solution of the weighted last-car recursion (Definition~\ref{def:weighted}) with weight $\omega(t)=t$. We consider a solution $W$ whose weight $\omega$ agrees with this weight up to some integer $k\ge 0$, i.e.\ $\omega(t)=t$ for $1\le t\le k$. The recursion is quadratic in $W$, but for $n\le 2k+1$ the deviation $c_n-W_n$ satisfies a linear relation with explicit coefficients, the \emph{single-chain formula}. We state it in two forms. The first gives $c_n-W_n$ in terms of the earlier values of $W$; for $n\le k$ its sum is empty, and it says that $W_n=c_n$.

\begin{proposition}\label{prop:forward}
Let $\omega=(\omega(t))_{t\ge 1}$ and $W=(W_n)_{n\ge 0}$ be sequences of real numbers such that $W$ is a solution of the weighted last-car recursion with weight $\omega$ (Definition~\ref{def:weighted}), and let $k\ge 0$ be an integer such that $\omega(t)=t$ for every integer $t$ with $1\le t\le k$. Then, for every integer $n$ with $0\le n\le 2k+1$,
\[
c_n-W_n=\sum_{t=k+1}^{n}c_{n-t}\Bigl(c_{t-1}-\frac{\omega(t)}{t}\,W_{t-1}\Bigr).
\]
\end{proposition}

\begin{proof}
By Definition~\ref{def:weighted} we have $W_0=1$ and
\[
(n+1)\,W_{n+1}=\sum_{j=0}^{n}\omega(j+1)\,W_j\,W_{n-j}\qquad(n\ge 0).
\]
We write
\[
\mathcal{W}=\sum_{n\ge 0}W_nX^n,\qquad U=\frac{\mathcal{W}}{C}=\sum_{i\ge 0}U_iX^i,\qquad L=\sum_{t\ge 1}L_tX^t,
\]
where $C=\sum_{n\ge 0}c_nX^n$ is the series of Section~\ref{sec:prelim}, and where
\[
L_t=c_{t-1}-\frac{\omega(t)}{t}\,W_{t-1}\qquad(t\ge 1).
\]
Thus $t\,L_t$ is the difference between the factor $t\,c_{t-1}$ of the recursion of Lemma~\ref{lem:crecursion} and the corresponding factor $\omega(t)\,W_{t-1}$ of the recursion for $W$.

We first show that
\[
W_n=c_n\quad(0\le n\le k)\qquad\text{and}\qquad L_t=0\quad(1\le t\le k).
\]
For $n<k$ every weight that occurs in the recursion for $W_{n+1}$ is $\omega(j+1)=j+1$, so that $W_0,\dots,W_k$ satisfy the recursion of Lemma~\ref{lem:crecursion}; since $W_0=c_0=1$, the first assertion follows by induction on $n$. The second follows from the first by the definition of $L_t$, since $\omega(t)=t$ for $1\le t\le k$.

Next we show that $U'=-L'\,U$. In terms of series, the recursion for $W$ reads
\[
\mathcal{W}'=\Bigl(\sum_{j\ge 0}\omega(j+1)\,W_j\,X^j\Bigr)\,\mathcal{W},
\]
and since $(XC)'=\sum_{j\ge 0}(j+1)\,c_jX^j$, the definition of $L_t$ gives
\[
L'=(XC)'-\sum_{j\ge 0}\omega(j+1)\,W_j\,X^j .
\]
Since $(1/C)'=-(1/C)\,(XC)'$ by \eqref{eq:prelim-inverse}, the product rule and the last two displays give
\begin{align*}
U'&=\frac1C\,\mathcal{W}'+\mathcal{W}\,\Bigl(\frac1C\Bigr)'\\
&=\frac1C\,\Bigl(\sum_{j\ge 0}\omega(j+1)\,W_j\,X^j\Bigr)\,\mathcal{W}-\frac1C\,(XC)'\,\mathcal{W}
=-L'\,U .
\end{align*}
Comparing the coefficients of $X^i$ on the two sides, we obtain
\[
(i+1)\,U_{i+1}=-\sum_{j=0}^{i}(j+1)\,L_{j+1}\,U_{i-j}\qquad(i\ge 0).
\]
Moreover $U_0=W_0/c_0=1$.

From this recursion we now determine the coefficients of $U$ up to degree $2k+1$. We claim that
\begin{equation}\label{eq:chain-linear}
U_i=0\quad(1\le i\le k)\qquad\text{and}\qquad U_i=-L_i\quad(1\le i\le 2k+1).
\end{equation}
With $i-1$ in place of $i$, the recursion reads
\[
i\,U_i=-\sum_{j=0}^{i-1}(j+1)\,L_{j+1}\,U_{i-1-j}\qquad(i\ge 1).
\]
For $1\le i\le k$ every coefficient $L_{j+1}$ on the right has $j+1\le k$ and vanishes, since $L_t=0$ for $1\le t\le k$, so that $U_i=0$. Now let $1\le i\le 2k+1$. The term with $j=i-1$ is $i\,L_i\,U_0=i\,L_i$. In a term with $j<i-1$, either $j+1\le k$, and then $L_{j+1}=0$ as before; or $j\ge k$, and then $1\le i-1-j\le k$, since $i\le 2k+1$, so that $U_{i-1-j}=0$ by the first part of the claim. Hence $i\,U_i=-i\,L_i$, which proves the claim. For $i=2k+2$ the argument no longer applies: the term with $j=k$ contains $L_{k+1}U_{k+1}=-L_{k+1}^2$, which need not vanish.

It remains to return from $U$ to $W$. Let $0\le n\le 2k+1$. Comparing the coefficients of $X^n$ in $\mathcal{W}=U\,C$ and using $U_0=1$ and the claim just proved, we find that
\[
W_n=\sum_{t=0}^{n}U_t\,c_{n-t}=c_n-\sum_{t=k+1}^{n}L_t\,c_{n-t},
\]
which is the assertion by the definition of $L_t$.
\end{proof}

The second form states the same relation for the coefficients of the quotient by $C$, in which $W$ is combined with the free kernel of Section~\ref{sec:prelim}.

\begin{proposition}\label{prop:quotientform}
Let $\omega$, $W$ and $k$ be as in Proposition~\ref{prop:forward}. Then, for every integer $n$ with $1\le n\le 2k+1$,
\[
W_n-\sum_{m=1}^{n}\varepsilon_m\,W_{n-m}=
\begin{cases}
0, & n\le k,\\[4pt]
\dfrac{\omega(n)}{n}\,W_{n-1}-c_{n-1}, & n>k.
\end{cases}
\]
\end{proposition}

\begin{proof}
We keep the notation $\mathcal{W}$, $U$ and $L$ of the proof of Proposition~\ref{prop:forward}. By Lemma~\ref{lem:quotient}, applied to $\mathcal{W}$, the coefficient of $X^n$ in $U=\mathcal{W}/C$ is
\[
U_n=W_n-\sum_{m=1}^{n}\varepsilon_m\,W_{n-m}.
\]
By the relations \eqref{eq:chain-linear} this is $0$ for $1\le n\le k$, and for $k<n\le 2k+1$ it is $-L_n=\frac{\omega(n)}{n}\,W_{n-1}-c_{n-1}$.
\end{proof}

\subsection{The level equation and the column equation}\label{subsec:chain-column}

We now apply the single-chain formula to the normalised counts of Section~\ref{sec:recursion}. Let $y\ge 1$ be a threshold. By Corollary~\ref{cor:normalised}, $z^{(y)}$ is a solution of the weighted last-car recursion with weight $\omega(t)=\min(t,y)$, which equals $t$ for $1\le t\le y$. Hence Propositions~\ref{prop:forward} and~\ref{prop:quotientform} apply to $W=z^{(y)}$ with $y$ in place of $k$, and the second of them gives, for every integer $n$ with $y\le n\le 2y+1$,
\begin{equation}\label{eq:chain-threshold}
z^{(y)}_n-\sum_{m=1}^{n}\varepsilon_m\,z^{(y)}_{n-m}=\frac{y}{n}\,z^{(y)}_{n-1}-c_{n-1}.
\end{equation}
Indeed, for $n>y$ this is the second case of the proposition with $\omega(n)=y$, and for $n=y$ both sides vanish, the right-hand side by Lemma~\ref{lem:freerange}.

This relation is linear in $z^{(y)}$, and the term $c_{n-1}$ on its right-hand side does not depend on $y$. It can therefore be subtracted for neighbouring thresholds, and the differences that arise are the normalised counts of the levels and the normalised steps of the columns. Recall from Section~\ref{sec:recursion} that
\begin{align*}
\beta^{(\ell)}_n&=z^{(\ell+1)}_n-z^{(\ell)}_n &&(\ell\ge 1,\ n\ge 0),\\
\beta^{(\ell)}_n&=0 &&(\ell\ge 1,\ 0\le n\le\ell),\\
\delta^{(k)}_n&=z^{(k+1)}_n-2\,z^{(k)}_n+z^{(k-1)}_n &&(k\ge 2,\ n\ge 0),\\
\delta^{(k)}_n&=0 &&(k\ge 2,\ 0\le n<k);
\end{align*}
here the third line is \eqref{eq:recursion-second} divided by $n!$.

For two neighbouring thresholds $\ell+1$ and $\ell$, the difference of the relations \eqref{eq:chain-threshold} is the \emph{level equation}, which determines $\beta^{(\ell)}_n$ row by row for $\ell+1\le n\le 2\ell+1$.

\begin{lemma}\label{lem:levelequation}
Let $\ell\ge 1$ and $n$ be integers with $\ell+1\le n\le 2\ell+1$. Then
\[
\beta^{(\ell)}_n=\frac{z^{(\ell+1)}_{n-1}}{n}+\frac{\ell}{n}\,\beta^{(\ell)}_{n-1}+\sum_{m=1}^{n-\ell-1}\varepsilon_m\,\beta^{(\ell)}_{n-m}.
\]
\end{lemma}

\begin{proof}
We subtract \eqref{eq:chain-threshold} for $y=\ell$ from the same relation for $y=\ell+1$; both apply in the stated range of $n$. The terms $c_{n-1}$ cancel, and since $\beta^{(\ell)}_m=z^{(\ell+1)}_m-z^{(\ell)}_m$ for every $m\ge 0$, we obtain
\[
\beta^{(\ell)}_n-\sum_{m=1}^{n}\varepsilon_m\,\beta^{(\ell)}_{n-m}
=\frac{\ell+1}{n}\,z^{(\ell+1)}_{n-1}-\frac{\ell}{n}\,z^{(\ell)}_{n-1}
=\frac{z^{(\ell+1)}_{n-1}}{n}+\frac{\ell}{n}\,\beta^{(\ell)}_{n-1}.
\]
The terms with $m\ge n-\ell$ vanish, because $\beta^{(\ell)}_{n-m}=0$ for $n-m\le\ell$, and the assertion follows.
\end{proof}

We turn to the step of a column $k\ge 2$. As recalled above, it vanishes in rows $n<k$, and Proposition~\ref{prop:forward} gives its value in row $k$.

\begin{lemma}\label{lem:initialstep}
For every integer $k\ge 2$,
\[
\delta^{(k)}_k=-\frac{a_k}{k}.
\]
\end{lemma}

\begin{proof}
We apply Proposition~\ref{prop:forward} to $W=z^{(k-1)}$, with $k-1$ in place of $k$ and with $n=k\le 2(k-1)+1$. The sum consists of the single term $t=k$, in which $\omega(k)=k-1$, and since $c_0=1$ and $z^{(k-1)}_{k-1}=c_{k-1}$ (Lemma~\ref{lem:freerange}), it gives
\[
c_k-z^{(k-1)}_k=c_{k-1}-\frac{k-1}{k}\,c_{k-1}=\frac{c_{k-1}}{k}.
\]
On the other hand, $\beta^{(k)}_k=0$, as recalled above, and $\beta^{(k-1)}_k=z^{(k)}_k-z^{(k-1)}_k=c_k-z^{(k-1)}_k$ by Lemma~\ref{lem:freerange}. Since $c_{k-1}=a_k$, we conclude that
\[
\delta^{(k)}_k=\beta^{(k)}_k-\beta^{(k-1)}_k=-\bigl(c_k-z^{(k-1)}_k\bigr)=-\frac{a_k}{k}.\qedhere
\]
\end{proof}

Multiplied by $k!$, the lemma reads $d_{k,k}=-k^{k-2}$, so that every column $k\ge 2$ begins with a negative step. From row $k+1$ on, the step satisfies a linear recursion as long as the relation \eqref{eq:chain-threshold} is available for the three thresholds $k+1$, $k$ and $k-1$, i.e.\ on the chain range $k+1\le n\le 2k-1$ of the column. This recursion is the column equation stated in Section~\ref{sec:intro}.

\begin{proposition}\label{prop:columnequation}
Let $k\ge 2$ and $n$ be integers with $k+1\le n\le 2k-1$. Then the normalised step of column $k$ satisfies
\[
\delta_n=\frac{k+1}{n}\,\delta_{n-1}+\frac{2}{n}\,\beta^{(k-1)}_{n-1}+\sum_{m=1}^{n-k}\varepsilon_m\,\delta_{n-m}.
\]
\end{proposition}

\begin{proof}
We argue as in the proof of Lemma~\ref{lem:levelequation}, now with three thresholds. For $k+1\le n\le 2k-1$ the relation \eqref{eq:chain-threshold} holds for each of the thresholds $y=k+1$, $k$ and $k-1\ge 1$. We take it for $y=k+1$, subtract twice the relation for $y=k$, and add the relation for $y=k-1$. The terms $c_{n-1}$ cancel, and we obtain
\begin{align*}
\delta_n-\sum_{m=1}^{n}\varepsilon_m\,\delta_{n-m}
&=\frac{k+1}{n}\,z^{(k+1)}_{n-1}-\frac{2k}{n}\,z^{(k)}_{n-1}+\frac{k-1}{n}\,z^{(k-1)}_{n-1}\\
&=\frac{k+1}{n}\,\delta_{n-1}+\frac{2}{n}\,\bigl(z^{(k)}_{n-1}-z^{(k-1)}_{n-1}\bigr)
=\frac{k+1}{n}\,\delta_{n-1}+\frac{2}{n}\,\beta^{(k-1)}_{n-1},
\end{align*}
where the first two equalities follow from the expression of the step as a second difference, $\delta_m=z^{(k+1)}_m-2\,z^{(k)}_m+z^{(k-1)}_m$, and the third from $\beta^{(k-1)}_{n-1}=z^{(k)}_{n-1}-z^{(k-1)}_{n-1}$. The terms with $m>n-k$ vanish, because $\delta_{n-m}=0$ for $n-m<k$, and the assertion follows.
\end{proof}

\section{The series of a threshold and three inequalities on all rows}\label{sec:series}
The column equation of Proposition~\ref{prop:columnequation} is confined to the chain range of a column. Beyond it no exact linear equation is available, and the proof of Theorem~\ref{thm:main} rests there on three inequalities that hold in every row: the column inequality (Proposition~\ref{prop:columninequality}), the one-tree bound (Proposition~\ref{prop:onetree}) and the diagonal lower bound (Proposition~\ref{prop:diagonallower}). To prove them we pass to the generating series of the $z^{(y)}_n$, for which the last-car recursion is a differential equation, and compare the series of adjacent thresholds.

We use the notation of Section~\ref{sec:prelim} for formal power series in $X$ with real coefficients, which we call series. For series $F$ and $G$ we write $F\succeq 0$ if every coefficient of $F$ is nonnegative, and $F\succeq G$, or $G\preceq F$, if $F-G\succeq 0$. Sums and products of series that are $\succeq 0$ are $\succeq 0$.

\subsection{The series of a threshold}\label{subsec:series-threshold}

For a threshold $y\ge 1$ let
\[
Z_y=\sum_{n\ge 0}z^{(y)}_n\,X^n ,
\]
and
\[
\Gamma_y=\sum_{j\ge 0}\min(j+1,y)\,z^{(y)}_j\,X^j .
\]
The constant term of $Z_y$ is $z^{(y)}_0=1$. With these two series the last-car recursion becomes a differential equation.

\begin{lemma}\label{lem:seriesode}
For every integer $y\ge 1$,
\[
Z_y'=\Gamma_y\,Z_y .
\]
\end{lemma}

\begin{proof}
The coefficients of $X^n$ in $Z_y'$ and in $\Gamma_y\,Z_y$ are $(n+1)\,z^{(y)}_{n+1}$ and $\sum_{j=0}^{n}\min(j+1,y)\,z^{(y)}_j\,z^{(y)}_{n-j}$, and these are equal for every $n\ge 0$ by Corollary~\ref{cor:normalised}.
\end{proof}

\subsection{The prime kernel}\label{subsec:series-kernel}

The proofs of the three inequalities divide by the series $Z_y$ of a threshold, and this subsection describes the coefficients of its inverse. Let $y\ge 1$ be a threshold. Since the constant term of $Z_y$ is $1$, the inverse $1/Z_y$ exists. We put
\[
K_y=1-\frac{1}{Z_y} ,
\]
and we call the sequence of the coefficients of this series, $\kappa^{(y)}_m=[X^m]K_y$ for $m\ge 0$, the \emph{prime kernel} of threshold $y$. By definition
\[
Z_y\,(1-K_y)=1 .
\]
The constant term of $1/Z_y$ is $1$, so that $\kappa^{(y)}_0=0$, and for every series $F=\sum_{n\ge 0}F_nX^n$ we have
\[
[X^n]\bigl(F\,(1-K_y)\bigr)=F_n-\sum_{m=1}^{n}\kappa^{(y)}_m\,F_{n-m}\qquad(n\ge 0).
\]
For $F=Z_y$ the product is $1$, and this formula becomes
\begin{equation}\label{eq:series-renewal}
z^{(y)}_n=\sum_{m=1}^{n}\kappa^{(y)}_m\,z^{(y)}_{n-m}\qquad(n\ge 1).
\end{equation}
This equation has the form of the renewal equation $c_n=\sum_{m=1}^{n}\varepsilon_m\,c_{n-m}$ of Corollary~\ref{cor:crenewal}, with $z^{(y)}$ in the place of $(c_n)$ and the prime kernel in the place of the free kernel.

\begin{lemma}\label{lem:kernelode}
For every integer $y\ge 1$,
\[
K_y'=(1-K_y)\,\Gamma_y .
\]
\end{lemma}

\begin{proof}
Differentiating the identity $Z_y\,(1-K_y)=1$ and using Lemma~\ref{lem:seriesode}, we get $Z_y\,K_y'=Z_y'\,(1-K_y)=\Gamma_y\,Z_y\,(1-K_y)=\Gamma_y$, and multiplying by $1-K_y$ gives the claim.
\end{proof}

\begin{lemma}\label{lem:kernelfree}
Let $y\ge 1$ and $m$ be integers with $1\le m\le y$. Then
\[
\kappa^{(y)}_m=\varepsilon_m .
\]
\end{lemma}

\begin{proof}
For $1\le n\le y$ Lemma~\ref{lem:freerange} turns the equation \eqref{eq:series-renewal} into $c_n=\sum_{i=1}^{n}\kappa^{(y)}_i\,c_{n-i}$. Subtracting the equation $c_n=\sum_{i=1}^{n}\varepsilon_i\,c_{n-i}$ of Corollary~\ref{cor:crenewal}, we obtain
\[
\sum_{i=1}^{n}\bigl(\kappa^{(y)}_i-\varepsilon_i\bigr)\,c_{n-i}=0\qquad(1\le n\le y),
\]
and since $c_0=1$ the assertion follows by induction on $m$.
\end{proof}

Beyond the free range the prime kernel satisfies a recursion with positive coefficients.

\begin{lemma}\label{lem:kernelrecursion}
Let $y\ge 2$ and $N\ge y$ be integers. Then
\[
N\,\kappa^{(y)}_N=\sum_{j=1}^{y-1}(y-j)\,a_j\,\kappa^{(y)}_{N-j} .
\]
\end{lemma}

\begin{proof}
Let $H_y=\sum_{i=0}^{y-2}(y-1-i)\,c_i\,X^i$. Since $z^{(y)}_i=c_i$ for $i\le y-2$ by Lemma~\ref{lem:freerange}, a comparison of coefficients shows that $\Gamma_y=y\,Z_y-H_y$. By Lemma~\ref{lem:kernelode}, and since $Z_y\,(1-K_y)=1$, it follows that
\[
K_y'=(1-K_y)\,\Gamma_y=y\,Z_y\,(1-K_y)-(1-K_y)\,H_y=y-(1-K_y)\,H_y .
\]
We compare the coefficients of $X^{N-1}$. On the left the coefficient is $N\,\kappa^{(y)}_N$. On the right the constant $y$ contributes nothing, because $N-1\ge 1$, and since $H_y$ has degree at most $y-2\le N-2$, only the coefficients $-\kappa^{(y)}_m$ of $1-K_y$ with $m\ge 1$ enter the product. Hence
\[
N\,\kappa^{(y)}_N=\sum_{i=0}^{y-2}(y-1-i)\,c_i\,\kappa^{(y)}_{N-1-i} ,
\]
which is the assertion, with $j=i+1$ and $c_{j-1}=a_j$.
\end{proof}

For example, for $y=2$ the recursion reads $N\,\kappa^{(2)}_N=\kappa^{(2)}_{N-1}$ for $N\ge 2$, because $a_1=1$, and with $\kappa^{(2)}_1=\varepsilon_1=1$ it gives $\kappa^{(2)}_N=1/N!$ for all $N\ge 1$. In particular $\kappa^{(2)}_3=1/6$, whereas $\varepsilon_3=2/3$, so that the range $m\le y$ of Lemma~\ref{lem:kernelfree} cannot be extended.

The recursion of Lemma~\ref{lem:kernelrecursion} gives the sign of the prime kernel. We state the result for the tilted form $\kappa^{(y)}_m\,x^m$ of the kernel at a tilt $x\ge 0$; for $x=1$ it says that $\kappa^{(y)}_m\ge 0$.

\begin{lemma}\label{lem:kernelnonneg}
Let $y\ge 2$ be an integer and $x\ge 0$ a real number. Then
\[
\kappa^{(y)}_m\,x^m\ge 0\qquad\text{for every integer } m\ge 0 .
\]
\end{lemma}

\begin{proof}
We use induction on $m$. We have $\kappa^{(y)}_0=0$, and $\kappa^{(y)}_m\,x^m=\varepsilon_m\,x^m\ge 0$ for $1\le m\le y$ by Lemma~\ref{lem:kernelfree}. For $m>y$ the recursion of Lemma~\ref{lem:kernelrecursion} with $N=m$, multiplied by $x^m$, reads
\[
m\,\kappa^{(y)}_m\,x^m=\sum_{j=1}^{y-1}(y-j)\,a_j\,x^j\cdot\kappa^{(y)}_{m-j}\,x^{m-j} ,
\]
and every term on the right is nonnegative, by the induction hypothesis and since $x\ge 0$.
\end{proof}

\subsection{The column inequality and the one-tree bound}\label{subsec:series-column}

Let $k\ge 2$ be an integer. Recall from Section~\ref{sec:recursion} that $\beta^{(\ell)}_n=z^{(\ell+1)}_n-z^{(\ell)}_n$ for $\ell\ge 1$, so that
\[
\delta^{(k)}_n=\beta^{(k)}_n-\beta^{(k-1)}_n=z^{(k+1)}_n-2\,z^{(k)}_n+z^{(k-1)}_n\qquad(n\ge 0).
\]
Thus the normalised steps of column $k$ are the coefficients of the series
\[
\mathsf{D}=Z_{k+1}-2\,Z_k+Z_{k-1} ,
\]
the second difference of $Z_y$ with respect to $y$ at $y=k$. The first two inequalities bound these steps from below in terms of a second sequence. For a column $k\ge 2$ let $(\xi_t)_{t\ge 0}$ be the sequence defined by $\xi_0=0$ and, for $t\ge 1$, by
\[
t\,\xi_t=\min(t,k+1)\,z^{(k+1)}_{t-1}-2\min(t,k)\,z^{(k)}_{t-1}+\min(t,k-1)\,z^{(k-1)}_{t-1} .
\]
We write $\xi^{(k)}_t$ for $\xi_t$ when the column has to be shown. We have $\xi_t=0$ for $t<k$: for $1\le t\le k-1$ the three minima in the definition are equal to $t$, and $z^{(k+1)}_{t-1}=z^{(k)}_{t-1}=z^{(k-1)}_{t-1}=c_{t-1}$ by Lemma~\ref{lem:freerange}. From index $k+1$ on, $\xi_t$ is expressed by the step and by level $k-1$, in the combination that occurs in the column equation.

\begin{lemma}\label{lem:curvatureidentity}
Let $k\ge 2$ and $n\ge k$ be integers. Then
\[
(n+1)\,\xi^{(k)}_{n+1}=(k+1)\,\delta^{(k)}_n+2\,\beta^{(k-1)}_n .
\]
\end{lemma}

\begin{proof}
For $t=n+1\ge k+1$ the three minima in the definition of $\xi_t$ are $k+1$, $k$ and $k-1$, so that the left-hand side is $(k+1)\,z^{(k+1)}_n-2k\,z^{(k)}_n+(k-1)\,z^{(k-1)}_n$. Since $\delta^{(k)}_n=z^{(k+1)}_n-2\,z^{(k)}_n+z^{(k-1)}_n$ and $\beta^{(k-1)}_n=z^{(k)}_n-z^{(k-1)}_n$, so is the right-hand side.
\end{proof}

The first inequality extends the column equation of Proposition~\ref{prop:columnequation} to all rows $n\ge k+1$: the free kernel is replaced by the prime kernel of threshold $k$, and the equation becomes an inequality that bounds the step from below. In the chain range $k+1\le n\le 2k-1$ it holds with equality: there $m\le n-k\le k-1$, so that $\kappa^{(k)}_m=\varepsilon_m$ by Lemma~\ref{lem:kernelfree}, and the inequality is the column equation multiplied by $n$. The proof compares the series of two adjacent thresholds. Applied to the thresholds $k$ and $k+1$ and to the thresholds $k-1$ and $k$, the comparison bounds $\mathsf{D}\,(1-K_k)$ from below, and the proposition is the coefficient of $X^n$ in this bound.

\begin{proposition}\label{prop:columninequality}
Let $k\ge 2$ and $n\ge k+1$ be integers. Then
\[
(k+1)\,\delta^{(k)}_{n-1}+2\,\beta^{(k-1)}_{n-1}\le n\,\Bigl(\delta^{(k)}_n-\sum_{m=1}^{n-k}\kappa^{(k)}_m\,\delta^{(k)}_{n-m}\Bigr).
\]
\end{proposition}

\begin{proof}
We write $\delta_j=\delta^{(k)}_j$ and $\xi_t=\xi^{(k)}_t$.

We begin with a positivity principle. For series $F$, $G$ and $S$,
\[
\text{if}\quad F'=G\,F+S,\quad G\succeq 0,\quad S\succeq 0\quad\text{and}\quad [X^0]F\ge 0,\quad\text{then}\quad F\succeq 0 .
\]
Indeed, in terms of the coefficients $F_n$, $G_n$ and $S_n$ of the three series the equation reads $(n+1)\,F_{n+1}=\sum_{j=0}^{n}G_j\,F_{n-j}+S_n$, and $F_n\ge 0$ follows by induction on $n$. For $G=0$ the principle says that a series with nonnegative constant term whose derivative is $\succeq 0$ is $\succeq 0$.

Next we introduce the series with which two adjacent thresholds are compared, and we show that it is $\succeq 0$. Let $y\ge 2$ be an integer, and let $\Omega_y$ be the series with constant term $0$ and derivative $\Omega_y'=\Gamma_y-\Gamma_{y-1}$. Its coefficients are
\begin{equation}\label{eq:series-log-coefficients}
[X^t]\Omega_y=\frac{1}{t}\Bigl(\min(t,y)\,z^{(y)}_{t-1}-\min(t,y-1)\,z^{(y-1)}_{t-1}\Bigr)\qquad(t\ge 1).
\end{equation}
Let further $\mathcal{R}_y=\sum_{j\ge y-1}z^{(y-1)}_jX^j$ be the sum of the terms of $Z_{y-1}$ of degree $j\ge y-1$. Then
\[
\Omega_y'=y\,(Z_y-Z_{y-1})+\mathcal{R}_y .
\]
Indeed, the coefficient of $X^j$ on the left is $\min(j+1,y)\,z^{(y)}_j-\min(j+1,y-1)\,z^{(y-1)}_j$. For $j\le y-2$ the two minima are both $j+1$, and $z^{(y)}_j=z^{(y-1)}_j=c_j$ by Lemma~\ref{lem:freerange}, so that the coefficients on both sides vanish. For $j\ge y-1$ the minima are $y$ and $y-1$, and $y\,z^{(y)}_j-(y-1)\,z^{(y-1)}_j=y\,\bigl(z^{(y)}_j-z^{(y-1)}_j\bigr)+z^{(y-1)}_j$. Recall that $0\le z^{(y-1)}_n\le z^{(y)}_n$ for all $n\ge 0$. The two series $y\,(Z_y-Z_{y-1})$ and $\mathcal{R}_y$ are therefore $\succeq 0$, and with the positivity principle for $G=0$ we obtain
\[
\Omega_y'\succeq 0\qquad\text{and}\qquad \Omega_y\succeq 0 .
\]

We now compare the thresholds $y-1$ and $y$. For $y\ge 2$ let
\[
\mathcal{Q}_y=Z_y\,(1-K_{y-1})\qquad\text{and}\qquad \Theta_y=1-Z_{y-1}\,(1-K_y) .
\]
By the definition of $K_y$ we have $\mathcal{Q}_y=Z_y/Z_{y-1}$ and $1-\Theta_y=Z_{y-1}/Z_y$. The constant term of $\mathcal{Q}_y$ is $1$ and that of $\Theta_y$ is $0$, and by the product rule and Lemmas~\ref{lem:seriesode} and~\ref{lem:kernelode},
\begin{align*}
\mathcal{Q}_y'&=\Gamma_y\,Z_y\,(1-K_{y-1})-Z_y\,(1-K_{y-1})\,\Gamma_{y-1}=\mathcal{Q}_y\,\Omega_y' ,\\
\Theta_y'&=-\Gamma_{y-1}\,Z_{y-1}\,(1-K_y)+Z_{y-1}\,(1-K_y)\,\Gamma_y=(1-\Theta_y)\,\Omega_y' .
\end{align*}
We show that
\begin{equation}\label{eq:series-comparison}
\Theta_y\preceq \Omega_y\preceq \mathcal{Q}_y-1\qquad(y\ge 2).
\end{equation}
Formally, the differential equations for $\mathcal{Q}_y$ and $\Theta_y$ say that $\mathcal{Q}_y=e^{\Omega_y}$ and $1-\Theta_y=e^{-\Omega_y}$, and the claim is the analogue for coefficients of the inequalities $1-e^{-s}\le s\le e^{s}-1$ for a real number $s$.

For the right-hand inequality let $Y=\mathcal{Q}_y-1-\Omega_y$. The constant term of $Y$ is $0$, and since $\mathcal{Q}_y'=\mathcal{Q}_y\,\Omega_y'$,
\[
Y'=\mathcal{Q}_y\,\Omega_y'-\Omega_y'=(Y+\Omega_y)\,\Omega_y'=\Omega_y'\,Y+\Omega_y\,\Omega_y' .
\]
Here $\Omega_y'\succeq 0$ and $\Omega_y\,\Omega_y'\succeq 0$, because $\Omega_y$ and $\Omega_y'$ are $\succeq 0$, and the positivity principle gives $Y\succeq 0$, i.e.\ $\Omega_y\preceq\mathcal{Q}_y-1$.

For the left-hand inequality we first show that $\Theta_y\succeq 0$. The equation $\Theta_y'=(1-\Theta_y)\,\Omega_y'$ does not show this: written as $\Theta_y'=-\Omega_y'\,\Theta_y+\Omega_y'$, it multiplies $\Theta_y$ by a series whose coefficients are nonpositive, and the positivity principle does not apply. We therefore rewrite its right-hand side. Since $Z_y\,(1-K_y)=1$, we have $\Theta_y=(Z_y-Z_{y-1})\,(1-K_y)$, hence $Z_y-Z_{y-1}=\Theta_y\,Z_y$ and
\[
(1-\Theta_y)\,(Z_y-Z_{y-1})=Z_{y-1}\,(1-K_y)\cdot\Theta_y\,Z_y=Z_{y-1}\,\Theta_y .
\]
Inserting the expression $y\,(Z_y-Z_{y-1})+\mathcal{R}_y$ for $\Omega_y'$, we obtain
\[
\Theta_y'=y\,(1-\Theta_y)\,(Z_y-Z_{y-1})+(1-\Theta_y)\,\mathcal{R}_y=\bigl(y\,Z_{y-1}-\mathcal{R}_y\bigr)\,\Theta_y+\mathcal{R}_y .
\]
In this linear equation the series that multiplies $\Theta_y$ is $\succeq 0$: the coefficient of $X^j$ in $y\,Z_{y-1}-\mathcal{R}_y$ is $y\,z^{(y-1)}_j$ for $j\le y-2$ and $(y-1)\,z^{(y-1)}_j$ for $j\ge y-1$. Since $\mathcal{R}_y\succeq 0$ as well and the constant term of $\Theta_y$ is $0$, the positivity principle gives $\Theta_y\succeq 0$. Now $\Omega_y-\Theta_y$ has constant term $0$, and since $\Theta_y'=(1-\Theta_y)\,\Omega_y'$ its derivative is
\[
\Omega_y'-(1-\Theta_y)\,\Omega_y'=\Theta_y\,\Omega_y' ,
\]
a product of the two series $\Theta_y\succeq 0$ and $\Omega_y'\succeq 0$. By the positivity principle with $G=0$ we get $\Omega_y-\Theta_y\succeq 0$, which completes the proof of \eqref{eq:series-comparison}.

We return to column $k$ and deduce the proposition. Since $Z_k\,(1-K_k)=1$, the definitions of $\mathcal{Q}_{k+1}$ and $\Theta_k$ give
\[
\mathsf{D}\,(1-K_k)=Z_{k+1}\,(1-K_k)-2+Z_{k-1}\,(1-K_k)=(\mathcal{Q}_{k+1}-1)-\Theta_k ,
\]
and the comparison \eqref{eq:series-comparison}, applied with $y=k+1$ to the first term and with $y=k$ to the second, gives
\begin{equation}\label{eq:series-column}
\mathsf{D}\,(1-K_k)\succeq \Omega_{k+1}-\Omega_k=\sum_{t\ge 0}\xi_t\,X^t .
\end{equation}
The equality holds by the formula \eqref{eq:series-log-coefficients} for the coefficients of $\Omega_y$ and by the definition of $\xi_t$. We compare the coefficients of $X^n$. Since $\kappa^{(k)}_0=0$, the coefficient on the left is $\delta_n-\sum_{m=1}^{n}\kappa^{(k)}_m\,\delta_{n-m}$, in which the terms with $n-m<k$ vanish, because $\delta_j=0$ for $j<k$ (Section~\ref{sec:recursion}). Hence
\[
\delta_n-\sum_{m=1}^{n-k}\kappa^{(k)}_m\,\delta_{n-m}\ge\xi_n .
\]
Since $n-1\ge k$, Lemma~\ref{lem:curvatureidentity} gives $n\,\xi_n=(k+1)\,\delta_{n-1}+2\,\beta^{(k-1)}_{n-1}$, and the assertion follows on multiplying by $n$.
\end{proof}

The second inequality bounds the step of column $k$ from below by the convolution of the sequence $(\xi_t)$ with the sequence $z^{(k)}$.

\begin{proposition}\label{prop:onetree}
Let $k\ge 2$ and $n\ge 0$ be integers. Then
\[
\sum_{t=0}^{n}\xi^{(k)}_t\,z^{(k)}_{n-t}\le\delta^{(k)}_n .
\]
\end{proposition}

\begin{proof}
We multiply the inequality \eqref{eq:series-column}, which was shown in the proof of Proposition~\ref{prop:columninequality}, by $Z_k\succeq 0$. Since $(1-K_k)\,Z_k=1$, this gives
\[
\mathsf{D}\succeq \Bigl(\sum_{t\ge 0}\xi^{(k)}_t\,X^t\Bigr)\,Z_k ,
\]
and the proposition is the coefficient of $X^n$ in this inequality.
\end{proof}

\subsection{The diagonal lower bound}\label{subsec:series-diagonal}

For an integer $k\ge 2$ let
\[
\gamma_k=\Bigl(1+\frac{1}{k}\Bigr)^{k-2} .
\]
The sequence $(\gamma_k)_{k\ge 2}$ is positive and nondecreasing: for $2\le i\le j$ Lemma~\ref{lem:powermono} and $1/j\le 1/i$ give
\[
\gamma_i\,\Bigl(1+\frac{1}{i}\Bigr)\le\gamma_j\,\Bigl(1+\frac{1}{j}\Bigr)\le\gamma_j\,\Bigl(1+\frac{1}{i}\Bigr).
\]

The third inequality compares two entries of the table that are neighbours along a diagonal; in terms of the counts it reads $\gamma_k\,n\,b_{n-1,k-1}\le b_{n,k}$. For $n=k+1$ it is an equality: since $\beta^{(k)}_k=\beta^{(k+1)}_{k+1}=0$, Lemma~\ref{lem:initialstep}, applied to columns $k$ and $k+1$, gives $\beta^{(k-1)}_k=a_k/k$ and $\beta^{(k)}_{k+1}=a_{k+1}/(k+1)$, and $a_{k+1}=(1+1/k)^{k-1}\,a_k$ by \eqref{eq:prelim-ratio}.

\begin{proposition}\label{prop:diagonallower}
Let $k\ge 2$ and $n\ge 1$ be integers. Then
\[
\gamma_k\,\beta^{(k-1)}_{n-1}\le\beta^{(k)}_n .
\]
\end{proposition}

\begin{proof}
Fix $k\ge 2$. Let $E_0=0$ and $E_n=\beta^{(k)}_n-\gamma_k\,\beta^{(k-1)}_{n-1}$ for $n\ge 1$. We show by induction on $n$ that $E_n\ge 0$, and we keep the notation of the proof of Proposition~\ref{prop:columninequality}. By the definition of the $\beta^{(\ell)}_n$, the $E_n$ are the coefficients of the series
\[
E=(Z_{k+1}-Z_k)-\gamma_k\,X\,(Z_k-Z_{k-1}) .
\]

We first divide $E$ by $Z_k$, that is, we multiply it by $1-K_k$. Since $Z_k\,(1-K_k)=1$, the definitions $\mathcal{Q}_{k+1}=Z_{k+1}\,(1-K_k)$ and $\Theta_k=1-Z_{k-1}\,(1-K_k)$ give $(Z_{k+1}-Z_k)\,(1-K_k)=\mathcal{Q}_{k+1}-1$ and $(Z_k-Z_{k-1})\,(1-K_k)=\Theta_k$. The comparison \eqref{eq:series-comparison}, applied with $y=k+1$ and with $y=k$, gives
\[
E\,(1-K_k)=(\mathcal{Q}_{k+1}-1)-\gamma_k\,X\,\Theta_k\succeq \Omega_{k+1}-\gamma_k\,X\,\Omega_k ,
\]
and, since $\kappa^{(k)}_0=0$, a comparison of the coefficients of $X^n$ shows that
\begin{equation}\label{eq:series-diagonal-division}
E_n-\sum_{m=1}^{n}\kappa^{(k)}_m\,E_{n-m}\ge[X^n]\bigl(\Omega_{k+1}-\gamma_k\,X\,\Omega_k\bigr)\qquad(n\ge 0).
\end{equation}

Next we bound the right-hand side of this inequality from below. We show that
\begin{equation}\label{eq:series-diagonal-step}
(i+1)\,[X^{i+1}]\bigl(\Omega_{k+1}-\gamma_k\,X\,\Omega_k\bigr)\ge k\,E_i\qquad(i\ge k).
\end{equation}
Let $i\ge k$. By the formula \eqref{eq:series-log-coefficients}, in which all the minima are now the thresholds themselves,
\begin{align*}
(i+1)\,[X^{i+1}]\Omega_{k+1}&=(k+1)\,z^{(k+1)}_i-k\,z^{(k)}_i ,\\
i\,[X^{i}]\Omega_k&=k\,z^{(k)}_{i-1}-(k-1)\,z^{(k-1)}_{i-1} .
\end{align*}
Inserting this and the definition of $E_i$, we find by a direct computation that
\begin{multline*}
(i+1)\,[X^{i+1}]\bigl(\Omega_{k+1}-\gamma_k\,X\,\Omega_k\bigr)-k\,E_i\\
=\Bigl(z^{(k+1)}_i-\gamma_k\,\frac{i+1}{i}\,z^{(k)}_{i-1}\Bigr)+\gamma_k\,\frac{i+1-k}{i}\,\Bigl(z^{(k)}_{i-1}-z^{(k-1)}_{i-1}\Bigr),
\end{multline*}
and we show that the two terms on the right are nonnegative. For the second this holds because $i\ge k$ and $z^{(k-1)}_{i-1}\le z^{(k)}_{i-1}$. For the first, Proposition~\ref{prop:diagonalmonotone} with threshold $k$ and index $i-1$ gives $z^{(k)}_{i-1}/c_{i-1}\le z^{(k+1)}_i/c_i$. Since $c_i=(1+1/i)^{i-1}\,c_{i-1}=\gamma_i\,\frac{i+1}{i}\,c_{i-1}$ by \eqref{eq:prelim-ratio}, it follows that
\[
z^{(k+1)}_i\ge\frac{c_i}{c_{i-1}}\,z^{(k)}_{i-1}=\gamma_i\,\frac{i+1}{i}\,z^{(k)}_{i-1}\ge\gamma_k\,\frac{i+1}{i}\,z^{(k)}_{i-1} ,
\]
where the last step holds because $\gamma_k\le\gamma_i$ for $2\le k\le i$. This proves \eqref{eq:series-diagonal-step}.

It remains to carry out the induction. For $n\le k$ we have $E_n=0$, because $\beta^{(\ell)}_j=0$ for $j\le\ell$ (Section~\ref{sec:recursion}). Let $i\ge k$, and suppose that $E_j\ge 0$ for all $j\le i$. By \eqref{eq:series-diagonal-division} for $n=i+1$ and by \eqref{eq:series-diagonal-step},
\[
E_{i+1}\ge\sum_{m=1}^{i+1}\kappa^{(k)}_m\,E_{i+1-m}+\frac{k}{i+1}\,E_i .
\]
Here $\kappa^{(k)}_m\ge 0$ by Lemma~\ref{lem:kernelnonneg} with $x=1$, and the $E_j$ on the right are nonnegative by the induction hypothesis. Hence $E_{i+1}\ge 0$.
\end{proof}

\section{The critical tilt}\label{sec:critical}
Recall from Section~\ref{sec:recursion} that a tilt is a real number $x$ and that the tilted form of a sequence $(s_n)$ is the sequence $(s_nx^n)$. The proof of the inequality $d_{n,k}\ge 0$ for $2k\le n$ (Proposition~\ref{prop:leftpart}) works, for each column $k\ge 2$, at one particular tilt, the critical tilt $x_k$. The main statement at this tilt is Proposition~\ref{prop:criticalmonotone}: for $k\ge 2$ the tilted counts $z^{(k)}_Nx_k^N$ of threshold $k$ never increase in $N$, so that they can serve as monotone weights in a summation by parts.

We begin with the definition. Since $c_j=a_{j+1}$, the ratio formula \eqref{eq:prelim-ratio} gives $c_{j+1}=\bigl(1+\frac1{j+1}\bigr)^{j}c_j$, and hence, for every real number $x$ and every integer $j\ge 0$,
\begin{equation}\label{eq:critical-ratio}
c_{j+1}\,x^{j+1}=x\Bigl(1+\frac{1}{j+1}\Bigr)^{j}\,c_j\,x^{j}.
\end{equation}
For an integer $k\ge 1$, the \emph{critical tilt} of column $k$ is the tilt at which the ratio $x\bigl(1+\frac1{j+1}\bigr)^{j}$ of two consecutive terms of the tilted sequence $(c_jx^j)$ equals $1$ for $j=k-1$, i.e.
\[
x_k=\Bigl(\frac{k}{k+1}\Bigr)^{k-1},\qquad\text{so that}\qquad x_k\Bigl(1+\frac1k\Bigr)^{k-1}=1.
\]
We have $0<x_k\le 1$, and $c_k\,x_k^{k}=c_{k-1}\,x_k^{k-1}$ by the choice of $x_k$. The first values are $x_1=1$, $x_2=\frac23$, $x_3=\frac{9}{16}$ and $x_4=\frac{64}{125}$, and $x_k>e^{-1}$ for every $k$, since $(1+1/k)^{k-1}<e$. The tilt is chosen with threshold $k$ in view. On the free range $N\le k$ the tilted counts of threshold $k$ are $z^{(k)}_Nx^N=c_Nx^N$ (Lemma~\ref{lem:freerange}), and for $x>x_k$ the count at $N=k$ is larger than the count at $N=k-1$; so $x_k$ is the largest tilt at which the tilted counts of threshold $k$ can be nonincreasing in $N$.

Here and in the rest of the paper, when a column $k$ is fixed, a hat denotes the tilted form of a sequence at the critical tilt $x_k$. Thus
\begin{gather*}
\hat c_n=c_n\,x_k^{n},\qquad \hat a_n=a_n\,x_k^{n},\qquad \hat\varepsilon_m=\varepsilon_m\,x_k^{m},\\
\hat z^{(y)}_N=z^{(y)}_N\,x_k^{N},\qquad \hat\beta^{(\ell)}_N=\beta^{(\ell)}_N\,x_k^{N},\qquad \hat\delta_n=\delta^{(k)}_n\,x_k^{n},
\end{gather*}
and $\hat\kappa_m=\kappa^{(k)}_m\,x_k^{m}$ for the prime kernel of threshold $k$ itself.

The numbers $\hat c_n$ do not increase up to index $k$ and do not decrease from index $k-1$ on, so that $\hat c_{k-1}=\hat c_k$ is their smallest value: for all integers $k\ge 1$ and $n\ge 0$,
\begin{equation}\label{eq:critical-updown}
\begin{aligned}
c_{n+1}\,x_k^{n+1}&\le c_n\,x_k^{n}&&\text{if } n+1\le k,\\
c_n\,x_k^{n}&\le c_{n+1}\,x_k^{n+1}&&\text{if } n+1\ge k.
\end{aligned}
\end{equation}
Indeed, by the identity \eqref{eq:critical-ratio} with $x=x_k$ and $j=n$ the two terms differ by the factor $x_k\bigl(1+\frac1{n+1}\bigr)^{n}$, and by Lemma~\ref{lem:powermono} and the definition of $x_k$ this factor is at most $x_k\bigl(1+\frac1k\bigr)^{k-1}=1$ if $n+1\le k$ and at least $1$ if $n+1\ge k$.

The tilted free kernel does not increase up to index $k$ either.

\begin{lemma}\label{lem:epsdown}
Let $k$ and $m$ be integers with $1\le m<k$. Then
\[
\varepsilon_{m+1}\,x_k^{m+1}\le\varepsilon_m\,x_k^{m}.
\]
\end{lemma}

\begin{proof}
For $m=1$ the claim is $\frac12x_k^2\le x_k$, since $\varepsilon_1=1$ and $\varepsilon_2=\frac12$, and it holds because $0<x_k\le 1$. Let $m\ge 2$. By the definitions of $a_{i+1}$ and $\varepsilon_{i+1}$,
\[
a_{i+1}=\Bigl(1+\frac1i\Bigr)^{i}\,\varepsilon_{i+1}\qquad(i\ge 1).
\]
Moreover, $\hat a_{m+1}\le\hat a_m$: indeed, $\hat a_{n+1}=x_k\,\hat c_n$ for every $n\ge 0$, and $\hat c_m\le\hat c_{m-1}$ by \eqref{eq:critical-updown}, since $m\le k$. Hence
\[
\bigl(1+\tfrac1m\bigr)^{m}\,\hat\varepsilon_{m+1}=\hat a_{m+1}\le\hat a_m=\bigl(1+\tfrac1{m-1}\bigr)^{m-1}\,\hat\varepsilon_m\le\bigl(1+\tfrac1m\bigr)^{m}\,\hat\varepsilon_m,
\]
where the two equalities are the identity for $i=m$ and for $i=m-1$, and the last inequality holds because $(1+1/i)^{i}$ is nondecreasing in $i\ge 1$ (Section~\ref{sec:prelim}). Dividing by $\bigl(1+\frac1m\bigr)^{m}$ gives the claim.
\end{proof}

We turn to the prime kernel $\kappa^{(y)}$ of a threshold $y\ge 2$, which was introduced in Section~\ref{sec:series}. For $1\le m\le y$ it is the free kernel (Lemma~\ref{lem:kernelfree}), and beyond this range it is no longer explicit. The next lemma and its proof show that it nevertheless inherits the monotonicity of the free kernel: if the tilted free kernel does not increase up to index $y$, then the tilted prime kernel of threshold $y$ does not increase at all, and beyond the threshold it decays at least like $1/m$. For $y=k$ and the critical tilt $x_k$ the hypothesis is Lemma~\ref{lem:epsdown}.

\begin{lemma}\label{lem:harmonic}
Let $y\ge 2$ be an integer and let $x\ge 0$ be a real number such that
\[
\varepsilon_{m+1}\,x^{m+1}\le\varepsilon_m\,x^{m}\qquad\text{for all integers $m$ with $1\le m<y$.}
\]
Then
\[
(m+1)\,\kappa^{(y)}_{m+1}\,x^{m+1}\le m\,\kappa^{(y)}_m\,x^{m}\qquad\text{for all integers $m\ge y$.}
\]
\end{lemma}

\begin{proof}
We prove the following two statements together, by induction on $m\ge 1$:
\begin{enumerate}
\item[\textup{(A$_m$)}] $\kappa^{(y)}_{m+1}\,x^{m+1}\le\kappa^{(y)}_m\,x^{m}$;
\item[\textup{(B$_m$)}] if $m\ge y$, then $(m+1)\,\kappa^{(y)}_{m+1}\,x^{m+1}\le m\,\kappa^{(y)}_m\,x^{m}$.
\end{enumerate}
The lemma is (B$_m$) for $m\ge y$. Let $m\ge 1$, and suppose that (A$_i$) holds for all $i$ with $1\le i<m$.

Suppose first that $m+1\le y$. Then (A$_m$) is the hypothesis of the lemma for this $m$, because at the indices $m$ and $m+1$ the prime kernel is the free kernel (Lemma~\ref{lem:kernelfree}).

Suppose now that $m\ge y$. Multiplying the recursion of Lemma~\ref{lem:kernelrecursion} by $x^{N}$, we obtain
\[
N\,\kappa^{(y)}_N\,x^{N}=\sum_{j=1}^{y-1}(y-j)\,a_j x^{j}\;\kappa^{(y)}_{N-j}\,x^{N-j}\qquad(N\ge y).
\]
We compare the right-hand sides for $N=m+1$ and for $N=m$. They have the same weights $(y-j)\,a_j x^{j}$, which are nonnegative. For $1\le j\le y-1$ we have $1\le m-j<m$, so (A$_i$) with $i=m-j$ gives $\kappa^{(y)}_{m+1-j}\,x^{m+1-j}\le\kappa^{(y)}_{m-j}\,x^{m-j}$. Hence each term of the sum for $N=m+1$ is at most the corresponding term of the sum for $N=m$, and (B$_m$) follows. Finally, $\kappa^{(y)}_{m+1}\,x^{m+1}\ge 0$ by Lemma~\ref{lem:kernelnonneg}, so that
\[
m\,\kappa^{(y)}_{m+1}\,x^{m+1}\le(m+1)\,\kappa^{(y)}_{m+1}\,x^{m+1}\le m\,\kappa^{(y)}_m\,x^{m},
\]
and division by $m$ gives (A$_m$). This completes the induction.
\end{proof}

The statements (A$_m$) of this induction are the monotonicity of the tilted prime kernel: under the hypothesis of Lemma~\ref{lem:harmonic}, that is, for an integer $y\ge 2$ and a real number $x\ge 0$ with $\varepsilon_{m+1}\,x^{m+1}\le\varepsilon_m\,x^{m}$ for $1\le m<y$, we have
\begin{equation}\label{eq:critical-kerneldown}
\kappa^{(y)}_{m+1}\,x^{m+1}\le\kappa^{(y)}_m\,x^{m}\qquad\text{for all integers $m\ge 1$.}
\end{equation}

We now come to the tilted counts $\hat z^{(k)}_N=z^{(k)}_N\,x_k^{N}$ of threshold $k$. Besides the first inequality of \eqref{eq:critical-updown}, the proof of their monotonicity (Proposition~\ref{prop:criticalmonotone} below) requires the inequality
\begin{equation}\label{eq:critical-masscondition}
\sigma_k\le 1,\qquad\text{where}\qquad\sigma_k=x_k\sum_{j=1}^{k}c_j\,x_k^{j}=\sum_{n=2}^{k+1}a_n\,x_k^{n},
\end{equation}
between explicit numbers; the two sums agree because $c_j=a_{j+1}$. We prove this inequality for $k\ge 12$ by estimating $a_n\,x_k^{n}$ with Stirling's formula (Lemma~\ref{lem:cayleysum}). This estimate does not reach the smallest columns, and for $2\le k\le 11$ we evaluate $\sigma_k$, in the proof of the proposition. The estimate uses the following bound for the excess of $x_k$ over $e^{-1}$.

\begin{lemma}\label{lem:tiltupper}
For every integer $k\ge 1$,
\[
e\,x_k\le\exp\Bigl(\frac{3}{2k}\Bigr).
\]
\end{lemma}

\begin{proof}
The elementary inequality $\log(1+t)\ge\frac{2t}{2+t}$ for $t\ge 0$, applied with $t=1/k$, gives $\log\bigl(1+\frac1k\bigr)\ge\frac{2}{2k+1}$. Multiplying by $k-1$, we obtain
\[
(k-1)\log\Bigl(1+\frac1k\Bigr)\ge\frac{2(k-1)}{2k+1}\ge 1-\frac{3}{2k},
\]
where the second inequality is checked directly. Hence $\bigl(1+\frac1k\bigr)^{k-1}\ge\exp\bigl(1-\frac{3}{2k}\bigr)$, and by the definition of $x_k$
\[
x_k=\Bigl(1+\frac1k\Bigr)^{-(k-1)}\le\exp\Bigl(-1+\frac{3}{2k}\Bigr).\qedhere
\]
\end{proof}

\begin{lemma}\label{lem:cayleysum}
For every integer $k\ge 12$,
\[
x_k\sum_{j=1}^{k}c_j\,x_k^{j}\le 1.
\]
\end{lemma}

\begin{proof}
Write $x=x_k$ and $\lambda=3/(2k)$. Since $c_j=a_{j+1}$, the left-hand side is $\sum_{n=2}^{k+1}a_nx^n$.

Let $2\le n\le k+1$. By the definition of $a_n$ and by the upper bound in \eqref{eq:prelim-stirling},
\[
a_n\,x^{n}=\frac{(e\,x)^n}{n}\cdot e^{-n}\,\frac{n^n}{n!}\le\frac{(e\,x)^n}{\sqrt{2\pi}\,n^{3/2}}.
\]
By Lemma~\ref{lem:tiltupper} we have $(e\,x)^n\le e^{\lambda n}$, and $0\le\lambda n\le\frac{3(k+1)}{2k}\le\frac{13}{8}$, because $n\le k+1$ and $k\ge 12$. Since $e^{t}\le 1+t+t^2$ for $0\le t\le 1.7$, as one easily checks, this gives with $t=\lambda n$
\[
\sqrt{2\pi}\;a_n\,x^{n}\le\frac{1+\lambda n+\lambda^2n^2}{n^{3/2}}=n^{-3/2}+\lambda\,n^{-1/2}+\lambda^2\,n^{1/2}\qquad(2\le n\le k+1).
\]

We sum this over $2\le n\le k+1$ and compare each of the three sums with an integral. Since $t^{-3/2}$ is convex, $t^{-1/2}$ decreases and $t^{1/2}$ increases, the value of these functions at $n$ is at most their mean over $[n-\frac12,n+\frac12]$, resp.\ over $[n-1,n]$, resp.\ over $[n,n+1]$. Therefore
\begin{align*}
\sum_{n=2}^{k+1}n^{-3/2}&\le\int_{3/2}^{k+3/2}t^{-3/2}\,dt=\frac{2}{\sqrt{3/2}}-\frac{2}{\sqrt{k+3/2}},\\
\sum_{n=2}^{k+1}n^{-1/2}&\le\int_{1}^{k+1}t^{-1/2}\,dt=2\sqrt{k+1}-2,\\
\sum_{n=2}^{k+1}n^{1/2}&\le\int_{2}^{k+2}t^{1/2}\,dt\le\frac23\,(k+2)^{3/2},
\end{align*}
and with $\lambda=3/(2k)$ we arrive at
\[
\sqrt{2\pi}\sum_{n=2}^{k+1}a_n\,x^{n}\le\frac{2}{\sqrt{3/2}}-\frac{2}{\sqrt{k+3/2}}+\frac{3\,(\sqrt{k+1}-1)}{k}+\frac{3\,(k+2)^{3/2}}{2k^2}.
\]

It remains to show that the right-hand side of this inequality is at most $\sqrt{2\pi}$. Put $r=\sqrt{k+2}$, so that $r\ge 3.74$, because $3.74^2<14\le k+2$. Since $\sqrt{k+1}\le r$ and $\sqrt{k+3/2}\le r$, the right-hand side is at most
\[
\frac{2}{\sqrt{3/2}}+\Bigl(-\frac{2}{r}+\frac{3(r-1)}{k}+\frac{3r^3}{2k^2}\Bigr).
\]
We claim that the expression in the bracket is at most $0.87$ for every $k\ge 12$. To prove this, we multiply the expression by $2rk^2>0$ and use the relation $r^2=k+2$ to remove the powers of $r$ higher than the first:
\begin{align*}
2rk^2\Bigl(-\frac{2}{r}+\frac{3(r-1)}{k}+\frac{3r^3}{2k^2}\Bigr)&=-4k^2+6\,(r^2-r)\,k+3r^4\\
&=-4k^2+6\,(k+2)\,k-6rk+3\,(k+2)^2\\
&=5k^2+24k+12-6rk.
\end{align*}
The claim is therefore equivalent to the inequality $5k^2+24k+12-6rk\le 1.74\,rk^2$, that is, to
\[
5k^2+24k+12\le rk\,(1.74\,k+6).
\]
Since $r\ge 3.74$, the right-hand side of this inequality is at least $3.74\,k\,(1.74\,k+6)=6.5076\,k^2+22.44\,k$. Subtracting the left-hand side from this lower bound leaves
\[
1.5076\,k^2-1.56\,k-12=k\,(1.5076\,k-1.56)-12\ge 12\cdot 16-12>0,
\]
where we have used that $k\ge 12$ and hence $1.5076\,k-1.56\ge 16$. This proves the claim.

Since $2/\sqrt{3/2}<1.6331$ and $\sqrt{2\pi}>2.5066$, we conclude that $\sqrt{2\pi}\sum_{n=2}^{k+1}a_n\,x^{n}\le 1.6331+0.87=2.5031<\sqrt{2\pi}$, and the lemma follows.
\end{proof}

In terms of the counts, the main statement of the section says that
\[
A_{N+1,k-1}\le(N+1)\,\Bigl(1+\frac1k\Bigr)^{k-1}A_{N,k-1}\qquad(N\ge 0,\ k\ge 2),
\]
and the factor $\bigl(1+\frac1k\bigr)^{k-1}$ cannot be replaced by a smaller one, since for $N=k-1$ both sides are equal to $(k+1)^{k-1}$.

\begin{proposition}\label{prop:criticalmonotone}
Let $k\ge 2$ be an integer. Then, for every integer $N\ge 0$,
\[
z^{(k)}_{N+1}\,x_k^{N+1}\le z^{(k)}_N\,x_k^{N}.
\]
\end{proposition}

\begin{proof}
We derive the proposition from a statement about an arbitrary solution of a tilted last-car recursion.

For the general statement, let $y\ge 1$ be an integer, let $x\ge 0$ be a real number, and let $(W_N)_{N\ge 0}$ be a solution of the weighted last-car recursion with weight $\omega(t)=x\min(t,y)$ in the sense of Definition~\ref{def:weighted}, i.e.\ $W_0=1$ and
\[
(N+1)\,W_{N+1}=x\sum_{i=0}^{N}\min(i+1,y)\,W_i\,W_{N-i}\qquad(N\ge 0).
\]
Suppose that
\begin{enumerate}
\item[\textup{(a)}] $W_{N+1}\le W_N$ for all $N<y$, and
\item[\textup{(b)}] $x\sum_{j=1}^{y}W_j\le 1$.
\end{enumerate}
We claim that $W_{N+1}\le W_N$ for all $N\ge 0$.

Since $W_0=1$ and $x\ge 0$, the recursion shows by induction on $N$ that $W_N\ge 0$ for all $N$.

Put $w_N=W_N-W_{N+1}$. For $N\ge y$ we have the identity
\[
(N+1)\,w_N=(1-xy)\,W_N+x\sum_{i=0}^{N-1}\min(i+1,y)\,W_i\,w_{N-1-i}.
\]
To see this, write $(N+1)\,w_N=W_N+N\,W_N-(N+1)\,W_{N+1}$ and insert the recursion, with $N-1$ in place of $N$ for the second term and as it stands for the third. In the latter the term with $i=N$ equals $xy\,W_N$, because $N+1>y$ and $W_0=1$, and the terms with $i\le N-1$ combine with those of the former, since $W_{N-1-i}-W_{N-i}=w_{N-1-i}$.

The first term on the right of this identity may be negative. We show that it is outweighed by those terms of the sum that contain the decreases $w_0,\dots,w_{y-1}$ on the free range, and prove in this way that $w_N\ge 0$, by induction on $N$. For $N<y$ this is condition (a). Let $N\ge y$, and suppose that $w_j\ge 0$ for all $j<N$. Then $W_0\ge W_1\ge\dots\ge W_N\ge 0$, and every term of the sum in the identity is nonnegative. Substituting $j=N-1-i$, we write this sum as $\sum_{j=0}^{N-1}\min(N-j,y)\,W_{N-1-j}\,w_j$, and we keep only the terms with $j\le y-1$; they are all present, because $y\le N$. For these terms
\[
\min(N-j,y)\ge y-j\qquad\text{and}\qquad W_{N-1-j}\ge W_N,
\]
the first because $N\ge y$. Hence
\[
\sum_{i=0}^{N-1}\min(i+1,y)\,W_i\,w_{N-1-i}\ge W_N\sum_{j=0}^{y-1}(y-j)\,w_j.
\]
Since $w_j=W_j-W_{j+1}$ and $W_0=1$, summation by parts gives
\[
\sum_{j=0}^{y-1}(y-j)\,w_j=y\,W_0-\sum_{j=1}^{y}W_j=y-\sum_{j=1}^{y}W_j.
\]
Inserting the last two displays into the identity for $(N+1)\,w_N$, where $x\ge 0$, we see that the terms $xy\,W_N$ cancel:
\begin{align*}
(N+1)\,w_N&\ge(1-xy)\,W_N+x\,W_N\Bigl(y-\sum_{j=1}^{y}W_j\Bigr)\\
&=W_N\Bigl(1-x\sum_{j=1}^{y}W_j\Bigr)\ge 0,
\end{align*}
by condition (b) and by $W_N\ge 0$. This completes the induction and proves the claim.

We now return to the tilted counts of threshold $k\ge 2$, and apply the claim with $y=k$ and $x=x_k$ to the sequence $W_N=z^{(k)}_N\,x_k^{N}$. By \eqref{eq:recursion-tilted} this sequence satisfies the recursion of the general statement, and $W_0=z^{(k)}_0=1$. On the free range $N\le k$ we have $W_N=c_N\,x_k^{N}=\hat c_N$ by Lemma~\ref{lem:freerange}. Condition (a) therefore reads $\hat c_{N+1}\le\hat c_N$ for $N<k$, which is the first inequality of \eqref{eq:critical-updown}, and condition (b) is the inequality $\sigma_k\le 1$ of \eqref{eq:critical-masscondition}. For $k\ge 12$ it is Lemma~\ref{lem:cayleysum}. For $2\le k\le 11$ it is checked directly: the number $\sigma_k=\sum_{n=2}^{k+1}a_n\,x_k^{n}$, with $a_n=n^{n-1}/n!$ and $x_k=(k/(k+1))^{k-1}$, is a sum of at most eleven rational numbers, and evaluating the ten numbers $\sigma_2,\dots,\sigma_{11}$ exactly we find that they decrease from $\sigma_2=\frac89$ to $\sigma_{11}=0.7791\ldots$, so that $\sigma_k<1$ in each of the ten cases. Hence (a) and (b) hold for every $k\ge 2$, and the claim gives $z^{(k)}_{N+1}\,x_k^{N+1}\le z^{(k)}_N\,x_k^{N}$ for all $N\ge 0$.
\end{proof}

Finally, we carry the diagonal lower bound of Proposition~\ref{prop:diagonallower} over to the critical tilt, where its constant becomes rational. Recall from Section~\ref{sec:recursion} that $\beta^{(\ell)}_N=z^{(\ell+1)}_N-z^{(\ell)}_N$ for $\ell\ge 1$; the corollary is thus an inequality between the tilted counts of thresholds $k-1$, $k$ and $k+1$.

\begin{corollary}\label{cor:diagonalcritical}
Let $k\ge 2$ and $n\ge 1$ be integers. Then
\[
k\,\beta^{(k-1)}_{n-1}\,x_k^{n-1}\le(k+1)\,\beta^{(k)}_n\,x_k^{n}.
\]
\end{corollary}

\begin{proof}
By Proposition~\ref{prop:diagonallower} we have $\gamma_k\,\beta^{(k-1)}_{n-1}\le\beta^{(k)}_n$, where $\gamma_k=\bigl(1+\frac1k\bigr)^{k-2}$, and the definition of $x_k$ can be written as $(k+1)\,x_k\,\gamma_k=k$. Hence
\[
k\,\beta^{(k-1)}_{n-1}\,x_k^{n-1}=(k+1)\,x_k^{n}\,\gamma_k\,\beta^{(k-1)}_{n-1}\le(k+1)\,x_k^{n}\,\beta^{(k)}_n.\qedhere
\]
\end{proof}

\section{Exact computations for the small rows and the small columns}\label{sec:tables}
In this section we establish the three facts in the proof of Theorem~\ref{thm:main} that are obtained by exact computation. The first concerns all rows $n\le 16$ (Lemma~\ref{lem:sixteenrows}); the second, Proposition~\ref{prop:nextcolumn} for the columns $9\le k\le 264$ (Lemma~\ref{lem:signpairs}); the third, Propositions~\ref{prop:chainrange} and~\ref{prop:leftpart} for the columns $2\le k\le 81$ (Lemma~\ref{lem:columntable}). These are the rows and the columns that our estimates do not reach.

The computations are carried out by computer and involve integers only, so that all of the arithmetic is exact. No parking function is enumerated: the numbers computed are cumulative counts $A_{n,\ell}$, or the differences $(n+1)^{n-1}-A_{n,\ell}$, and they are obtained from a recursion in $n$, either the last-car recursion (Proposition~\ref{prop:lastcar}) or the single-chain formula (Proposition~\ref{prop:forward}). Each of the three facts is then a finite number of comparisons between integers that are formed from these numbers. The second and the third computation depend on one number for each column, $\zeta(k)$ and $n_0(k)$, which locates the rows that are examined. We list these numbers in full, in Definitions~\ref{def:zeta} and~\ref{def:startrow}, and the proof of each lemma gives the recursion, the rows up to which it is run and the comparisons that are made, so that each computation can be repeated from what is written here.

\subsection{Rows up to $16$}\label{subsec:tables-rows}

The inequality $d_{n,k}\ge 0$ for $2k\le n$ is proved for the rows $n\ge 17$ only (Proposition~\ref{prop:leftpart}), and in smaller rows it can fail. We therefore compute all rows $n\le 16$. Each of them turns out to be unimodal, and the following lemma states this in the form needed for Proposition~\ref{prop:descentascent}.

\begin{lemma}\label{lem:sixteenrows}
Let $n$, $i$ and $j$ be integers with $0\le n\le 16$ and $1\le i<j\le n-1$. Then row $n$ does not have both a strict descent at $i$ and a strict ascent at $j$.
\end{lemma}

\begin{proof}
For $0\le\ell\le 15$ and $0\le m\le 16$ we compute $A_{m,\ell}$, starting from $A_{0,\ell}=1$, by the last-car recursion of Proposition~\ref{prop:lastcar},
\[
A_{m+1,\ell}=\sum_{j=0}^{m}\min(j+1,\ell+1)\binom{m}{j}A_{j,\ell}\,A_{m-j,\ell}.
\]
The entries of row $n$ are then $b_{n,0}=A_{n,0}$ and $b_{n,\ell}=A_{n,\ell}-A_{n,\ell-1}$ for $1\le\ell\le n-1$. We find that the rows with $1\le n\le 11$ are
\[
\begin{array}{r|l}
n & b_{n,0},\ b_{n,1},\ \dots,\ b_{n,n-1}\\ \hline
1 & 1\\
2 & 2,\ 1\\
3 & 6,\ 7,\ 3\\
4 & 24,\ 51,\ 34,\ 16\\
5 & 120,\ 421,\ 377,\ 253,\ 125\\
6 & 720,\ 3963,\ 4594,\ 3688,\ 2546,\ 1296\\
7 & 5040,\ 42253,\ 62145,\ 57398,\ 46142,\ 32359,\ 16807\\
8 & 40320,\ 505515,\ 929856,\ 979430,\ 865970,\ 702292,\ 497442,\ 262144\\
9 & 362880,\ 6724381,\ 15298809,\ 18289811,\ 17520519,\\
 & 15455851,\ 12587507,\ 8977273,\ 4782969\\
10 & 3628800,\ 98618763,\ 275075236,\ 371706836,\ 385131546,\\
 & 359973796,\ 318092120,\ 259627152,\ 186093442,\ 100000000\\
11 & 39916800,\ 1582715773,\ 5373200117,\ 8177419578,\\
 & 9160454598,\ 9008899440,\ 8390168112,\ 7411908040,\\
 & 6056786384,\ 4357947691,\ 2357947691
\end{array}
\]
and that the rows with $12\le n\le 16$ are
\[
\begin{array}{r|l}
n & b_{n,0},\ b_{n,1},\ \dots,\ b_{n,n-1}\\ \hline
12 & 479001600,\ 27612565995,\ 113407725802,\ 193804329368,\\
 & 234483247248,\ 242766518158,\ 234882851708,\ 218509882050,\\
 & 192828376414,\ 157676318044,\ 113792213426,\ 61917364224\\
13 & 6227020800,\ 520631327581,\ 2573638729605,\ 4926465506021,\\
 & 6428493655749,\ 7014734252870,\ 7025563491106,\ 6770961415170,\\
 & 6285825855182,\ 5539649733253,\ 4531383858509,\ 3278177135413,\\
 & 1792160394037\\
14 & 87178291200,\ 10554164679243,\ 62521925061642,\\
 & 133777256144940,\ 188000138187750,\ 216295074328842,\\
 & 224520420787986,\ 222211003267142,\ 213687649761928,\\
 & 197900551632256,\ 174165050709924,\ 142490660510968,\\
 & 103290082620258,\ 56693912375296\\
15 & 1307674368000,\ 228975516609853,\ 1619488339989969,\\
 & 3866365396216526,\ 5844430081201758,\ 7086949855748155,\\
 & 7637297774350499,\ 7747121093904120,\ 7638524284820008,\\
 & 7325291684300523,\ 6767645976331523,\ 5947866448755034,\\
 & 4866510228104930,\ 3533624614867663,\ 1946195068359375\\
16 & 20922789888000,\ 5294731892093355,\ 44569555672124464,\\
 & 118532872437716726,\ 192529294913883426,\ 245884738565040108,\\
 & 275378436591129714,\ 286698620154512400,\ 288176016598878112,\\
 & 283408731647719814,\ 271000724330122940,\ 249788273565590526,\\
 & 219249309938712778,\ 179389782285766308,\ 130443446088709186,\\
 & 72057594037927936
\end{array}
\]
In each of these sixteen rows the entries do not decrease up to the largest entry of the row and do not increase after it. Hence a strict descent at an index $i$ stands to the right of the largest entry and is not followed by a strict ascent at an index $j>i$. This is the assertion of the lemma.
\end{proof}

\subsection{Two signs for each column from $9$ to $264$}\label{subsec:tables-signs}

For $9\le k\le 264$, Proposition~\ref{prop:nextcolumn} is derived in Section~\ref{sec:next} from Proposition~\ref{prop:chainrange} and from two signs: column $k$ is positive in some row of its chain range, and column $k+1$ is negative one row earlier. The two rows are $k+\zeta(k)+1$ and $k+\zeta(k)$, with $\zeta(k)$ defined as follows.

\begin{definition}\label{def:zeta}
For an integer $k$ with $9\le k\le 264$, the number $\zeta(k)$ is given by the following table: $\zeta(k)$ is the value $\zeta$ in whose range of columns $k$ lies.
\[
\begin{array}{c|c@{\qquad\quad}c|c@{\qquad\quad}c|c}
\zeta & k & \zeta & k & \zeta & k\\ \hline
7 & 9 & 17 & 67\text{--}74 & 27 & 163\text{--}174\\
8 & 10\text{--}14 & 18 & 75\text{--}83 & 28 & 175\text{--}185\\
9 & 15\text{--}20 & 19 & 84\text{--}91 & 29 & 186\text{--}197\\
10 & 21\text{--}25 & 20 & 92\text{--}101 & 30 & 198\text{--}210\\
11 & 26\text{--}31 & 21 & 102\text{--}110 & 31 & 211\text{--}222\\
12 & 32\text{--}37 & 22 & 111\text{--}120 & 32 & 223\text{--}235\\
13 & 38\text{--}44 & 23 & 121\text{--}130 & 33 & 236\text{--}248\\
14 & 45\text{--}51 & 24 & 131\text{--}141 & 34 & 249\text{--}261\\
15 & 52\text{--}58 & 25 & 142\text{--}151 & 35 & 262\text{--}264\\
16 & 59\text{--}66 & 26 & 152\text{--}162 & &
\end{array}
\]
\end{definition}

\begin{lemma}\label{lem:signpairs}
Let $k$ be an integer with $9\le k\le 264$, and let $\zeta=\zeta(k)$. Then
\begin{enumerate}
\item $\zeta+2\le k$;
\item $b_{k+\zeta+1,\,k-1}<b_{k+\zeta+1,\,k}$;
\item $b_{k+\zeta,\,k+1}<b_{k+\zeta,\,k}$.
\end{enumerate}
\end{lemma}

\begin{proof}
Part~(1) is read from the table of Definition~\ref{def:zeta}, in each line of which $\zeta+2$ is at most the first column of the range.

Write $p_n=(n+1)^{n-1}$ for $n\ge 0$, so that $p_0=1$ and $p_n=n!\,c_n$, and for a threshold $y\ge 1$ and $n\ge 0$ let
\[
U^{(y)}_n=p_n-A_{n,y-1}.
\]
Recall that $A_{n,\ell}$ is nondecreasing in $\ell$ and equals $p_n$ for $n\le\ell+1$. Hence $U^{(y)}_n$ is a nonnegative integer, the number of parking functions of length $n$ in which some displacement is at least $y$, and $U^{(y)}_n=0$ for $n\le y$. We claim that, for $y\ge 1$ and $n\le 2y+1$,
\begin{equation}\label{eq:tables-forward}
U^{(y)}_n=\sum_{t=y+1}^{n}\binom{n}{t}\,p_{n-t}\,\bigl((t-y)\,p_{t-1}+y\,U^{(y)}_{t-1}\bigr).
\end{equation}
Indeed, by Corollary~\ref{cor:normalised} the sequence $z^{(y)}$ is a solution of the weighted last-car recursion with weight $\omega(t)=\min(t,y)$, so that Proposition~\ref{prop:forward} applies to $W=z^{(y)}$ with $y$ in place of $k$. Since $\omega(t)=y$ for $t\ge y+1$, it gives, for $n\le 2y+1$,
\[
c_n-z^{(y)}_n=\sum_{t=y+1}^{n}c_{n-t}\Bigl(c_{t-1}-\frac{y}{t}\,z^{(y)}_{t-1}\Bigr).
\]
Multiplying this identity by $n!$ and using $c_m=p_m/m!$, $z^{(y)}_m=A_{m,y-1}/m!$ and $n!/((n-t)!\,(t-1)!)=t\binom{n}{t}$, we obtain the claim, because $t\,p_{t-1}-y\,A_{t-1,y-1}=(t-y)\,p_{t-1}+y\,U^{(y)}_{t-1}$.

For $n\ge y+1$ the right-hand side of the recursion \eqref{eq:tables-forward} contains, besides binomial coefficients and the $p_t$, only the $U^{(y)}_t$ with $y\le t\le n-1$, and $U^{(y)}_y=0$. So the recursion gives $U^{(y)}_{y+1},U^{(y)}_{y+2},\dots$ one after another, up to row $2y+1$, and no row of the free range has to be computed.

The steps are expressed by these numbers: since $A_{n,\ell}=p_n-U^{(\ell+1)}_n$, the identity \eqref{eq:recursion-second} gives
\[
d_{n,j}=2\,U^{(j)}_n-U^{(j+1)}_n-U^{(j-1)}_n\qquad(j\ge 2,\ n\ge 0).
\]
With $m=k+\zeta$ we have in particular
\begin{align*}
d_{k+\zeta+1,\,k}&=2\,U^{(k)}_{m+1}-U^{(k+1)}_{m+1}-U^{(k-1)}_{m+1},\\
d_{k+\zeta,\,k+1}&=2\,U^{(k+1)}_{m}-U^{(k+2)}_{m}-U^{(k)}_{m}.
\end{align*}
Rows $m$ and $m+1$ lie in the range $n\le 2y+1$ of the recursion \eqref{eq:tables-forward} for each of the four thresholds $y$ that occur: the strongest of the conditions is $m+1\le 2(k-1)+1$, for threshold $k-1$, and this is part~(1). For each of the $256$ columns we compute the six numbers on the right-hand sides by this recursion. The largest row that occurs is row $300$, for $k=264$, and the integers that occur have up to about $740$ decimal digits. For $k=9$, for instance, where $\zeta=7$ and the two rows are $17$ and $16$, we find that
\begin{align*}
d_{17,9}&=18882703175184212>0,\\
d_{16,10}&=-12408007317596874<0.
\end{align*}
The second of these numbers agrees with the proof of Lemma~\ref{lem:sixteenrows}: by definition $d_{16,10}=b_{16,10}-b_{16,9}$, and the two entries of row $16$ displayed there, which are computed by the last-car recursion, have this difference. In the same way we find in each of the $256$ columns that $d_{k+\zeta+1,k}>0$ and that $d_{k+\zeta,k+1}<0$. These are parts~(2) and~(3).
\end{proof}

By part~(1) and since $\zeta(k)\ge 7$, the rows $k+\zeta(k)+1$ and $k+\zeta(k)$ lie in the chain range of column $k$, resp.\ of column $k+1$.

\subsection{Columns $2$ to $81$}\label{subsec:tables-columns}

For $2\le k\le 81$ two statements are supplied by computation: Proposition~\ref{prop:chainrange}, which concerns the finitely many rows $k+1,\dots,2k-1$ of the chain range, and the inequality $d_{n,k}\ge 0$ of Proposition~\ref{prop:leftpart}, which concerns all rows $n\ge\max(2k,17)$. For the latter, column $k$ is computed up to a row $n_0(k)$, and at this row we verify a condition from which the propagation theorem (Proposition~\ref{prop:propagation}) derives the inequality in all later rows. These rows are the following.

\begin{definition}\label{def:startrow}
For an integer $k$ with $2\le k\le 81$, the number $n_0(k)$ is given by the following table.
\[
\begin{array}{c|cccccccccc}
k & 2 & 3 & 4 & 5 & 6 & 7 & 8 & 9 & 10 & 11\\
n_0(k) & 7 & 10 & 13 & 15 & 18 & 20 & 23 & 25 & 28 & 30\\ \hline
k & 12 & 13 & 14 & 15 & 16 & 17 & 18 & 19 & 20 & 21\\
n_0(k) & 32 & 35 & 37 & 39 & 41 & 44 & 46 & 48 & 50 & 52\\ \hline
k & 22 & 23 & 24 & 25 & 26 & 27 & 28 & 29 & 30 & 31\\
n_0(k) & 55 & 57 & 59 & 61 & 63 & 66 & 68 & 70 & 72 & 74\\ \hline
k & 32 & 33 & 34 & 35 & 36 & 37 & 38 & 39 & 40 & 41\\
n_0(k) & 76 & 78 & 80 & 82 & 85 & 87 & 89 & 91 & 93 & 95\\ \hline
k & 42 & 43 & 44 & 45 & 46 & 47 & 48 & 49 & 50 & 51\\
n_0(k) & 97 & 99 & 101 & 103 & 105 & 107 & 109 & 111 & 114 & 116\\ \hline
k & 52 & 53 & 54 & 55 & 56 & 57 & 58 & 59 & 60 & 61\\
n_0(k) & 118 & 120 & 122 & 124 & 126 & 128 & 130 & 132 & 134 & 136\\ \hline
k & 62 & 63 & 64 & 65 & 66 & 67 & 68 & 69 & 70 & 71\\
n_0(k) & 138 & 140 & 142 & 144 & 146 & 148 & 150 & 152 & 154 & 156\\ \hline
k & 72 & 73 & 74 & 75 & 76 & 77 & 78 & 79 & 80 & 81\\
n_0(k) & 158 & 160 & 162 & 164 & 166 & 168 & 170 & 171 & 173 & 175
\end{array}
\]
\end{definition}

Recall that $x_k=(k/(k+1))^{k-1}$ is the critical tilt of column $k$ (Section~\ref{sec:critical}) and that the sequence $(\xi^{(k)}_t)_{t\ge 0}$ of column $k$ is defined in Section~\ref{sec:series}.

\begin{lemma}\label{lem:columntable}
Let $k$ be an integer with $2\le k\le 81$, and let $n_0=n_0(k)$. Then
\begin{enumerate}
\item $2k\le n_0$;
\item for every integer $n$ with $k+1\le n\le 2k-2$: if $b_{n,k-1}<b_{n,k}$, then $b_{n+1,k-1}<b_{n+1,k}$;
\item for every integer $n$ with $2k\le n$, $17\le n$ and $n<n_0$: $b_{n,k-1}\le b_{n,k}$;
\item for every integer $s$ with $k\le s\le n_0$: $\sum_{t=s}^{n_0}\xi^{(k)}_t\,x_k^{t}\ge 0$.
\end{enumerate}
\end{lemma}

\begin{proof}
Part~(1) is read from the table of Definition~\ref{def:startrow}, in which $n_0(k)\ge 2k+3$ for every $k$. We write $\xi_t=\xi^{(k)}_t$.

For $\ell=k$, $k-1$ and $k-2$ we compute $A_{n,\ell}$ for $0\le n\le n_0-1$, starting from $A_{0,\ell}=1$, by the last-car recursion of Proposition~\ref{prop:lastcar}, as in the proof of Lemma~\ref{lem:sixteenrows}. By \eqref{eq:recursion-second} these numbers give the steps
\[
d_{n,k}=A_{n,k}-2A_{n,k-1}+A_{n,k-2}\qquad(0\le n\le n_0-1).
\]
They also give the numbers $\Xi_t=t!\,\xi_t$ for $1\le t\le n_0$: by the definition of $\xi_t$ and by $z^{(y)}_{t-1}=A_{t-1,y-1}/(t-1)!$, these are the integers
\[
\Xi_t=\min(t,k+1)\,A_{t-1,k}-2\min(t,k)\,A_{t-1,k-1}+\min(t,k-1)\,A_{t-1,k-2}.
\]
Parts~(2) and~(3) are statements about the signs of the steps in rows $k+1,\dots,2k-1$, resp.\ in rows $\max(2k,17),\dots,n_0-1$. No row beyond $n_0-1$ occurs in them, because $2k-1\le n_0-1$ by part~(1).

For part~(4) we clear the denominators. For $k\le s\le n_0$ let
\[
\mathsf{S}_s=\sum_{t=s}^{n_0}\Xi_t\,\frac{n_0!}{t!}\,(k+1)^{(k-1)(n_0-t)}\,k^{(k-1)(t-s)}.
\]
These numbers are integers, and they are computed one after another for $s=n_0,n_0-1,\dots,k$, from $\mathsf{S}_{n_0+1}=0$ and
\[
\mathsf{S}_s=\Xi_s\,\frac{n_0!}{s!}\,(k+1)^{(k-1)(n_0-s)}+k^{k-1}\,\mathsf{S}_{s+1}.
\]
Since $\Xi_t\,n_0!/t!=n_0!\,\xi_t$ and $x_k=k^{k-1}/(k+1)^{k-1}$, the definition of $\mathsf{S}_s$ reads
\[
\mathsf{S}_s=\frac{n_0!\,(k+1)^{(k-1)n_0}}{k^{(k-1)s}}\,\sum_{t=s}^{n_0}\xi_t\,x_k^{t}.
\]
The factor in front of the sum is positive, so part~(4) says that $\mathsf{S}_s\ge 0$ for every $s$ with $k\le s\le n_0$.

We carry out this computation for each of the eighty columns. The largest row for which the numbers $A_{n,\ell}$ are computed is row $174$, for $k=81$, and in that column the integers $\mathsf{S}_s$ have up to about $14700$ decimal digits. For $k=2$, for instance, where $n_0=7$ and no row occurs in parts~(2) and~(3), the values of $\Xi_t$ for $1\le t\le 7$ and of $\mathsf{S}_t$ for $2\le t\le 7$ are
\[
\begin{array}{c|rrrrrrr}
t & 1 & 2 & 3 & 4 & 5 & 6 & 7\\ \hline
\Xi_t & 0 & -1 & -1 & 2 & 51 & 710 & 9819\\
\mathsf{S}_t & & 3912 & 308136 & 188088 & 88374 & 34548 & 9819
\end{array}
\]
and the six numbers $\mathsf{S}_t$ are positive. In each of the eighty columns we find the following. For every $n$ with $k+1\le n\le 2k-2$ we have $d_{n,k}\le 0$ or $d_{n+1,k}>0$, which is part~(2). For every $n$ with $\max(2k,17)\le n\le n_0-1$ we have $d_{n,k}\ge 0$, which is part~(3). And $\mathsf{S}_s\ge 0$ for every $s$ with $k\le s\le n_0$, which is part~(4).
\end{proof}

\section{The crossing criterion}\label{sec:crossing}
For the columns $k\ge 265$, every positive row from which the later arguments start is obtained from one statement, the crossing criterion (Proposition~\ref{prop:criterion}). For a column $k\ge 2$ and an integer $R$ with $1\le R\le k-1$ it names two inequalities between finite sums that depend on $k$ and $R$ only, and it asserts that, when both hold, column $k$ is positive in some row $j$ with $k<j\le k+R$. The sums are replaced by bounds in closed form in Section~\ref{sec:strip}, where the criterion is applied with $R$ of order $k^{2/3}$. The sums are formed at the diagonal tilt of the column, which we introduce first.

\subsection{The diagonal tilt}\label{subsec:crossing-tilt}

Let $k\ge 2$ be an integer. The \emph{diagonal tilt} of column $k$ is the positive number
\[
\vartheta_k=\frac{1}{\gamma_k}=\Bigl(\frac{k}{k+1}\Bigr)^{k-2},
\]
where $\gamma_k=(1+1/k)^{k-2}$ is the factor in the diagonal lower bound, Proposition~\ref{prop:diagonallower}; recall from Section~\ref{sec:series} that $\gamma_k$ is nondecreasing in $k$. When a column $k$ is fixed, a tilde denotes the tilted form of a sequence at this tilt:
\[
\tilde\varepsilon_m=\varepsilon_m\vartheta_k^m,\qquad
\tilde\beta^{(\ell)}_j=\beta^{(\ell)}_j\vartheta_k^j,\qquad
\tilde\delta_n=\delta^{(k)}_n\vartheta_k^n .
\]
Since $\vartheta_k>0$, each $\tilde\delta_n$ has the sign of $d_{n,k}$. The diagonal tilt is not the critical tilt $x_k$ of Section~\ref{sec:critical}; the two are related by $x_k=\frac{k}{k+1}\,\vartheta_k$.

The tilt is chosen so that $\vartheta_k\gamma_k=1$, which removes the factor $\gamma_k$ from the diagonal lower bound. Indeed, this bound says that $\gamma_k\,\beta^{(k-1)}_n\le\beta^{(k)}_{n+1}$ for every $n\ge 0$, and $\vartheta_k^{n}=\vartheta_k^{n+1}\gamma_k$. Hence, whenever column $k$ is nonpositive in row $n+1$, i.e.\ whenever $\beta^{(k)}_{n+1}\le\beta^{(k-1)}_{n+1}$, we have
\begin{equation}\label{eq:crossing-monotone}
\tilde\beta^{(k-1)}_n=\vartheta_k^{n+1}\gamma_k\,\beta^{(k-1)}_n\le\vartheta_k^{n+1}\beta^{(k)}_{n+1}\le\vartheta_k^{n+1}\beta^{(k-1)}_{n+1}=\tilde\beta^{(k-1)}_{n+1}.
\end{equation}
Thus the tilted counts of level $k-1$ do not decrease as long as column $k$ has not turned positive. Since $\gamma_i$ is nondecreasing in $i$, we also have
\[
\vartheta_k\gamma_i\ge\vartheta_k\gamma_k=1\qquad\text{for } i\ge k .
\]

The following lemma locates the tilt relative to $e^{-1}$. For a column $k$ we put
\[
\bar a=\frac{5}{2k},\qquad \underline a=\frac{5}{2k}-\frac{4}{3k^2};
\]
the proof of the lemma shows that these two numbers enclose $1+\log\vartheta_k$.

\begin{lemma}\label{lem:tiltbounds}
Let $k\ge 2$ be an integer. Then the diagonal tilt of column $k$ satisfies
\[
e^{-1}\le\vartheta_k\le e^{-1}e^{\bar a}.
\]
\end{lemma}

\begin{proof}
We show that
\[
\underline a\le 1+\log\vartheta_k\le\bar a ,
\]
which gives the lemma because $\underline a=(15k-8)/(6k^2)\ge 0$. By Taylor's formula with remainder,
\[
u-\frac{u^2}{2}\le\log(1+u)\le u-\frac{u^2}{2}+\frac{u^3}{3}
\]
for every real number $u\ge 0$. We apply this with $u=1/k$ to
\[
1+\log\vartheta_k=1-(k-2)\,\log\Bigl(1+\frac1k\Bigr).
\]
Expanding the products, we find that
\[
\frac{5}{2k}-\frac{4}{3k^2}+\frac{2}{3k^3}\le 1+\log\vartheta_k\le\frac{5}{2k}-\frac{1}{k^2},
\]
where the left-hand side is at least $\underline a$ and the right-hand side is at most $\bar a$.
\end{proof}

\subsection{Statement and proof of the criterion}\label{subsec:crossing-criterion}

We fix a column $k\ge 2$ and introduce the notation in which the criterion is stated. For an integer $M\ge 0$ we set
\[
\chi_M=1-\vartheta_k-\sum_{m=1}^{M}\tilde\varepsilon_m .
\]
Since $\varepsilon_m\ge 0$ and $\varepsilon_1=1$, we have $\tilde\varepsilon_m\ge 0$ and $\tilde\varepsilon_1=\vartheta_k$; hence
\[
\chi_0=1-\vartheta_k,\qquad \chi_1=1-2\vartheta_k,\qquad \chi_0\ge\chi_1\ge\chi_2\ge\cdots .
\]
For an integer $R\ge 0$ let
\[
S_R=\sum_{i=0}^{R-1}\frac{(i\,\vartheta_k)^i}{i!},
\]
in which the term with $i=0$ is $1$. Since $m\,\varepsilon_m=(m-1)^{m-1}/(m-1)!$ for $m\ge 1$, we have
\[
\sum_{m=1}^{R}m\,\tilde\varepsilon_m=\vartheta_kS_R ,
\]
so that $\vartheta_kS_R$ is the first moment of the tilted kernel up to index $R$.

The remaining quantities are built from the numbers
\[
\pi_j=\frac{\vartheta_k^j\,a_j}{j}\qquad(j\ge 1),
\]
which are positive. The ratio formula \eqref{eq:prelim-ratio} and the definition of $\gamma_i$ give
\[
a_{i+1}=\Bigl(1+\frac1i\Bigr)^{i-1}a_i=\gamma_i\,\frac{i+1}{i}\,a_i\qquad(i\ge 2).
\]
Hence, since $\vartheta_k\gamma_j\ge\vartheta_k\gamma_k=1$ for $j\ge k$,
\[
\pi_{j+1}=\vartheta_k\gamma_j\,\pi_j\quad(j\ge 2),\qquad \pi_k=\pi_{k+1}\le\pi_{k+2}\le\pi_{k+3}\le\cdots .
\]
Thus the numbers $\pi_n$ do not decrease from $n=k$ on.
For an integer $i$ let
\[
\Pi_i=\sum_{n=k}^{i}\pi_n
\]
be the partial sums of the numbers $\pi_n$ from index $k$ on, so that $\Pi_i=0$ for $i<k$. If these numbers were constant, the quotient $\Pi_{j+1}/\pi_{j+1}$ would be the number $j-k+2$ of terms of $\Pi_{j+1}$. For $j\ge k$ we denote by $\mathsf e_j$ the amount by which the quotient falls short of this number,
\[
\mathsf e_j=(j-k+2)-\frac{\Pi_{j+1}}{\pi_{j+1}}=\sum_{n=k}^{j+1}\Bigl(1-\frac{\pi_n}{\pi_{j+1}}\Bigr).
\]
Since $\pi_n\le\pi_{j+1}$ for $k\le n\le j+1$, every term of the last sum is nonnegative, so that $\mathsf e_j\ge 0$, and $\mathsf e_k=0$ because $\pi_k=\pi_{k+1}$. We put
\[
\mathsf E=\mathsf E(k,R)=\sum_{j=k}^{k+R-2}\frac{\mathsf e_j}{j+1},
\]
which is nonnegative and vanishes for $R\le 1$.

Finally, for $i\ge 2$ we set
\[
\phi_i=1-\frac{i}{(i+1)\,\vartheta_k\gamma_i}=1-\frac{\gamma_k\,c_{i-1}}{c_i};
\]
the two expressions agree because $c_i/c_{i-1}=a_{i+1}/a_i=\gamma_i\,(i+1)/i$ and $\gamma_k=1/\vartheta_k$. For $i\ge k$ we have $0<\phi_i<1$, since $\vartheta_k\gamma_i\ge 1$. We also put
\[
\Sigma_\phi=\Sigma_\phi(k,R)=\sum_{i=k+1}^{k+R-1}(i-k+1)\,\phi_i ,
\]
which is likewise nonnegative and vanishes for $R\le 1$.

The proof of the criterion argues by contradiction. If column $k$ is nonpositive in rows $k+1,\dots,k+R$, the tilted counts $\tilde\beta^{(k-1)}_j$ of level $k-1$ are largest in row $k+R$, and the level equation at the diagonal tilt bounds this largest count by multiples of one and the same term of the equation, from above when it is read in row $k+R$ and from below when it is summed over rows $k,\dots,k+R$; the lower bound contains a correction, which the proof controls with the column equation. When the two bounds are combined, the sizes of this term and of the largest count, which are unknown, cancel, and an inequality remains that contradicts the second of the two inequalities that the criterion assumes. In this notation the criterion reads as follows.

\begin{proposition}\label{prop:criterion}
Let $k\ge 2$ and $R$ be integers with $1\le R\le k-1$, and put $N=k+R$. Suppose that the quantities $\chi_R$, $S_R$, $\mathsf e_{N-1}$, $\mathsf E=\mathsf E(k,R)$ and $\Sigma_\phi=\Sigma_\phi(k,R)$ of column $k$ satisfy
\[
\chi_R\ge 0
\]
and
\begin{equation}\label{eq:crossing-hypothesis}
\Bigl(\chi_R+\frac{\vartheta_k\,(R+1)}{N}\Bigr)\bigl(\mathsf e_{N-1}+\Sigma_\phi\bigr)<\vartheta_k\,\Bigl(\frac{R(R+1)}{N}-1-S_R-\mathsf E\Bigr).
\end{equation}
Then there is an integer $j$ with $k<j\le N$ and $d_{j,k}>0$.
\end{proposition}

\begin{proof}
We write $\delta_j=\delta^{(k)}_j$ and $\lambda=\chi_R+\vartheta_k(R+1)/N$, the first factor on the left of \eqref{eq:crossing-hypothesis}. We assume that the conclusion fails, i.e.
\[
\delta_j\le 0\qquad\text{for all } k<j\le N,
\]
and we derive
\begin{equation}\label{eq:crossing-opposite}
\vartheta_k\,\Bigl(\frac{R(R+1)}{N}-1-S_R-\mathsf E\Bigr)\le\lambda\,\bigl(\mathsf e_{N-1}+\Sigma_\phi\bigr),
\end{equation}
which contradicts the second hypothesis.

We first record two consequences of $\chi_R\ge 0$. Since $R\ge 1$ and the numbers $\chi_M$ do not increase with $M$, we have
\[
\chi_M\ge\chi_R\ge 0\quad(0\le M\le R)\qquad\text{and}\qquad \vartheta_k\le\tfrac12 ,
\]
the latter because $1-2\vartheta_k=\chi_1\ge\chi_R$. In particular $\lambda>0$.

Next we normalise by the number of all parking functions of length $n$, which, we recall, is $c_n\,n!$. For $n\ge 0$ we put
\[
F_n=\frac{z^{(k)}_n}{c_n},\qquad F^-_n=\frac{z^{(k-1)}_n}{c_n},\qquad P_n=\frac{\beta^{(k-1)}_n}{c_n}.
\]
These are the proportions, among all parking functions of length $n$, of those whose maximum displacement is at most $k-1$, at most $k-2$, and exactly $k-1$. Thus $P_n\ge 0$ and
\[
F_n=F^-_n+P_n .
\]
By Lemma~\ref{lem:freerange}, $z^{(y)}_n=c_n$ for $n\le y$, and therefore
\[
F_n=1\quad(n\le k),\qquad F^-_n=1\ \text{ and }\ P_n=0\quad(n\le k-1).
\]
For $i\ge 1$ let $J_i=F_i-F^-_{i-1}$; this is the increase of the cumulative proportion when the length and the bound on the displacements are both raised by one. We have $J_i\ge 0$ by Proposition~\ref{prop:diagonalmonotone}, applied with threshold $k-1$, and $J_k=0$ because $F_k=F^-_{k-1}=1$. Since $F_n=F^-_n+P_n$, it follows that
\[
F_n-F_{n+1}=P_n-J_{n+1}\qquad(n\ge 0).
\]

We multiply the level equation of Lemma~\ref{lem:levelequation} for level $k-1$ by $\vartheta_k^n$. Since $a_n=c_{n-1}$, its first term becomes $\vartheta_k^n\,z^{(k)}_{n-1}/n=\pi_nF_{n-1}$, and the equation reads
\begin{equation}\label{eq:crossing-level}
\tilde\beta^{(k-1)}_n=\pi_n\,F_{n-1}+\frac{(k-1)\,\vartheta_k}{n}\,\tilde\beta^{(k-1)}_{n-1}+\sum_{m=1}^{n-k}\tilde\varepsilon_m\,\tilde\beta^{(k-1)}_{n-m}
\end{equation}
for $k\le n\le 2k-1$, where $\tilde\beta^{(k-1)}_n=0$ for $n\le k-1$ because $P_n=0$ for these $n$. In the same way the column equation of Proposition~\ref{prop:columnequation}, whose source term is $\frac{2}{n}\,\beta^{(k-1)}_{n-1}$, becomes
\begin{equation}\label{eq:crossing-step}
\tilde\delta_n=2\,\pi_n\,P_{n-1}+\frac{(k+1)\,\vartheta_k}{n}\,\tilde\delta_{n-1}+\sum_{m=1}^{n-k}\tilde\varepsilon_m\,\tilde\delta_{n-m}
\end{equation}
for $k+1\le n\le 2k-1$, and the initial value $\delta_k=-a_k/k$ of Lemma~\ref{lem:initialstep} becomes $\tilde\delta_k=-\pi_k$. Since $R\le k-1$, we have $N\le 2k-1$, so that the tilted level equation holds in rows $k,\dots,N$ and the tilted column equation in rows $k+1,\dots,N$.

Our assumption implies that the tilted counts of level $k-1$ are largest in row $N$. Indeed, for $k\le n<N$, column $k$ is nonpositive in row $n+1$, and the chain of inequalities \eqref{eq:crossing-monotone} gives $\tilde\beta^{(k-1)}_n\le\tilde\beta^{(k-1)}_{n+1}$. Since the tilted counts are nonnegative and vanish for $j\le k-1$, it follows that
\[
0\le\tilde\beta^{(k-1)}_j\le\tilde\beta^{(k-1)}_N\qquad\text{for every } j\le N .
\]

An upper bound for $\tilde\beta^{(k-1)}_N$ now follows from the equation of row $N$. In the tilted level equation \eqref{eq:crossing-level} with $n=N$ the coefficients are nonnegative, and, as we have just seen, every tilted count on the right is at most $\tilde\beta^{(k-1)}_N$. Hence
\[
\tilde\beta^{(k-1)}_N\le\pi_N\,F_{N-1}+\Bigl(\frac{(k-1)\,\vartheta_k}{N}+\sum_{m=1}^{R}\tilde\varepsilon_m\Bigr)\,\tilde\beta^{(k-1)}_N .
\]
By the definition of $\chi_R$ and since $N-(k-1)=R+1$, the bracket equals $1-\lambda$. Therefore
\begin{equation}\label{eq:crossing-top}
\lambda\,\tilde\beta^{(k-1)}_N\le\pi_N\,F_{N-1}.
\end{equation}

The lower bound for $\tilde\beta^{(k-1)}_N$ comes from an identity that involves all rows $k,\dots,N$. We claim that, without any assumption on signs,
\begin{multline}\label{eq:crossing-summed}
F_{N-1}\,\Pi_N=\tilde\beta^{(k-1)}_N+\Bigl(1-\vartheta_k-\frac{(k-1)\,\vartheta_k}{N}\Bigr)\,\tilde\beta^{(k-1)}_{N-1}\\
+\sum_{j=k}^{N-2}\Bigl(\chi_{N-j}+\frac{\vartheta_k\,\mathsf e_j}{j+1}\Bigr)\,\tilde\beta^{(k-1)}_j+\sum_{i=k}^{N-1}\Pi_i\,J_i .
\end{multline}
To see this, we add the tilted level equations \eqref{eq:crossing-level} for $n=k,\dots,N$ and collect the coefficient of each $\tilde\beta^{(k-1)}_j$ on the right. Since $\tilde\beta^{(k-1)}_{k-1}=0$, this gives
\[
\sum_{n=k}^{N}\pi_n\,F_{n-1}=\tilde\beta^{(k-1)}_N+\sum_{j=k}^{N-1}\Bigl(1-\frac{(k-1)\,\vartheta_k}{j+1}-\sum_{m=1}^{N-j}\tilde\varepsilon_m\Bigr)\,\tilde\beta^{(k-1)}_j .
\]
We sum the left-hand side by parts: since $F_{i-1}-F_i=P_{i-1}-J_i$,
\[
\sum_{n=k}^{N}\pi_n\,F_{n-1}=F_{N-1}\,\Pi_N+\sum_{i=k}^{N-1}\Pi_i\,(P_{i-1}-J_i).
\]
The terms with $P$ are multiples of the tilted counts: $P_{k-1}=0$, and for $j\ge k$ we have $\pi_{j+1}P_j=\vartheta_k\,\tilde\beta^{(k-1)}_j/(j+1)$, again by $a_{j+1}=c_j$, and $\Pi_{j+1}=\pi_{j+1}\,(j-k+2-\mathsf e_j)$ by the definition of $\mathsf e_j$. So
\[
\sum_{i=k}^{N-1}\Pi_i\,P_{i-1}=\sum_{j=k}^{N-2}\frac{\vartheta_k\,(j-k+2-\mathsf e_j)}{j+1}\,\tilde\beta^{(k-1)}_j .
\]
Equating the two expressions for $\sum_n\pi_nF_{n-1}$, we find that the coefficient of $\tilde\beta^{(k-1)}_j$ is $\chi_{N-j}+\vartheta_k\,\mathsf e_j/(j+1)$ for $k\le j\le N-2$, because $(k-1)+(j-k+2)=j+1$, and $1-\vartheta_k-(k-1)\vartheta_k/N$ for $j=N-1$, because $\tilde\varepsilon_1=\vartheta_k$. This is the claimed identity.

Besides multiples of the tilted counts, the right-hand side of this identity contains the numbers $J_i$, and we bound them next: for every $i\ge k$,
\[
J_i\le\phi_i\,(1-F^-_{i-1}).
\]
For the proof we compare the parking functions that the two cumulative counts leave out. Since $\beta^{(\ell)}_n=z^{(\ell+1)}_n-z^{(\ell)}_n$ for $\ell\ge 1$, and since $z^{(n)}_n=c_n$ for $n\ge 1$ by Lemma~\ref{lem:freerange}, we have the telescoping sums
\[
c_{i-1}-z^{(k-1)}_{i-1}=\sum_{s=0}^{i-k-1}\beta^{(k-1+s)}_{i-1},\qquad c_i-z^{(k)}_i=\sum_{s=0}^{i-k-1}\beta^{(k+s)}_i .
\]
The diagonal lower bound, Proposition~\ref{prop:diagonallower}, in column $k+s$ gives $\gamma_{k+s}\,\beta^{(k-1+s)}_{i-1}\le\beta^{(k+s)}_i$. Here $\gamma_k\le\gamma_{k+s}$, and the normalised counts are nonnegative. Summing over $s$, we obtain
\[
\gamma_k\,\bigl(c_{i-1}-z^{(k-1)}_{i-1}\bigr)\le c_i-z^{(k)}_i .
\]
Dividing by $c_i$, we obtain $(1-\phi_i)\,(1-F^-_{i-1})\le 1-F_i$, because $\gamma_k\,c_{i-1}/c_i=1-\phi_i$, and hence $J_i=(1-F^-_{i-1})-(1-F_i)\le\phi_i\,(1-F^-_{i-1})$.

The identity \eqref{eq:crossing-summed} also contains the proportion $F_{N-1}$, and the bound for $J_i$ contains the proportions $F^-_{i-1}$. We bound them by means of the column equation, through the inequality
\[
\sum_{j=k}^{N-1}P_j\le\frac12 ,
\]
which says that the proportions of parking functions with maximum displacement exactly $k-1$, summed over rows $k,\dots,N-1$, are at most $\frac12$ as long as column $k$ has not turned positive. To prove it, we add the tilted column equations \eqref{eq:crossing-step} for $n=k+1,\dots,N$ and collect the terms as before:
\[
2\sum_{n=k+1}^{N}\pi_n\,P_{n-1}=\tilde\delta_N-\tilde\delta_k+\sum_{i=k}^{N-1}\Bigl(1-\frac{(k+1)\,\vartheta_k}{i+1}-\sum_{m=1}^{N-i}\tilde\varepsilon_m\Bigr)\,\tilde\delta_i .
\]
For $k\le i\le N-1$ the coefficient in the last sum is at least $1-\vartheta_k-\sum_{m=1}^{R}\tilde\varepsilon_m=\chi_R\ge 0$, because $(k+1)/(i+1)\le 1$ and $N-i\le R$. The values $\tilde\delta_i$ with $k\le i\le N$ are nonpositive, by $\tilde\delta_k=-\pi_k$ and by our assumption. So the last sum and $\tilde\delta_N$ are nonpositive, whence $2\sum_{n=k+1}^{N}\pi_nP_{n-1}\le-\tilde\delta_k=\pi_k$. Since $\pi_n\ge\pi_k$ for $n\ge k$ and $P_{n-1}\ge 0$, the left-hand side is at least $2\pi_k\sum_{j=k}^{N-1}P_j$, and the asserted inequality follows.

Since $F_n-F_{n+1}=P_n-J_{n+1}$ and $J_{n+1}\ge 0$, we have $F_n-F_{n+1}\le P_n$, and $F_k=1$; hence $1-F_n\le\sum_{j=k}^{n-1}P_j$ and, because $F_n=F^-_n+P_n$, also $1-F^-_n=1-F_n+P_n\le\sum_{j=k}^{n}P_j$ for $n\ge k$. Since $\sum_{j=k}^{N-1}P_j\le\frac12$, we obtain
\[
F_{N-1}\ge\frac12\qquad\text{and}\qquad 1-F^-_n\le\frac12\quad(k\le n\le N-1).
\]

We now combine these estimates, beginning with an upper bound for the right-hand side of the identity \eqref{eq:crossing-summed}. The coefficients of the tilted counts there are nonnegative: by the two consequences of $\chi_R\ge 0$ recorded at the beginning, $1-\vartheta_k-(k-1)\vartheta_k/N\ge 1-2\vartheta_k\ge 0$ and $\chi_{N-j}\ge 0$ for $k\le j\le N-2$, and $\mathsf e_j\ge 0$. Since the tilted counts are largest in row $N$, we may replace every $\tilde\beta^{(k-1)}_j$ by $\tilde\beta^{(k-1)}_N$, and the terms that contain the tilted counts are together at most $L\,\tilde\beta^{(k-1)}_N$, where
\[
L=1+\Bigl(1-\vartheta_k-\frac{(k-1)\,\vartheta_k}{N}\Bigr)+\sum_{M=2}^{R}\chi_M+\vartheta_k\,\mathsf E
\]
is nonnegative. In the sum $\sum_{i=k}^{N-1}\Pi_iJ_i$ the summand with $i=k$ vanishes because $J_k=0$. For $k<i<N$ the bound $J_i\le\phi_i\,(1-F^-_{i-1})$ and the inequality $1-F^-_{i-1}\le\frac12$ give $0\le J_i\le\phi_i/2$, and $\Pi_i\le(i-k+1)\,\pi_i\le(i-k+1)\,\pi_N$ because the numbers $\pi_n$ do not decrease from $n=k$ on. Hence $\sum_{i=k}^{N-1}\Pi_iJ_i\le\frac12\pi_N\Sigma_\phi$, and altogether
\[
F_{N-1}\,\Pi_N\le L\,\tilde\beta^{(k-1)}_N+\frac{\pi_N}{2}\,\Sigma_\phi .
\]
Now we eliminate $\tilde\beta^{(k-1)}_N$ and $F_{N-1}$. By the definition of $\mathsf e_{N-1}$ we have $\Pi_N=\pi_N\,(R+1-\mathsf e_{N-1})$. We multiply the last display by $\lambda$, use the upper bound \eqref{eq:crossing-top} in the form $L\,\lambda\,\tilde\beta^{(k-1)}_N\le L\,\pi_NF_{N-1}$, and divide by $\pi_N$:
\[
F_{N-1}\,\bigl[(R+1-\mathsf e_{N-1})\,\lambda-L\bigr]\le\frac{\lambda\,\Sigma_\phi}{2}.
\]
If the bracket is not positive, it is at most $\lambda\Sigma_\phi$, because $\lambda\Sigma_\phi\ge 0$. If it is positive, then $F_{N-1}\ge\frac12$ shows that half of it is at most $\frac12\lambda\Sigma_\phi$. In both cases
\begin{equation}\label{eq:crossing-eliminated}
(R+1)\,\lambda-L\le\lambda\,\bigl(\mathsf e_{N-1}+\Sigma_\phi\bigr).
\end{equation}
It remains to evaluate the left-hand side. For $M\le R$ we have $\chi_M-\chi_R=\sum_{m=M+1}^{R}\tilde\varepsilon_m$, and in the sum of these differences over $0\le M\le R$ the number $\tilde\varepsilon_m$ occurs $m$ times; hence, since $\sum_{m=1}^{R}m\,\tilde\varepsilon_m=\vartheta_kS_R$,
\[
\sum_{M=0}^{R}\chi_M=(R+1)\,\chi_R+\vartheta_k\,S_R .
\]
Removing $\chi_0=1-\vartheta_k$ and $\chi_1=1-2\vartheta_k$ and inserting the result into the definition of $L$, we get
\[
L=(R+1)\,\chi_R+\vartheta_k\,S_R+2\vartheta_k-\frac{(k-1)\,\vartheta_k}{N}+\vartheta_k\,\mathsf E ,
\]
and, since $(R+1)^2+(k-1)-2N=R(R+1)-N$,
\[
(R+1)\,\lambda-L=\vartheta_k\,\Bigl(\frac{R(R+1)}{N}-1-S_R-\mathsf E\Bigr).
\]
Inserted into \eqref{eq:crossing-eliminated}, this is the inequality \eqref{eq:crossing-opposite}, which contradicts the second hypothesis of the proposition. Hence our assumption is false, and there is an integer $j$ with $k<j\le N$ and $d_{j,k}>0$.
\end{proof}

\begin{remark}
Since $\chi_R\le\chi_1=1-2\vartheta_k$, the first hypothesis of the criterion requires $\vartheta_k\le\frac12$. Since $\vartheta_k=1/\gamma_k$ does not increase with $k$ and $\vartheta_7=(7/8)^5=16807/32768>\frac12$, the criterion says nothing about columns $k\le 7$. For a given column the same hypothesis limits how large $R$ may be taken: at tilt $e^{-1}$ the number $\chi_M$ would be the positive tail $\sum_{m>M}\varepsilon_me^{-m}$, by Lemma~\ref{lem:epsmass}, and since $\vartheta_k\ge e^{-1}$ by Lemma~\ref{lem:tiltbounds}, it is smaller than this tail by an excess that grows with $M$. In the second hypothesis $R(R+1)/N$ is of order $R^2/k$, while $S_R$ is at most of order $\sqrt R$ and the other terms are corrections; so this hypothesis can hold once $R^{3/2}$ exceeds a multiple of $k$.
\end{remark}

\section{The crossing row}\label{sec:strip}
The hypotheses of the crossing criterion, Proposition~\ref{prop:criterion}, are inequalities between finite sums whose number of terms grows with $R$. We now replace each of these sums by a bound in closed form and verify the resulting inequalities for $R$ of order $k^{2/3}$. The outcome is that every column $k\ge 265$ is positive in some row $j>k$ with $25\,(j-k-1)^3\le 121\,k^2$ (Proposition~\ref{prop:strip}) and, together with the computed signs for the columns below $265$, that every column $k\ge 82$ is positive in some row $k+r$ with $1\le r\le 27k/100$ (Corollary~\ref{cor:earlypositive}). The proofs of Propositions~\ref{prop:chainrange} and~\ref{prop:leftpart} start from such an early positive row, and the first bound confines the rows before the crossing to the range that the proof of Proposition~\ref{prop:nextcolumn} treats.

\subsection{The crossing criterion in closed form}\label{subsec:strip-closed}

Let $k\ge 2$ and $R\ge 1$ be integers, and put $N=k+R$. We use the quantities of column $k$ that are introduced in Section~\ref{sec:crossing}: the diagonal tilt $\vartheta_k=1/\gamma_k$, the numbers $\chi_R$ and $S_R$, the $\pi_j$ with their partial sums $\Pi_i$, the $\mathsf e_j$ with their sum $\mathsf E$, the $\phi_i$ with their sum $\Sigma_\phi$, and the numbers $\bar a=5/(2k)$ and $\underline a=5/(2k)-4/(3k^2)$; the proof of Proposition~\ref{prop:closedcriterion} recalls each definition where it is used. Recall also that $c_\ast=1/\sqrt{2\pi}$. For $R\ge 2$ we put
\begin{equation}\label{eq:strip-upper}
\begin{aligned}
\bar S&=1+c_\ast\Bigl[2\sqrt{R-1}-1+\frac{2\bar a}{3}\,R\sqrt R+\frac25\cdot\frac{633}{1000}\,\bar a^{2}R^{2}\sqrt R\Bigr],\\
\bar{\mathsf E}&=\frac{5\,(R-2)(R-1)\,R\,(R+1)}{24\,k^{3}},\\
\bar H&=(R-1)-k\,\log\frac{k+R+1}{k+2},\\
\bar G&=\frac{R^{2}-1}{2}-(k-1)(R-1)+k\,(k-1)\,\log\frac{k+R}{k+1},
\end{aligned}
\end{equation}
and
\begin{equation}\label{eq:strip-deficit}
\begin{aligned}
\underline t&=\frac{2e^{-1}c_\ast\bigl(1-\frac{1}{12R}\bigr)}{\sqrt{R+1}}-e^{-1}\,\bar a\,\Bigl[e^{\bar a}+1+c_\ast\bigl(2\sqrt{R-1}-1\bigr)\Bigr]\\
&\qquad-e^{-1}\,\frac{633}{1000}\,\bar a^{2}\,\Bigl[1+c_\ast\Bigl(\frac23\,R\sqrt R+2\sqrt{R-1}-1\Bigr)\Bigr],\\
\bar t&=\frac{2e^{-1}c_\ast}{\sqrt{R-1}}-e^{-1}\,\underline a\,\Bigl[2+c_\ast\Bigl(2\sqrt R-2-\frac14\Bigr)\Bigr].
\end{aligned}
\end{equation}
We shall show that $\bar S$ bounds $S_R$, that $\bar{\mathsf E}$ bounds $\mathsf E$ and that $\bar H+\frac5k\bar G$ bounds $\mathsf e_{N-1}+\Sigma_\phi$, all from above, and that $\underline t\le\chi_R\le\bar t$. We also write
\begin{equation}\label{eq:strip-three}
\mathsf{A}=e\,\bar t+e^{\bar a}\,\frac{R+1}{N},\qquad\mathsf{B}=\bar H+\frac5k\,\bar G,\qquad\mathsf{C}=\frac{R(R+1)}{N}-1-\bar S-\bar{\mathsf E}.
\end{equation}
With these bounds the criterion takes the following form, in which the last hypothesis is the inequality $\mathsf{A}\,\mathsf{B}<\mathsf{C}$ multiplied by $e^{-1}$.

\begin{proposition}\label{prop:closedcriterion}
Let $k\ge 2$ and $R$ be integers with $2\le R\le k-1$, and put $N=k+R$. Let $\bar S$, $\bar{\mathsf E}$, $\bar H$, $\bar G$, $\underline t$ and $\bar t$ be the numbers defined in \eqref{eq:strip-upper} and \eqref{eq:strip-deficit} for these $k$ and $R$. Suppose that
\[
\bar a\,R\le\frac{6679}{10000},\qquad\underline t\ge 0,
\]
and
\begin{equation}\label{eq:strip-closed}
\Bigl(\bar t+e^{-1}e^{\bar a}\,\frac{R+1}{N}\Bigr)\Bigl(\bar H+\frac5k\,\bar G\Bigr)<e^{-1}\Bigl(\frac{R(R+1)}{N}-1-\bar S-\bar{\mathsf E}\Bigr).
\end{equation}
Then there is an integer $j$ with $k<j\le N$ and $d_{j,k}>0$.
\end{proposition}

\begin{proof}
By Proposition~\ref{prop:criterion} it suffices to show that $\chi_R\ge 0$ and that
\begin{equation}\label{eq:strip-criterion}
\Bigl(\chi_R+\frac{\vartheta_k\,(R+1)}{N}\Bigr)\bigl(\mathsf e_{N-1}+\Sigma_\phi\bigr)<\vartheta_k\Bigl(\frac{R(R+1)}{N}-1-S_R-\mathsf E\Bigr).
\end{equation}
We first prove the bounds
\begin{equation}\label{eq:strip-five}
\begin{gathered}
S_R\le\bar S,\qquad\mathsf E\le\bar{\mathsf E},\\
0\le\mathsf e_{N-1}+\Sigma_\phi\le\bar H+\frac5k\,\bar G,\qquad\underline t\le\chi_R\le\bar t.
\end{gathered}
\end{equation}
We write
\[
\sigma=1+\log\vartheta_k,\qquad\text{so that}\qquad\vartheta_k=e^{\sigma-1},
\]
and we use throughout that for real $t\ge 0$ the Taylor polynomial of $\log(1+t)$ of even (resp.\ odd) degree bounds it from below (resp.\ above).

The bounds that we have listed rest on two estimates for the tilt, which we prove first. By Lemma~\ref{lem:tiltbounds} we have $e^{-1}\le\vartheta_k\le e^{-1}e^{\bar a}$, and its proof shows more precisely that $\underline a\le 1+\log\vartheta_k\le\bar a$. Since $\underline a=(15k-8)/(6k^2)\ge 0$, this gives
\[
0\le\underline a\le\sigma\le\bar a.
\]
We also need a bound for the products $\vartheta_k\gamma_i$: for every integer $i\ge k$,
\begin{equation}\label{eq:strip-loggamma}
0\le\log(\vartheta_k\gamma_i)\le\frac{5\,(i-k)}{2\,k\,i}.
\end{equation}
The lower bound holds because $\vartheta_k\gamma_k=1$ and the sequence $(\gamma_i)_{i\ge 2}$ is nondecreasing (Section~\ref{sec:series}). For the upper bound we write $\log(\vartheta_k\gamma_i)=\log\gamma_i-\log\gamma_k$ as the sum of the differences $\log\gamma_{m+1}-\log\gamma_m$ over $k\le m<i$. For $m\ge 2$ we bound $\log\gamma_{m+1}=(m-1)\log(1+\frac{1}{m+1})$ from above by the Taylor polynomial of degree $3$ and $\log\gamma_m=(m-2)\log(1+\frac1m)$ from below by the Taylor polynomial of degree $4$. A direct computation with these two bounds gives
\begin{align*}
\log\gamma_{m+1}-\log\gamma_m&\le\frac{5}{2m(m+1)}-\frac{29m^4+21m^3+m^2+7m+6}{12\,m^4(m+1)^3}\\
&\le\frac52\Bigl(\frac1m-\frac1{m+1}\Bigr),
\end{align*}
and summing over $k\le m<i$ we obtain the asserted upper bound.

We turn to the first of the bounds, the inequality $S_R\le\bar S$. Recall that $S_R=\sum_{i=0}^{R-1}(i\vartheta_k)^i/i!$. Since $\vartheta_k=e^{\sigma-1}$, the term with index $i$ is $e^{\sigma i}\cdot e^{-i}i^i/i!$, the Poisson mass at $i$ multiplied by $e^{\sigma i}$. To bound the exponential we use the elementary inequality
\[
e^s\le 1+s+\frac{633}{1000}\,s^2\qquad\Bigl(0\le s\le\frac{6679}{10000}\Bigr),
\]
which holds because the factor $\sum_{j\ge 0}s^j/(j+2)!$ of $s^2$ in the exponential series is smaller than $0.6328$ on this interval. For a real number $y$ with $0\le y\le R$ we have $0\le\sigma y\le\bar ay\le\bar aR\le\frac{6679}{10000}$, because $0\le\sigma\le\bar a$ and by the hypothesis on $\bar aR$, so that the inequality for $e^s$ gives
\begin{equation}\label{eq:strip-quad}
e^{\sigma y}\le e^{\bar ay}\le 1+\bar a\,y+\frac{633}{1000}\,\bar a^{2}y^{2}\qquad(0\le y\le R).
\end{equation}
The Poisson mass at $0$ is $1$, and the Poisson mass at $i\ge 1$ is at most $c_\ast/\sqrt i$ by \eqref{eq:prelim-stirling}. Therefore
\begin{equation}\label{eq:strip-poisson}
\begin{aligned}
\sum_{i=0}^{R-1}\frac{e^{-i}i^i}{i!}&\le 1+c_\ast\sum_{i=1}^{R-1}\frac{1}{\sqrt i}\le 1+c_\ast\bigl(2\sqrt{R-1}-1\bigr),\\
\sum_{i=0}^{R-1}i\,\frac{e^{-i}i^i}{i!}&\le c_\ast\sum_{i=1}^{R-1}\sqrt i\le\frac23\,c_\ast\,R\sqrt R,\\
\sum_{i=0}^{R-1}i^2\,\frac{e^{-i}i^i}{i!}&\le c_\ast\sum_{i=1}^{R-1}i\sqrt i\le\frac25\,c_\ast\,R^2\sqrt R,
\end{aligned}
\end{equation}
where the second inequality in each line is a comparison of the sum with an integral. By the bound for $e^{\sigma y}$ with $y=i$, the sum $S_R$ is at most the first of these three sums, plus $\bar a$ times the second, plus $\frac{633}{1000}\bar a^2$ times the third, and the three bounds add up to $\bar S$.

Next we prove that $\mathsf E\le\bar{\mathsf E}$. Recall that $\pi_j=\vartheta_k^j\,a_j/j$ and $\Pi_i=\pi_k+\dots+\pi_i$, that $\mathsf e_j=(j-k+2)-\Pi_{j+1}/\pi_{j+1}$, and that $\mathsf E=\sum_{j=k}^{N-2}\mathsf e_j/(j+1)$. The numbers $\pi_j$ are positive, and for $n\ge 2$,
\[
\pi_{n+1}=\vartheta_k\gamma_n\,\pi_n,
\]
because $a_{n+1}=(1+\frac1n)^{n-1}a_n=\gamma_n\,\frac{n+1}{n}\,a_n$ by \eqref{eq:prelim-ratio}. Since $\vartheta_k\gamma_k=1$, we get $\pi_{k+1}=\pi_k$ and hence $\mathsf e_k=0$. For $j\ge k$, the relations $\Pi_{j+2}=\Pi_{j+1}+\pi_{j+2}$ and $\pi_{j+2}=\vartheta_k\gamma_{j+1}\,\pi_{j+1}$ give
\[
\mathsf e_{j+1}=(j-k+2)\Bigl(1-\frac{1}{\vartheta_k\gamma_{j+1}}\Bigr)+\frac{\mathsf e_j}{\vartheta_k\gamma_{j+1}}.
\]
Since $\vartheta_k\gamma_{j+1}\ge\vartheta_k\gamma_k=1$, it follows by induction that $\mathsf e_j\ge 0$ for all $j\ge k$, and then, with $1-1/t\le\log t$ and $1/t\le 1$ for real $t\ge 1$, that
\begin{equation}\label{eq:strip-step}
\mathsf e_{j+1}\le(j-k+2)\,\log(\vartheta_k\gamma_{j+1})+\mathsf e_j\qquad(j\ge k).
\end{equation}
We claim that $\mathsf e_{k+m}\le\frac{5}{6k^2}\,m(m+1)(m+2)$ for every integer $m\ge 0$. This is clear for $m=0$, and the induction step follows from the last display with $j=k+m$ and from the bound \eqref{eq:strip-loggamma} with $i=k+m+1$, where $1/i\le 1/k$:
\begin{align*}
\mathsf e_{k+m+1}&\le(m+2)\,\frac{5\,(m+1)}{2k^2}+\frac{5}{6k^2}\,m(m+1)(m+2)\\
&=\frac{5}{6k^2}\,(m+1)(m+2)(m+3).
\end{align*}
Finally $\mathsf e_j/(j+1)\le\mathsf e_j/k$ for $j\ge k$, and $\sum_{m=0}^{R-2}m(m+1)(m+2)=\frac14(R-2)(R-1)R(R+1)$. Hence
\[
\mathsf E=\sum_{j=k}^{N-2}\frac{\mathsf e_j}{j+1}\le\frac{5}{6k^3}\sum_{m=0}^{R-2}m(m+1)(m+2)=\bar{\mathsf E}.
\]

We come to the two bounds for $\mathsf e_{N-1}+\Sigma_\phi$. Recall that $\phi_i=1-i/((i+1)\,\vartheta_k\gamma_i)$ and $\Sigma_\phi=\sum_{i=k+1}^{N-1}(i-k+1)\,\phi_i$. We have just seen that $\mathsf e_j\ge 0$ for all $j\ge k$; in particular $\mathsf e_{N-1}\ge 0$. Moreover $\phi_i\ge 0$ for $i\ge k$ because $\vartheta_k\gamma_i\ge 1$, so that $\Sigma_\phi\ge 0$. For the upper bound we write $\bar H(k,m)$ and $\bar G(k,m)$ for the expressions that define $\bar H$ and $\bar G$, with an integer $m\ge 1$ in place of $R$, so that $\bar H(k,1)=\bar G(k,1)=0$. The increments of these expressions dominate two families of summands: for every integer $m\ge 1$,
\[
\begin{aligned}
\frac{m+1}{k+m+1}&\le\bar H(k,m+1)-\bar H(k,m),\\
\frac{m(m+1)}{k+m}&\le\bar G(k,m+1)-\bar G(k,m).
\end{aligned}
\]
Both inequalities follow by a direct computation from the bounds for $\log(1+t)$: the first increment is $1-k\log\bigl(1+\frac{1}{k+m+1}\bigr)\ge 1-\frac{k}{k+m+1}$, and the second is $m+\frac32-k+k(k-1)\log\bigl(1+\frac{1}{k+m}\bigr)$, which by the Taylor polynomial of degree $2$ exceeds $\frac{m(m+1)}{k+m}=m+1-k+\frac{k(k-1)}{k+m}$ by at least $\frac12-\frac{k(k-1)}{2(k+m)^2}\ge 0$. Now let $1\le m\le R-1$. By the bound \eqref{eq:strip-loggamma} with $i=k+m$ and by the inequality for the increment of $\bar G$,
\[
(m+1)\,\log(\vartheta_k\gamma_{k+m})\le\frac{5}{2k}\cdot\frac{m(m+1)}{k+m}\le\frac{5}{2k}\bigl(\bar G(k,m+1)-\bar G(k,m)\bigr).
\]
We use this twice. First, the inequality \eqref{eq:strip-step} with $j=k+m-1$ gives $\mathsf e_{k+m}-\mathsf e_{k+m-1}\le\frac{5}{2k}\bigl(\bar G(k,m+1)-\bar G(k,m)\bigr)$, and summing over $1\le m\le R-1$, with $\mathsf e_k=0$ and $\bar G(k,1)=0$, we get $\mathsf e_{N-1}\le\frac{5}{2k}\,\bar G$. Second,
\[
\phi_i=\frac{1}{i+1}+\frac{i}{i+1}\Bigl(1-\frac{1}{\vartheta_k\gamma_i}\Bigr)\le\frac{1}{i+1}+\log(\vartheta_k\gamma_i)
\]
for $i\ge k$, so that the term of $\Sigma_\phi$ with index $i=k+m$ satisfies
\[
(m+1)\,\phi_{k+m}\le\bigl(\bar H(k,m+1)-\bar H(k,m)\bigr)+\frac{5}{2k}\bigl(\bar G(k,m+1)-\bar G(k,m)\bigr).
\]
Summing over $1\le m\le R-1$ we get $\Sigma_\phi\le\bar H+\frac{5}{2k}\,\bar G$. The two bounds add up to $\mathsf e_{N-1}+\Sigma_\phi\le\bar H+\frac5k\,\bar G$.

Finally, we prove the bounds $\underline t\le\chi_R\le\bar t$. Recall that $\chi_R=1-\vartheta_k-\sum_{m=1}^{R}\varepsilon_m\vartheta_k^m$. Since $\vartheta_k^m=e^{\sigma m}e^{-m}$ and $\sum_{m\ge 1}e^{-m}\varepsilon_m=1-e^{-1}$ by Lemma~\ref{lem:epsmass},
\[
\chi_R=\sum_{m>R}e^{-m}\varepsilon_m-e^{-1}\bigl(e^{\sigma}-1\bigr)-\sum_{m=1}^{R}\bigl(e^{\sigma m}-1\bigr)\,e^{-m}\varepsilon_m.
\]
Thus $\chi_R$ is the tail of the series minus the excess that is caused by $\sigma>0$, i.e.\ minus the sum of the two terms that are subtracted in this formula, and we bound the tail and the excess from both sides. The terms are Poisson masses: since $m\,\varepsilon_m=(m-1)^{m-1}/(m-1)!$,
\[
e^{-m}\varepsilon_m=\frac{e^{-1}}{m}\cdot\frac{e^{-(m-1)}(m-1)^{m-1}}{(m-1)!}\qquad(m\ge 1).
\]

For the tail we use that for real $y>0$ the integral of $t^{-3/2}$ over $[y,y+1]$ lies between the values of the integrand at the two ends,
\[
\frac{1}{(y+1)\sqrt{y+1}}\le 2\Bigl(\frac{1}{\sqrt y}-\frac{1}{\sqrt{y+1}}\Bigr)\le\frac{1}{y\sqrt y}.
\]
Let $m\ge R+1$, so that $m-1\ge R\ge 2$. By the expression of $e^{-m}\varepsilon_m$ through the Poisson mass at $m-1$ and the upper bound in \eqref{eq:prelim-stirling},
\[
e^{-m}\varepsilon_m\le\frac{e^{-1}c_\ast}{m\sqrt{m-1}}\le\frac{e^{-1}c_\ast}{(m-1)\sqrt{m-1}}\le 2e^{-1}c_\ast\Bigl(\frac{1}{\sqrt{m-2}}-\frac{1}{\sqrt{m-1}}\Bigr).
\]
By the same expression and the lower bound in \eqref{eq:prelim-stirling}, in which the factor $1-\frac{1}{12(m-1)}$ is at least $1-\frac{1}{12R}$,
\begin{align*}
e^{-m}\varepsilon_m&\ge\frac{e^{-1}c_\ast\bigl(1-\frac{1}{12R}\bigr)}{m\sqrt{m-1}}\ge\frac{e^{-1}c_\ast\bigl(1-\frac{1}{12R}\bigr)}{m\sqrt m}\\
&\ge 2e^{-1}c_\ast\Bigl(1-\frac{1}{12R}\Bigr)\Bigl(\frac{1}{\sqrt m}-\frac{1}{\sqrt{m+1}}\Bigr).
\end{align*}
Both bounds telescope when they are summed over $m\ge R+1$, and we obtain
\[
\frac{2e^{-1}c_\ast\bigl(1-\frac{1}{12R}\bigr)}{\sqrt{R+1}}\le\sum_{m>R}e^{-m}\varepsilon_m\le\frac{2e^{-1}c_\ast}{\sqrt{R-1}}.
\]

For the lower bound of $\chi_R$ we bound the excess from above. From $1-\sigma\le e^{-\sigma}$ we get $e^{\sigma}-1\le\sigma e^{\sigma}\le\bar a\,e^{\bar a}$. For $1\le m\le R$, the bound \eqref{eq:strip-quad} with $y=m$ and the expression of $e^{-m}\varepsilon_m$ through the Poisson mass at $m-1$ give
\[
\bigl(e^{\sigma m}-1\bigr)\,e^{-m}\varepsilon_m\le e^{-1}\Bigl(\bar a+\frac{633}{1000}\,\bar a^{2}\,m\Bigr)\frac{e^{-(m-1)}(m-1)^{m-1}}{(m-1)!}.
\]
We sum over $1\le m\le R$, write $i=m-1$, and use the bounds \eqref{eq:strip-poisson}: the sum of the Poisson masses at $i<R$ is at most $1+c_\ast(2\sqrt{R-1}-1)$, and the sum of these masses with weights $i+1$ is at most $1+c_\ast\bigl(\frac23R\sqrt R+2\sqrt{R-1}-1\bigr)$. Subtracting the resulting bounds for the excess from the lower bound for the tail, we obtain $\underline t\le\chi_R$.

For the upper bound of $\chi_R$ we bound the excess from below by its part of the first order: $e^{\sigma}-1\ge\sigma$ and $e^{\sigma m}-1\ge\sigma m$, so that, since $m\,e^{-m}\varepsilon_m$ is $e^{-1}$ times the Poisson mass at $m-1$, the excess is at least $e^{-1}\sigma\bigl(1+\sum_{i=0}^{R-1}e^{-i}i^i/i!\bigr)$. By the lower bound in \eqref{eq:prelim-stirling},
\[
\sum_{i=0}^{R-1}\frac{e^{-i}i^i}{i!}\ge 1+c_\ast\sum_{i=1}^{R-1}\Bigl(\frac{1}{\sqrt i}-\frac{1}{12\,i\sqrt i}\Bigr)\ge 1+c_\ast\Bigl(2\sqrt R-2-\frac14\Bigr),
\]
since $\sum_{i=1}^{n}i^{-1/2}\ge 2\sqrt{n+1}-2$ and $\sum_{i=1}^{n}i^{-3/2}\le 3$ for $n\ge 1$, again by comparison with integrals. As $\sigma\ge\underline a\ge 0$, the excess is at least $e^{-1}\underline a\,\bigl[2+c_\ast(2\sqrt R-2-\frac14)\bigr]$, and subtracting this from the upper bound for the tail we obtain $\chi_R\le\bar t$.

It remains to deduce the two conditions from the bounds \eqref{eq:strip-five}. Since $\underline t\le\chi_R$ and $\underline t\ge 0$ by hypothesis, we have $\chi_R\ge 0$. By the bounds for $\chi_R$ and for $\mathsf e_{N-1}+\Sigma_\phi$ and by $0<\vartheta_k\le e^{-1}e^{\bar a}$, the two factors on the left of the inequality \eqref{eq:strip-criterion} that is to be proved are nonnegative, the first is at most $e^{-1}\mathsf{A}$, and the second is at most $\mathsf{B}$. In particular $e^{-1}\mathsf{A}\,\mathsf{B}\ge 0$, so that the last hypothesis of the proposition, which reads $e^{-1}\mathsf{A}\,\mathsf{B}<e^{-1}\mathsf{C}$, implies $\mathsf{C}>0$. Since $\vartheta_k\ge e^{-1}$, $S_R\le\bar S$ and $\mathsf E\le\bar{\mathsf E}$, we obtain
\begin{align*}
\Bigl(\chi_R+\frac{\vartheta_k\,(R+1)}{N}\Bigr)\bigl(\mathsf e_{N-1}+\Sigma_\phi\bigr)
&\le e^{-1}\mathsf{A}\,\mathsf{B}<e^{-1}\mathsf{C}\\
&\le\vartheta_k\,\mathsf{C}\le\vartheta_k\Bigl(\frac{R(R+1)}{N}-1-S_R-\mathsf E\Bigr),
\end{align*}
which is the inequality to be proved. This completes the proof.
\end{proof}

\subsection{The choice of $R$}\label{subsec:strip-scale}

We shall apply Proposition~\ref{prop:closedcriterion} to a column $k\ge 265$, and the difficulty lies in the choice of $R$, because the hypotheses pull in opposite directions. The inequality \eqref{eq:strip-closed} needs $R$ large: on its right-hand side the term $R(R+1)/N$, which is of order $R^2/k$, has to exceed $\bar S$, which is of order $\sqrt R$, and this requires $R$ of order at least $k^{2/3}$. The other two hypotheses need $R$ small: $\bar aR$ grows with $R$, and in $\underline t$ the first term decreases with $R$ while the terms that are subtracted grow. For $k=265$ a numerical evaluation of the closed forms shows that the three hypotheses hold together exactly for $64\le R\le 70$.

Let $k\ge 265$ be an integer. We take for $R$ the least positive integer with $121\,k^2<25\,R^3$, which exists because $121\,k^2<25\,k^3$; it is $\lfloor(11k/5)^{2/3}\rfloor+1$ and equals $70$ for $k=265$. By the minimality of $R$,
\[
25\,(R-1)^3\le 121\,k^2<25\,R^3.
\]
To verify the hypotheses for all $k\ge 265$ at once, we work with the real number $u>0$ defined by
\[
k\,u^3=\frac{5}{11},\qquad\text{so that}\qquad 121\,k^2u^6=25.
\]
Multiplying the two inequalities for $R$ by $u^6$ and taking cube roots, we get $1<Ru^2\le 1+u^2$. Since $k\ge 265$, we have $265\,u^3\le\frac{5}{11}<265\cdot(0.119705)^3$, hence $u<0.119705<\frac{3}{25}$. In particular $R\ge 2$, because $Ru^2>1>u^2$. Since $\frac1k=\frac{11}{5}u^3$, the numbers $\bar a=5/(2k)$ and $\underline a=5/(2k)-4/(3k^2)$ are $\bar a=\frac{11}{2}u^3$ and $\underline a=\frac{11}{2}u^3-\frac{484}{75}u^6$.

Thus $R$ is about $u^{-2}$, and the square roots of $R-1$, $R$ and $R+1$, which occur in the closed forms $\bar S$, $\underline t$ and $\bar t$, are about $u^{-1}$. More precisely, the bounds for $Ru^2$ give
\[
1-u^2<(R-1)\,u^2\le 1,\qquad 1<Ru^2\le 1+u^2,\qquad(R+1)\,u^2\le 1+2u^2,
\]
and since $1+u^2\le(1+\frac{u^2}{2})^2$ and $1+2u^2\le(1+u^2)^2$, taking square roots we get $\sqrt{R-1}\,u\le 1$, $1<\sqrt R\,u\le 1+\frac{u^2}{2}$ and $\sqrt{R+1}\,u\le 1+u^2$. For $\sqrt{R-1}$ we also need a bound from below. Since $(R-1)\,u^2>1-u^2>0$, we have
\begin{align*}
(R-1)\,u^2\,\Bigl(1+\frac{51}{100}\,u^2\Bigr)^2&>(1-u^2)\Bigl(1+\frac{51}{50}\,u^2\Bigr)\\
&=1+\frac{u^2}{50}\,(1-51\,u^2)\ge 1,
\end{align*}
where the last step holds because $u<\frac{3}{25}$ gives $u^2<\frac{9}{625}<\frac{1}{51}$. Therefore $1\le\sqrt{R-1}\,u\,\bigl(1+\frac{51}{100}u^2\bigr)$. Finally, $k\ge 265$ gives $0\le\bar a\le\frac{5}{530}=\frac{1}{106}$, and $e^s\le\frac{1}{1-s}$ for $0\le s<1$; together with the enclosure $0.39894\le c_\ast\le 0.39895$ of Section~\ref{sec:prelim} this gives $1\le e^{\bar a}\le\frac{106}{105}$ and $\frac{19947}{50000}\le c_\ast\le\frac{7979}{20000}$.

For later reference we collect the relations that have been obtained:
\begin{equation}\label{eq:strip-relations}
\begin{gathered}
k\,u^3=\frac{5}{11},\qquad 0<u<0.119705<\frac{3}{25},\qquad 1<Ru^2\le 1+u^2,\\
\bar a=\frac{11}{2}\,u^3,\qquad\underline a=\frac{11}{2}\,u^3-\frac{484}{75}\,u^6,\qquad 1\le e^{\bar a}\le\frac{106}{105},\\
\sqrt{R-1}\,u\le 1\le\sqrt{R-1}\,u\,\Bigl(1+\frac{51}{100}\,u^2\Bigr),\qquad 1<\sqrt R\,u\le 1+\frac{u^2}{2},\\
\sqrt{R+1}\,u\le 1+u^2,\qquad\frac{19947}{50000}\le c_\ast\le\frac{7979}{20000}.
\end{gathered}
\end{equation}

By these relations, each closed form is a power of $u$ times a polynomial in $u$, up to errors that can be controlled. To leading order, $e\,\underline t$ is $c_\ast u\,(2-11u)$, and the numbers $\mathsf{A}$, $\mathsf{B}$ and $\mathsf{C}$ of \eqref{eq:strip-three} are
\[
\Bigl(2c_\ast+\frac{11}{5}\Bigr)u,\qquad\frac{11}{10\,u},\qquad\frac{1}{u}\Bigl(\frac{11}{5}-2c_\ast\Bigr).
\]
So $\mathsf{A}\,\mathsf{B}$ stays near $3.3$ while $\mathsf{C}$ grows like $1.40/u$, i.e.\ like $k^{1/3}$, and both $\underline t\ge 0$ and $\mathsf{A}\,\mathsf{B}<\mathsf{C}$ hold when $k$ is large. The work is therefore at the columns just above $265$, where $u$ is close to $\frac{3}{25}$ and the terms of first order in $u$ are not small.

It remains to prove, for every $k\ge 265$ and for this $R$, the two inequalities
\begin{equation}\label{eq:strip-two}
\underline t\ge 0\qquad\text{and}\qquad\mathsf{A}\,\mathsf{B}<\mathsf{C}
\end{equation}
themselves, keeping track of these terms. This is done in the next two subsections, in which $k\ge 265$ is an integer, $R$ and $u$ are the numbers just defined, and the relations \eqref{eq:strip-relations} are used throughout. We write
\[
\mathfrak{c}_-=\frac{19947}{50000},\qquad\mathfrak{c}_+=\frac{7979}{20000}
\]
for the two rational bounds of $c_\ast$ in \eqref{eq:strip-relations}, so that $0<\mathfrak{c}_-\le c_\ast\le\mathfrak{c}_+<1$.

\subsection{The inequality $\underline t\ge 0$}\label{subsec:strip-deficit}

By the definition \eqref{eq:strip-deficit} of $\underline t$,
\begin{multline}\label{eq:strip-deficitformula}
e\,\underline t=\frac{2c_\ast\bigl(1-\frac{1}{12R}\bigr)}{\sqrt{R+1}}-\bar a\,\bigl(e^{\bar a}+1+c_\ast\,(2\sqrt{R-1}-1)\bigr)\\
-\frac{633}{1000}\,\bar a^{2}\Bigl(1+c_\ast\Bigl(\frac23\,R\sqrt R+2\sqrt{R-1}-1\Bigr)\Bigr).
\end{multline}
We bound the three terms on the right-hand side by polynomials in $u$ and $c_\ast$. Note that $u^2\le\frac{9}{625}$.

For the first term, $Ru^2>1$ and $\sqrt{R+1}\,u\le 1+u^2$ give $1-\frac{1}{12R}\ge 1-\frac{u^2}{12}\ge 0$ and $\frac{1}{\sqrt{R+1}}\ge\frac{u}{1+u^2}\ge u\,(1-u^2)\ge 0$, so that
\[
\frac{2c_\ast\bigl(1-\frac{1}{12R}\bigr)}{\sqrt{R+1}}\ge 2c_\ast\,u\Bigl(1-\frac{u^2}{12}\Bigr)(1-u^2).
\]
For the second term, $\bar a=\frac{11}{2}u^3$, $e^{\bar a}\le\frac{106}{105}$ and $\sqrt{R-1}\,u\le 1$ give
\[
\bar a\,\bigl(e^{\bar a}+1+c_\ast\,(2\sqrt{R-1}-1)\bigr)\le\frac{11}{2}\,u^3\Bigl(\frac{211}{105}-c_\ast\Bigr)+11\,c_\ast\,u^2.
\]
For the third term, $\bar a^2=\frac{121}{4}u^6$, $R\sqrt R\,u^3=Ru^2\cdot\sqrt R\,u\le(1+u^2)\bigl(1+\frac{u^2}{2}\bigr)$ and $\sqrt{R-1}\,u\le 1$ give
\begin{align*}
&\bar a^{2}\Bigl(1+c_\ast\Bigl(\frac23\,R\,\sqrt R+2\sqrt{R-1}-1\Bigr)\Bigr)\\
&\qquad\le\frac{121}{4}\Bigl[(1-c_\ast)\,u^6+\frac23\,c_\ast\,u^3\,(1+u^2)\Bigl(1+\frac{u^2}{2}\Bigr)+2c_\ast\,u^5\Bigr].
\end{align*}
We subtract the second bound and $\frac{633}{1000}$ times the third from the first and expand. The right-hand side of \eqref{eq:strip-deficitformula} is at least
\begin{multline*}
u\,\Bigl[c_\ast\,\Bigl(2-11u-\frac{56593}{6000}\,u^2-\frac{687337}{12000}\,u^4+\frac{76593}{4000}\,u^5-\frac{25531}{4000}\,u^6\Bigr)\\
-\frac{2321}{210}\,u^2-\frac{76593}{4000}\,u^5\Bigr].
\end{multline*}
For $0\le u\le\frac{3}{25}$ the polynomial that multiplies $c_\ast$ is positive: since $u^2\le\frac{9}{625}$, it is at least $2-11u-10.3\,u^2\ge 0.5$. Since $c_\ast\ge\mathfrak{c}_-$, the bracket is therefore at least its value with $\mathfrak{c}_-$ in the place of $c_\ast$. This value is a polynomial in $u$ with constant term $2\mathfrak{c}_-=0.79788$ and with no term in $u^3$. Since $0<\mathfrak{c}_-<1$, its coefficients of $u$, $u^2$, $u^4$, $u^5$ and $u^6$ are negative; here the two terms with $u^5$ combine to $-\frac{76593}{4000}\,(1-\mathfrak{c}_-)\,u^5$. On $[0,\frac{3}{25}]$ the polynomial is therefore at least its value at $u=\frac{3}{25}$, and a direct evaluation shows that this value is $0.0529\ldots$. Hence $e\,\underline t\ge 0$, which is the first of the inequalities \eqref{eq:strip-two}.

\subsection{The inequality $\mathsf{A}\,\mathsf{B}<\mathsf{C}$}\label{subsec:strip-main}

By their definitions \eqref{eq:strip-three}, \eqref{eq:strip-upper} and \eqref{eq:strip-deficit}, and since $N=k+R$, the numbers $\mathsf{A}$, $\mathsf{B}$ and $\mathsf{C}$ are
\begin{align*}
\mathsf{A}&=\frac{2c_\ast}{\sqrt{R-1}}-\underline a\,\Bigl(2+c_\ast\,\Bigl(2\sqrt R-2-\frac14\Bigr)\Bigr)+\frac{e^{\bar a}\,(R+1)}{k+R},\\
\mathsf{B}&=(R-1)-k\log\frac{k+R+1}{k+2}\\
&\qquad+\frac5k\Bigl(\frac{R^2-1}{2}-(k-1)(R-1)+k(k-1)\log\frac{k+R}{k+1}\Bigr),\\
\mathsf{C}&=\frac{R\,(R+1)}{k+R}-1-\Bigl(1+c_\ast\,\Bigl(2\sqrt{R-1}-1+\frac{2\bar a}{3}\,R\sqrt R+\frac25\cdot\frac{633}{1000}\,\bar a^{2}R^{2}\sqrt R\Bigr)\Bigr)\\
&\qquad-\frac{5\,(R-2)(R-1)\,R\,(R+1)}{24\,k^{3}}.
\end{align*}
We write
\[
h=R\,u^2,\qquad w=\frac{11}{5}\,u,
\]
so that $1<h\le 1+u^2$ and $0<w\le\frac{33}{125}$. The relation $ku^3=\frac{5}{11}$, i.e.\ $ku^2w=1$, serves to eliminate $k$:
\[
(k+R)\,u^3=\frac{5}{11}+hu,\qquad\frac5k=11\,u^3,\qquad k^3u^9=\frac{125}{1331}.
\]
We define three polynomials in $u$,
\begin{align*}
\mathsf{A}^{*}(u)&=2\mathfrak{c}_+\Bigl(1+\frac{51}{100}\,u^2\Bigr)-\Bigl(2-\frac94\,\mathfrak{c}_+\Bigr)\Bigl(\frac{11}{2}\,u^2-\frac{484}{75}\,u^5\Bigr)\\
&\qquad-2\mathfrak{c}_+\Bigl(\frac{11}{2}\,u-\frac{484}{75}\,u^4\Bigr)+\frac{106}{105}\cdot\frac{11}{5}\,(1+2u^2)(1-w+w^2),\\
\mathsf{B}^{*}(u)&=\frac{11}{10}+\frac{484}{75}\,u+\frac{26}{25}\,u^2,\\
\mathsf{C}_{*}(u)&=\frac{11}{5}\,(1-w)\Bigl(1+\frac{146}{121}\,w^2\Bigr)-2u-\frac{1331}{600}\,u^2\,(1+u^2)^2\\
&\qquad-\mathfrak{c}_+\Bigl[2-u+\Bigl(\frac{11}{3}\,u+\frac{121}{10}\cdot\frac{633}{1000}\,u^2\,(1+u^2)\Bigr)(1+u^2)\Bigl(1+\frac{u^2}{2}\Bigr)\Bigr],
\end{align*}
and prove successively
\[
\begin{gathered}
0\le\mathsf{A}\le u\,\mathsf{A}^{*}(u),\qquad u\,\mathsf{B}\le\mathsf{B}^{*}(u),\\
\mathsf{C}_{*}(u)\le u\,\mathsf{C},\qquad u\,\mathsf{A}^{*}(u)\,\mathsf{B}^{*}(u)<\mathsf{C}_{*}(u).
\end{gathered}
\]
Since $\mathsf{B}^{*}(u)\ge 0$, these four statements give
\[
u\,\mathsf{A}\,\mathsf{B}\le\mathsf{A}\,\mathsf{B}^{*}(u)\le u\,\mathsf{A}^{*}(u)\,\mathsf{B}^{*}(u)<\mathsf{C}_{*}(u)\le u\,\mathsf{C},
\]
and the second of the inequalities \eqref{eq:strip-two} follows on division by $u$.

We begin with the bounds for $\mathsf{A}$, and treat its three terms one by one. From $\sqrt{R-1}\,u\le 1$ and $1\le\sqrt{R-1}\,u\,(1+\frac{51}{100}u^2)$ we get
\[
2c_\ast\,u\le\frac{2c_\ast}{\sqrt{R-1}}\le 2c_\ast\,u\,\Bigl(1+\frac{51}{100}\,u^2\Bigr).
\]
For the second term note that $\underline a=u\,\bigl(\frac{11}{2}u^2-\frac{484}{75}u^5\bigr)$ and $0\le\underline a\le\frac{11}{2}u^3$. From $\sqrt R\,u\ge 1$ we get
\[
\underline a\,\Bigl(2+c_\ast\,\Bigl(2\sqrt R-2-\frac14\Bigr)\Bigr)\ge\underline a\,\Bigl(2-\frac94\,c_\ast\Bigr)+2c_\ast\,\Bigl(\frac{11}{2}\,u^2-\frac{484}{75}\,u^5\Bigr),
\]
and from $\sqrt R\,u\le 1+\frac{u^2}{2}$ we get
\[
\underline a\,\Bigl(2+c_\ast\,\Bigl(2\sqrt R-2-\frac14\Bigr)\Bigr)\le\frac{11}{2}\,u^3\,(2+2c_\ast\,\sqrt R)\le 11\,u^3+11\,c_\ast\,u^2\Bigl(1+\frac{u^2}{2}\Bigr).
\]
For the third term we have $\frac{R+1}{k+R}=u\,(h+u^2)/(\frac{5}{11}+hu)$, and
\[
h+u^2\le\frac{11}{5}\,(1+2u^2)(1-w+w^2)\Bigl(\frac{5}{11}+hu\Bigr).
\]
Indeed, since $h\ge 1$ and $\frac{11}{5}\bigl(\frac{5}{11}+u\bigr)=1+w$, the right-hand side is at least $(1+\nobreak 2u^2)(1-w+w^2)(1+\nobreak w)=(1+2u^2)(1+w^3)$, and this is at least $1+2u^2\ge h+u^2$. Since $e^{\bar a}\le\frac{106}{105}$, it follows that
\[
0\le\frac{e^{\bar a}\,(R+1)}{k+R}\le\frac{106}{105}\cdot\frac{11}{5}\,u\,(1+2u^2)(1-w+w^2).
\]
The lower bounds of the first and the third term and the upper bound of the second give
\[
\mathsf{A}\ge u\,\Bigl[c_\ast\,\Bigl(2-11u-\frac{11}{2}\,u^3\Bigr)-11\,u^2\Bigr],
\]
and for $c_\ast\ge\mathfrak{c}_-$ and $0<u\le\frac{3}{25}$ the bracket is at least $0.39\cdot 0.67-0.16>0$. The upper bounds of the first and the third term and the lower bound of the second give an upper bound for $\mathsf{A}$ that is affine in $c_\ast$, and in which the coefficient of $c_\ast$ is
\[
u\,\Bigl[2\Bigl(1+\frac{51}{100}\,u^2\Bigr)+\frac94\Bigl(\frac{11}{2}\,u^2-\frac{484}{75}\,u^5\Bigr)-2\Bigl(\frac{11}{2}\,u-\frac{484}{75}\,u^4\Bigr)\Bigr]\ge u\,(2-11u)>0.
\]
Since $c_\ast\le\mathfrak{c}_+$, this upper bound is therefore at most its value with $\mathfrak{c}_+$ in the place of $c_\ast$, which is $u\,\mathsf{A}^{*}(u)$.

We turn to the bound for $\mathsf{B}$. As in the proof of Proposition~\ref{prop:closedcriterion}, we bound $\log(1+t)$ for $t\ge 0$ from below by its Taylor polynomial of degree $4$ and from above by that of degree $5$. Put $\tau=\frac{R-1}{k+2}$ and $\tau'=\frac{R-1}{k+1}$. Then $0\le\tau\le\tau'\le w$, the last inequality because $(k+1)\,\tau'u^2=h-u^2\le 1=ku^2w\le(k+1)\,u^2w$. Since $\frac{k+R+1}{k+2}=1+\tau$ and $R-1=(k+2)\,\tau$, the Taylor polynomial of degree $4$ gives
\[
(R-1)-k\log\frac{k+R+1}{k+2}=2\tau+k\bigl(\tau-\log(1+\tau)\bigr)\le 2\tau+k\Bigl(\frac{\tau^2}{2}-\frac{\tau^3}{3}+\frac{\tau^4}{4}\Bigr).
\]
The right-hand side does not decrease in $\tau\ge 0$, as its derivative shows. We may therefore replace $\tau$ by $w$. Multiplying by $u$ and using $ku^2w=1$, we obtain
\begin{align*}
u\,\Bigl[(R-1)-k\log\frac{k+R+1}{k+2}\Bigr]&\le 2uw+\frac{w}{2u}-\frac{w^2}{3u}+\frac{w^3}{4u}\\
&=\frac{11}{10}-\frac{121}{75}\,u+\frac{3531}{500}\,u^2.
\end{align*}
In the same way $\frac{k+R}{k+1}=1+\tau'$ and $R-1=(k+1)\,\tau'$, hence $\frac{R^2-1}{2}=(k+1)\,\tau'+\frac{(k+1)^2}{2}\,\tau'^2$. Since $k(k-1)\ge 0$, the Taylor polynomial of degree $5$ gives
\begin{multline*}
\frac{R^2-1}{2}-(k-1)(R-1)+k(k-1)\log\frac{k+R}{k+1}\\
\le 2\tau'+\frac{3k+1}{2}\,\tau'^2+k(k-1)\Bigl(\frac{\tau'^3}{3}-\frac{\tau'^4}{4}+\frac{\tau'^5}{5}\Bigr).
\end{multline*}
The right-hand side does not decrease in $\tau'\ge 0$ either, so we may replace $\tau'$ by $w$, and we may then omit the term $-k\bigl(\frac{w^3}{3}-\frac{w^4}{4}+\frac{w^5}{5}\bigr)\le 0$. Multiplying by $\frac{5u}{k}=11\,u^4$ and using $ku^2w=1$, we obtain
\begin{multline*}
\frac{5u}{k}\,\Bigl[\frac{R^2-1}{2}-(k-1)(R-1)+k(k-1)\log\frac{k+R}{k+1}\Bigr]\\
\le\frac{121}{15}\,u-\frac{1331}{100}\,u^2+\frac{74657}{1250}\,u^3+\frac{242}{5}\,u^5+\frac{1331}{50}\,u^6.
\end{multline*}
We have thus bounded $u$ times each of the two lines in the definition of $\mathsf{B}$ by a polynomial in $u$. The sum of the two polynomials is
\[
\mathsf{B}^{*}(u)-u^2\,\Bigl(\frac{911}{125}-\frac{74657}{1250}\,u-\frac{242}{5}\,u^3-\frac{1331}{50}\,u^4\Bigr),
\]
and for $0\le u\le\frac{3}{25}$ the last bracket is positive: it decreases in $u$, and at $u=\frac{3}{25}$ it is larger than $0.03$. Hence $u\,\mathsf{B}\le\mathsf{B}^{*}(u)$.

For the bound on $\mathsf{C}$ we treat the four terms of $u\,\mathsf{C}$ one by one. First, $u\,\frac{R(R+1)}{k+R}=h\,(h+u^2)/(\frac{5}{11}+hu)$, and
\[
h\,(h+u^2)\ge\lambda\,\Bigl(\frac{5}{11}+hu\Bigr),\qquad\text{where}\quad\lambda=\frac{11}{5}\,(1-w)\Bigl(1+\frac{146}{121}\,w^2\Bigr).
\]
For $h=1$ the right-hand side is $(1-w^2)\bigl(1+\frac{146}{121}w^2\bigr)=1+u^2-\frac{146}{121}w^4\le 1+u^2$, because $\frac{25}{121}w^2=u^2$. For $h>1$ the difference $h\,(h+u^2)-\lambda\,(\frac{5}{11}+hu)$ of the two sides exceeds its value at $h=1$ by $(h-1)(h+1+u^2-\lambda u)$, which is positive because $\lambda u\le\frac{11}{5}u<1$. Hence the first term of $u\,\mathsf{C}$ is at least $\lambda$. Second, the term $-1$ of $\mathsf{C}$ contributes $-u$. Third, with $\bar a=\frac{11}{2}u^3$,
\begin{multline*}
u\,\Bigl(1+c_\ast\,\Bigl(2\sqrt{R-1}-1+\frac{2\bar a}{3}\,R\,\sqrt R+\frac25\cdot\frac{633}{1000}\,\bar a^{2}R^{2}\sqrt R\Bigr)\Bigr)\\
=u+c_\ast\,\Bigl(2\sqrt{R-1}\,u-u+\frac{11}{3}\,u\,h\,(\sqrt R\,u)+\frac{121}{10}\cdot\frac{633}{1000}\,u^2h^2\,(\sqrt R\,u)\Bigr),
\end{multline*}
and by $\sqrt{R-1}\,u\le 1$, $h\le 1+u^2$, $\sqrt R\,u\le 1+\frac{u^2}{2}$ and $c_\ast\le\mathfrak{c}_+$ this is at most $u$ plus $\mathfrak{c}_+$ times the square bracket in the definition of $\mathsf{C}_{*}(u)$. Fourth,
\begin{align*}
u\cdot\frac{5\,(R-2)(R-1)\,R\,(R+1)}{24\,k^3}&=\frac{1331}{600}\,u^2\,(h-2u^2)(h-u^2)\,h\,(h+u^2)\\
&\le\frac{1331}{600}\,u^2\,(1+u^2)^2,
\end{align*}
because $(h-2u^2)(h+u^2)\le h\,(h-u^2)$ and $0\le h\,(h-u^2)\le 1+u^2$. Subtracting the last three bounds from the first, we obtain $\mathsf{C}_{*}(u)\le u\,\mathsf{C}$.

We are left with the inequality between the polynomials. Since $\mathfrak{c}_+$ is a rational number, the difference $\mathsf{C}_{*}(u)-u\,\mathsf{A}^{*}(u)\,\mathsf{B}^{*}(u)$ is a polynomial of degree $8$ in $u$ with rational coefficients. Expanding it in exact arithmetic, we find that it is
\begin{multline*}
\frac{14021}{10000}-\frac{7857223}{700000}\,u-\frac{21491467261}{12600000000}\,u^2+\frac{39724085447}{2520000000}\,u^3\\
-\frac{149002190689}{2800000000}\,u^4+\frac{72772583951}{3150000000}\,u^5-\frac{124571939481}{700000000}\,u^6\\
-\frac{14494390471}{196875000}\,u^7-\frac{10711569649}{1200000000}\,u^8,
\end{multline*}
which, with its coefficients rounded, is
\begin{multline*}
1.4021-11.22\,u-1.706\,u^2+15.76\,u^3-53.22\,u^4\\
+23.10\,u^5-178.0\,u^6-73.62\,u^7-8.926\,u^8 .
\end{multline*}
Its constant term is $\frac{11}{5}-2\mathfrak{c}_+$, its coefficients of $u^3$ and $u^5$ are positive, and its other six coefficients are negative. For $0\le u\le\frac{3}{25}$ the polynomial is therefore at least its constant term plus the values at $u=\frac{3}{25}$ of the six terms with a negative coefficient. These six values are $-1.34695\ldots$, $-0.02456\ldots$, $-0.01103\ldots$, $-0.00053\ldots$, $-0.00002\ldots$ and a number of absolute value less than $10^{-6}$, and they add up to more than $-1.38315$. Hence
\[
\mathsf{C}_{*}(u)-u\,\mathsf{A}^{*}(u)\,\mathsf{B}^{*}(u)>1.4021-1.38315>0.018\qquad\Bigl(0\le u\le\frac{3}{25}\Bigr).
\]
This proves the last of the four statements, and with it the inequality $\mathsf{A}\,\mathsf{B}<\mathsf{C}$.

\subsection{The crossing in the strip}\label{subsec:strip-row}

The \emph{strip} of a column $k\ge 2$ consists of the rows $k+r$ with $r\ge 1$ and $25\,r^3\le 121\,k^2$, i.e.\ with $1\le r\le(11k/5)^{2/3}$. We now show that a column $k\ge 265$ is positive in a row of its strip or in the first row beyond it.

\begin{proposition}\label{prop:strip}
Let $k\ge 265$ be an integer. Then there is an integer $j$ with
\[
k<j,\qquad 25\,(j-k-1)^3\le 121\,k^2,\qquad d_{j,k}>0.
\]
\end{proposition}

\begin{proof}
Let $R$ be the least positive integer with $121\,k^2<25\,R^3$ and let $u>0$ be the real number with $k\,u^3=\frac{5}{11}$, as in Section~\ref{subsec:strip-scale}, so that the relations \eqref{eq:strip-relations} hold. We apply Proposition~\ref{prop:closedcriterion} to $k$ and $R$.

First, $2\le R\le k-1$. We have seen that $R\ge 2$, and $R+1<k$ because $Ru^2\le 1+u^2$ and $u<\frac{3}{25}$ give
\[
(R+1)\,u^3=u\cdot Ru^2+u^3\le u+2u^3<\frac{5}{11}=k\,u^3.
\]
Second, $\bar a\,R\le\frac{6679}{10000}$, because $\bar a=\frac{11}{2}u^3$, $Ru^2\le 1+u^2$ and $u<0.119705$ give
\begin{align*}
\bar a\,R=\frac{11}{2}\,u\cdot Ru^2&\le\frac{11}{2}\,u\,(1+u^2)\\
&<\frac{11}{2}\cdot 0.119705\cdot\bigl(1+(0.119705)^2\bigr)<0.66782<\frac{6679}{10000}.
\end{align*}

Third, $\underline t\ge 0$, and fourth, the inequality \eqref{eq:strip-closed}, which is the inequality $\mathsf{A}\,\mathsf{B}<\mathsf{C}$ multiplied by $e^{-1}$: these are the two inequalities \eqref{eq:strip-two}, which were proved in Sections~\ref{subsec:strip-deficit} and~\ref{subsec:strip-main}.

Proposition~\ref{prop:closedcriterion} now gives an integer $j$ with $k<j\le k+R$ and $d_{j,k}>0$, and $25\,(j-k-1)^3\le 25\,(R-1)^3\le 121\,k^2$ by the minimality of $R$.
\end{proof}

Consequently, if $k\ge 265$, $n>k$ and column $k$ is nonpositive in all rows $k+1,\dots,n$, then row $n$ lies in the strip of the column: the integer $j$ of the proposition is larger than $n$, so that $n-k\le j-k-1$.

\begin{remark}
The bound of Proposition~\ref{prop:strip} is not sharp. For $k=265$ a computation with the last-car recursion shows that the crossing row is row $k+36$, while the proposition gives a positive row $j\le k+70$. The criterion behind the proposition replaces the tilted counts of level $k-1$ in the rows before the crossing by the largest of them, and it is not designed to find the exact row.
\end{remark}

Together with the computed signs of Lemma~\ref{lem:signpairs} for the columns below $265$, the proposition gives a bound for the crossing row of every column $k\ge 82$.

\begin{corollary}\label{cor:earlypositive}
For every integer $k\ge 82$ there is an integer $r$ with
\[
1\le r,\qquad 100\,r\le 27\,k,\qquad d_{k+r,k}>0.
\]
\end{corollary}

\begin{proof}
Let first $82\le k\le 264$, and let $\zeta=\zeta(k)$ be the number of Definition~\ref{def:zeta}. By Lemma~\ref{lem:signpairs} we have $b_{k+\zeta+1,k-1}<b_{k+\zeta+1,k}$, i.e.\ $d_{k+r,k}>0$ for $r=\zeta+1\ge 1$. It remains to check that $100\,(\zeta(k)+1)\le 27\,k$. The table of Definition~\ref{def:zeta} shows that $\zeta$ does not decrease and that $\zeta(92)=20$, $\zeta(133)=24$ and $\zeta(264)=35$. Hence $100\,(\zeta(k)+1)$ is at most $2100$ for $82\le k\le 92$, at most $2500$ for $93\le k\le 133$ and at most $3600$ for $134\le k\le 264$, and the claim follows from
\[
2100\le 27\cdot 82,\qquad 2500\le 27\cdot 93,\qquad 3600\le 27\cdot 134.
\]

Let now $k\ge 265$. Proposition~\ref{prop:strip} gives an integer $j>k$ with $d_{j,k}>0$ and $25\,s^3\le 121\,k^2$, where $s=j-k-1\ge 0$. We take $r=j-k=s+1$. If $s$ were larger than $0.266\,k$, we would have $25\,s^3>25\cdot(0.266)^3\cdot 265\,k^2>121\,k^2$, because $k\ge 265$, $(0.266)^3\cdot 265>4.98$ and $\frac{121}{25}=4.84$. Hence $s\le 0.266\,k$, and $100\,r\le 26.6\,k+100\le 27\,k$, because $0.4\,k\ge 106$.
\end{proof}

\section{Persistence on the near rows}\label{sec:near}
We turn to persistence: a column that has turned positive in its chain range stays positive there (Proposition~\ref{prop:chainrange}). On the near rows this is Proposition~\ref{prop:nearrows}: if a column $k\ge 42$ is nonpositive in rows $k,\dots,q-1$ and positive in rows $q,\dots,n-1$, where $k<q<n\le k+0.47k$, then it is positive in row $n$ as well. The rows beyond $k+0.47k$ need a different argument, which is given in Sections~\ref{sec:renewal} and~\ref{sec:far}.

\subsection{A lemma on two rows}\label{subsec:near-twopoint}

The argument is contained in the following lemma about real sequences; in the proof of Proposition~\ref{prop:nearrows} it is applied to the column equation at the diagonal tilt. The lemma concerns a sequence $\eta$ that satisfies, at a row $j$, an equation of the form
\[
\eta_j=\sigma_j\,U_{j-1}+\psi_j\,\eta_{j-1}+\sum_{m=1}^{j-k}p_m\,\eta_{j-m},
\]
which is the form of the column equation (cf.\ Proposition~\ref{prop:columnequation}). In the first term $U$ is a nonnegative sequence, and this term acts as a source at row $j$. The sum, a convolution of $\eta$ with the kernel $p$, feeds the values of $\eta$ in rows $k,\dots,j-1$ back into row $j$, the value in row $j-m$ with weight $p_m>0$. If $\eta$ is nonpositive in rows $k,\dots,q-1$ and positive in rows $q,\dots,n-1$, the feedback at row $n$ still contains the rows before $q$, with values that are nonpositive and about whose size, row by row, nothing is known, so that the positivity of the source does not decide the sign of $\eta_n$. The lemma uses the equation at rows $q=k+r$ and $n=q+t$ only. For a log-convex kernel the weight of every row before $q$ changes between these two rows by at most the factor $p_{r+t}/p_r$ by which the weight of row $k$ changes, and so the feedback of the rows before $q$ at row $n$ is controlled by their feedback at row $q$, where the source outweighed it.

\begin{lemma}\label{lem:tworows}
Let $k\ge 0$, $r\ge 1$ and $t\ge 1$ be integers, and put $q=k+r$ and $n=q+t$. Let $\eta=(\eta_j)_{j\ge 0}$, $U=(U_j)_{j\ge 0}$ and $p=(p_m)_{m\ge 0}$ be sequences of real numbers, and let $\sigma_q$, $\sigma_n$, $\psi_q$ and $\psi_n$ be real numbers with $\sigma_q\ge 0$, $\psi_q\ge 0$ and $\psi_n>0$. Suppose that the following hold.
\begin{enumerate}
\item (equation at the two rows) We have
\begin{align*}
\eta_q&=\sigma_q\,U_{q-1}+\psi_q\,\eta_{q-1}+\sum_{m=1}^{r}p_m\,\eta_{q-m},\\
\eta_n&=\sigma_n\,U_{n-1}+\psi_n\,\eta_{n-1}+\sum_{m=1}^{r+t}p_m\,\eta_{n-m}.
\end{align*}
\item (initial value) $\eta_k=-U_k$.
\item (diagonal bound) $U_{j-1}-U_j\le\eta_j$ for every $j$ with $k+1\le j\le n-1$.
\item (nonnegativity) $U_j\ge 0$ for every $j\ge 0$.
\item (kernel) $p_m>0$ for $1\le m\le r+t$, and $p_{m+1}\le p_m$ for $1\le m\le r+t-1$.
\item (cross inequality) $p_{m+t}\,p_r\le p_{r+t}\,p_m$ for $1\le m\le r$.
\item (two-point condition) $p_{r+t}\,\sigma_q\le p_r\,\sigma_n$.
\item (signs) $\eta_j\le 0$ for $k\le j\le q-1$, and $\eta_j>0$ for $q\le j\le n-1$.
\end{enumerate}
Then $\eta_n>0$.
\end{lemma}

\begin{proof}
Suppose, for a contradiction, that $\eta_n\le 0$. Let
\[
E=\sum_{j=q}^{n-1}\eta_j\qquad\text{and}\qquad H=\sum_{j=k}^{q-1}\eta_j ,
\]
so that $E>0$ and $H\le 0$ by~(8). Let further
\[
L=\sum_{m=1}^{t}p_m\,\eta_{n-m},\qquad B_n=\sum_{m=1}^{r}p_{m+t}\,\eta_{q-m},\qquad B_q=\sum_{m=1}^{r}p_m\,\eta_{q-m}.
\]
Thus $L$ is the feedback at row $n$ of rows $q,\dots,n-1$, and $B_n$ (resp.\ $B_q$) is the feedback of rows $k,\dots,q-1$ at row $n$ (resp.\ at row $q$). Splitting the sum in the second equation of~(1) at $m=t$, we may write the two equations as
\[
\eta_n=\sigma_n\,U_{n-1}+\psi_n\,\eta_{n-1}+L+B_n,\qquad \eta_q=\sigma_q\,U_{q-1}+\psi_q\,\eta_{q-1}+B_q.
\]

We need three comparisons of these sums:
\[
L\ge p_t\,E,\qquad B_n\ge p_{t+1}\,H,\qquad p_{r+t}\,B_q\le p_r\,B_n.
\]
The first two hold because the kernel does not increase, by~(5): we have $p_m\ge p_t$ and $\eta_{n-m}>0$ for $1\le m\le t$, and $p_{m+t}\le p_{t+1}$ and $\eta_{q-m}\le 0$ for $1\le m\le r$. The third is the cross inequality~(6), multiplied by $\eta_{q-m}\le 0$ and summed over $1\le m\le r$.

We also need two consequences of the diagonal bound~(3):
\[
U_{q-1}-U_{n-1}\le E\qquad\text{and}\qquad -U_{q-1}\le H.
\]
The first is the sum of~(3) over $q\le j\le n-1$, which telescopes. In the same way the sum of~(3) over $k+1\le j\le q-1$ is $U_k-U_{q-1}\le H-\eta_k$, and since $\eta_k=-U_k$ by~(2), this is the second. Adding the two inequalities we get
\[
-(E+H)\le U_{n-1}.
\]

We first bound $\sigma_n$ from above. Since $\eta_n\le 0$ and $\psi_n\,\eta_{n-1}>0$, the equation at row $n$, in the form $\eta_n=\sigma_n\,U_{n-1}+\psi_n\,\eta_{n-1}+L+B_n$, gives
\[
\sigma_n\,U_{n-1}+L<-B_n.
\]
Hence
\begin{align*}
\sigma_n\,U_{n-1}&<-B_n-L\le -p_{t+1}\,H-p_t\,E\\
&\le -p_{t+1}\,(E+H)\le p_{t+1}\,U_{n-1},
\end{align*}
by the comparisons $L\ge p_t\,E$ and $B_n\ge p_{t+1}\,H$, by $p_{t+1}\le p_t$, and because $-(E+H)\le U_{n-1}$. This excludes $U_{n-1}=0$, so that $U_{n-1}>0$ by~(4), and
\[
\sigma_n<p_{t+1}.
\]

A lower bound for $\sigma_n$ comes from row $q$. There we have $\eta_q>0$, $\eta_{q-1}\le 0$ and $\psi_q\ge 0$, so the equation at row $q$, in the form $\eta_q=\sigma_q\,U_{q-1}+\psi_q\,\eta_{q-1}+B_q$, gives $-B_q<\sigma_q\,U_{q-1}$. We transfer this inequality to row $n$. By the comparison $p_{r+t}\,B_q\le p_r\,B_n$ and the inequality just obtained,
\[
-p_r\,B_n\le -p_{r+t}\,B_q<p_{r+t}\,\sigma_q\,U_{q-1},
\]
and by the bound $U_{q-1}-U_{n-1}\le E$ and the two-point condition~(7),
\[
p_{r+t}\,\sigma_q\,U_{q-1}\le p_{r+t}\,\sigma_q\,(U_{n-1}+E)\le p_r\,\sigma_n\,(U_{n-1}+E).
\] Dividing by $p_r$, we obtain $-B_n<\sigma_n\,(U_{n-1}+E)$. Since $\sigma_n\,U_{n-1}+L<-B_n$, as we saw in the first step, this gives $L<\sigma_n\,E$, and with $L\ge p_t\,E$ it gives $p_t\,E<\sigma_n\,E$. Since $E>0$, we conclude that
\[
p_t<\sigma_n.
\]

The two bounds for $\sigma_n$ and~(5) give $p_t<\sigma_n<p_{t+1}\le p_t$, which is impossible. Hence $\eta_n>0$.
\end{proof}

\subsection{The free kernel at the diagonal tilt}\label{subsec:near-kernel}

Beyond its log-convexity, the application of Lemma~\ref{lem:tworows} requires of the kernel $p$ two monotonicity properties, which the free kernel at the diagonal tilt has on an initial range of indices. Recall from Section~\ref{sec:prelim} that the free kernel $\varepsilon$ satisfies $\varepsilon_{m+1}=g_m\,\varepsilon_m$ for $m\ge 1$, where $g_1=\frac12$ and $g_m=\frac{m}{m+1}\bigl(1+\frac{1}{m-1}\bigr)^{m-1}$ for $m\ge 2$, and from Section~\ref{sec:crossing} that $\tilde\varepsilon_m=\varepsilon_m\,\vartheta_k^{m}$, where $\vartheta_k=1/\gamma_k$ is the diagonal tilt of column $k\ge 2$ and $\gamma_k=(1+1/k)^{k-2}$.

The two properties are that $\tilde\varepsilon_{m+1}\le\tilde\varepsilon_m$ as long as $m+1\le\frac35k$, and, for $k\ge 42$, that $(k+\rho)\,\tilde\varepsilon_\rho$ does not increase in $\rho$ on the shorter range $1\le\rho\le 0.47k$; the second gives the two-point condition of Lemma~\ref{lem:tworows}. We state the first for $\varepsilon$ itself, as a bound for its growth in one step by $\gamma_k$; multiplied by $\vartheta_k^{m+1}$, it takes the form just given.

\begin{lemma}\label{lem:epsgrowth}
Let $k\ge 2$ and $m\ge 1$ be integers with $5\,(m+1)\le 3\,k$. Then
\[
\varepsilon_{m+1}\le\gamma_k\,\varepsilon_m .
\]
\end{lemma}

\begin{proof}
Since $\varepsilon_{m+1}=g_m\,\varepsilon_m$ and $\varepsilon_m\ge 0$, it suffices to show that $g_m\le\gamma_k$. For $m=1$ this holds because $g_1=\frac12$ and $\gamma_k=(1+1/k)^{k-2}\ge 1$. For $m\ge 2$ we compare the two numbers through the bounds
\begin{align}
g_m&\le\exp\Bigl(1-\frac{3}{2\,(m+1)}\Bigr)&&(m\ge 2),\label{eq:near-ratio}\\
\gamma_k&\ge\exp\Bigl(1-\frac{5}{2k}\Bigr)&&(k\ge 2).\label{eq:near-gamma}
\end{align}
The second of these is the upper bound $\vartheta_k\le e^{-1}e^{5/(2k)}$ of Lemma~\ref{lem:tiltbounds}, because $\gamma_k=1/\vartheta_k$. To prove the first, let $m\ge 2$ and put $i=m-1$, so that $i\ge 1$ and
\[
\log g_m=i\,\log\Bigl(1+\frac1i\Bigr)-\log\Bigl(1+\frac{1}{i+1}\Bigr).
\]
By the elementary inequalities $u-\frac{u^2}{2}\le\log(1+u)\le u-\frac{u^2}{2}+\frac{u^3}{3}$ for $u\ge 0$, of which we use the upper one at $u=1/i$ and the lower one at $u=1/(i+1)$,
\[
\log g_m\le 1-\frac{1}{2i}+\frac{1}{3i^2}-\frac{1}{i+1}+\frac{1}{2\,(i+1)^2}\le 1-\frac{3}{2\,(i+2)} ,
\]
where the second inequality holds for every $i\ge 1$: a direct computation shows that it is equivalent to $7i^3+4i^2\ge 4i+4$. Since $i+2=m+1$, this is the bound for $g_m$.

Finally, the hypothesis $5\,(m+1)\le 3\,k$ says that $\frac{5}{2k}\le\frac{3}{2\,(m+1)}$, so that for $m\ge 2$ the two bounds give $g_m\le\gamma_k$.
\end{proof}

\begin{lemma}\label{lem:tworowscondition}
Let $k\ge 42$, $r$ and $s$ be integers with $1\le r\le s$ and $100\,s\le 47\,k$. Then the free kernel at the diagonal tilt of column $k$ satisfies
\[
(k+s)\,\tilde\varepsilon_s\le(k+r)\,\tilde\varepsilon_r .
\]
\end{lemma}

\begin{proof}
It suffices to prove that, for every integer $\rho\ge 1$ with $100\,(\rho+1)\le 47\,k$,
\[
(k+\rho+1)\,\tilde\varepsilon_{\rho+1}\le(k+\rho)\,\tilde\varepsilon_\rho.
\]
Indeed, the lemma follows by applying this inequality for $\rho=r,r+1,\dots,s-1$.

Let $\rho$ be such an integer. Since $\varepsilon_{\rho+1}=g_\rho\,\varepsilon_\rho$ and $\vartheta_k\gamma_k=1$, the inequality to be proved is
\[
(k+\rho+1)\,g_\rho\le(k+\rho)\,\gamma_k,
\]
multiplied by $\varepsilon_\rho\,\vartheta_k^{\rho+1}\ge 0$. For $\rho=1$ the displayed inequality holds because $g_1=\frac12$ and $\gamma_k\ge 1$.

Let $\rho\ge 2$. We first show that
\[
\frac{1}{k+\rho}\le\frac{3}{2\,(\rho+1)}-\frac{5}{2k}.
\]
A direct computation shows that this inequality is equivalent to $5\rho^2+(4k+5)\,\rho+7k\le 3k^2$. The left-hand side of the latter increases with $\rho$, and $\rho\le 0.47k-1$ by the hypothesis on $\rho$. At $\rho=0.47k-1$ it equals $2.9845\,k^2+0.65\,k$, which is at most $3k^2$ because $0.0155\,k\ge 0.651$ for $k\ge 42$.

By the inequality $1+y\le e^y$ for real $y$, by the inequality just proved, and by the bounds \eqref{eq:near-ratio} for $g_\rho$ and \eqref{eq:near-gamma} for $\gamma_k$ from the proof of Lemma~\ref{lem:epsgrowth}, we obtain
\begin{align*}
\Bigl(1+\frac{1}{k+\rho}\Bigr)\,g_\rho&\le\exp\Bigl(\frac{1}{k+\rho}\Bigr)\,g_\rho
\le\exp\Bigl(\frac{3}{2\,(\rho+1)}-\frac{5}{2k}\Bigr)\exp\Bigl(1-\frac{3}{2\,(\rho+1)}\Bigr)\\
&=\exp\Bigl(1-\frac{5}{2k}\Bigr)\le\gamma_k .
\end{align*}
Multiplied by $k+\rho$, this is the inequality $(k+\rho+1)\,g_\rho\le(k+\rho)\,\gamma_k$ that remained to be proved.
\end{proof}

\subsection{Rows up to $k+0.47k$}\label{subsec:near-rows}

We now apply Lemma~\ref{lem:tworows} to a column at its diagonal tilt. In the following proposition $q$ is the crossing row of column $k$; the case $t=1$ is included, so that the proposition applies from the first row after the crossing row on.

\begin{proposition}\label{prop:nearrows}
Let $k\ge 42$, $r\ge 1$ and $t\ge 1$ be integers with $100\,(r+t)\le 47\,k$, and put $q=k+r$ and $n=q+t$. Suppose that
\begin{enumerate}
\item $d_{j,k}\le 0$ for every $j$ with $k\le j<q$, and
\item $d_{j,k}>0$ for every $j$ with $q\le j<n$.
\end{enumerate}
Then $d_{n,k}>0$.
\end{proposition}

\begin{proof}
Let $\delta_j=d_{j,k}/j!$ be the normalised step of column $k$, and let $\vartheta_k=1/\gamma_k$ be its diagonal tilt. We apply Lemma~\ref{lem:tworows}, with the given integers $k$, $r$ and $t$, to the sequences
\[
\eta_j=\tilde\delta_j=\delta_j\,\vartheta_k^{j},\qquad U_j=\tilde\beta^{(k-1)}_j=\beta^{(k-1)}_j\,\vartheta_k^{j},\qquad p_m=\tilde\varepsilon_m=\varepsilon_m\,\vartheta_k^{m}
\]
and to the numbers
\[
\sigma_q=\frac{2\vartheta_k}{q},\qquad\sigma_n=\frac{2\vartheta_k}{n},\qquad\psi_q=\frac{(k+1)\,\vartheta_k}{q},\qquad\psi_n=\frac{(k+1)\,\vartheta_k}{n},
\]
all of which are positive. Each $\eta_j$ has the sign of $d_{j,k}$, so that hypothesis~(8) of the lemma consists of the two assumptions of the proposition, and its conclusion $\eta_n>0$ is the claim. It remains to verify hypotheses (1) to~(7) of the lemma.

For~(1), note that rows $q$ and $n$ lie in the chain range $k+1\le j\le 2k-1$ of column $k$, because $r+t\le 0.47k<k$. For a row $j$ of the chain range, the equation of Proposition~\ref{prop:columnequation}, multiplied by $\vartheta_k^{j}$, reads
\[
\tilde\delta_j=\frac{2\vartheta_k}{j}\,\tilde\beta^{(k-1)}_{j-1}+\frac{(k+1)\,\vartheta_k}{j}\,\tilde\delta_{j-1}+\sum_{m=1}^{j-k}\tilde\varepsilon_m\,\tilde\delta_{j-m},
\]
and for $j=q$ and $j=n$ these are the two equations of~(1).

Hypothesis~(2) holds because $\beta^{(k)}_k=0$ (Section~\ref{sec:recursion}), so that $\delta_k=-\beta^{(k-1)}_k$. For~(3), Proposition~\ref{prop:diagonallower} gives $\gamma_k\,\beta^{(k-1)}_{j-1}\le\beta^{(k)}_j$ for every $j\ge 1$. Multiplying by $\vartheta_k^{j}$, which removes the factor $\gamma_k$ since $\vartheta_k\gamma_k=1$, and subtracting $U_j$ on both sides, we get
\[
U_{j-1}-U_j\le\bigl(\beta^{(k)}_j-\beta^{(k-1)}_j\bigr)\,\vartheta_k^{j}=\eta_j.
\]
Hypothesis~(4) holds because $\beta^{(k-1)}_j=b_{j,k-1}/j!\ge 0$.

For~(5), the $p_m$ with $m\ge 1$ are positive because the $\varepsilon_m$ are. For $1\le m\le r+t-1$ we have $5\,(m+1)\le 5\,(r+t)\le 3k$, because $100\,(r+t)\le 47\,k$, so that Lemma~\ref{lem:epsgrowth} gives $\varepsilon_{m+1}\le\gamma_k\,\varepsilon_m$; multiplied by $\vartheta_k^{m+1}$, this is $p_{m+1}\le p_m$.

For~(6), the sequence $(\varepsilon_m)_{m\ge 1}$ is log-convex (Lemma~\ref{lem:freelogconvex}), so that the ratios $\varepsilon_{m+1}/\varepsilon_m$ do not decrease in $m$. Hence $\varepsilon_{m+t}/\varepsilon_m\le\varepsilon_{r+t}/\varepsilon_r$ for $1\le m\le r$, since each side is a product of $t$ consecutive ratios, and multiplying by $\varepsilon_m\,\varepsilon_r\,\vartheta_k^{m+t+r}$ we obtain $p_{m+t}\,p_r\le p_{r+t}\,p_m$.

Finally, the two-point condition~(7), multiplied by $qn/(2\vartheta_k)$, becomes $(k+r+t)\,\tilde\varepsilon_{r+t}\le(k+r)\,\tilde\varepsilon_r$. This is Lemma~\ref{lem:tworowscondition} with $s=r+t$, whose hypotheses $k\ge 42$ and $100\,s\le 47\,k$ are those of the proposition. Thus Lemma~\ref{lem:tworows} applies, and it gives $\eta_n>0$, i.e.\ $d_{n,k}>0$.
\end{proof}

\section{A comparison with renewal sequences}\label{sec:renewal}
On the far rows of a column, the lower bound for the step with which Section~\ref{sec:far} works contains terms whose coefficients may be negative. They are bounded there by a comparison which we prove here for arbitrary sequences: for a sequence $v$ about which we know only that $\sum_{i=1}^{s}f_i\,v_{s-i}\le v_s$ for $s\le\rho$, where $f$ is a log-convex kernel, a weighted sum of terms $f_mv_{\rho-m}$ is at most a multiple of $v_\rho$ that does not depend on $v$ (Proposition~\ref{prop:extremecolumn}). We then make this multiple explicit for the kernel to which the comparison is applied.

\subsection{Renewal sequences of log-convex kernels}\label{subsec:renewal-general}

The sequence with which $v$ is compared is the renewal sequence of the kernel.

\begin{definition}\label{def:renewal}
Let $f=(f_i)_{i\ge 0}$ and $u=(u_n)_{n\ge 0}$ be sequences of real numbers. We say that $u$ is \emph{the renewal sequence of the kernel $f$} if
\[
u_0=1\qquad\text{and}\qquad u_n=\sum_{i=1}^{n}f_i\,u_{n-i}\quad\text{for every } n\ge 1.
\]
\end{definition}

The second condition is the \emph{renewal equation} at index $n$. The term $f_0$ does not occur in it, and the two conditions determine $u$ from $f$ term by term: $u_1=f_1$, $u_2=f_1^2+f_2$, and so on. For example, $(c_n)$ is the renewal sequence of the free kernel $(\varepsilon_m)$, by Corollary~\ref{cor:crenewal} and since $c_0=1$. We shall use repeatedly that the renewal sequence $u$ of a kernel $f$ with positive terms is positive:
\begin{equation}\label{eq:renewal-positive}
\text{if } f_i>0 \text{ for every } i\ge 1,\quad\text{then } u_n>0 \text{ for every } n\ge 0,
\end{equation}
as the renewal equation shows by induction on $n$.

For a kernel $f$ and a sequence $v=(v_s)_{s\ge 0}$ of real numbers, the numbers
\[
\iota_s=v_s-\sum_{i=1}^{s}f_i\,v_{s-i}\qquad(s\ge 0)
\]
are called the \emph{innovations} of $v$ with respect to $f$. Thus $\iota_s$ is what $v_s$ contains beyond the feedback of the earlier values through the kernel, and $\iota_0=v_0$. For an integer $\rho\ge 0$ we say that $v$ \emph{has nonnegative innovations with respect to $f$ up to $\rho$} if $\iota_s\ge 0$ for every $s\le\rho$, i.e.\ if
\[
\sum_{i=1}^{s}f_i\,v_{s-i}\le v_s\qquad\text{for every integer } s \text{ with } 0\le s\le\rho .
\]
In these terms the renewal sequence of $f$ is the sequence whose innovations are $\iota_0=1$ and $\iota_s=0$ for $s\ge 1$.

\begin{lemma}\label{lem:convolutionform}
Let $f$, $G$, $v$ and $\iota$ be sequences of real numbers, indexed by the integers $s\ge 0$. Suppose that $G$ is the renewal sequence of the kernel $f$ (Definition~\ref{def:renewal}), and that
\[
v_s=\iota_s+\sum_{i=1}^{s}f_i\,v_{s-i}\qquad\text{for every } s\ge 0 .
\]
Then
\[
v_s=\sum_{t=0}^{s}G_{s-t}\,\iota_t\qquad\text{for every } s\ge 0 .
\]
\end{lemma}

\begin{proof}
We argue by induction on $s$ and assume the formula for all indices smaller than $s$. Inserting it for $v_{s-i}$ with $1\le i\le s$ and interchanging the two summations, we obtain
\[
\sum_{i=1}^{s}f_i\,v_{s-i}=\sum_{t=0}^{s-1}\Bigl(\sum_{i=1}^{s-t}f_i\,G_{s-t-i}\Bigr)\iota_t=\sum_{t=0}^{s-1}G_{s-t}\,\iota_t,
\]
where the inner sum equals $G_{s-t}$ by the renewal equation at index $s-t\ge 1$. Adding $\iota_s=G_0\,\iota_s$ gives $v_s=\sum_{t=0}^{s}G_{s-t}\,\iota_t$.
\end{proof}

We use the log-convexity of a sequence of positive numbers in cross-multiplied form, and we need it also for a sequence of which only an initial segment is known to be log-convex. Let $(q_n)_{n\ge 0}$ be a sequence of positive real numbers, and let $N\ge 0$ be an integer such that $q_{n+1}^2\le q_n\,q_{n+2}$ for every $n<N$. Then
\begin{equation}\label{eq:renewal-cross}
q_{i+1}\,q_j\le q_i\,q_{j+1}\qquad(0\le i\le j\le N).
\end{equation}
Indeed, the hypothesis says that $q_{n+1}/q_n$ does not decrease for $0\le n\le N$, and the inequality compares these ratios at $n=i$ and $n=j$. If $q_{n+1}^2\le q_n\,q_{n+2}$ holds for every $n\ge 0$, then moreover
\begin{equation}\label{eq:renewal-shift}
q_{i+m}\,q_j\le q_i\,q_{j+m}\qquad(0\le i\le j,\ m\ge 0),
\end{equation}
because $q_{i+m}/q_i$ is the product of the $m$ ratios $q_{i+t+1}/q_{i+t}$ with $0\le t<m$, and each of them is at most the corresponding ratio $q_{j+t+1}/q_{j+t}$.

The comparison rests on the classical theorem that log-convexity passes from a kernel to its renewal sequence. For a kernel of sum $1$ this is Theorem~1 of de~Bruijn and Erd\H{o}s \cite[pp.~161--162]{debruijn1953}, who state it with strict inequalities and remark that it holds with inequalities that are not strict as well. We need it for kernels whose sum is not $1$, such as the kernel of Section~\ref{subsec:renewal-frozen}. The proof in \cite{debruijn1953} does not use that the sum is $1$, and since the theorem is not stated without this condition there, we repeat the proof in our notation.

\begin{proposition}\label{prop:logconvexrenewal}
Let $f=(f_i)_{i\ge 0}$ and $u=(u_n)_{n\ge 0}$ be sequences of real numbers such that $u$ is the renewal sequence of the kernel $f$ (Definition~\ref{def:renewal}). Suppose that $f_i>0$ for every $i\ge 1$ and that $f_{j+2}^2\le f_{j+1}\,f_{j+3}$ for every $j\ge 0$, i.e.\ that the sequence $f_1,f_2,f_3,\dots$ is log-convex. Then
\[
u_{n+1}^2\le u_n\,u_{n+2}\qquad\text{for every } n\ge 0 .
\]
\end{proposition}

\begin{proof}
By \eqref{eq:renewal-positive} all $u_n$ are positive. We prove the inequality by induction on $n$.

For $n=0$ the claim is $f_1^2\le f_1^2+f_2$, since $u_1=f_1$ and $u_2=f_1^2+f_2$.

Let $n\ge 1$, and suppose that $u_{j+1}^2\le u_j\,u_{j+2}$ for every $j<n$. We apply the cross-multiplied form \eqref{eq:renewal-cross} of log-convexity twice. For the sequence $q_j=f_{j+1}$, which is positive and log-convex by the hypothesis on the kernel, it gives
\[
f_{i+1}\,f_{n+1}\le f_i\,f_{n+2}\qquad(1\le i\le n).
\]
For $u$ and $N=n$, where the hypothesis on the sequence is the induction hypothesis, it gives
\[
u_{n-i+1}\,u_n\le u_{n-i}\,u_{n+1}\qquad(1\le i\le n).
\]
Hence both factors of every term of the sum
\[
S=\sum_{i=1}^{n}\bigl(f_i\,f_{n+2}-f_{i+1}\,f_{n+1}\bigr)\bigl(u_{n-i}\,u_{n+1}-u_{n-i+1}\,u_n\bigr)
\]
are nonnegative, and $S\ge 0$.

On the other hand, when the product is multiplied out, four sums over $i$ arise, and the renewal equation at indices $n$, $n+1$ and $n+2$ evaluates them:
\begin{align*}
\sum_{i=1}^{n}f_i\,u_{n-i}&=u_n, & \sum_{i=1}^{n}f_i\,u_{n-i+1}&=u_{n+1}-f_{n+1},\\
\sum_{i=1}^{n}f_{i+1}\,u_{n-i}&=u_{n+1}-f_1u_n, & \sum_{i=1}^{n}f_{i+1}\,u_{n-i+1}&=u_{n+2}-f_1u_{n+1}-f_{n+2}.
\end{align*}
Inserting these values, we find that
\[
S=f_{n+1}\bigl(u_n\,u_{n+2}-u_{n+1}^2\bigr),
\]
since all other terms cancel in pairs. As $S\ge 0$ and $f_{n+1}>0$, we conclude that $u_{n+1}^2\le u_n\,u_{n+2}$.
\end{proof}

We now come to the main result of this subsection. Let $f$ be a kernel as in Proposition~\ref{prop:logconvexrenewal}, let $G$ be its renewal sequence, and let $v$ have nonnegative innovations with respect to $f$ up to $\rho$. For an integer $M\ge 1$ and weights $\theta_m$, the sum $\sum_{m=M}^{\rho}\theta_mf_mv_{\rho-m}$ is the part with $m\ge M$ of the feedback $\sum_{i=1}^{\rho}f_iv_{\rho-i}$ into $v_\rho$, taken with these weights. The following proposition says that, if the weights are nonnegative and do not decrease, this part is largest in proportion to $v_\rho$ when $v=G$, that is, for the sequence whose only innovation occurs at index $0$.

\begin{proposition}\label{prop:extremecolumn}
Let $f$, $G$, $v$ and $\theta$ be sequences of real numbers, indexed by the integers $n\ge 0$, and let $\rho\ge 0$ and $M\ge 1$ be integers. Assume that
\begin{enumerate}
\item $f_i>0$ and $f_{i+1}^2\le f_i\,f_{i+2}$ for every $i\ge 1$;
\item $G$ is the renewal sequence of the kernel $f$ (Definition~\ref{def:renewal});
\item $\theta_M\ge 0$, and $\theta_m\le\theta_{m'}$ whenever $M\le m\le m'\le\rho$;
\item $\sum_{i=1}^{s}f_i\,v_{s-i}\le v_s$ for every integer $s$ with $0\le s\le\rho$.
\end{enumerate}
Then
\[
\Bigl(\sum_{m=M}^{\rho}\theta_m\,f_m\,v_{\rho-m}\Bigr)\,G_\rho\le\Bigl(\sum_{m=M}^{\rho}\theta_m\,f_m\,G_{\rho-m}\Bigr)\,v_\rho .
\]
\end{proposition}

\begin{proof}
By~(1) and~(2), the sequence $G$ is positive, by \eqref{eq:renewal-positive}, and it satisfies $G_{n+1}^2\le G_n\,G_{n+2}$ for every $n\ge 0$, by Proposition~\ref{prop:logconvexrenewal}. Hence the inequality \eqref{eq:renewal-shift}, applied to $G$ with indices $s-m\le\rho-m$ and shift $m$, gives
\[
G_s\,G_{\rho-m}\le G_{s-m}\,G_\rho\qquad(0\le m\le s\le\rho).
\]
For $0\le s\le\rho$ let
\[
\mathsf{F}_s=\sum_{m=M}^{s}\theta_m\,f_m\,G_{s-m},
\]
so that $\mathsf{F}_s=0$ for $s<M$ and the sum on the right-hand side of the claim is $\mathsf{F}_\rho$. By~(3) we have $\theta_m\ge\theta_M\ge 0$ for $M\le m\le\rho$, and therefore $\mathsf{F}_s\ge 0$ for every $s$.

We begin with the inequality in which the weights $\theta_m$ are left out: for $M\le s\le\rho$ we have
\begin{equation}\label{eq:renewal-tails}
\Bigl(\sum_{m=M}^{s}f_m\,G_{s-m}\Bigr)\,G_\rho\le\Bigl(\sum_{m=M}^{\rho}f_m\,G_{\rho-m}\Bigr)\,G_s .
\end{equation}
Indeed, by the renewal equation at indices $s$ and $\rho$, the two sums in brackets equal $G_s-\sum_{m=1}^{M-1}f_mG_{s-m}$ and $G_\rho-\sum_{m=1}^{M-1}f_mG_{\rho-m}$, so that the inequality to be shown is equivalent to
\[
\sum_{m=1}^{M-1}f_m\,G_s\,G_{\rho-m}\le\sum_{m=1}^{M-1}f_m\,G_{s-m}\,G_\rho ,
\]
which holds term by term, because $G_s\,G_{\rho-m}\le G_{s-m}\,G_\rho$ for $0\le m\le s$, as we saw above.

Next we show that the same holds with the weights, i.e.\ that
\begin{equation}\label{eq:renewal-weighted}
\mathsf{F}_s\,G_\rho\le\mathsf{F}_\rho\,G_s\qquad(0\le s\le\rho).
\end{equation}
For $s<M$ the left-hand side is $0$ and the right-hand side is nonnegative. Let $M\le s\le\rho$. We split $\mathsf{F}_\rho$ into the terms with $m\le s$ and those with $m>s$, which gives
\begin{align*}
\mathsf{F}_\rho\,G_s-\mathsf{F}_s\,G_\rho
&=\sum_{m=s+1}^{\rho}\theta_m\,f_m\,G_{\rho-m}\,G_s\\
&\quad-\sum_{m=M}^{s}\theta_m\,f_m\bigl(G_{s-m}\,G_\rho-G_s\,G_{\rho-m}\bigr).
\end{align*}
In the second sum we have $\theta_m\le\theta_s$ for $M\le m\le s$ by~(3), and every difference in brackets is nonnegative, again because $G_s\,G_{\rho-m}\le G_{s-m}\,G_\rho$ for $0\le m\le s$. Hence
\begin{align*}
\sum_{m=M}^{s}\theta_m\,f_m\bigl(G_{s-m}\,G_\rho-G_s\,G_{\rho-m}\bigr)
&\le\theta_s\sum_{m=M}^{s}f_m\bigl(G_{s-m}\,G_\rho-G_s\,G_{\rho-m}\bigr)\\
&\le\theta_s\sum_{m=s+1}^{\rho}f_m\,G_{\rho-m}\,G_s ,
\end{align*}
where the second inequality is \eqref{eq:renewal-tails}, with the terms $M\le m\le s$ of its right-hand side moved to the left. In the first sum we have $\theta_m\ge\theta_s$ for $s<m\le\rho$, again by~(3), so that
\[
\sum_{m=s+1}^{\rho}\theta_m\,f_m\,G_{\rho-m}\,G_s\ge\theta_s\sum_{m=s+1}^{\rho}f_m\,G_{\rho-m}\,G_s .
\]
The last two displays together show that $\mathsf{F}_\rho\,G_s-\mathsf{F}_s\,G_\rho\ge 0$.

It remains to pass from $G$ to $v$. Let $\iota_s=v_s-\sum_{i=1}^{s}f_iv_{s-i}$ for $s\ge 0$ be the innovations of $v$. By~(4) we have $\iota_s\ge 0$ for $s\le\rho$, and by Lemma~\ref{lem:convolutionform} we have $v_s=\sum_{t=0}^{s}G_{s-t}\,\iota_t$ for every $s\ge 0$. Inserting this formula for $v_{\rho-m}$ and interchanging the two summations, we obtain
\[
\sum_{m=M}^{\rho}\theta_m\,f_m\,v_{\rho-m}=\sum_{t=0}^{\rho}\iota_t\sum_{m=M}^{\rho-t}\theta_m\,f_m\,G_{\rho-t-m}=\sum_{t=0}^{\rho}\iota_t\,\mathsf{F}_{\rho-t} .
\]
Multiplying by $G_\rho$ and applying the inequality \eqref{eq:renewal-weighted} with $s=\rho-t$ to each term, where $\iota_t\ge 0$, we obtain
\[
\Bigl(\sum_{m=M}^{\rho}\theta_m\,f_m\,v_{\rho-m}\Bigr)\,G_\rho=\sum_{t=0}^{\rho}\iota_t\,\mathsf{F}_{\rho-t}\,G_\rho\le\mathsf{F}_\rho\sum_{t=0}^{\rho}G_{\rho-t}\,\iota_t=\mathsf{F}_\rho\,v_\rho .\qedhere
\]
\end{proof}

\subsection{An explicit kernel and its renewal sequence}\label{subsec:renewal-frozen}

We now turn to the kernel to which Section~\ref{sec:far} applies Proposition~\ref{prop:extremecolumn}. The \emph{frozen kernel} is the sequence
\[
f_0=0,\qquad f_1=\tfrac32\,e^{-1},\qquad f_i=e^{-i}\,\varepsilon_i\quad(i\ge 2),
\]
where $(\varepsilon_i)$ is the free kernel of Section~\ref{sec:prelim}. Since $\varepsilon_0=0$ and $\varepsilon_1=1$, this can be written as
\[
f_i=e^{-i}\,\varepsilon_i+\begin{cases}\tfrac12\,e^{-1}&\text{if } i=1,\\ 0&\text{otherwise}\end{cases}\qquad(i\ge 0).
\]
The factor $e^{-i}$ is there because $\sum_{m}e^{-m}\varepsilon_m=1-e^{-1}$ by Lemma~\ref{lem:epsmass}, so that the frozen kernel has sum $1-\frac12e^{-1}<1$; the summand $\frac12e^{-1}$ at $i=1$ takes the place of a coefficient of the level equation (Lemma~\ref{lem:levelequation}) that is at least $\frac12$ in the rows to which Section~\ref{sec:far} applies the bound. From here to the end of the section $f$ denotes the frozen kernel. We have $f_i>0$ for every $i\ge 1$, because $\varepsilon_i>0$ for $i\ge 1$.

The renewal sequence of $f$ is given by a closed formula, from which the lower bound of Lemma~\ref{lem:frozenlower} will be obtained, and we therefore define a sequence $G$ by this formula. For integers $j\ge 0$ and $s\ge 0$ let
\[
h_{j,s}=[X^s]\,\bigl(X^jC^{j+1}\bigr)
\]
be the coefficient of $X^s$ in the power series $X^jC^{j+1}$, where $C=\sum_{n\ge 0}c_nX^n$ is the series of Section~\ref{sec:prelim}, and for $s\ge 0$ let
\begin{equation}\label{eq:renewal-resolvent}
H_s=\sum_{j=0}^{s}2^{-j}\,h_{j,s},\qquad G_s=e^{-s}H_s=e^{-s}\sum_{j=0}^{s}2^{-j}\,[X^s]\,X^jC^{j+1} .
\end{equation}
From here to the end of the section $G$ denotes the sequence $(G_s)_{s\ge 0}$ defined in this way. Since $X^jC^{j+1}$ is divisible by $X^j$, and since $[X^s]C=c_s$, we have
\[
h_{j,s}=0\quad\text{for } s<j,\qquad\text{and}\qquad h_{0,s}=c_s .
\]
In particular the sum over $j$ that defines $H_s$ may be extended to any range $0\le j\le N$ with $N\ge s$ without changing its value. The first values are $G_0=1$, $G_1=\frac32e^{-1}$ and $G_2=\frac{11}{4}e^{-2}$.

The formula for $G_s$ is found from generating series. By Lemma~\ref{lem:quotient} with $F=1$ we have $1/C=1-\sum_{m\ge 1}\varepsilon_mX^m$. Before the scaling by $e^{-s}$, the renewal equation of the frozen kernel reads $H_s=\sum_{m=1}^{s}\varepsilon_mH_{s-m}+\frac12H_{s-1}$ for $s\ge 1$, with $H_0=1$; it says that the series $\sum_sH_sX^s$ is the inverse of $1/C-X/2$, and
\[
\frac{1}{1/C-X/2}=\frac{C}{1-XC/2}=\sum_{j\ge 0}\Bigl(\frac{X}{2}\Bigr)^{j}C^{j+1} ,
\]
where only the terms with $j\le s$ contribute to the coefficient of $X^s$. In the proof of the following lemma we verify the renewal equation for the coefficients directly.

\begin{lemma}\label{lem:frozenrenewal}
The sequence $G$ defined by \eqref{eq:renewal-resolvent} is the renewal sequence of the frozen kernel $f$ in the sense of Definition~\ref{def:renewal}: we have $G_0=1$ and, for every integer $n\ge 1$,
\[
G_n=\sum_{i=1}^{n}f_i\,G_{n-i} .
\]
\end{lemma}

\begin{proof}
Let $h_{j,s}$ and $H_s$ be as in the definition of $G$. Since $h_{0,s}=c_s$, we have $H_0=h_{0,0}=c_0=1$, and hence $G_0=1$.

Let $j\ge 0$ and $s\ge 1$. Lemma~\ref{lem:quotient}, applied to $F=X^jC^{j+1}$, gives
\[
[X^s]\Bigl(\frac{1}{C}\,X^jC^{j+1}\Bigr)=h_{j,s}-\sum_{m=1}^{s}\varepsilon_m\,h_{j,s-m} .
\]
The series on the left is $X^jC^{j}$. Its coefficient of $X^s$ is $0$ for $j=0$, because $s\ge 1$, and it is $h_{j-1,s-1}$ for $j\ge 1$, because $X^jC^{j}=X\cdot X^{j-1}C^{j}$. Thus
\[
h_{0,s}=\sum_{m=1}^{s}\varepsilon_m\,h_{0,s-m}\qquad\text{and}\qquad h_{j,s}=\sum_{m=1}^{s}\varepsilon_m\,h_{j,s-m}+h_{j-1,s-1}\quad(j\ge 1).
\]
We multiply the identity for $j$ by $2^{-j}$ and sum over $0\le j\le s$. On the left we obtain $H_s$. On the right, since $h_{j,s-m}=0$ for $j>s-m$, we have $\sum_{j=0}^{s}2^{-j}h_{j,s-m}=H_{s-m}$ for $1\le m\le s$; the remaining terms add up to
\[
\sum_{j=1}^{s}2^{-j}\,h_{j-1,s-1}=\frac12\sum_{j=0}^{s-1}2^{-j}\,h_{j,s-1}=\frac12\,H_{s-1} .
\]
Hence
\[
H_s=\sum_{m=1}^{s}\varepsilon_m\,H_{s-m}+\tfrac12\,H_{s-1}\qquad(s\ge 1).
\]
Multiplied by $e^{-s}$, this equation is $G_s=\sum_{i=1}^{s}f_i\,G_{s-i}$, because $f_1=e^{-1}\varepsilon_1+\frac12e^{-1}$ and $f_i=e^{-i}\varepsilon_i$ for $i\ge 2$.
\end{proof}

We turn to the two estimates for $G$ through which the bound of Proposition~\ref{prop:extremecolumn} becomes explicit: an upper bound for the partial sums of $G$, and a lower bound of order $s^{-3/2}$ for a single term $G_s$. The lower bound is needed because in that proposition $G_\rho$ is a factor on the side of the sum that is to be bounded.

\begin{lemma}\label{lem:frozenmass}
For every integer $N\ge 0$, the renewal sequence $G$ of the frozen kernel satisfies
\[
\sum_{s=0}^{N-1}G_s\le 2e .
\]
\end{lemma}

\begin{proof}
We apply Lemma~\ref{lem:convolutionform} to the constant sequence $1$. For this we first bound the partial sums of the frozen kernel $f$: for every integer $s\ge 0$,
\[
\sum_{i=1}^{s}f_i\le 1-\tfrac12\,e^{-1} .
\]
Indeed, by the definition of the frozen kernel the sum on the left is at most $\frac12e^{-1}+\sum_{i=1}^{s}e^{-i}\varepsilon_i$, and the partial sums of the series of the nonnegative numbers $e^{-i}\varepsilon_i$ are at most its sum, which is $1-e^{-1}$ by Lemma~\ref{lem:epsmass}.

Without loss of generality, $N\ge 1$. The constant sequence $v_s=1$ satisfies $v_s=\iota_s+\sum_{i=1}^{s}f_iv_{s-i}$ for every $s\ge 0$, with $\iota_s=1-\sum_{i=1}^{s}f_i$. Since $G$ is the renewal sequence of $f$ by Lemma~\ref{lem:frozenrenewal}, Lemma~\ref{lem:convolutionform} applies, and at index $N-1$ it gives
\[
1=\sum_{t=0}^{N-1}G_{N-1-t}\,\iota_t .
\]
Here $\iota_t\ge\frac12e^{-1}$ by the bound for the partial sums of $f$, and $G_{N-1-t}>0$ by \eqref{eq:renewal-positive}, since $f_i>0$ for $i\ge 1$. Hence
\[
1\ge\tfrac12\,e^{-1}\sum_{t=0}^{N-1}G_{N-1-t}=\tfrac12\,e^{-1}\sum_{s=0}^{N-1}G_s ,
\]
which is the claim.
\end{proof}

\begin{lemma}\label{lem:frozenlower}
For every integer $s\ge 6$, the renewal sequence $G$ of the frozen kernel satisfies
\[
G_s\ge\frac{e}{s+1}\cdot\frac{1-\dfrac{1}{12(s+1)}}{\sqrt{2\pi(s+1)}}\cdot\Bigl(4-\frac{21}{s+1}\Bigr).
\]
\end{lemma}

\begin{proof}
Let $h_{j,s}$ and $H_s$ be as in the definition \eqref{eq:renewal-resolvent} of $G$, so that $G_s=e^{-s}H_s$. Recall that the weights of Lemma~\ref{lem:abel} are $w_0(x)=1$ and $w_n(x)=x\,(x+n)^{n-1}$ for $n\ge 1$.

We first bring the identity of Abel into convolution form: for every integer $n\ge 1$ and all real numbers $x$ and $y$,
\[
\sum_{i=0}^{n}\binom{n}{i}\,w_i(x)\,w_{n-i}(y)=(x+y)\,(x+y+n)^{n-1} .
\]
To see this, put $L=\sum_{i=0}^{n-1}\binom{n}{i}\,w_i(x)\,(y+n-i)^{n-1-i}$. Since $w_{n-i}(y)=y\,(y+n-i)^{n-1-i}$ for $i\le n-1$ and $w_0(y)=1$, the left-hand side of the identity to be proved is $yL+w_n(x)$. On the other hand, Lemma~\ref{lem:abel} with $y+n$ in the place of $y$ gives
\[
\sum_{i=0}^{n}\binom{n}{i}\,w_i(x)\,(y+n-i)^{n-i}=(x+y+n)^{n} .
\]
The term with $i=n$ on the left is again $w_n(x)$. In the terms with $i\le n-1$ we write $(y+n-i)^{n-i}=\bigl(y+(n-i)\bigr)(y+n-i)^{n-1-i}$ and use that $(n-i)\binom{n}{i}=n\binom{n-1}{i}$. The identity then reads
\[
yL+n\sum_{i=0}^{n-1}\binom{n-1}{i}\,w_i(x)\,(y+n-i)^{n-1-i}+w_n(x)=(x+y+n)^{n} ,
\]
and the sum in the middle equals $(x+y+n)^{n-1}$ by Lemma~\ref{lem:abel} for $n-1$, again with $y+n$ in the place of $y$. Therefore
\[
yL+w_n(x)=(x+y+n)^{n}-n\,(x+y+n)^{n-1}=(x+y)\,(x+y+n)^{n-1} .
\]

With the convolution form we can compute the coefficients $h_{j,s}$. We claim that, for all integers $j\ge 0$ and $n\ge 0$,
\[
[X^n]\,C^{j+1}=\frac{w_n(j+1)}{n!} .
\]
The proof is by induction on $j$. For $j=0$ the claim is the identity $c_n=w_n(1)/n!$ of Section~\ref{sec:prelim}. If the claim holds for $j$, then the formula for the coefficients of a product gives
\[
[X^n]\,C^{j+2}=\sum_{i=0}^{n}\,[X^i]\,C^{j+1}\cdot c_{n-i}=\frac{1}{n!}\sum_{i=0}^{n}\binom{n}{i}\,w_i(j+1)\,w_{n-i}(1) .
\]
For $n\ge 1$ this is $(j+2)\,(j+2+n)^{n-1}/n!=w_n(j+2)/n!$ by the convolution form of the identity of Abel with $x=j+1$ and $y=1$, and for $n=0$ it is $1=w_0(j+2)$. This proves the claim.

For $j\le s$ we have $h_{j,s}=[X^s]\,(X^jC^{j+1})=[X^{s-j}]\,C^{j+1}$, and the formula just proved gives
\[
h_{j,s}=\frac{w_{s-j}(j+1)}{(s-j)!}=\frac{(j+1)\,(s+1)^{s-j-1}}{(s-j)!}\qquad(0\le j\le s),
\]
where for $j=s$ both expressions are $1$. Dividing by $c_s=(s+1)^{s-1}/s!$ and writing $s!/(s-j)!=\prod_{i=1}^{j}(s+1-i)$, we obtain the closed form
\[
h_{j,s}=c_s\,(j+1)\prod_{i=1}^{j}\Bigl(1-\frac{i}{s+1}\Bigr)\qquad(0\le j\le s).
\]

We now bound the $h_{j,s}$ from below. For real numbers $t_1,\dots,t_j$ in $[0,1]$ we have $\prod_{i=1}^{j}(1-t_i)\ge 1-\sum_{i=1}^{j}t_i$, by a routine induction on $j$. Applied to $t_i=i/(s+1)$ with $1\le i\le j$, whose sum is $j(j+1)/(2(s+1))$, this turns the closed form into
\[
h_{j,s}\ge c_s\,(j+1)\Bigl(1-\frac{j(j+1)}{2(s+1)}\Bigr)\qquad(0\le j\le s).
\]

For the summation over $j$ we use the two formulas
\[
\sum_{j=0}^{s}\frac{j+1}{2^j}=4-\frac{s+3}{2^s},\qquad\sum_{j=0}^{s}\frac{j\,(j+1)^2}{2^j}=40-\frac{s^3+8s^2+27s+40}{2^s} ,
\]
which hold for every integer $s\ge 0$ and are easily proved by induction on $s$. Multiplying the lower bound for $h_{j,s}$ by $2^{-j}$ and summing over $0\le j\le s$, we obtain
\[
H_s\ge c_s\Bigl(4-\frac{s+3}{2^s}\Bigr)-\frac{c_s}{2(s+1)}\Bigl(40-\frac{s^3+8s^2+27s+40}{2^s}\Bigr) .
\]
A direct computation shows that the right-hand side equals
\[
c_s\Bigl(4-\frac{21}{s+1}\Bigr)+\frac{c_s}{s+1}\Bigl(1+\frac{s^3+6s^2+19s+34}{2^{s+1}}\Bigr) ,
\]
and the second term of this expression is positive. Hence
\[
H_s\ge c_s\Bigl(4-\frac{21}{s+1}\Bigr) .
\]

Finally, we insert the lower Stirling bound for $c_s$. Since $c_s=(s+1)^{s-1}/s!=(s+1)^{s}/(s+1)!$, the lower bound in \eqref{eq:prelim-stirling} at $n=s+1$ gives
\[
e^{-s}c_s=\frac{e}{s+1}\cdot e^{-(s+1)}\,\frac{(s+1)^{s+1}}{(s+1)!}\ge\frac{e}{s+1}\cdot\frac{1-\dfrac{1}{12(s+1)}}{\sqrt{2\pi(s+1)}} .
\]
Since $s\ge 6$, the factor $4-\frac{21}{s+1}$ is positive. So we may multiply the inequality $H_s\ge c_s\,\bigl(4-\frac{21}{s+1}\bigr)$ by $e^{-s}$ and insert the last display, which gives
\[
G_s=e^{-s}H_s\ge e^{-s}c_s\Bigl(4-\frac{21}{s+1}\Bigr)\ge\frac{e}{s+1}\cdot\frac{1-\dfrac{1}{12(s+1)}}{\sqrt{2\pi(s+1)}}\cdot\Bigl(4-\frac{21}{s+1}\Bigr) .\qedhere
\]
\end{proof}

\section{Persistence on the far rows and the single crossing}\label{sec:far}
In this section we prove Proposition~\ref{prop:chainrange}, one of the three propositions about columns from which Theorem~\ref{thm:main} is derived: a column that is positive in a row of its chain range is positive in the next row of that range. What remains to be treated are the far rows, those beyond row $k+0.47k$, which Proposition~\ref{prop:nearrows} does not reach. This is done in Proposition~\ref{prop:farrows} for the columns $k\ge 82$. Together with Proposition~\ref{prop:nearrows} it gives the single crossing of every column $k\ge 82$ (Proposition~\ref{prop:onecrossing}), and with the computed columns $k\le 81$ it completes the proof of Proposition~\ref{prop:chainrange}.

\subsection{Auxiliary results}\label{subsec:far-lemmas}

We begin with three auxiliary results for the proof of Proposition~\ref{prop:farrows}. The first is a lemma about arbitrary sequences. From an inequality of the form of the column equation, which is hypothesis~(3) below, and a lower bound of the form of Proposition~\ref{prop:diagonallower}, which is hypothesis~(4), it derives a lower bound for one term of a sequence $\eta$ in terms of a second sequence $u$ alone, provided that a run of terms of $\eta$ before it is nonnegative.

\begin{lemma}\label{lem:windowbound}
Let $(h_m)_{m\ge 0}$, $(u_j)_{j\ge 0}$ and $(\eta_j)_{j\ge 0}$ be sequences of real numbers, let $\gamma>0$ be a real number, and let $k\ge 1$, $n$ and $\tau$ be integers with $k+1\le n$ and $2\le\tau\le n-k$. Assume that
\begin{enumerate}
\item $h_m>0$ for every $m\ge 1$;
\item $u_{k-1}=0$;
\item $\dfrac2n\,u_{n-1}+\dfrac{k+1}{n}\,\eta_{n-1}+\displaystyle\sum_{m=1}^{n-k}h_m\,\eta_{n-m}\le\eta_n$;
\item $\gamma\,u_{j-1}-u_j\le\eta_j$ for every $j$ with $k\le j\le n-1$;
\item $h_{m+1}\le\gamma\,h_m$ for every $m$ with $1\le m<\tau$;
\item $\eta_j\ge 0$ for every $j$ with $n-\tau+1\le j\le n-1$.
\end{enumerate}
Then
\[
\Bigl(\frac2n-\frac{h_\tau}{\gamma^{\tau-1}}\Bigr)u_{n-1}+\sum_{m=\tau}^{n-k-1}\bigl(\gamma\,h_m-h_{m+1}\bigr)\,u_{n-1-m}\le\eta_n .
\]
\end{lemma}

\begin{proof}
Put $\psi=h_\tau/\gamma^{\tau-1}$, which is positive by~(1). We first note that
\[
\psi\,\gamma^{m-1}\le h_m\qquad(1\le m\le\tau).
\]
For $m=\tau$ this is an equality, and for smaller $m$ it follows from~(5) by downward induction on $m$. Since $\tau\ge 2$, we have $\eta_{n-1}\ge 0$ by~(6), and~(3) gives
\begin{equation}\label{eq:far-reduced}
\frac2n\,u_{n-1}+\sum_{m=1}^{n-k}h_m\,\eta_{n-m}\le\eta_n .
\end{equation}
For $1\le m\le n-k$, hypothesis~(4) at $j=n-m$ reads $\gamma\,u_{n-1-m}-u_{n-m}\le\eta_{n-m}$.

Let $1\le m\le\tau-1$. Then $\eta_{n-m}\ge 0$ by~(6), and the bound $\psi\,\gamma^{m-1}\le h_m$ and~(4) give
\[
h_m\,\eta_{n-m}\ge\psi\gamma^{m-1}\,\eta_{n-m}\ge\psi\,\bigl(\gamma^{m}\,u_{n-1-m}-\gamma^{m-1}\,u_{n-m}\bigr).
\]
The right-hand sides telescope, and since $\psi\gamma^{\tau-1}=h_\tau$ we obtain
\[
\sum_{m=1}^{\tau-1}h_m\,\eta_{n-m}\ge h_\tau\,u_{n-\tau}-\psi\,u_{n-1}.
\]

Let now $\tau\le m\le n-k$. Here we use only that $h_m>0$, so that $h_m\,\eta_{n-m}\ge h_m\,(\gamma\,u_{n-1-m}-u_{n-m})$ by~(4). Summing over these $m$ and shifting the index by one in the terms $h_m\,u_{n-m}$, we get
\begin{align*}
\sum_{m=\tau}^{n-k}h_m\,\eta_{n-m}&\ge\sum_{m=\tau}^{n-k}\gamma\,h_m\,u_{n-1-m}-\sum_{m=\tau-1}^{n-k-1}h_{m+1}\,u_{n-1-m}\\
&=-h_\tau\,u_{n-\tau}+\sum_{m=\tau}^{n-k-1}\bigl(\gamma\,h_m-h_{m+1}\bigr)\,u_{n-1-m}+\gamma\,h_{n-k}\,u_{k-1},
\end{align*}
where the last term vanishes by~(2). Adding the two bounds, in which the terms $h_\tau\,u_{n-\tau}$ cancel, and inserting the result into the inequality \eqref{eq:far-reduced} that we derived from~(3), we obtain the claim.
\end{proof}

The second result is an upper bound for $\frac N2\,\varepsilon_{t+1}\vartheta_k^{\,t}$ when $t$ is close to a fixed multiple $\sigma k$ of $k$ and $N$ is at most a multiple $Lk$ of $k$; recall that $\vartheta_k=1/\gamma_k$ is the diagonal tilt of column $k$ (Section~\ref{sec:crossing}). Let $k\ge 15$ and $t\ge 1$ be integers, and let $\sigma>0$, $\lambda>0$, $L$ and $N$ be real numbers with
\[
\sigma\,(k+3)\le\lambda\,(\sigma k+1),\qquad\sigma k\le t\le\sigma k+1,\qquad 0\le N\le L\,k .
\]
Then
\begin{equation}\label{eq:far-loss}
\frac N2\;\varepsilon_{t+1}\,\vartheta_k^{\,t}\le\frac{\lambda\,L\,e^{5\sigma/2}}{2\sigma\sqrt{2\pi\sigma k}} .
\end{equation}
To see this, write $\psi=\varepsilon_{t+1}\vartheta_k^{t}$. We have $(t+1)\,\varepsilon_{t+1}=t^t/t!$, and $\vartheta_k\le e^{-1}e^{5/(2k)}$ by Lemma~\ref{lem:tiltbounds}. Hence the upper bound in \eqref{eq:prelim-stirling} gives
\[
(t+1)\,\psi=\frac{t^t}{t!}\,\vartheta_k^{t}\le e^{-t}\,\frac{t^t}{t!}\,e^{5t/(2k)}\le\frac{e^{5t/(2k)}}{\sqrt{2\pi t}}.
\]
Since $t\le\sigma k+1$, we have $e^{5t/(2k)}\le e^{5\sigma/2}\,e^{5/(2k)}$. Moreover $e^{5/(2k)}\le 1+\frac3k$, by the elementary inequality $e^{u}\le 1+\frac65u$ for $0\le u\le\frac16$, which applies to $u=5/(2k)$ because $k\ge 15$. Thus
\[
(t+1)\,\sqrt{2\pi t}\;\psi\le e^{5\sigma/2}\Bigl(1+\frac3k\Bigr).
\]
Since $\sigma k\le t$, we have $\bigl(1+\frac3k\bigr)\,\sigma k=\sigma\,(k+3)\le\lambda\,(\sigma k+1)\le\lambda\,(t+1)$, so that $\sqrt{2\pi t}\;\sigma k\,\psi\le\lambda\,e^{5\sigma/2}$, and with $0\le N\le Lk$ this gives
\[
\frac N2\,\psi\cdot 2\sigma\sqrt{2\pi t}\le\lambda\,L\,e^{5\sigma/2}.
\]
Since $t\ge\sigma k$, the bound \eqref{eq:far-loss} follows.

The third result, again a lemma, concerns the ratios $g_m$ of the free kernel (Section~\ref{sec:prelim}), for which $\varepsilon_{m+1}=g_m\,\varepsilon_m$, and the frozen kernel $f$ of Section~\ref{sec:renewal}, for which $f_1=\frac32e^{-1}$ and $f_i=e^{-i}\varepsilon_i$ for $i\ge 2$. From here on $f$ always denotes the frozen kernel and $G$ its renewal sequence (Lemma~\ref{lem:frozenrenewal}).

\begin{lemma}\label{lem:farweights}
Let $k$ and $m$ be integers with $m\ge 49$, $m+1\le k$ and $3k\le 5\,(m+1)$. Then the frozen kernel $f$ satisfies
\[
(g_m-\gamma_k)\,f_m\le\frac{105}{100}\cdot\frac{c_\ast}{k^2\sqrt k} .
\]
\end{lemma}

\begin{proof}
Write $s=m+1$, so that $50\le s\le k$ and $3k\le 5s$.

We begin with $g_m-\gamma_k$. Since $g_m=\frac{m}{m+1}\bigl(1+\frac{1}{m-1}\bigr)^{m-1}$, we have $\log g_m=(m-1)\log\bigl(1+\frac{1}{m-1}\bigr)-\log\bigl(1+\frac1m\bigr)$. By the Taylor bounds $u-\frac{u^2}{2}\le\log(1+u)\le u-\frac{u^2}{2}+\frac{u^3}{3}$ for $u\ge 0$, used at $u=\frac1m$ and at $u=\frac{1}{m-1}$,
\begin{align*}
\log g_m&\le 1-\frac{1}{2(m-1)}+\frac{1}{3(m-1)^2}-\frac1m+\frac{1}{2m^2}\\
&=1-\frac{3}{2(m+1)}-\frac{7m^3-17m^2+9m-3}{6\,(m-1)^2\,m^2\,(m+1)}\le 1-\frac{3}{2(m+1)},
\end{align*}
where the equality is a direct computation and the last fraction is clearly positive. Hence $g_m\le\exp\bigl(1-\frac{3}{2s}\bigr)$. On the other hand, $\gamma_k=1/\vartheta_k\ge\exp\bigl(1-\frac{5}{2k}\bigr)$ by Lemma~\ref{lem:tiltbounds}. Put
\[
\alpha=\frac{5}{2k}-\frac{3}{2s}=\frac{5s-3k}{2ks},
\]
which is nonnegative because $3k\le 5s$. Then $\gamma_k\ge\exp\bigl(1-\frac{3}{2s}\bigr)\,e^{-\alpha}\ge\exp\bigl(1-\frac{3}{2s}\bigr)\,(1-\alpha)$, and therefore
\[
g_m-\gamma_k\le\exp\Bigl(1-\frac{3}{2s}\Bigr)\,\alpha\le e\,\alpha .
\]

Next we bound $f_m=e^{-m}\varepsilon_m$. Since $m\,\varepsilon_m=(m-1)^{m-1}/(m-1)!$, the upper bound in \eqref{eq:prelim-stirling} at $m-1$ gives
\[
f_m=\frac{e^{-1}}{m}\cdot e^{-(m-1)}\,\frac{(m-1)^{m-1}}{(m-1)!}\le\frac{e^{-1}}{m\,\sqrt{2\pi\,(m-1)}}.
\]

Since $f_m>0$ and $\alpha\ge 0$, the two bounds give
\[
(g_m-\gamma_k)\,f_m\le e\,\alpha\,f_m\le\frac{\alpha}{m\,\sqrt{2\pi\,(m-1)}}=\frac{c_\ast\,(5s-3k)}{2ks\,(s-1)\sqrt{s-2}}.
\]
The right-hand side is at most $\frac{105}{100}\,c_\ast\,k^{-5/2}$, as the lemma claims, if and only if
\[
(5s-3k)\,k^{3/2}\le\frac{105}{100}\cdot 2s\,(s-1)\sqrt{s-2}.
\]
This follows from two inequalities. The first is $(5s-3k)\,k^{3/2}\le 2s^{5/2}$, which holds because
\[
2s^{5/2}-(5s-3k)\,k^{3/2}=\bigl(\sqrt s-\sqrt k\bigr)^2\,\bigl(2s^{3/2}+4s\,k^{1/2}+6s^{1/2}k+3k^{3/2}\bigr)\ge 0,
\]
as one sees by expanding the product. The second is $s^{3/2}\le\frac{105}{100}\,(s-1)\sqrt{s-2}$ for $s\ge 50$; after squaring it is the polynomial inequality $s^3\le 1.1025\,(s-1)^2(s-2)$, which is easy to check. Together they give the required inequality.
\end{proof}

\subsection{Rows beyond $k+0.47k$}\label{subsec:far-rows}

Proposition~\ref{prop:farrows} below says that a column $k\ge 82$ is positive in a far row $k+r$ if it is positive in the rows before it, from some row $k+r_q$ with $r_q\le k/3$ on. The comparison with the crossing row, on which Section~\ref{sec:near} rests, does not reach beyond row $k+0.47k$, and the proof uses the column equation at row $k+r$ alone. If column $k$ is positive in a window, that is, in a run of rows before this row, Lemma~\ref{lem:windowbound} bounds the step in row $k+r$ from below by a combination of values of level $k-1$; for a far row there is a window whose length is a fixed fraction of $k$, because $3\,r_q\le k$. The difficulty is that for $r>\frac35k$ some coefficients of this combination may be negative. The values of level $k-1$ that carry them are known only to satisfy the level equation, and this is the situation for which the comparison of Section~\ref{sec:renewal} is made: relative to the value of level $k-1$ in row $k+r-1$, it bounds these terms by an explicit quantity, which is bounded but is not small when $r$ is close to $k$. The proof therefore ends with a comparison of explicit constants, and this comparison holds for $k\ge 82$.

\begin{proposition}\label{prop:farrows}
Let $k\ge 82$, $r$ and $r_q$ be nonnegative integers with
\[
47\,k<100\,r,\qquad r+1\le k,\qquad 3\,r_q\le k .
\]
Assume that $d_{j,k}>0$ for every integer $j$ with $k+r_q\le j<k+r$. Then $d_{k+r,k}>0$.
\end{proposition}

\begin{proof}
Put $n=k+r$, and write $\delta_j=\delta^{(k)}_j=d_{j,k}/j!$, which has the sign of $d_{j,k}$. Since $\frac{47}{100}k<r\le k-1$, we have $k+1\le n\le 2k-1$, so that row $n$ lies in the chain range of column $k$. We have to show that $\delta_n>0$.

We begin with a lower bound for $\delta_n$. Let $t_0\ge 1$ be an integer with $t_0+1\le r$ and $5\,(t_0+1)\le 3k$, suppose that $\delta_j\ge 0$ for $n-t_0\le j\le n-1$, and put $\psi=\varepsilon_{t_0+1}\vartheta_k^{t_0}$. We claim that
\begin{equation}\label{eq:far-window}
\Bigl(\frac2n-\psi\Bigr)\beta^{(k-1)}_{n-1}+\sum_{m=t_0+1}^{r-1}\bigl(\gamma_k\varepsilon_m-\varepsilon_{m+1}\bigr)\,\beta^{(k-1)}_{n-1-m}\le\delta_n .
\end{equation}
We apply Lemma~\ref{lem:windowbound} with $h_m=\varepsilon_m$, $u_j=\beta^{(k-1)}_j$, $\eta_j=\delta_j$, $\gamma=\gamma_k$ and $\tau=t_0+1$; then $2\le\tau\le r=n-k$, and $h_\tau/\gamma^{\tau-1}=\psi$ because $\vartheta_k=1/\gamma_k$. Hypotheses~(1) and~(2) hold because $\varepsilon_m>0$ for $m\ge 1$ and $\beta^{(k-1)}_{k-1}=0$, and~(3) holds with equality, by the column equation at row $n$ (Proposition~\ref{prop:columnequation}). Hypothesis~(4) is the diagonal lower bound: Proposition~\ref{prop:diagonallower} gives $\gamma_k\,\beta^{(k-1)}_{j-1}\le\beta^{(k)}_j$, and subtracting $\beta^{(k-1)}_j$ from both sides we obtain
\[
\gamma_k\,\beta^{(k-1)}_{j-1}-\beta^{(k-1)}_{j}\le\delta_j\qquad(j\ge 1).
\]
Hypothesis~(5) holds by Lemma~\ref{lem:epsgrowth}, because $5\,(m+1)\le 5\,(t_0+1)\le 3k$ for $1\le m<\tau$, and~(6) is the assumption on rows $n-t_0,\dots,n-1$. The conclusion of the lemma is the claimed bound.

Next we choose $t_0$, the length of the window. We distinguish three ranges of $r$, which together contain every $r$ with $\frac{47}{100}k<r\le k-1$:
\[
\textup{(a)}\ \ \tfrac{47}{100}k<r\le\tfrac35k,\qquad\textup{(b)}\ \ \tfrac35k<r\le\tfrac45k,\qquad\textup{(c)}\ \ \tfrac45k<r\le k-1 .
\]
Let $\sigma=\frac{41}{300}$ in range~(a), $\sigma=\frac{4}{15}$ in range~(b) and $\sigma=\frac{7}{15}$ in range~(c), so that $r>\bigl(\sigma+\frac13\bigr)k$ in each case, and let
\[
t_0=\lfloor\sigma k\rfloor+1,\qquad\text{so that}\qquad\sigma k\le t_0\le\sigma k+1 .
\]
Since $3\,r_q\le k$, the integer $r-r_q$ is larger than $\sigma k$, and hence $r-r_q\ge t_0$. So rows $n-t_0,\dots,n-1$ are among rows $k+r_q,\dots,k+r-1$, in which column $k$ is positive by assumption. Moreover, $t_0+1\le\sigma k+2<\bigl(\sigma+\frac13\bigr)k<r$ and $5\,(t_0+1)\le\frac73k+10\le 3k$, because $\sigma\le\frac{7}{15}$ and $k\ge 15$. Hence the lower bound \eqref{eq:far-window} holds with this $t_0$. From now on $t_0$ and $\psi=\varepsilon_{t_0+1}\vartheta_k^{t_0}$ have this meaning.

We also need that the numbers $\beta^{(k-1)}_j$ of level $k-1$ are positive in rows $k,\dots,2k-1$. By Lemma~\ref{lem:levelequation}, applied to level $k-1$,
\begin{equation}\label{eq:far-level}
\beta^{(k-1)}_j=\frac{z^{(k)}_{j-1}}{j}+\frac{k-1}{j}\,\beta^{(k-1)}_{j-1}+\sum_{m=1}^{j-k}\varepsilon_m\,\beta^{(k-1)}_{j-m}\qquad(k\le j\le 2k-1).
\end{equation}
All terms on the right are nonnegative. For $j=k$ the first term is $z^{(k)}_{k-1}/k=c_{k-1}/k>0$, by Lemma~\ref{lem:freerange}, and for $k<j\le 2k-1$ the second term is positive if $\beta^{(k-1)}_{j-1}>0$. By induction on $j$ we get
\[
\beta^{(k-1)}_j>0\qquad(k\le j\le 2k-1),
\]
and in particular $\beta^{(k-1)}_{n-1}>0$.

Suppose first that $r\le\frac35k$, which is range~(a). Every index $m$ of the sum in \eqref{eq:far-window} satisfies $5\,(m+1)\le 5r\le 3k$, so that $\gamma_k\varepsilon_m-\varepsilon_{m+1}\ge 0$ by Lemma~\ref{lem:epsgrowth}. Since the $\beta^{(k-1)}_{n-1-m}$ are nonnegative, the sum is nonnegative, and we obtain
\[
\Bigl(\frac2n-\psi\Bigr)\beta^{(k-1)}_{n-1}\le\delta_n .
\]
Since $\beta^{(k-1)}_{n-1}>0$, it therefore suffices to show that $\frac n2\,\psi<1$; this is done at the end of the proof.

Suppose now that $r>\frac35k$, so that $r$ lies in range~(b) or~(c); some of the coefficients $\gamma_k\varepsilon_m-\varepsilon_{m+1}$ of the lower bound may then be negative. Put
\[
\rho=r-1,\qquad M=\Bigl\lfloor\frac35k\Bigr\rfloor,\qquad B_k=\frac{105}{100}\cdot\frac{c_\ast}{k^2\sqrt k}.
\]
Then $49\le M\le\rho\le k-2$, because $k\ge 82$ and the integer $r$ satisfies $\frac35k<r\le k-1$. We show that
\begin{equation}\label{eq:far-structural}
\Bigl[\Bigl(\frac2n-\psi\Bigr)G_\rho-B_k\sum_{i=0}^{\rho-M}G_i\Bigr]\,\beta^{(k-1)}_{n-1}\le\delta_n\,G_\rho .
\end{equation}

We begin by separating the terms of the lower bound \eqref{eq:far-window} whose coefficients may be negative. Its sum runs over $t_0+1\le m\le\rho$, and $t_0+1\le\sigma k+2\le\frac35k-1<M$, because $\sigma\le\frac{7}{15}$ and $k\ge 23$. For $t_0+1\le m\le M-1$ we have $5\,(m+1)\le 5M\le 3k$, so the coefficient $\gamma_k\varepsilon_m-\varepsilon_{m+1}$ is nonnegative by Lemma~\ref{lem:epsgrowth}, and we omit these terms. For $M\le m\le\rho$ let
\[
\theta_m=\max(g_m-\gamma_k,\,0).
\]
Since $\varepsilon_{m+1}=g_m\,\varepsilon_m$, we have $\gamma_k\varepsilon_m-\varepsilon_{m+1}=(\gamma_k-g_m)\,\varepsilon_m\ge-\theta_m\,\varepsilon_m$. As the $\beta^{(k-1)}_{n-1-m}$ are nonnegative, the lower bound gives
\begin{equation}\label{eq:far-split}
\Bigl(\frac2n-\psi\Bigr)\beta^{(k-1)}_{n-1}-\sum_{m=M}^{\rho}\theta_m\,\varepsilon_m\,\beta^{(k-1)}_{n-1-m}\le\delta_n .
\end{equation}

Next we pass to the scale $e^{-1}$. For an integer $s\ge 0$ let $v_s=e^{-s}\beta^{(k-1)}_{k+s}$. We claim that
\[
\sum_{i=1}^{s}f_i\,v_{s-i}\le v_s\qquad(0\le s\le k-2),
\]
that is, in the terms of Section~\ref{sec:renewal}, that $v$ has nonnegative innovations with respect to $f$ up to $k-2$. For $s=0$ it says that $v_0\ge 0$. Let $1\le s\le k-2$. In the level equation \eqref{eq:far-level} at $j=k+s$ the first term is nonnegative, and $\frac{k-1}{k+s}\ge\frac12$ because $s\le k-2$. Hence
\[
\beta^{(k-1)}_{k+s}\ge\frac12\,\beta^{(k-1)}_{k+s-1}+\sum_{m=1}^{s}\varepsilon_m\,\beta^{(k-1)}_{k+s-m}.
\]
Multiplying by $e^{-s}$, we obtain
\[
v_s\ge\frac12\,e^{-1}\,v_{s-1}+\sum_{m=1}^{s}e^{-m}\varepsilon_m\,v_{s-m}=\sum_{i=1}^{s}f_i\,v_{s-i},
\]
because $f_1=e^{-1}\varepsilon_1+\frac12e^{-1}$ and $f_i=e^{-i}\varepsilon_i$ for $i\ge 2$.

We now apply Proposition~\ref{prop:extremecolumn} to $f$, $G$, $v$ and the weights $\theta_m$, with $\rho$ and $M$ as above. Its hypotheses hold. We have $f_i>0$ for $i\ge 1$, and $f$ is log-convex: for $i\ge 2$ the inequality $f_{i+1}^2\le f_i\,f_{i+2}$ is the inequality $\varepsilon_{i+1}^2\le\varepsilon_i\,\varepsilon_{i+2}$ of Lemma~\ref{lem:freelogconvex}, multiplied by $e^{-2(i+1)}$, and for $i=1$ it follows in the same way from $\varepsilon_2^2\le\varepsilon_1\,\varepsilon_3$, because $f_1\ge e^{-1}\varepsilon_1$. By Lemma~\ref{lem:frozenrenewal}, $G$ is the renewal sequence of $f$. The weights satisfy $\theta_M\ge 0$ and $\theta_m\le\theta_{m'}$ for $M\le m\le m'\le\rho$, because $g_m=\varepsilon_{m+1}/\varepsilon_m$ does not decrease with $m$, again by Lemma~\ref{lem:freelogconvex}. Finally, $v$ has nonnegative innovations with respect to $f$ up to $\rho$, since $\rho\le k-2$. The proposition gives
\[
\Bigl(\sum_{m=M}^{\rho}\theta_m\,f_m\,v_{\rho-m}\Bigr)\,G_\rho\le\Bigl(\sum_{m=M}^{\rho}\theta_m\,f_m\,G_{\rho-m}\Bigr)\,v_\rho .
\]
Since $k+\rho=n-1$, we have $f_m\,v_{\rho-m}=e^{-\rho}\,\varepsilon_m\,\beta^{(k-1)}_{n-1-m}$ for $M\le m\le\rho$ and $v_\rho=e^{-\rho}\beta^{(k-1)}_{n-1}$. Dividing by $e^{-\rho}$, we obtain
\[
\Bigl(\sum_{m=M}^{\rho}\theta_m\,\varepsilon_m\,\beta^{(k-1)}_{n-1-m}\Bigr)\,G_\rho\le\Bigl(\sum_{m=M}^{\rho}\theta_m\,f_m\,G_{\rho-m}\Bigr)\,\beta^{(k-1)}_{n-1}.
\]

It remains to bound the weights $\theta_m\,f_m$. For $M\le m\le\rho$ Lemma~\ref{lem:farweights} applies, because $m\ge M\ge 49$, $m+1\le r\le k$ and $m+1>\frac35k$. Hence $(g_m-\gamma_k)\,f_m\le B_k$, and therefore $\theta_m\,f_m\le B_k$. Clearly $G_s>0$ for every $s\ge 0$, since $G$ is the renewal sequence of a kernel with positive terms. Thus
\[
\sum_{m=M}^{\rho}\theta_m\,f_m\,G_{\rho-m}\le B_k\sum_{m=M}^{\rho}G_{\rho-m}=B_k\sum_{i=0}^{\rho-M}G_i .
\]
Multiplying \eqref{eq:far-split} by $G_\rho$ and inserting the last two displays, we obtain \eqref{eq:far-structural}.

We now reduce the claim to two numerical inequalities. Suppose that $p_1$ and $p_2$ are real numbers with $p_1+p_2<1$ such that
\begin{equation}\label{eq:far-closing}
\frac n2\,\psi\le p_1\qquad\text{and}\qquad\frac n2\,B_k\sum_{i=0}^{\rho-M}G_i\le p_2\,G_\rho .
\end{equation}
Then
\[
\frac n2\,\Bigl[\Bigl(\frac2n-\psi\Bigr)G_\rho-B_k\sum_{i=0}^{\rho-M}G_i\Bigr]\ge(1-p_1-p_2)\,G_\rho>0,
\]
so the bracket in \eqref{eq:far-structural} is positive. Since $\beta^{(k-1)}_{n-1}>0$ and $G_\rho>0$, it follows that $\delta_n>0$. So it suffices to find such $p_1$ and $p_2$.

The first of the two numerical inequalities will come from the bound \eqref{eq:far-loss}. For the second we use the following sufficient condition. For a real number $x\ge 6$ let
\[
\mathcal{T}(x)=\Bigl(1-\frac{1}{12x}\Bigr)\Bigl(4-\frac{21}{x}\Bigr).
\]
Both factors are positive and increasing on this range, so $\mathcal{T}$ is increasing. Suppose that $S$, $w>0$, $\underline r\ge 6$ and $p_2\ge 0$ are real numbers with
\begin{equation}\label{eq:far-four}
\sum_{i=0}^{\rho-M}G_i\le S\,e,\qquad r\le w^2k,\qquad r\ge\underline r,\qquad\frac{105}{100}\,S\,(1+w^2)\,w^3\le 2\,p_2\,\mathcal{T}(\underline r).
\end{equation}
Then the second of the two numerical inequalities holds. Indeed, Lemma~\ref{lem:frozenlower}, which applies at $s=\rho$ because $\rho\ge 6$, reads $G_\rho\ge e\,c_\ast\,r^{-3/2}\,\mathcal{T}(r)$, and $\mathcal{T}(r)\ge\mathcal{T}(\underline r)$. From $r\le w^2k$ we get $(k+r)\,r^{3/2}\le(1+w^2)\,w^3\,k^{5/2}$. Hence
\begin{align*}
\frac n2\,B_k\sum_{i=0}^{\rho-M}G_i&\le\frac{k+r}{2}\cdot\frac{105}{100}\cdot\frac{c_\ast}{k^{5/2}}\cdot S\,e\le\frac{105}{100}\cdot\frac{S\,(1+w^2)\,w^3}{2}\cdot\frac{e\,c_\ast}{r^{3/2}}\\
&\le p_2\,\mathcal{T}(\underline r)\,\frac{e\,c_\ast}{r^{3/2}}\le p_2\,G_\rho .
\end{align*}

Before we choose the constants, we sharpen the bound for the sum of the $G_i$ in one case. By Lemma~\ref{lem:frozenmass} the condition $\sum_{i=0}^{\rho-M}G_i\le S\,e$ holds with $S=2$. In range~(c) this is not enough when $82\le k\le 109$, and we show that for these $k$ and for $r>\frac45k$
\[
\sum_{i=0}^{\rho-M}G_i\le 1.747\,e .
\]
The number of terms of the sum is
\[
J=\rho-M+1=r-\Bigl\lfloor\frac35k\Bigr\rfloor,
\]
and $\frac k5<J<\frac25k$, because $\frac45k<r\le k-1$. As $82\le k\le 109$, this gives $17\le J\le 43$. For $i\ge 17$, Lemma~\ref{lem:frozenlower} and the monotonicity of $\mathcal{T}$ give
\[
G_i\ge e\,c_\ast\,(i+1)^{-3/2}\,\mathcal{T}(i+1)\ge e\,c_\ast\,\mathcal{T}(18)\,(i+1)^{-3/2}\ge 1.12\,e\,(i+1)^{-3/2},
\]
because $\mathcal{T}(18)=\frac{215}{216}\cdot\frac{17}{6}$, $c_\ast\ge 0.39894$ and $0.39894\cdot\frac{215}{216}\cdot\frac{17}{6}=1.1250\ldots$. Let $J'=16\,(J+1)-1$, so that $\sqrt{J'+1}=4\sqrt{J+1}$. Comparing the sum of the $(i+1)^{-3/2}$ with an integral, we find that
\[
\sum_{i=J}^{J'-1}G_i\ge 1.12\,e\int_{J+1}^{J'+1}x^{-3/2}\,dx=\frac{1.68\,e}{\sqrt{J+1}}\ge 0.253\,e,
\]
where the last step holds because $J+1\le 44\le 6.64^2$ and $1.68/6.64=0.2530\ldots$. Since $\sum_{i=0}^{J'-1}G_i\le 2e$ by Lemma~\ref{lem:frozenmass}, subtracting gives
\[
\sum_{i=0}^{\rho-M}G_i=\sum_{i=0}^{J-1}G_i\le(2-0.253)\,e=1.747\,e .
\]

It remains to choose the constants and to verify the numerical inequalities, in all three ranges. We apply the bound \eqref{eq:far-loss} with $t=t_0$ and $N=n$, and for $r>\frac35k$ the condition \eqref{eq:far-four}, with the following constants.
\[
\begin{array}{l|cccc|cccc}
\text{range} & \sigma & L & \lambda & p_1 & S & w & \underline r & p_2\\ \hline
\textup{(a)}\rule{0pt}{2.6ex} & \frac{41}{300} & \frac85 & 1 & 0.99 & & & & \\[2pt]
\textup{(b)} & \frac{4}{15} & \frac95 & 1 & 0.57 & 2 & 0.8945 & 50 & 0.385\\[2pt]
\textup{(c)},\ k\ge 110 & \frac{7}{15} & 2 & 1.011 & 0.395 & 2 & 1 & 89 & 0.56\\[2pt]
\textup{(c)},\ k\le 109 & \frac{7}{15} & 2 & 1.011 & 0.455 & 1.747 & 1 & 66 & 0.5
\end{array}
\]
The three conditions under which \eqref{eq:far-loss} holds are satisfied. For $\lambda=1$ the first holds because $3\sigma\le 1$, and for $\sigma=\frac{7}{15}$ and $\lambda=1.011$ it reads $0.389\le 0.011\cdot\frac{7}{15}\,k$, which is true for $k\ge 76$. The second holds by the choice of $t_0$, and the third because $n=k+r\le Lk$ in each range. The square of the right-hand side of \eqref{eq:far-loss} is $\lambda^2L^2e^{5\sigma}/(8\pi\,\sigma^3k)$, so that this right-hand side is at most $p_1$ as soon as
\[
k\ge\frac{\lambda^2L^2\,e^{5\sigma}}{8\pi\,\sigma^3\,p_1^2} .
\]
A direct evaluation gives for the number on the right the values $80.63\ldots$, $79.37\ldots$, $105.79\ldots$ and $79.73\ldots$ in the four lines of the table, and we have $k\ge 82$ in the first, second and fourth line and $k\ge 110$ in the third. So $\frac n2\,\psi\le p_1$ in each line. In range~(a) this is $\frac n2\,\psi\le 0.99<1$, which, as we have seen, completes the proof in this range.

Let now $r>\frac35k$. In the last three lines of the table $p_1+p_2=0.955<1$, and the first inequality in \eqref{eq:far-closing} has just been shown; so it remains to verify the four conditions \eqref{eq:far-four}, which give the second. The first condition holds with $S=2$ by Lemma~\ref{lem:frozenmass}, and in the last line with $S=1.747$ by the sharper bound shown above for $82\le k\le 109$. The second condition holds in range~(b) because $r\le\frac45k$ and $0.8945^2=0.80013\ldots$, and in range~(c) because $r\le k$. For the third we have $r>\frac35k\ge 49.2$ in range~(b), and in range~(c) we have $r>\frac45k\ge 88$ for $k\ge 110$ and $r>\frac45k\ge 65.6$ for $k\ge 82$; since $r$ is an integer, this gives $r\ge\underline r$ in each line. The last condition is again checked by direct evaluation: its two sides are $2.7056\ldots$ and $2.7520\ldots$ in range~(b), they are $4.2$ and $4.2117\ldots$ in range~(c) for $k\ge 110$, and they are $3.6687$ and $3.6771\ldots$ in range~(c) for $k\le 109$.
This completes the proof.
\end{proof}

\subsection{The single crossing and the chain range}\label{subsec:far-crossing}

We now combine Proposition~\ref{prop:farrows} with Proposition~\ref{prop:nearrows}: a column $k\ge 82$ whose crossing row $k+r_q$ lies in the first third of its chain range stays positive up to row $2k-1$.

\begin{proposition}\label{prop:staypositive}
Let $k\ge 82$ and $r_q\ge 1$ be integers with $3\,r_q\le k$. Assume that $d_{j,k}\le 0$ for every integer $j$ with $k\le j<k+r_q$, and that $d_{k+r_q,k}>0$. Then $d_{n,k}>0$ for every integer $n$ with $k+r_q\le n\le 2k-1$.
\end{proposition}

\begin{proof}
We prove the claim by induction on $n$. For $n=k+r_q$ it is the second assumption. Let $k+r_q<n\le 2k-1$, and suppose that $d_{j,k}>0$ for every $j$ with $k+r_q\le j<n$. Write $n=k+r$ and $r=r_q+t$, so that $t\ge 1$ and $r\le k-1$.

If $100\,r>47\,k$, then Proposition~\ref{prop:farrows} applies to $k$, $r$ and $r_q$, its assumption on rows $k+r_q,\dots,k+r-1$ being the induction hypothesis, and it gives $d_{n,k}>0$.

If $100\,r\le 47\,k$, we apply Proposition~\ref{prop:nearrows} with $r_q$ in the place of $r$, so that its rows $q$ and $n$ are $k+r_q$ and $k+r_q+t=n$. Its hypothesis~(1) is the first assumption, and its hypothesis~(2) is the induction hypothesis. Hence $d_{n,k}>0$ in this case as well, which completes the induction.
\end{proof}

Together with Corollary~\ref{cor:earlypositive}, Proposition~\ref{prop:staypositive} gives the single crossing: along a column $k\ge 82$, from row $k$ to row $2k-1$, the step changes its sign exactly once, from nonpositive to positive.

\begin{proposition}\label{prop:onecrossing}
For every integer $k\ge 82$ there is an integer $q$ such that
\begin{enumerate}
\item $k+1\le q$ and $100\,(q-k)\le 27\,k$;
\item $d_{j,k}\le 0$ for every integer $j$ with $k\le j<q$;
\item $d_{n,k}>0$ for every integer $n$ with $q\le n\le 2k-1$.
\end{enumerate}
\end{proposition}

\begin{proof}
By Corollary~\ref{cor:earlypositive} there is an integer $r\ge 1$ with $100\,r\le 27\,k$ and $d_{k+r,k}>0$. Let $q$ be the least row $j\ge k+1$ with $d_{j,k}>0$. Then $k+1\le q\le k+r$, which gives~(1). For $k+1\le j<q$ we have $d_{j,k}\le 0$ by the choice of $q$, and $d_{k,k}<0$, because $\delta^{(k)}_k=-a_k/k$ by Lemma~\ref{lem:initialstep}; this is~(2). Finally, $r_q=q-k$ satisfies $r_q\ge 1$ and $3\,r_q\le k$ by~(1), so that~(2), $d_{q,k}>0$ and Proposition~\ref{prop:staypositive} give~(3).
\end{proof}

We can now prove Proposition~\ref{prop:chainrange}, for all columns $k\ge 2$.

\begin{proof}[Proof of Proposition~\ref{prop:chainrange}]
Let $k\ge 2$ and $n$ be integers with $k+1\le n\le 2k-2$. In terms of the step $d_{n,k}=b_{n,k}-b_{n,k-1}$, we have to show that $d_{n,k}>0$ implies $d_{n+1,k}>0$. For $2\le k\le 81$ this is part~(2) of Lemma~\ref{lem:columntable}. Let $k\ge 82$, and let $q$ be as in Proposition~\ref{prop:onecrossing}. Suppose that $d_{n,k}>0$. Since $n\ge k$ and $d_{j,k}\le 0$ for $k\le j<q$, we have $q\le n$. Then $q\le n+1\le 2k-1$, and part~(3) of Proposition~\ref{prop:onecrossing} gives $d_{n+1,k}>0$.
\end{proof}

\section{Property K}\label{sec:propk}
This section and Section~\ref{sec:layer} prepare the proof of Proposition~\ref{prop:leftpart} for the columns $k\ge 82$ (Proposition~\ref{prop:largecolumns}). In a row beyond the chain range of such a column only inequalities are available, and into them enter the steps of the column in the rows before its crossing row, which are nonpositive and about whose size nothing is known. The proof in Section~\ref{sec:left} outweighs these steps by means of one inequality, Proposition~\ref{prop:debt}. To describe it, we fix a column $k$ and, as in Section~\ref{sec:critical}, let a hat denote the tilted form of a sequence at the critical tilt $x_k$, so that $\hat c_n=c_n\,x_k^{n}$, $\hat z^{(y)}_N=z^{(y)}_N\,x_k^{N}$ and $\hat\delta_n=\delta^{(k)}_n\,x_k^{n}$; this notation is used throughout the section. The proposition says that, for $k\ge 7$ and $J\ge k$, the sum of the tilted steps $\hat\delta_n$ over rows $n=k,\dots,J$, taken with certain positive weights, is bounded from below by the negative of a quantity that consists of three nonnegative terms belonging to rows $J-1$ and $J$. This quantity is called the weighted debt of the column at row $J$. It is defined, together with the weights, in Section~\ref{sec:layer}, where the proposition is proved.

The proof of Proposition~\ref{prop:debt} is a summation by parts. Recall from Section~\ref{sec:recursion} that the step of a column $k\ge 2$ is a second difference in the threshold, taken at one length; in tilted form,
\[
\hat\delta_n=\hat z^{(k+1)}_n-2\,\hat z^{(k)}_n+\hat z^{(k-1)}_n .
\]
When the tilted steps of rows $k,\dots,J$ are added with their weights, the terms regroup along the diagonals, on which the threshold and the length are raised together: apart from terms at the two ends of the range, the sum consists of the expressions
\[
\hat z^{(k+1)}_{N+1}-2\,\hat z^{(k)}_N+\hat z^{(k-1)}_{N-1}\qquad(k\le N\le J-1),
\]
with each of the three tilted counts multiplied by the weight of its row. The bound therefore rests on the nonnegativity of these expressions, which is a convexity of the tilted counts of thresholds $k-1$, $k$ and $k+1$ along the diagonal (Corollary~\ref{cor:weightedconvexity}). Section~\ref{sec:layer} proves it first for the lengths $k\le N\le 2k-2$, where the three weights are equal, and extends it from there to all larger lengths. For the first of these two steps one would argue by induction on the last-car recursion, but convexity does not pass through that induction. We therefore strengthen it to a property of a triple of real numbers, property~K (Definition~\ref{def:propertyk}), which does, and we prove property~K for the explicit triples that the recursion produces at the three thresholds: the numbers $\hat c_n$ at the top, sums of products of two of them, and the direct term. The induction itself is carried out in Section~\ref{sec:layer}.

\subsection{The definition and the closure under factors}\label{subsec:propk-closure}

The convexity in question compares the tilted counts of thresholds $y=k-1$, $k$ and $k+1$ at lengths $y+\rho$, with the same integer $\rho\ge 0$ for the three thresholds, and it is proved by induction on $\rho$. The induction rests on the last-car recursion in the form
\[
\hat z^{(y)}_{y+\rho}=\sum_{i=0}^{\rho-2}x_k\,\hat c_i\,\frac{y+i+1}{y+\rho}\,\hat z^{(y)}_{y+\rho-1-i}+\mathcal S_\rho(y)\qquad(1\le\rho\le y-1),
\]
which is derived in Section~\ref{sec:layer}; here $\mathcal S_\rho(y)$ is an explicit expression in the $\hat c_n$, the direct term of Section~\ref{subsec:propk-direct}. The tilted counts on the right-hand side belong to smaller values of $\rho$. They are multiplied by $x_k\,\hat c_i\ge 0$, which is the same for the three thresholds, and by $(y+i+1)/(y+\rho)$, which increases with the threshold. Now the convexity $2w_0\le w_-+w_+$ of a triple $(w_-,w_0,w_+)$ passes to sums and to nonnegative multiples, but it does not pass to products with a factor that increases with the threshold: the triple $(2,1,0)$ is convex, and for positive numbers $F_-<F_0<F_+$ the triple $(2F_-,F_0,0)$ is not, since $2F_0>2F_-$. We therefore work with a family of inequalities, one for each real number $\mu\ge 0$, of which convexity is the member $\mu=0$.

\begin{definition}\label{def:propertyk}
Let $k\ne 0$, $w_-$, $w_0$ and $w_+$ be real numbers. The triple $(w_-,w_0,w_+)$ \emph{has property~K at $k$} if
\begin{equation}\label{eq:propk-def}
2\,e^{\mu/k}\,w_0\le e^{\mu}\,w_++e^{-\mu}\,w_-\qquad\text{for every real number }\mu\ge 0.
\end{equation}
\end{definition}

In all our applications $k$ is the index of a column, and $w_-$, $w_0$, $w_+$ are the values of one quantity at thresholds $k-1$, $k$, $k+1$, in this order. For $\mu=0$ the inequality \eqref{eq:propk-def} is the convexity $2w_0\le w_-+w_+$ of the three values. Since the inequality is linear in the triple, property~K at $k$ is preserved when two triples are added entry by entry and when a triple is multiplied by a nonnegative number.

What we prove about factors that depend on the threshold comes from one computation. Let $k>0$, let $(w_-,w_0,w_+)$ have property~K at $k$, and let $L$, $M$ and $\mu$ be real numbers with $\mu+\frac{L-M}{2}\ge 0$. Then
\begin{equation}\label{eq:propk-shift}
2\,e^{\mu/k}\exp\Bigl(\frac{(k+1)L+(k-1)M}{2k}\Bigr)\,w_0\le e^{\mu}\,\bigl(e^{L}w_+\bigr)+e^{-\mu}\,\bigl(e^{M}w_-\bigr).
\end{equation}
Indeed, this is the inequality of property~K for $\mu'=\mu+\frac{L-M}{2}\ge 0$, multiplied by $e^{(L+M)/2}$, since $\frac{L+M}{2}+\frac{\mu'}{k}=\frac{\mu}{k}+\frac{(k+1)L+(k-1)M}{2k}$. Thus multiplying the upper value by $e^{L}$ and the lower value by $e^{M}$ has the same effect on the inequality of property~K as moving $\mu$ to $\mu+\frac{L-M}{2}$ and multiplying the middle value by the second exponential on the left-hand side of the last display. If $L\ge M$, the condition $\mu+\frac{L-M}{2}\ge 0$ holds for every $\mu\ge 0$: a factor that does not decrease with the threshold turns the inequality for $\mu$ into the inequality for a larger value of $\mu$, so that the family as a whole passes through the induction although its single members do not. If $L<M$, the computation is available only for $\mu\ge\frac{M-L}{2}$, because property~K says nothing about negative values of $\mu$.

The factor $e^{\mu/k}$ on the middle value in the definition of property~K is chosen for the factors of the recursion. With it, a positive nondecreasing function $F$ of the threshold is an admissible factor as soon as $y\log F(y)$ is convex at $k-1$, $k$, $k+1$, since the second exponential in \eqref{eq:propk-shift} is then at least $1$; this holds for the functions $(y+p)/(y+q)$ with $0\le p\le q$, as the proof of Lemma~\ref{lem:kratio} shows. Without it the condition would be the log-convexity of $F$, and these functions are log-concave.

The two lemmas that follow apply \eqref{eq:propk-shift} to the factors $(y+p)/(y+q)$ with $0\le p\le q$, which are those of the last-car recursion, and to the factors $(y+s)(y+t)/y$; for $t=0$ the latter are the factors $y+s$, through which Section~\ref{sec:layer} passes from property~K to a convexity with weights. The direct term is not covered by these lemmas, and Proposition~\ref{prop:kdirect} establishes its property~K by a separate estimate.

\begin{lemma}\label{lem:kratio}
Let $k>1$ be a real number, let $(w_-,w_0,w_+)$ be a triple of real numbers that has property~K at $k$ and satisfies $w_0\ge 0$, and let $p$ and $q$ be real numbers with $0\le p\le q$. Then the triple
\[
\Bigl(\frac{k-1+p}{k-1+q}\,w_-,\ \ \frac{k+p}{k+q}\,w_0,\ \ \frac{k+1+p}{k+1+q}\,w_+\Bigr)
\]
has property~K at $k$.
\end{lemma}

\begin{proof}
We first show that, for any positive real numbers $F_-$, $F_0$, $F_+$ with
\begin{equation}\label{eq:propk-factor}
F_-\le F_+\qquad\text{and}\qquad(k+1)\log\frac{F_+}{F_0}+(k-1)\log\frac{F_-}{F_0}\ge 0,
\end{equation}
the triple $(F_-w_-,\,F_0w_0,\,F_+w_+)$ has property~K at $k$. Put $L=\log(F_+/F_0)$ and $M=\log(F_-/F_0)$. Then $L\ge M$, so that \eqref{eq:propk-shift} holds for every $\mu\ge 0$. The second exponential on its left-hand side is at least $1$, because $(k+1)L+(k-1)M\ge 0$ by the second condition, and since $w_0\ge 0$, that side is at least $2e^{\mu/k}w_0$. Multiplying by $F_0$, we obtain the inequality of property~K for the triple $(F_-w_-,F_0w_0,F_+w_+)$.

The second condition in \eqref{eq:propk-factor} follows from a criterion without logarithms: with $u=F_+/F_0-1$ and $v=F_0/F_--1$, it holds as soon as
\[
u\ge 0\qquad\text{and}\qquad(k-1)\,v\le(k+1)\,\frac{2u}{u+2}.
\]
Indeed, the elementary inequalities $\log(1+u)\ge\frac{2u}{u+2}$ for $u\ge 0$ and $\log r\ge 1-\frac1r$ for $r=F_-/F_0$ give, since $k>1$,
\[
(k+1)\log\frac{F_+}{F_0}+(k-1)\log\frac{F_-}{F_0}\ge(k+1)\,\frac{2u}{u+2}-(k-1)\,v\ge 0.
\]

Now let $F(y)=(y+p)/(y+q)$, and put $F_-=F(k-1)$, $F_0=F(k)$ and $F_+=F(k+1)$. These numbers are positive, and $F_-\le F_+$ because $p\le q$. A direct computation gives
\[
u=\frac{F_+}{F_0}-1=\frac{q-p}{(k+p)(k+1+q)}\ge 0,\qquad v=\frac{F_0}{F_-}-1=\frac{q-p}{(k+q)(k-1+p)},
\]
and, clearing the positive denominators, one checks that the inequality $(k-1)\,v\le(k+1)\,\frac{2u}{u+2}$ is equivalent to
\[
(q-p)\,\bigl[p\,(4q+3k+1)+q\,(k-1)\bigr]\ge 0,
\]
which holds because $0\le p\le q$ and $k>1$. Thus the criterion applies, both conditions in \eqref{eq:propk-factor} hold, and the lemma follows.
\end{proof}

\begin{lemma}\label{lem:kproduct}
Let $k>1$ be a real number, let $(w_-,w_0,w_+)$ be a triple of real numbers that has property~K at $k$ and satisfies $w_0\ge 0$, and let $s$ and $t$ be nonnegative real numbers with $st\le k^2-1$. Then the triple
\[
\Bigl(\frac{(k-1+s)(k-1+t)}{k-1}\,w_-,\ \ \frac{(k+s)(k+t)}{k}\,w_0,\ \ \frac{(k+1+s)(k+1+t)}{k+1}\,w_+\Bigr)
\]
has property~K at $k$.
\end{lemma}

\begin{proof}
Let $F(y)=(y+\nobreak s)(y+\nobreak t)/y=y+s+t+st/y$, and put $F_-=F(k-1)$, $F_0=F(k)$ and $F_+=F(k+1)$. These numbers are positive. As was shown in the proof of Lemma~\ref{lem:kratio}, it suffices to verify that $F_-\le F_+$, that $u\ge 0$ and that $(k-1)\,v\le(k+1)\,\frac{2u}{u+2}$, where again $u=F_+/F_0-1$ and $v=F_0/F_--1$.

First, $F_+-F_-=2-\frac{2st}{k^2-1}\ge 0$, because $st\le k^2-1$. Second, a direct computation gives
\begin{gather*}
u=\frac{F_+}{F_0}-1=\frac{k(k+1)-st}{(k+1)(k+s)(k+t)},\\
v=\frac{F_0}{F_-}-1=\frac{k(k-1)-st}{k\,(k-1+s)(k-1+t)},
\end{gather*}
and $u\ge 0$ because $st\le k^2-1\le k(k+1)$. Third, clearing the positive denominators, one checks that the inequality $(k-1)\,v\le(k+1)\,\frac{2u}{u+2}$ is equivalent to the nonnegativity of
\[
k^2(k-1)^2(k+1)+4k^2(k^2-1)\,(s+t)+st\,\bigl(10k^3+2k^2-4k-(3k+1)\,st\bigr).
\]
The first two terms are nonnegative, and so is the third, since $st\ge 0$ and, by $st\le k^2-1$,
\begin{align*}
10k^3+2k^2-4k-(3k+1)\,st&\ge 10k^3+2k^2-4k-(3k+1)(k^2-1)\\
&=7k^3+k^2-k+1>0.
\end{align*}
Thus the three conditions hold, and the lemma follows.
\end{proof}

\subsection{The numbers $\hat c_n$}\label{subsec:propk-cayley}

From here on $k$ is an integer, the index of the fixed column. The numbers $\hat c_n=c_n\,x_k^{n}$ are positive. We first prove the inequality of property~K at $k$ for triples of consecutive ones: for three of them whose indices do not exceed $k$, as the input for Section~\ref{subsec:propk-pairs}, and for the triple $(\hat c_{k-1},\hat c_k,\hat c_{k+1})$ at the top, for which we also bound from below the amount by which the inequality holds.

Both statements are read off from the ratios $\hat c_{n+1}/\hat c_n$. For a real number $t>0$ let
\[
\varphi(t)=(t-1)\,\log\frac{t+1}{t}.
\]
For an integer $n\ge 1$ we have $e^{\varphi(n)}=(1+1/n)^{n-1}$, which by \eqref{eq:prelim-ratio} is the ratio $a_{n+1}/a_n=c_n/c_{n-1}$; and $x_k=e^{-\varphi(k)}$ by the definition of the critical tilt. Hence, for every integer $n\ge 0$,
\begin{equation}\label{eq:propk-ratio}
\hat c_{n+1}=x_k\,\frac{c_{n+1}}{c_n}\,\hat c_n=e^{\varphi(n+1)-\varphi(k)}\,\hat c_n.
\end{equation}
In particular, $\hat c_k=\hat c_{k-1}$. By Lemma~\ref{lem:powermono}, $\varphi$ does not decrease on the positive integers.

We shall need comparisons between the increments of $\varphi$. For these we write the logarithm in the definition of $\varphi$ as an integral,
\begin{equation}\label{eq:propk-logintegral}
\log\frac{t+1}{t}=\int_0^1\frac{dr}{t+r}\qquad(t>0),
\end{equation}
so that a linear combination of such logarithms, taken at several points $t$, is the integral over $0\le r\le 1$ of a rational function of $r$. The first instance is a lower bound for one increment: for every integer $n\ge 2$,
\begin{equation}\label{eq:propk-increment}
\varphi(n)-\varphi(n-1)\ge\frac{3}{2n(n+1)}.
\end{equation}
Indeed, the difference $\varphi(n)-\varphi(n-1)$ is the integral over $0\le r\le 1$ of
\[
\frac{n-1}{n+r}-\frac{n-2}{n-1+r}=\frac{1+r}{(n+r)(n-1+r)}\ge\frac{1+r}{n(n+1)},
\]
and $\int_0^1(1+r)\,dr=\frac32$.

Finally, for $k\ge 2$ put
\[
\Upsilon_k=(k+1)\log\frac{k+1}{k}-(k-1)\log\frac{k}{k-1}.
\]
Then, for all real numbers $\mu$, $\alpha$ and $\beta$,
\begin{equation}\label{eq:propk-amgm}
e^{\mu+\alpha}+e^{-\mu+\beta}\ge 2\exp\Bigl(\frac{\mu}{k}+\frac{(k+1)\,\alpha+(k-1)\,\beta-\Upsilon_k}{2k}\Bigr).
\end{equation}
This is the inequality between the weighted arithmetic and geometric means of $\frac{k}{k+1}\,e^{\mu+\alpha}$ and $\frac{k}{k-1}\,e^{-\mu+\beta}$, with weights $\frac{k+1}{2k}$ and $\frac{k-1}{2k}$.

\begin{lemma}\label{lem:tiltedconvexity}
Let $k\ge 3$ be an integer. Then, for every integer $j\ge 0$ with $j+2\le k$ and every real number $\mu$,
\[
2\,e^{\mu/k}\,\hat c_{j+1}\le e^{\mu}\,\hat c_{j+2}+e^{-\mu}\,\hat c_{j}.
\]
\end{lemma}

\begin{proof}
Put $n=j+1$, so that $1\le n\le k-1$. By the ratio formula \eqref{eq:propk-ratio} we have $\hat c_{j+2}=e^{\alpha}\,\hat c_{j+1}$ and $\hat c_{j}=e^{\beta}\,\hat c_{j+1}$ with
\[
\alpha=\varphi(n+1)-\varphi(k),\qquad\beta=\varphi(k)-\varphi(n).
\]
Since $\hat c_{j+1}>0$, the claim is that $e^{\mu+\alpha}+e^{-\mu+\beta}\ge 2e^{\mu/k}$, and by the inequality \eqref{eq:propk-amgm} between the means it suffices to show that $(k+1)\,\alpha+(k-1)\,\beta\ge\Upsilon_k$. In terms of
\[
\Phi(m)=(k+1)\,\varphi(m+1)-(k-1)\,\varphi(m)-2\,\varphi(k)\qquad(m\ge 1),
\]
this is the inequality
\[
\Phi(n)\ge\Upsilon_k\qquad(1\le n\le k-1).
\]
We show first that $\Phi(m+1)\le\Phi(m)$ for $1\le m\le k-2$, so that $\Phi(n)\ge\Phi(k-1)$.

This rests on a comparison of two consecutive increments of $\varphi$: for every integer $i\ge 2$,
\[
(i+2)\,\bigl(\varphi(i+1)-\varphi(i)\bigr)\le i\,\bigl(\varphi(i)-\varphi(i-1)\bigr).
\]
Indeed, by the definition of $\varphi$ and the integral formula \eqref{eq:propk-logintegral}, the right-hand side minus the left-hand side is the integral over $0\le r\le 1$ of
\[
-\frac{i(i-2)}{i-1+r}+\frac{2(i^2-1)}{i+r}-\frac{i(i+2)}{i+1+r}=\frac{2\,(1-r^2)}{(i-1+r)(i+r)(i+1+r)}\ge 0.
\]
Now let $1\le m\le k-2$. Then
\begin{align*}
\Phi(m)-\Phi(m+1)&=(k-1)\bigl(\varphi(m+1)-\varphi(m)\bigr)-(k+1)\bigl(\varphi(m+2)-\varphi(m+1)\bigr)\\
&\ge\Bigl(k-1-\frac{(k+1)(m+1)}{m+3}\Bigr)\bigl(\varphi(m+1)-\varphi(m)\bigr)\\
&=\frac{2\,(k-m-2)}{m+3}\,\bigl(\varphi(m+1)-\varphi(m)\bigr)\ge 0.
\end{align*}
Here the first inequality is the comparison of increments for $i=m+1$, and the last one holds because $m+2\le k$ and because $\varphi(m+1)-\varphi(m)\ge 0$, as $\varphi$ does not decrease on the positive integers.

It remains to show that $\Phi(k-1)\ge\Upsilon_k$. We note that $\Phi(k-1)=(k-1)\bigl(\varphi(k)-\varphi(k-1)\bigr)$, so that, by the definitions of $\varphi$ and of $\Upsilon_k$,
\[
\Phi(k-1)-\Upsilon_k=(k-3)\,\Bigl(k\log\frac{k+1}{k}-(k-1)\log\frac{k}{k-1}\Bigr).
\]
Here $k-3\ge 0$, and the second factor is nonnegative, because the sequence $(1+1/i)^{i}$ is nondecreasing in $i\ge 1$ (see Section~\ref{sec:prelim}). Hence $\Phi(n)\ge\Phi(k-1)\ge\Upsilon_k$, which proves the lemma.
\end{proof}

The condition $j+2\le k$ excludes the triple at the top, $(\hat c_{k-1},\hat c_k,\hat c_{k+1})$. For this triple the ratio formula \eqref{eq:propk-ratio} gives $\hat c_{k-1}=\hat c_k$ and $\hat c_{k+1}=e^{\lambda_k}\,\hat c_k$, where
\[
\lambda_k=\varphi(k+1)-\varphi(k)\ge\frac{3}{2(k+1)(k+2)}
\]
by the bound \eqref{eq:propk-increment} for one increment. Thus, for every real number $\mu$,
\begin{equation}\label{eq:propk-surplus}
e^{\mu}\,\hat c_{k+1}+e^{-\mu}\,\hat c_{k-1}-2\,e^{\mu/k}\,\hat c_k=\hat c_k\,\bigl(e^{\mu+\lambda_k}+e^{-\mu}-2\,e^{\mu/k}\bigr).
\end{equation}
The left-hand side of this identity is the surplus of the triple at the top, the amount by which it exceeds what property~K at $k$ requires for $\mu$. The proof of Proposition~\ref{prop:kdirect} needs the following lower bound for it.

\begin{lemma}\label{lem:topsurplus}
Let $k\ge 3$ be an integer. Then, for every real number $\mu$,
\[
\hat c_k\,e^{\mu/k}\cdot\frac1k\cdot\frac{k^2-2k-1}{2k(k+1)(k+2)}\le e^{\mu}\,\hat c_{k+1}+e^{-\mu}\,\hat c_{k-1}-2\,e^{\mu/k}\,\hat c_k.
\]
\end{lemma}

\begin{proof}
By the inequality \eqref{eq:propk-amgm} with $\alpha=\lambda_k$ and $\beta=0$, and by the inequality $e^{t}\ge 1+t$,
\[
e^{\mu+\lambda_k}+e^{-\mu}\ge 2\,e^{\mu/k}\exp\Bigl(\frac{(k+1)\,\lambda_k-\Upsilon_k}{2k}\Bigr)\ge 2\,e^{\mu/k}\Bigl(1+\frac{(k+1)\,\lambda_k-\Upsilon_k}{2k}\Bigr).
\]
In view of the identity \eqref{eq:propk-surplus} and of $\hat c_k>0$ it therefore suffices to show that
\[
(k+1)\,\lambda_k-\Upsilon_k\ge\frac{k^2-2k-1}{2k(k+1)(k+2)}.
\]
By the definitions of $\varphi$ and of $\Upsilon_k$ and by the integral formula \eqref{eq:propk-logintegral}, the left-hand side is the integral over $0\le r\le 1$ of
\[
\frac{k-1}{k-1+r}-\frac{k(k+1)}{k+r}+\frac{k(k+1)}{k+1+r}=\frac{r\,\bigl(k^2-2k-1+(k-1)\,r\bigr)}{(k-1+r)(k+r)(k+1+r)}.
\]
Since $k\ge 3$, we have $k^2-2k-1>0$, and for $0\le r\le 1$ the denominator is at most $k(k+1)(k+2)$. So the integrand is at least $\frac{(k^2-2k-1)\,r}{k(k+1)(k+2)}$, and since $\int_0^1r\,dr=\frac12$, the required inequality follows.
\end{proof}

For $k\ge 3$ the left-hand side of the inequality of Lemma~\ref{lem:topsurplus} is nonnegative, so that $2e^{\mu/k}\,\hat c_k\le e^{\mu}\,\hat c_{k+1}+e^{-\mu}\,\hat c_{k-1}$ for every real number $\mu$. Thus the triple $(\hat c_{k-1},\hat c_k,\hat c_{k+1})$ at the top has property~K at $k$ for every $k\ge 3$. By Lemma~\ref{lem:freerange}, it is the triple $\bigl(\hat z^{(k-1)}_{k-1},\hat z^{(k)}_{k},\hat z^{(k+1)}_{k+1}\bigr)$ of the tilted counts of the three thresholds, with which the induction of Section~\ref{sec:layer} starts.

\subsection{Self-convolutions of the numbers $\hat c_n$}\label{subsec:propk-pairs}

For integers $\rho\ge 0$ and $j\ge 0$ let
\[
Q_\rho(j)=\sum_{a=0}^{j}\hat c_{\rho+a}\,\hat c_{\rho+j-a},
\]
a sum of products of two of the numbers $\hat c_n$. The sums of this kind that occur in the direct terms of thresholds $k-1$, $k$ and $k+1$ (Section~\ref{subsec:propk-direct}) are $Q_\rho(n)$, $Q_\rho(n+1)$ and $Q_\rho(n+2)$ with $\rho+n+2=k$, and we show that they form a triple with property~K at $k$.

For a sequence $f=(f(a))_{a\ge 0}$ of real numbers and an integer $j\ge 0$ we write $(f*f)(j)=\sum_{a=0}^{j}f(a)\,f(j-a)$ for the self-convolution of $f$ at length $j$. Thus $Q_\rho(j)$ is the self-convolution at length $j$ of the shifted sequence $a\mapsto\hat c_{\rho+a}$. Under the conditions of the following lemma the inequality of property~K passes from the consecutive values of a sequence to its self-convolutions. In the proof, the expression whose sign is to be shown is sorted by one of the two indices of the convolutions; this produces inner terms, which are nonnegative by the inequality of property~K for consecutive values, and two boundary terms, whose sum may be negative. If the parameter $\mu$ of property~K satisfies $\cosh\mu\ge e^{\mu/k}$, a crude lower bound for the inner terms suffices; otherwise the proof compares the sequence $e^{\mu a}f(a)$ with a solution of a three-term recurrence, whose contribution cancels the boundary terms.

\begin{lemma}\label{lem:kconvolution}
Let $k$ and $n$ be integers with $n\ge 0$ and $n+2\le k$, and let $f=(f(a))_{a\ge 0}$ be a sequence of real numbers that satisfies the following conditions:
\begin{enumerate}
\item $f(a)>0$ for $0\le a\le n+2$;
\item $f(n+1)\le f(a)$ for $0\le a\le n$;
\item $f(n+2)=f(n+1)$;
\item $2\,e^{\mu/k}\,f(j+1)\le e^{\mu}\,f(j+2)+e^{-\mu}\,f(j)$ for every real number $\mu\ge 0$ and every integer $j$ with $0\le j\le n$.
\end{enumerate}
Then the triple
\[
\bigl((f*f)(n),\ (f*f)(n+1),\ (f*f)(n+2)\bigr)
\]
has property~K at $k$.
\end{lemma}

\begin{proof}
Fix a real number $\mu\ge 0$. We have to show that
\begin{equation}\label{eq:propk-convgoal}
e^{\mu}\,(f*f)(n+2)+e^{-\mu}\,(f*f)(n)-2\,e^{\mu/k}\,(f*f)(n+1)\ge 0.
\end{equation}
Sorting the left-hand side by one of the two indices of the convolutions gives the following identity, which holds for every sequence $f$ and all real numbers $A$, $B$ and $D$:
\begin{multline}\label{eq:propk-sorted}
A\,(f*f)(n+2)+B\,(f*f)(n)-D\,(f*f)(n+1)\\
=\sum_{j=0}^{n}f(n-j)\,\bigl(A\,f(j+2)+B\,f(j)-D\,f(j+1)\bigr)\\
+f(n+1)\,\bigl(A\,f(1)-D\,f(0)\bigr)+A\,f(n+2)\,f(0).
\end{multline}
It is obtained by separating the term with $a=0$ of $(f*f)(n+1)$ and the terms with $a=0$ and $a=1$ of $(f*f)(n+2)$, and by shifting the index of summation in the remaining terms. We call the two terms in its last line the boundary terms.

Suppose first that $\cosh\mu\ge e^{\mu/k}$. We use the identity \eqref{eq:propk-sorted} with $A=e^{\mu}$, $B=e^{-\mu}$ and $D=2e^{\mu/k}$. By condition~(4) every bracket in the sum over $j$ is nonnegative, and by condition~(2) we have $f(n-j)\ge f(n+1)$ for $0\le j\le n$. So the sum does not increase when each factor $f(n-j)$ is replaced by $f(n+1)$. Let $S=\sum_{a=0}^{n+1}f(a)$. The sums of $f(j)$, of $f(j+1)$ and of $f(j+2)$ over $0\le j\le n$ are then $S-f(n+1)$, $S-f(0)$ and $S-f(0)-f(1)+f(n+2)$. Inserting them and using condition~(3), we find that the terms that contain $f(0)$ or $f(1)$ cancel against the two boundary terms of the identity, and that the left-hand side of \eqref{eq:propk-convgoal} is at least
\[
2\,f(n+1)\,\Bigl(\bigl(\cosh\mu-e^{\mu/k}\bigr)\,S+\sinh\mu\cdot f(n+1)\Bigr).
\]
This is nonnegative by condition~(1), because $\mu\ge 0$ and $\cosh\mu\ge e^{\mu/k}$, which proves the claim in this case.

Suppose now that $\cosh\mu<e^{\mu/k}$. Then $\mu>0$, and there is a unique real number $\theta>0$ with $\cosh\theta=e^{\mu/k}$; since $\cosh\mu<\cosh\theta$, we have $\mu<\theta$.

Put $h(a)=e^{\mu a}f(a)$ for $a\ge 0$. Then $(h*h)(j)=e^{\mu j}\,(f*f)(j)$ for every $j$, so that the left-hand side of \eqref{eq:propk-convgoal}, multiplied by $e^{\mu(n+1)}$, is
\begin{equation}\label{eq:propk-tilted}
(h*h)(n+2)+(h*h)(n)-2\cosh\theta\cdot(h*h)(n+1),
\end{equation}
and we have to show that this number is nonnegative. Condition~(4) for the present $\mu$, multiplied by $e^{\mu(j+1)}$, says that
\[
r_j=h(j+2)+h(j)-2\cosh\theta\cdot h(j+1)\ge 0\qquad(0\le j\le n).
\]

For an integer $y$ let
\[
H(y)=\sinh\bigl((n+2-y)\,\theta\bigr)\,h(n+1)-\sinh\bigl((n+1-y)\,\theta\bigr)\,h(n+2).
\]
By the addition theorem $\sinh(x+\theta)+\sinh(x-\theta)=2\cosh\theta\,\sinh x$, the sequence $H$ satisfies the recurrence
\[
H(y+1)+H(y-1)=2\cosh\theta\cdot H(y)
\]
for every integer $y$, and $H(n+1)=\sinh\theta\cdot h(n+1)$ and $H(n+2)=\sinh\theta\cdot h(n+2)$. Let
\[
E(j)=\sinh\theta\cdot h(j)-H(j)\qquad(0\le j\le n+2).
\]
Then $E(n+1)=E(n+2)=0$, and by the recurrence of $H$ and since $r_j\ge 0$,
\[
E(j+2)+E(j)-2\cosh\theta\cdot E(j+1)=\sinh\theta\cdot r_j\ge 0\qquad(0\le j\le n).
\]
It follows by downward induction on $j$ that $0\le E(j+1)\le E(j)$ for $0\le j\le n+1$. For $j=n+1$ both numbers are $0$, and if $0\le E(j+2)\le E(j+1)$ for some $0\le j\le n$, then
\[
E(j)\ge 2\cosh\theta\cdot E(j+1)-E(j+2)\ge 2\,E(j+1)-E(j+2)\ge E(j+1)\ge 0,
\]
where we used $\cosh\theta\ge 1$ and $E(j+1)\ge 0$. In particular, $E(j)\ge 0$ for $0\le j\le n$.

We now apply the identity \eqref{eq:propk-sorted} to $h$, with $A=B=1$ and $D=2\cosh\theta$, and multiply by $\sinh\theta>0$. In the term of index $j$ we write $\sinh\theta\cdot h(n-j)=E(n-j)+H(n-j)$. The part $E(n-j)\,r_j$ is nonnegative. For the other part let
\[
Y_j=H(n+1-j)\,h(j+1)-H(n-j)\,h(j)\qquad(0\le j\le n+1).
\]
By the recurrence of $H$ at $y=n-j$ we have $H(n-j)\,r_j=Y_{j+1}-Y_j$, so that $\sum_{j=0}^{n}H(n-j)\,r_j=Y_{n+1}-Y_0$. Moreover, the recurrence at $y=n+1$ gives $H(n)=\sinh\theta\,\bigl(2\cosh\theta\cdot h(n+1)-h(n+2)\bigr)$, and hence
\[
Y_0=\sinh\theta\,\Bigl(h(n+1)\,\bigl(h(1)-2\cosh\theta\cdot h(0)\bigr)+h(n+2)\,h(0)\Bigr).
\]
This is the product of $\sinh\theta$ and the two boundary terms of the identity, which therefore cancel against $-Y_0$. Altogether, the product of $\sinh\theta$ and the number \eqref{eq:propk-tilted} is at least $Y_{n+1}=H(0)\,h(n+2)-H(-1)\,h(n+1)$, and by the definition of $H$,
\begin{multline*}
Y_{n+1}=2\,h(n+1)\,h(n+2)\,\sinh\bigl((n+2)\,\theta\bigr)\\
-h(n+2)^2\,\sinh\bigl((n+1)\,\theta\bigr)-h(n+1)^2\,\sinh\bigl((n+3)\,\theta\bigr).
\end{multline*}

It remains to show that $Y_{n+1}\ge 0$. By condition~(3) we have $h(n+\nobreak 2)=e^{\mu}\,h(n+\nobreak 1)$. With $R=e^{\mu}$ and $m=n+1$ we get $Y_{n+1}=h(n+1)^2\,P$, where
\[
P=2R\,\sinh\bigl((m+1)\,\theta\bigr)-R^2\,\sinh(m\theta)-\sinh\bigl((m+2)\,\theta\bigr).
\]
Put $t=m\theta/2\ge 0$. By the addition theorems,
\[
P=2\,\bigl(\sinh(t+\theta)-R\,\sinh t\bigr)\,\bigl(R\,\cosh t-\cosh(t+\theta)\bigr),
\]
as one sees by multiplying out the right-hand side. The first factor is nonnegative. Indeed, $\mu<\theta$ gives $R\le e^{\theta}=\cosh\theta+\sinh\theta$, and $0\le\sinh t\le\cosh t$, so that
\[
R\,\sinh t\le\sinh t\,\cosh\theta+\cosh t\,\sinh\theta=\sinh(t+\theta).
\]
For the second factor we use the elementary inequality $\sinh x\le x\cosh x$ for $x\ge 0$ at $x=t$ and at $x=\theta$, where $t\theta=m\theta^2/2$; it gives
\[
\cosh(t+\theta)=\cosh t\,\cosh\theta+\sinh t\,\sinh\theta\le\cosh t\,\cosh\theta\,\Bigl(1+\frac{m\theta^2}{2}\Bigr).
\]
By Bernoulli's inequality and since $1+\theta^2/2\le\cosh\theta$,
\[
1+\frac{m\theta^2}{2}\le\Bigl(1+\frac{\theta^2}{2}\Bigr)^{m}\le(\cosh\theta)^{m},
\]
and $(\cosh\theta)^{m+1}\le(\cosh\theta)^{k}=e^{\mu}=R$, because $\cosh\theta\ge 1$ and $m+1=n+2\le k$, and by the choice of $\theta$. Hence $\cosh(t+\theta)\le R\,\cosh t$, so that the second factor is nonnegative as well. Thus $P\ge 0$ and $Y_{n+1}\ge 0$, so that the number \eqref{eq:propk-tilted} is nonnegative. This proves the claim in the second case as well.
\end{proof}

\begin{lemma}\label{lem:kpairsums}
Let $k\ge 3$, $\rho\ge 0$ and $n\ge 0$ be integers with $\rho+n+2=k$. Then the triple
\[
\bigl(Q_\rho(n),\ Q_\rho(n+1),\ Q_\rho(n+2)\bigr)
\]
has property~K at $k$.
\end{lemma}

\begin{proof}
We apply Lemma~\ref{lem:kconvolution} with $f(a)=\hat c_{\rho+a}$, for which $(f*f)(j)=Q_\rho(j)$. Condition~(1) holds because the $\hat c_n$ are positive, and condition~(2) because $f(n+1)=\hat c_{k-1}$ and the $\hat c_n$ do not increase up to index $k$, by \eqref{eq:critical-updown}. Condition~(3) is the equality $\hat c_k=\hat c_{k-1}$, which the ratio formula \eqref{eq:propk-ratio} gives, and condition~(4) is the inequality of Lemma~\ref{lem:tiltedconvexity} at index $\rho+j$, where $\rho+j+2\le k$ for $0\le j\le n$.
\end{proof}

\subsection{The direct term}\label{subsec:propk-direct}

For integers $\rho$ and $y$ with $1\le\rho\le y-1$, the \emph{direct term} is
\[
\mathcal S_\rho(y)=x_k\,\hat c_{\rho-1}\,\hat c_y+\frac{x_k}{2}\cdot\frac{y+\rho+1}{y+\rho}\,Q_\rho(y-\rho-1).
\]
It is the part of $\hat z^{(y)}_{y+\rho}$ that comes from the terms of the last-car recursion in which both factors lie in the free range of threshold $y$, as Section~\ref{sec:layer} shows. The main result of this section, Proposition~\ref{prop:kdirect} below, is that for $k\ge 7$ the direct terms of thresholds $k-1$, $k$ and $k+1$ form a triple with property~K at $k$. All three are formed with one and the same tilt $x_k$, the critical tilt of the middle threshold.

Both ingredients of the direct term have property~K: the triple $(\hat c_{k-1},\hat c_k,\hat c_{k+1})$ by Lemma~\ref{lem:topsurplus}, and the triple of the three sums $Q_\rho(y-\rho-1)$ by Lemma~\ref{lem:kpairsums}. If these sums carried the factor $\frac{x_k}{2}$ alone, the proposition would follow, since property~K is preserved by nonnegative multiples and by sums. But they are multiplied by $\frac{y+\rho+1}{y+\rho}$, which \emph{decreases} as $y$ grows. In terms of the computation \eqref{eq:propk-shift} this is the case $L<M$: the factor moves $\mu$ downward, by an amount of order $(k+\rho)^{-2}$, and for $\mu$ below this amount the family of inequalities of property~K has no member to which one could move. So the second term of $\mathcal S_\rho$ may fail the inequality of property~K, and the amount by which it fails has to be made up by the first term, which satisfies its inequality with the surplus \eqref{eq:propk-surplus}. The proof compares the two amounts for each $\mu\ge 0$, in one way for $\mu\ge\frac{1}{k(k+2)}$, where that computation applies to the sums $Q_\rho$, and in another for smaller $\mu$.

\begin{proposition}\label{prop:kdirect}
Let $k\ge 7$ and $\rho$ be integers with $1\le\rho\le k-2$. Then the triple
\[
\bigl(\mathcal S_\rho(k-1),\ \mathcal S_\rho(k),\ \mathcal S_\rho(k+1)\bigr)
\]
of direct terms has property~K at $k$.
\end{proposition}

\begin{proof}
Let $N=k+\rho$ and $n=k-\rho-2$, so that $k+1\le N\le 2k-2$ and $n\ge 0$. We write
\[
Q_-=Q_\rho(n),\qquad Q_0=Q_\rho(n+1),\qquad Q_+=Q_\rho(n+2)
\]
for the sums $Q_\rho(y-\rho-1)$ with $y=k-1$, $k$, $k+1$, and
\[
F_-=\frac{N}{N-1},\qquad F_0=\frac{N+1}{N},\qquad F_+=\frac{N+2}{N+1}
\]
for the values of $\frac{y+\rho+1}{y+\rho}$ at $y=k-1$, $k$, $k+1$. The numbers $Q_-$, $Q_0$ and $Q_+$ are positive, and since $\rho+n+2=k$, the triple $(Q_-,Q_0,Q_+)$ has property~K at $k$ by Lemma~\ref{lem:kpairsums}. In this notation $\mathcal S_\rho(k-1)=x_k\,\hat c_{\rho-1}\,\hat c_{k-1}+\frac{x_k}{2}\,F_-\,Q_-$, and similarly for thresholds $k$ and $k+1$. Fix a real number $\mu\ge 0$, and put
\begin{align*}
\mathrm{Sur}&=e^{\mu}\,\hat c_{k+1}+e^{-\mu}\,\hat c_{k-1}-2\,e^{\mu/k}\,\hat c_k,\\
\mathrm{Def}&=e^{\mu}\,F_+\,Q_++e^{-\mu}\,F_-\,Q_--2\,e^{\mu/k}\,F_0\,Q_0.
\end{align*}
Thus $\mathrm{Sur}$ is the surplus \eqref{eq:propk-surplus} of the triple at the top, and $\mathrm{Def}$ is the corresponding quantity for the second terms of the $\mathcal S_\rho(y)$; it may be negative, and $-\mathrm{Def}$ is then the defect of these terms, the amount by which they fail the inequality of property~K. We have
\[
e^{\mu}\,\mathcal S_\rho(k+1)+e^{-\mu}\,\mathcal S_\rho(k-1)-2\,e^{\mu/k}\,\mathcal S_\rho(k)=x_k\,\hat c_{\rho-1}\,\mathrm{Sur}+\frac{x_k}{2}\,\mathrm{Def},
\]
and we have to show that the right-hand side is nonnegative.

We first reduce the claim to a lower bound for each of $\mathrm{Sur}$ and $\mathrm{Def}$. Suppose that there is a real number $\tau\ge 0$ such that
\begin{equation}\label{eq:propk-tau}
\mathrm{Sur}\ge\hat c_k\,\tau\qquad\text{and}\qquad\mathrm{Def}\ge-2\,\hat c_\rho\,\hat c_k\,\tau.
\end{equation}
Then $x_k\,\hat c_{\rho-1}\,\mathrm{Sur}+\frac{x_k}{2}\,\mathrm{Def}$ is at least $x_k\,\hat c_k\,\tau\,(\hat c_{\rho-1}-\hat c_\rho)$, and this is nonnegative, because $\hat c_\rho\le\hat c_{\rho-1}$ by \eqref{eq:critical-updown}. So it suffices to find such a $\tau$. We do this separately for $\mu\le\mu_0$ and for $\mu\ge\mu_0$, where $\mu_0=\frac{1}{k(k+2)}$.

Consider first the case $\mu\le\mu_0$. Here we take $\tau=\frac{e^{\mu}}{N(N+1)}$.

We begin with $\mathrm{Def}$. We have $F_0-F_+=\frac{1}{N(N+1)}$ and $F_--F_0=\frac{1}{N(N-1)}\ge\frac{1}{N(N+1)}$. In the decomposition
\[
\mathrm{Def}=F_0\,\bigl(e^{\mu}Q_++e^{-\mu}Q_--2\,e^{\mu/k}Q_0\bigr)-(F_0-F_+)\,e^{\mu}Q_++(F_--F_0)\,e^{-\mu}Q_-
\]
the first term is nonnegative by property~K of the triple $(Q_-,Q_0,Q_+)$, the second is $-\frac{1}{N(N+1)}\,e^{\mu}Q_+$, and the third is at least $\frac{1}{N(N+1)}\,e^{-\mu}Q_-$. Hence
\[
\mathrm{Def}\ge-\frac{1}{N(N+1)}\,\bigl(e^{\mu}\,Q_+-e^{-\mu}\,Q_-\bigr).
\]
The sum $Q_+$ has two more terms than $Q_-$, and we claim that
\[
e^{\mu}\,Q_+-e^{-\mu}\,Q_-\le 2\,e^{\mu}\,\hat c_\rho\,\hat c_k.
\]
For the proof we note first that for $\mu\le\mu_0$ the numbers $\hat c_i$ below the top fall at each step by a factor of at least $e^{\mu}$:
\[
e^{\mu}\,\hat c_{i+1}\le\hat c_i\qquad(0\le i\le k-2).
\]
Indeed, $\hat c_i=e^{\varphi(k)-\varphi(i+1)}\,\hat c_{i+1}$ by the ratio formula \eqref{eq:propk-ratio}, and
\[
\varphi(k)-\varphi(i+1)\ge\varphi(k)-\varphi(k-1)\ge\frac{3}{2k(k+1)}\ge\frac{1}{k(k+2)}\ge\mu;
\]
the first inequality holds because $i+1\le k-1$ and $\varphi$ does not decrease on the positive integers, and the second is the bound \eqref{eq:propk-increment} for $n=k$. Now in $Q_+=\sum_{a=0}^{n+2}\hat c_{\rho+a}\,\hat c_{\rho+n+2-a}$ the terms with $a=0$ and $a=n+2$ are both equal to $\hat c_\rho\,\hat c_k$, because $\rho+n+2=k$. The other terms are $\hat c_{\rho+b+1}\,\hat c_{\rho+(n-b)+1}$ for $0\le b\le n$. Since $\rho+b$ and $\rho+(n-b)$ are at most $k-2$, the inequality $e^{\mu}\,\hat c_{i+1}\le\hat c_i$ applies to both factors and gives
\[
e^{\mu}\,\hat c_{\rho+b+1}\,\hat c_{\rho+(n-b)+1}\le e^{-\mu}\,\hat c_{\rho+b}\,\hat c_{\rho+(n-b)}.
\]
The right-hand side is $e^{-\mu}$ times the term with $a=b$ of $Q_-$. Summing over $0\le b\le n$, we obtain $e^{\mu}\,\bigl(Q_+-2\,\hat c_\rho\,\hat c_k\bigr)\le e^{-\mu}\,Q_-$, which is the claim. Inserted into the lower bound for $\mathrm{Def}$ obtained above, it gives $\mathrm{Def}\ge-2\,\hat c_\rho\,\hat c_k\,\tau$, since $\tau=\frac{e^{\mu}}{N(N+1)}$.

We turn to $\mathrm{Sur}$. For these $\mu$ the bound of Lemma~\ref{lem:topsurplus} does not suffice: at $\mu=0$ it is about $\hat c_k/\bigl(2(k+1)(k+2)\bigr)$, whereas $\hat c_k\,\tau$ is then $\hat c_k/\bigl((k+1)(k+2)\bigr)$ for $\rho=1$. We use another bound. Recall that $\mathrm{Sur}$ equals $\hat c_k$ times $e^{\mu+\lambda_k}+e^{-\mu}-2\,e^{\mu/k}$. By the inequality $e^{t}\ge 1+t$ for $t=\mu+\lambda_k$ and for $t=-\mu$, and by $\lambda_k\ge\frac{3}{2(k+1)(k+2)}$,
\[
\mathrm{Sur}\ge\hat c_k\,\bigl(2+\lambda_k-2\,e^{\mu/k}\bigr)\ge\hat c_k\,\Bigl(\frac{3}{2(k+1)(k+2)}-2\,\bigl(e^{\mu/k}-1\bigr)\Bigr).
\]
It remains to show that the bracket on the right is at least $\tau$. Put $U=\frac{1}{(k+1)(k+2)}$. We claim that
\[
\tau\le\frac{63}{62}\,U\qquad\text{and}\qquad e^{\mu/k}-1\le\frac{9}{55}\,U;
\]
that bracket is then at least $\bigl(\frac32-\frac{18}{55}\bigr)\,U=\frac{129}{110}\,U$, which is larger than $\frac{63}{62}\,U\ge\tau$. For the first bound, $\frac{1}{N(N+1)}\le U$ because $N\ge k+1$, and $e^{\mu}\le\frac{1}{1-\mu}\le\frac{63}{62}$ because $\mu\le\mu_0\le\frac1{63}$ for $k\ge 7$. For the second, $k\ge 7$ gives $\frac{\mu}{k}\le\frac{1}{k^2(k+2)}=\frac{k+1}{k^2}\,U\le\frac{8}{49}\,U$ and $\frac{\mu}{k}\le\frac{1}{441}$, so that $e^{\mu/k}-1\le\frac{\mu/k}{1-\mu/k}\le\frac{441}{440}\cdot\frac{8}{49}\,U=\frac{9}{55}\,U$. Thus $\mathrm{Sur}\ge\hat c_k\,\tau$, and both bounds \eqref{eq:propk-tau} hold in the case $\mu\le\mu_0$.

Now let $\mu\ge\mu_0$. In this case we take
\[
\tau=e^{\mu/k}\cdot\frac{N+1}{N}\cdot(2k-N)\cdot\frac{N-k+1}{N^2(N+2)\,k},
\]
which is nonnegative because $k+1\le N\le 2k-2$.

We begin again with $\mathrm{Def}$. Let $L=\log(F_+/F_0)$ and $M=\log(F_-/F_0)$, where
\[
\frac{F_+}{F_0}=1-\frac{1}{(N+1)^2}\qquad\text{and}\qquad\frac{F_-}{F_0}=1+\frac{1}{N^2-1}.
\]
The inequalities $1-\frac1r\le\log r\le r-1$ for $r>0$ give
\[
L\ge-\frac{1}{N(N+2)},\qquad\frac{1}{N^2}\le M\le\frac{1}{(N-1)(N+1)}.
\]
Hence $\frac{M-L}{2}\le\frac{1}{(N-1)(N+1)}\le\frac{1}{k(k+2)}\le\mu$, where we used $N(N+2)\ge(N-1)(N+1)$ and $N\ge k+1$. So $\mu+\frac{L-M}{2}\ge 0$, and the inequality \eqref{eq:propk-shift} applies to the triple $(Q_-,Q_0,Q_+)$. Multiplied by $F_0$, it reads
\[
2\,e^{\mu/k}\exp\Bigl(\frac{(k+1)L+(k-1)M}{2k}\Bigr)\,F_0\,Q_0\le e^{\mu}\,F_+\,Q_++e^{-\mu}\,F_-\,Q_-.
\]
By the bounds for $L$ and $M$,
\[
(k+1)\,L+(k-1)\,M\ge-\frac{k+1}{N(N+2)}+\frac{k-1}{N^2}=-\frac{2\,(N-k+1)}{N^2(N+2)},
\]
so that the exponential is at least $1-\frac{N-k+1}{N^2(N+2)\,k}$, by $e^{t}\ge 1+t$. Since $Q_0>0$, it follows that
\[
\mathrm{Def}\ge-2\,e^{\mu/k}\,F_0\cdot\frac{N-k+1}{N^2(N+2)\,k}\cdot Q_0.
\]
Finally, the sum $Q_0=\sum_{b=0}^{n+1}\hat c_{\rho+b}\,\hat c_{\rho+n+1-b}$ has $n+2=2k-N$ terms, and each of them is at most $\hat c_\rho\,\hat c_{\rho+n+1}=\hat c_\rho\,\hat c_{k-1}=\hat c_\rho\,\hat c_k$. Indeed, since $c_m=a_{m+1}$, Lemma~\ref{lem:cayleylogconvex} with $u=\rho+1$, $v=\rho+n+2-b$ and $s=b$ says that $c_{\rho+b}\,c_{\rho+n+1-b}\le c_\rho\,c_{\rho+n+1}$, and the tilt multiplies both sides by the same power of $x_k$. Thus $Q_0\le(2k-N)\,\hat c_\rho\,\hat c_k$, and we obtain $\mathrm{Def}\ge-2\,\hat c_\rho\,\hat c_k\,\tau$.

For $\mathrm{Sur}$ we use Lemma~\ref{lem:topsurplus}, by which
\[
\mathrm{Sur}\ge\hat c_k\,e^{\mu/k}\cdot\frac1k\cdot\frac{k^2-2k-1}{2k(k+1)(k+2)}.
\]
After division by $\hat c_k\,e^{\mu/k}/k$, the inequality $\mathrm{Sur}\ge\hat c_k\,\tau$ therefore follows from
\[
\frac{(N+1)\,(2k-N)\,(N-k+1)}{N^3\,(N+2)}\le\frac{k^2-2k-1}{2k(k+1)(k+2)}.
\]
To prove this inequality, note that $2k-N$ and $N-k+1$ are positive and add up to $k+1$, so that $4\,(2k-N)(N-k+1)\le(k+1)^2$; and that $(N+1)\,(k+1)^3\le(N+2)\,N^3$, because $N\ge k+1$. Hence the left-hand side of the displayed inequality is at most
\[
\frac{(N+1)\,(k+1)^2}{4\,N^3\,(N+2)}\le\frac{1}{4(k+1)}.
\]
Finally, $\frac{1}{4(k+1)}\le\frac{k^2-2k-1}{2k(k+1)(k+2)}$, because this inequality is equivalent to $k^2-6k-2\ge 0$, which holds since $k\ge 7$. Thus both bounds hold in this case as well.

We have thus found, for every $\mu\ge 0$, a real number $\tau\ge 0$ with the two bounds \eqref{eq:propk-tau}. By the reduction at the beginning of the proof, $x_k\,\hat c_{\rho-1}\,\mathrm{Sur}+\frac{x_k}{2}\,\mathrm{Def}$ is therefore nonnegative for every $\mu\ge 0$, that is, the triple of the direct terms has property~K at $k$.
\end{proof}

\section{Convexity along the diagonal and the weighted debt}\label{sec:layer}
With property~K available for the triples of Section~\ref{sec:propk}, an induction along the last-car recursion gives property~K, and with it convexity along the diagonal, for the tilted counts of thresholds $k-1$, $k$ and $k+1$ in the first layer; a second induction extends the convexity, in a weighted form, to all lengths (Corollary~\ref{cor:weightedconvexity}). Summation by parts then yields the bound of the weighted debt, Proposition~\ref{prop:debt}: for a column $k\ge 7$, a weighted sum of the tilted steps of the column over rows $k,\dots,J$ is bounded from below by three terms that belong to rows $J-1$ and $J$. In the proof of Proposition~\ref{prop:leftpart} (Section~\ref{sec:left}) this is the bound by which the rows before the crossing are outweighed. As in Section~\ref{sec:critical}, a hat denotes the tilted form of a sequence at the critical tilt $x_k$ of the column under consideration; thus $\hat z^{(y)}_N=z^{(y)}_N\,x_k^{N}$, $\hat c_n=c_n\,x_k^{n}$, $\hat\beta^{(\ell)}_N=\beta^{(\ell)}_N\,x_k^{N}$ and $\hat\delta_n=\delta^{(k)}_n\,x_k^{n}$.

\subsection{The last-car recursion in the first layer of a threshold}\label{subsec:layer-split}

Let $k\ge 1$ be an integer and let $y\ge 1$ be a threshold. By \eqref{eq:recursion-tilted} with $x=x_k$, the tilted counts of threshold $y$ satisfy
\begin{equation}\label{eq:layer-recursion}
(N+1)\,\hat z^{(y)}_{N+1}=x_k\sum_{j=0}^{N}\min(j+1,y)\,\hat z^{(y)}_j\,\hat z^{(y)}_{N-j}
\end{equation}
for every integer $N\ge 0$. Recall from Section~\ref{sec:recursion} that the lengths $n\le y$ form the free range of threshold $y$ and the lengths $y+1\le n\le 2y+1$ its first layer. A factor $\hat z^{(y)}_i$ whose length $i$ lies in the free range is one of the numbers $\hat c_i$: by Lemma~\ref{lem:freerange},
\[
\hat z^{(y)}_i=\hat c_i\qquad(0\le i\le y).
\]

A length of the first layer is of the form $y+P+1$ with $0\le P\le y$. In the recursion with $N=y+P$ the lengths of the two factors add up to $y+P\le 2y$, so that at most one of the two factors lies outside the free range. Let
\[
\mathcal D(y,P)=\sum_{a=0}^{y-P}\min(P+a+1,y)\,\hat c_{P+a}\,\hat c_{y-a}
\]
be the sum of the terms of the recursion at length $y+P+1$ in which both factors lie in the free range. Sorting the other terms brings the recursion into the form that supports an induction on the amount by which the length exceeds the threshold.

\begin{lemma}\label{lem:layersplit}
Let $k\ge 1$, $y\ge 1$ and $P$ be integers with $0\le P\le y$. Then
\[
(y+P+1)\,\hat z^{(y)}_{y+P+1}=x_k\Bigl(\,\sum_{i=0}^{P-1}(y+i+1)\,\hat c_i\,\hat z^{(y)}_{y+P-i}+\mathcal D(y,P)\Bigr).
\]
\end{lemma}

\begin{proof}
We split the sum in the recursion \eqref{eq:layer-recursion} with $N=y+P$ into the three ranges $0\le j<P$, $P\le j\le y$ and $y<j\le y+P$. In the first range we write $j=i$; then $\min(j+1,y)=i+1$, and the first factor is $\hat c_i$, because its length $i$ lies in the free range. In the third range we write $j=y+P-i$ with $0\le i<P$; then $\min(j+1,y)=y$, and the second factor is $\hat c_i$. For each $i$ the two terms contain the same product, so that together the two ranges give
\[
\sum_{i=0}^{P-1}\bigl((i+1)+y\bigr)\,\hat c_i\,\hat z^{(y)}_{y+P-i}.
\]
In the second range we write $j=P+a$ with $0\le a\le y-P$. Both factors then lie in the free range, and the terms of this range are those of $\mathcal D(y,P)$.
\end{proof}

\subsection{Property K and convexity for $k\le N\le 2k-2$}\label{subsec:layer-first}

Fix a column $k\ge 2$. For an integer $\rho\ge 0$ the three tilted counts
\[
\hat z^{(k-1)}_{k-1+\rho},\qquad \hat z^{(k)}_{k+\rho},\qquad \hat z^{(k+1)}_{k+1+\rho}
\]
lie on a diagonal: from one to the next the threshold and the length are both raised by one, so that each length exceeds its threshold by $\rho$. All three are tilted by the same number $x_k$, the critical tilt of the middle threshold. We show by induction on $\rho$, along the recursion of Lemma~\ref{lem:layersplit}, that they form a triple with property~K (Definition~\ref{def:propertyk}).

\begin{proposition}\label{prop:klayer}
Let $k\ge 7$ and $\rho$ be integers with $0\le\rho\le k-2$. Then the triple
\[
\bigl(\hat z^{(k-1)}_{k-1+\rho},\ \hat z^{(k)}_{k+\rho},\ \hat z^{(k+1)}_{k+1+\rho}\bigr)
\]
has property~K at $k$.
\end{proposition}

\begin{proof}
Before the induction on $\rho$ we bring the recursion of Lemma~\ref{lem:layersplit}, for one threshold $y$ at a time, into a form in which $\mathcal D(y,P)$ is replaced by the direct term of Section~\ref{sec:propk}. Recall from Section~\ref{sec:propk} that $Q_s(j)=\sum_{a=0}^{j}\hat c_{s+a}\,\hat c_{s+j-a}$ and that the direct term is
\[
\mathcal S_\rho(y)=x_k\,\hat c_{\rho-1}\,\hat c_y+\frac{x_k}{2}\cdot\frac{y+\rho+1}{y+\rho}\,Q_\rho(y-\rho-1).
\]

We first evaluate $\mathcal D(y,P)$. For integers $y$ and $P$ with $P\ge 0$ and $P+2\le y$, a direct computation with the sum that defines $\mathcal D(y,P)$ in Section~\ref{subsec:layer-split} gives
\[
\mathcal D(y,P)=(y+P+1)\,\hat c_P\,\hat c_y+\frac{y+P+2}{2}\,Q_{P+1}(y-P-2).
\]
Here the first term on the right comes from the terms with $a=0$ and $a=y-P$, which are distinct because $P+2\le y$ and have the weights $P+1$ and $y$. The products of the other terms add up to $Q_{P+1}(y-P-2)$, and their weights $P+a+1$ may be replaced by $\frac12(y+P+2)$, the mean of the weights of the two terms with indices $a$ and $y-P-a$, which contain the same product.

The recursion of Lemma~\ref{lem:layersplit}, divided by the length, now takes the form announced above: for integers $y$ and $\rho$ with $1\le\rho\le y-1$,
\begin{equation}\label{eq:layer-divided}
\hat z^{(y)}_{y+\rho}=\sum_{i=0}^{\rho-2}x_k\,\hat c_i\,\frac{y+i+1}{y+\rho}\,\hat z^{(y)}_{y+\rho-1-i}+\mathcal S_\rho(y).
\end{equation}
Indeed, in the identity of Lemma~\ref{lem:layersplit} with $P=\rho-1$ we insert the value of $\mathcal D(y,\rho-1)$ just obtained, which applies because $\rho+1\le y$, and divide by $y+\rho$; by the definition of $\mathcal S_\rho(y)$ the result is the stated identity.

We turn to the induction on $\rho$. Recall that property~K at $k$ is preserved when triples are added entry by entry and when a triple is multiplied by a nonnegative number.

Let $\rho=0$. The three lengths lie in the free range of their thresholds, so that the triple is $(\hat c_{k-1},\hat c_k,\hat c_{k+1})$ by Lemma~\ref{lem:freerange}, and it has property~K at $k$ by Lemma~\ref{lem:topsurplus}, whose left-hand side is nonnegative for $k\ge 3$.

Let now $1\le\rho\le k-2$, and suppose that the proposition holds for every smaller value of $\rho$. Since $\rho\le k-2$, the recursion \eqref{eq:layer-divided} holds for $y=k-1$, $k$ and $k+1$, with the same range of $i$ and the same factors $x_k\,\hat c_i$. Hence the triple of the proposition is the sum of the following triples: the triple $\bigl(\mathcal S_\rho(k-1),\ \mathcal S_\rho(k),\ \mathcal S_\rho(k+1)\bigr)$ of the direct terms; and, for each $i$ with $0\le i\le\rho-2$, the triple with the entries
\[
x_k\,\hat c_i\,\frac{y+i+1}{y+\rho}\,\hat z^{(y)}_{y+\rho'}\qquad(y=k-1,\ k,\ k+1),
\]
where $\rho'=\rho-1-i$. It suffices to show that each of these triples has property~K at $k$.

The first of them has property~K at $k$ by Proposition~\ref{prop:kdirect}, since $k\ge 7$ and $1\le\rho\le k-2$. For the others, fix $i$ with $0\le i\le\rho-2$. Then $1\le\rho'\le\rho-1$, so that by the induction hypothesis the numbers $\hat z^{(k-1)}_{k-1+\rho'}$, $\hat z^{(k)}_{k+\rho'}$ and $\hat z^{(k+1)}_{k+1+\rho'}$ form a triple with property~K at $k$, and its middle entry is nonnegative. By Lemma~\ref{lem:kratio} with $p=i+1$ and $q=\rho$, the triple keeps property~K at $k$ when its entries are multiplied by the values of $(y+i+1)/(y+\rho)$ at $y=k-1$, $k$ and $k+1$. Multiplying further by $x_k\,\hat c_i>0$, we obtain the triple that belongs to $i$ in the sum above. This completes the induction.
\end{proof}

We record two consequences of the proposition: the convexity of the tilted counts along the diagonal, and the form of it that enters the induction of the next subsection, in which the three counts carry weights that increase by one from each to the next.

\begin{corollary}\label{cor:layerconvex}
Let $k\ge 7$ and $N$ be integers with $k\le N\le 2k-2$. Then
\[
\hat z^{(k+1)}_{N+1}-2\,\hat z^{(k)}_N+\hat z^{(k-1)}_{N-1}\ge 0.
\]
\end{corollary}

\begin{proof}
This is the case $\mu=0$ of Proposition~\ref{prop:klayer} with $\rho=N-k$.
\end{proof}

In terms of the counts, after division by $x_k^{N-1}$ and multiplication by $(N+1)!$, the corollary says that
\[
x_k^{2}\,A_{N+1,k}-2\,(N+1)\,x_k\,A_{N,k-1}+N(N+1)\,A_{N-1,k-2}\ge 0
\]
for $k\ge 7$ and $k\le N\le 2k-2$.

For the second consequence we write, for a column $k\ge 2$, a real number $\lambda$ and an integer $N\ge 1$,
\begin{equation}\label{eq:layer-weighted}
\Delta_\lambda(N)=(\lambda+1)\,\hat z^{(k+1)}_{N+1}-2\lambda\,\hat z^{(k)}_N+(\lambda-1)\,\hat z^{(k-1)}_{N-1}.
\end{equation}

\begin{corollary}\label{cor:layerweighted}
Let $k\ge 7$ and $N$ be integers with $k\le N\le 2k-2$, and let $\lambda\ge k$ be a real number. Then
\[
(\lambda+1)\,\hat z^{(k+1)}_{N+1}-2\lambda\,\hat z^{(k)}_N+(\lambda-1)\,\hat z^{(k-1)}_{N-1}\ge 0.
\]
\end{corollary}

\begin{proof}
Apply Lemma~\ref{lem:kproduct} with $s=\lambda-k\ge 0$ and $t=0$ to the triple of Proposition~\ref{prop:klayer} with $\rho=N-k$, whose middle entry $\hat z^{(k)}_N$ is nonnegative: the lemma multiplies the three entries by $\lambda-1$, $\lambda$ and $\lambda+1$, and the case $\mu=0$ of property~K for the resulting triple is the assertion.
\end{proof}

\subsection{Weighted convexity for $N\ge 2k-1$}\label{subsec:layer-beyond}

For $N\ge 2k-1$ we prove the convexity along the diagonal in the form that the recursion \eqref{eq:layer-recursion} produces, in which each tilted count is multiplied by its length. We first collect the facts about the tilted counts of the three thresholds of a column $k\ge 2$ that the remaining proofs of this section use.

Recall from Section~\ref{sec:recursion} that $0\le z^{(y)}_i\le z^{(y')}_i$ for thresholds $y\le y'$ and that $\beta^{(\ell)}_i=z^{(\ell+1)}_i-z^{(\ell)}_i$ for $\ell\ge 1$. Since $x_k>0$, the tilted counts of the three thresholds at one length are ordered,
\[
0\le\hat z^{(k-1)}_i\le\hat z^{(k)}_i\le\hat z^{(k+1)}_i\qquad(i\ge 0),
\]
and their differences are the tilted counts of levels $k-1$ and $k$,
\[
\hat\beta^{(k-1)}_i=\hat z^{(k)}_i-\hat z^{(k-1)}_i,\qquad \hat\beta^{(k)}_i=\hat z^{(k+1)}_i-\hat z^{(k)}_i\qquad(i\ge 0),
\]
which are nonnegative. Since $\delta^{(k)}_i=\beta^{(k)}_i-\beta^{(k-1)}_i$, the tilted step of the column is
\[
\hat\delta_i=\hat\beta^{(k)}_i-\hat\beta^{(k-1)}_i=\hat z^{(k+1)}_i-2\,\hat z^{(k)}_i+\hat z^{(k-1)}_i\qquad(i\ge 0).
\]
By Lemma~\ref{lem:freerange} the tilted counts of the three thresholds agree as long as the length lies in the free range of all three:
\[
\hat z^{(k-1)}_i=\hat z^{(k)}_i=\hat z^{(k+1)}_i=\hat c_i\qquad(0\le i\le k-1).
\]

Next, the tilted counts do not decrease along a diagonal from length $k-1$ on: for every threshold $y\ge 1$,
\begin{equation}\label{eq:layer-diagonal}
\hat z^{(y)}_i\le\hat z^{(y+1)}_{i+1}\qquad(i\ge k-1).
\end{equation}
Indeed, $\hat z^{(y)}_i=(z^{(y)}_i/c_i)\,\hat c_i$, and here $z^{(y)}_i/c_i\le z^{(y+1)}_{i+1}/c_{i+1}$ by Proposition~\ref{prop:diagonalmonotone} and $\hat c_i\le\hat c_{i+1}$ by \eqref{eq:critical-updown}, since $i+1\ge k$.

Finally, we need three inequalities between the numbers $\hat c_n$ near index $k$:
\begin{equation}\label{eq:layer-cayley}
\hat c_{k-2}\ge\hat c_{k-1},\qquad \hat c_{k-1}\le\hat c_k,\qquad (k-1)\,\hat c_{k-2}\le k\,\hat c_{k-1}.
\end{equation}
The first two are instances of \eqref{eq:critical-updown}. For the third, the identity \eqref{eq:critical-ratio} with $x=x_k$ and $j=k-2$ gives $\hat c_{k-1}=x_k\bigl(1+\frac{1}{k-1}\bigr)^{k-2}\hat c_{k-2}$, and hence
\begin{align*}
k\,\hat c_{k-1}&=(k-1)\,x_k\Bigl(1+\frac{1}{k-1}\Bigr)^{k-1}\hat c_{k-2}\\
&\ge(k-1)\,x_k\Bigl(1+\frac{1}{k}\Bigr)^{k-1}\hat c_{k-2}=(k-1)\,\hat c_{k-2},
\end{align*}
where the last step is the equation $x_k\,(1+1/k)^{k-1}=1$ that defines the critical tilt.

The following proposition is proved by induction on $N$. By the recursion \eqref{eq:layer-recursion} for the three thresholds, the left-hand side of its inequality, divided by $x_k$, is a combination of products of two tilted counts of smaller lengths. In the proof, the products in which one factor has length at most $k-2$ combine to numbers $\Delta_\lambda(M)$ with $M<N$, which Corollary~\ref{cor:layerweighted} and the induction hypothesis show to be nonnegative, and the other products are controlled by Corollary~\ref{cor:diagonalcritical}.

\begin{proposition}\label{prop:beyondlayer}
Let $k\ge 7$ and $N$ be integers with $N\ge 2k-1$. Then
\[
(N+1)\,\hat z^{(k+1)}_{N+1}-2N\,\hat z^{(k)}_N+(N-1)\,\hat z^{(k-1)}_{N-1}\ge 0.
\]
\end{proposition}

\begin{proof}
In the notation \eqref{eq:layer-weighted} the assertion reads $\Delta_N(N)\ge 0$, and we prove it by induction on $N$.

Let $N\ge 2k-1$, and suppose that $\Delta_M(M)\ge 0$ for every integer $M$ with $2k-1\le M<N$. We write $N=2k-1+r$ with $r\ge 0$.

We first show that the weighted second differences at smaller lengths are nonnegative: for integers $M$ and real numbers $\lambda$,
\begin{equation}\label{eq:layer-smaller}
\Delta_\lambda(M)\ge 0\qquad(k\le M\le N-1,\ k\le\lambda\le 2k-1),
\end{equation}
and moreover $\Delta_k(k-1)\ge 0$.

For $M\le 2k-2$ the first of these inequalities is Corollary~\ref{cor:layerweighted}. Let $2k-1\le M\le N-1$. Comparing the coefficients of the three tilted counts in the definition of $\Delta_\lambda(M)$, one sees that
\[
M\,\Delta_\lambda(M)=\lambda\,\Delta_M(M)+(M-\lambda)\,\bigl(\hat z^{(k+1)}_{M+1}-\hat z^{(k-1)}_{M-1}\bigr).
\]
Here $\Delta_M(M)\ge 0$ by the induction hypothesis, $0\le\lambda\le M$, and $\hat z^{(k-1)}_{M-1}\le\hat z^{(k)}_M\le\hat z^{(k+1)}_{M+1}$ by the monotonicity \eqref{eq:layer-diagonal} along the diagonal. Hence $\Delta_\lambda(M)\ge 0$. As for $\Delta_k(k-1)$, the lengths of its three tilted counts lie in the free range of their thresholds, so that
\[
\Delta_k(k-1)=(k+1)\,\hat c_k-2k\,\hat c_{k-1}+(k-1)\,\hat c_{k-2},
\]
and this is nonnegative by the first two inequalities in \eqref{eq:layer-cayley}.

We now express $\Delta_N(N)$ through tilted counts of smaller lengths. Applied to thresholds $k+1$, $k$ and $k-1$, at lengths $N+1$, $N$ and $N-1$, the recursion gives $(N+1)\,\hat z^{(k+1)}_{N+1}=x_k\,\Sigma^{+}$, $N\,\hat z^{(k)}_N=x_k\,\Sigma^{0}$ and $(N-1)\,\hat z^{(k-1)}_{N-1}=x_k\,\Sigma^{-}$, where
\begin{align*}
\Sigma^{+}&=\sum_{i=0}^{N}\min(i+1,k+1)\,\hat z^{(k+1)}_i\,\hat z^{(k+1)}_{N-i},\\
\Sigma^{0}&=\sum_{i=0}^{N-1}\min(i+1,k)\,\hat z^{(k)}_i\,\hat z^{(k)}_{N-1-i},\\
\Sigma^{-}&=\sum_{i=0}^{N-2}\min(i+1,k-1)\,\hat z^{(k-1)}_i\,\hat z^{(k-1)}_{N-2-i}.
\end{align*}
Thus $\Delta_N(N)=x_k\,(\Sigma^{+}-2\,\Sigma^{0}+\Sigma^{-})$, and it suffices to show that $\Sigma^{+}-2\,\Sigma^{0}+\Sigma^{-}\ge 0$.

To this end we sort the terms into three kinds. In each of the three sums we separate the terms in which the first factor has length at most $k-2$, the terms in which the second factor has length at most $k-2$, and the terms in which both factors have length at least $k-1$. Since $N\ge 2k-1$, the lengths of the two factors add up to at least $N-2\ge 2k-3$, so that no term is of the first and of the second kind at once.

Consider a factor of length $\ell\le k-2$. This length lies in the free range of the three thresholds, so that the factor is $\hat c_\ell$ in each of the three sums, and the other factor is $\hat z^{(k+1)}_{N-\ell}$, $\hat z^{(k)}_{N-1-\ell}$ or $\hat z^{(k-1)}_{N-2-\ell}$. When it is the first factor, the weight is $\ell+1$ in each sum. When it is the second factor, the length of the first factor is at least the threshold of the sum, because $N-\ell\ge k+1$, and the weight is that threshold, i.e.\ $k+1$, $k$ or $k-1$. The terms of the first two kinds therefore contribute to $\Sigma^{+}-2\,\Sigma^{0}+\Sigma^{-}$ the amount
\begin{align*}
&\sum_{\ell=0}^{k-2}\hat c_\ell\,\Bigl((\ell+k+2)\,\hat z^{(k+1)}_{N-\ell}-2\,(\ell+k+1)\,\hat z^{(k)}_{N-1-\ell}+(\ell+k)\,\hat z^{(k-1)}_{N-2-\ell}\Bigr)\\
&\qquad=\sum_{\ell=0}^{k-2}\hat c_\ell\,\Delta_{\ell+k+1}(N-1-\ell).
\end{align*}

In a term of the third kind the first factor has a length $i$ with $k-1\le i\le N-k+1$ in $\Sigma^{+}$, with $k-1\le i\le N-k$ in $\Sigma^{0}$, and with $k-1\le i\le N-k-1$ in $\Sigma^{-}$. The weight is the threshold of the sum, with one exception: in the term $i=k-1$ of $\Sigma^{+}$ the weight is $k$, and the first factor of this term is $\hat z^{(k+1)}_{k-1}=\hat c_{k-1}$. We index the other terms by the length of the second factor and put $m_t=N-k-t$ for $0\le t\le r$, so that $m_0=N-k$ and $m_r=k-1$. Then the terms of the third kind are
\begin{align*}
&k\,\hat c_{k-1}\,\hat z^{(k+1)}_{N-k+1}\quad\text{and}\quad(k+1)\,\hat z^{(k+1)}_{k+t}\,\hat z^{(k+1)}_{m_t}\quad(0\le t\le r)&&\text{in }\Sigma^{+},\\
&k\,\hat z^{(k)}_{k-1+t}\,\hat z^{(k)}_{m_t}\quad(0\le t\le r)&&\text{in }\Sigma^{0},\\
&(k-1)\,\hat z^{(k-1)}_{k-2+t}\,\hat z^{(k-1)}_{m_t}\quad(1\le t\le r)&&\text{in }\Sigma^{-}.
\end{align*}
To have the same range of $t$ in the three sums, we add to the last family the term with $t=0$, which is $(k-1)\,\hat c_{k-2}\,\hat z^{(k-1)}_{N-k}$ because $\hat z^{(k-1)}_{k-2}=\hat c_{k-2}$, and subtract it again. With
\[
E_t=(k+1)\,\hat z^{(k+1)}_{k+t}\,\hat z^{(k+1)}_{m_t}-2k\,\hat z^{(k)}_{k-1+t}\,\hat z^{(k)}_{m_t}+(k-1)\,\hat z^{(k-1)}_{k-2+t}\,\hat z^{(k-1)}_{m_t}
\]
for $0\le t\le r$, the terms of the third kind contribute to $\Sigma^{+}-2\,\Sigma^{0}+\Sigma^{-}$ the amount
\begin{equation}\label{eq:layer-middle}
\sum_{t=0}^{r}E_t+k\,\hat c_{k-1}\,\hat z^{(k+1)}_{N-k+1}-(k-1)\,\hat c_{k-2}\,\hat z^{(k-1)}_{N-k},
\end{equation}
and $\Sigma^{+}-2\,\Sigma^{0}+\Sigma^{-}$ is the sum of this amount and of the amount that the terms of the first two kinds contribute.

The amount that the terms of the first two kinds contribute is the sum of the numbers $\hat c_\ell\,\Delta_{\ell+k+1}(N-1-\ell)$, and it is nonnegative. Indeed, for $0\le\ell\le k-2$ the weight $\ell+k+1$ lies between $k+1$ and $2k-1$, and the length $N-1-\ell$ lies between $N-k+1\ge k$ and $N-1$, so that $\Delta_{\ell+k+1}(N-1-\ell)\ge 0$ by \eqref{eq:layer-smaller}.

It remains to show that the amount \eqref{eq:layer-middle}, which the terms of the third kind contribute, is nonnegative. Let $0\le t\le r$. Since the tilted counts of levels $k-1$ and $k$ are the differences of the tilted counts of the three thresholds, we have $\hat z^{(k+1)}_{m_t}=\hat z^{(k)}_{m_t}+\hat\beta^{(k)}_{m_t}$ and $\hat z^{(k-1)}_{m_t}=\hat z^{(k)}_{m_t}-\hat\beta^{(k-1)}_{m_t}$. Inserting this into $E_t$, we find that the terms with $\hat z^{(k)}_{m_t}$ have the coefficient $\Delta_k(k-1+t)$, so that
\begin{align*}
E_t&=\hat z^{(k)}_{m_t}\,\Delta_k(k-1+t)\\
&\quad+(k+1)\,\hat z^{(k+1)}_{k+t}\,\hat\beta^{(k)}_{m_t}-(k-1)\,\hat z^{(k-1)}_{k-2+t}\,\hat\beta^{(k-1)}_{m_t}.
\end{align*}
The first term on the right is nonnegative: for $t=0$ because $\Delta_k(k-1)\ge 0$, as we showed at the beginning of the proof, and for $t\ge 1$ by \eqref{eq:layer-smaller}, since then $k\le k-1+t\le N-k$. In the term that contains $\hat\beta^{(k)}_{m_t}$ we use Corollary~\ref{cor:diagonalcritical} with $n=m_t\ge 1$, which gives $(k+1)\,\hat\beta^{(k)}_{m_t}\ge k\,\hat\beta^{(k-1)}_{m_t-1}$. We obtain
\[
E_t\ge k\,\hat z^{(k+1)}_{k+t}\,\hat\beta^{(k-1)}_{m_t-1}-(k-1)\,\hat z^{(k-1)}_{k-2+t}\,\hat\beta^{(k-1)}_{m_t}\qquad(0\le t\le r).
\]
We sum this over $t$. Among the added terms, the one with $t=r$ vanishes, because $\hat\beta^{(k-1)}_{m_r-1}=\hat\beta^{(k-1)}_{k-2}=0$: the tilted counts of thresholds $k-1$ and $k$ agree at length $k-2$. Among the subtracted terms, the one with $t=0$ is $(k-1)\,\hat c_{k-2}\,\hat\beta^{(k-1)}_{N-k}$, and in the others we replace $t$ by $t+1$ and use that $m_{t+1}=m_t-1$. Therefore
\begin{align*}
\sum_{t=0}^{r}E_t&\ge\sum_{t=0}^{r-1}\hat\beta^{(k-1)}_{m_t-1}\,\Bigl(k\,\hat z^{(k+1)}_{k+t}-(k-1)\,\hat z^{(k-1)}_{k-1+t}\Bigr)\\
&\quad-(k-1)\,\hat c_{k-2}\,\hat\beta^{(k-1)}_{N-k}.
\end{align*}
Every term of the sum on the right is nonnegative, since $\hat\beta^{(k-1)}_{m_t-1}\ge 0$ and $\hat z^{(k-1)}_{k-1+t}\le\hat z^{(k)}_{k+t}\le\hat z^{(k+1)}_{k+t}$; here the first inequality is again the monotonicity along the diagonal, and the second holds because the tilted counts of the three thresholds at one length are ordered. Hence
\[
\sum_{t=0}^{r}E_t\ge-(k-1)\,\hat c_{k-2}\,\hat\beta^{(k-1)}_{N-k}.
\]
Adding the two further terms of the amount \eqref{eq:layer-middle} and using that $\hat\beta^{(k-1)}_{N-k}+\hat z^{(k-1)}_{N-k}=\hat z^{(k)}_{N-k}$, we see that this amount is at least
\[
k\,\hat c_{k-1}\,\hat z^{(k+1)}_{N-k+1}-(k-1)\,\hat c_{k-2}\,\hat z^{(k)}_{N-k}.
\]
This number is nonnegative, because $\hat z^{(k+1)}_{N-k+1}\ge\hat z^{(k)}_{N-k}\ge 0$, once more by the monotonicity along the diagonal, and $k\,\hat c_{k-1}\ge(k-1)\,\hat c_{k-2}$ by the last inequality in \eqref{eq:layer-cayley}.

Thus the two amounts are both nonnegative, so that $\Sigma^{+}-2\,\Sigma^{0}+\Sigma^{-}\ge 0$, and hence $\Delta_N(N)=x_k\,(\Sigma^{+}-2\,\Sigma^{0}+\Sigma^{-})\ge 0$. This completes the induction.
\end{proof}

\subsection{The weights $\Psi(N)$ and the weighted debt}\label{subsec:layer-debt}

For a column $k$ and an integer $N\ge 0$ let
\[
\Psi(N)=\max(N,\,2k-1)
\]
be the weight of row $N$. It depends on the column, which is fixed wherever $\Psi$ occurs. The weight is constant, $\Psi(N)=2k-1$, for $N\le 2k-1$, and $\Psi(N)=N$ for $N\ge 2k-1$; the constant part is chosen for the rows before the crossing row, in which Section~\ref{sec:left} compares the weights with values of the tilted prime kernel of threshold $k$.

\begin{corollary}\label{cor:weightedconvexity}
Let $k\ge 7$ and $N$ be integers with $N\ge k$. Then
\[
\Psi(N+1)\,\hat z^{(k+1)}_{N+1}-2\,\Psi(N)\,\hat z^{(k)}_N+\Psi(N-1)\,\hat z^{(k-1)}_{N-1}\ge 0.
\]
\end{corollary}

\begin{proof}
We use the notation \eqref{eq:layer-weighted}. If $k\le N\le 2k-2$, the three weights are equal to $2k-1$, and the left-hand side is $2k-1$ times $\hat z^{(k+1)}_{N+1}-2\,\hat z^{(k)}_N+\hat z^{(k-1)}_{N-1}$, which is nonnegative by Corollary~\ref{cor:layerconvex}. If $N=2k-1$, the weights are $N+1$, $N$ and $\Psi(N-1)=N$, and the left-hand side is $\Delta_N(N)+\hat z^{(k-1)}_{N-1}$, where $\Delta_N(N)\ge 0$ by Proposition~\ref{prop:beyondlayer} and $\hat z^{(k-1)}_{N-1}\ge 0$. If $N\ge 2k$, the left-hand side is $\Delta_N(N)$, which is nonnegative by the same proposition.
\end{proof}

We now sum by parts. A column is nonpositive in the rows before its crossing row, and nothing is known about the size of its steps in these rows one by one. What can be bounded from below is a weighted sum of the steps, and this without any assumption on their signs, so that the bound applies wherever the crossing row lies. Indeed, the tilted step $\hat\delta_n=\hat z^{(k+1)}_n-2\,\hat z^{(k)}_n+\hat z^{(k-1)}_n$ is a second difference in the threshold, taken at one length, and when the tilted steps of rows $k,\dots,J$ are added with the weights $\Psi(n)$, their terms regroup along the diagonals into the expressions of Corollary~\ref{cor:weightedconvexity}, which are nonnegative, up to terms at the two ends of the range. What is left at the upper end, in rows $J-1$ and $J$, is estimated by the following quantity. For a column $k\ge 2$ and a row $J\ge k$, the \emph{weighted debt} of the column at row $J$ is
\[
\mathfrak D_J=\Psi(J)\,\hat\beta^{(k-1)}_J+\Psi(J-1)\,\hat\beta^{(k-1)}_{J-1}+\bigl(\Psi(J)-\Psi(J-1)\bigr)\,\hat z^{(k)}_J .
\]
Its three terms are nonnegative and belong to rows $J-1$ and $J$; the coefficient $\Psi(J)-\Psi(J-1)$ is $0$ for $J\le 2k-1$ and $1$ for $J\ge 2k$.

\begin{proposition}\label{prop:debt}
Let $k\ge 7$ and $J$ be integers with $J\ge k$. Then
\[
\sum_{n=k}^{J}\Psi(n)\,\hat\delta_n\ge-\mathfrak D_J .
\]
\end{proposition}

\begin{proof}
For an integer $N\ge k$ let
\[
L_N=\Psi(N+1)\,\hat z^{(k+1)}_{N+1}-2\,\Psi(N)\,\hat z^{(k)}_N+\Psi(N-1)\,\hat z^{(k-1)}_{N-1}
\]
denote the left-hand side of the inequality of Corollary~\ref{cor:weightedconvexity}. Summation by parts gives
\begin{multline*}
\sum_{n=k}^{J}\Psi(n)\,\hat\delta_n=\sum_{N=k}^{J-1}L_N+\Psi(k)\,\hat z^{(k+1)}_k-\Psi(k-1)\,\hat z^{(k-1)}_{k-1}\\
-2\,\Psi(J)\,\hat z^{(k)}_J+\Psi(J)\,\hat z^{(k-1)}_J+\Psi(J-1)\,\hat z^{(k-1)}_{J-1}.
\end{multline*}
Indeed, since $\hat\delta_n=\hat z^{(k+1)}_n-2\,\hat z^{(k)}_n+\hat z^{(k-1)}_n$, the left-hand side consists of the terms $\Psi(n)\,\hat z^{(k+1)}_n$, $-2\,\Psi(n)\,\hat z^{(k)}_n$ and $\Psi(n)\,\hat z^{(k-1)}_n$ with $k\le n\le J$, whereas in the sum of the $L_N$ these three families have the ranges $k+1\le n\le J$, $k\le n\le J-1$ and $k-1\le n\le J-2$. The two sums differ by the five terms at the ends of these ranges that the identity displays.

Every $L_N$ is nonnegative by Corollary~\ref{cor:weightedconvexity}. Next, $\Psi(k)=\Psi(k-1)=2k-1$, and the two tilted counts that carry these weights have lengths in the free range of their thresholds, so that
\[
\Psi(k)\,\hat z^{(k+1)}_k-\Psi(k-1)\,\hat z^{(k-1)}_{k-1}=(2k-1)\,(\hat c_k-\hat c_{k-1})\ge 0
\]
by \eqref{eq:critical-updown}. Dropping these nonnegative terms, we get
\[
\sum_{n=k}^{J}\Psi(n)\,\hat\delta_n\ge-2\,\Psi(J)\,\hat z^{(k)}_J+\Psi(J)\,\hat z^{(k-1)}_J+\Psi(J-1)\,\hat z^{(k-1)}_{J-1}.
\]
In terms of the tilted counts of level $k-1$ we have $\hat z^{(k-1)}_J=\hat z^{(k)}_J-\hat\beta^{(k-1)}_J$ and $\hat z^{(k-1)}_{J-1}=\hat z^{(k)}_{J-1}-\hat\beta^{(k-1)}_{J-1}$, so that the right-hand side is
\[
-\Psi(J)\,\hat\beta^{(k-1)}_J-\Psi(J-1)\,\hat\beta^{(k-1)}_{J-1}+\Psi(J-1)\,\hat z^{(k)}_{J-1}-\Psi(J)\,\hat z^{(k)}_J .
\]
Finally, the tilted counts of threshold $k$ do not increase with the length: $\hat z^{(k)}_J\le\hat z^{(k)}_{J-1}$ by Proposition~\ref{prop:criticalmonotone}. Hence
\[
\Psi(J-1)\,\hat z^{(k)}_{J-1}-\Psi(J)\,\hat z^{(k)}_J\ge-\bigl(\Psi(J)-\Psi(J-1)\bigr)\,\hat z^{(k)}_J,
\]
and with the two preceding displays this gives $\sum_{n=k}^{J}\Psi(n)\,\hat\delta_n\ge-\mathfrak D_J$.
\end{proof}

\section{The left half of a row}\label{sec:left}
In this section we prove Proposition~\ref{prop:leftpart}: for all integers $n\ge 17$ and $k\ge 1$ with $2k\le n$ we have $b_{n,k-1}\le b_{n,k}$, so that a row $n\ge 17$ has no strict descent at an index $i$ with $2i\le n$. For $k\ge 82$ we argue by induction on the row, starting from the single crossing of Proposition~\ref{prop:onecrossing}; for $2\le k\le 81$ we use the computation of Lemma~\ref{lem:columntable}; and column $k=1$ is treated directly.

\subsection{The tilted prime kernel at the critical tilt}\label{subsec:left-kernel}

Throughout this subsection $k\ge 2$ is a fixed integer and, as in Section~\ref{sec:critical}, a hat denotes the tilted form of a sequence at the critical tilt $x_k$; in particular, $\hat\kappa_m=\kappa^{(k)}_m\,x_k^{m}$ is the tilted prime kernel of threshold $k$. We prove the estimates for $\hat\kappa$ on which the proof of Proposition~\ref{prop:largecolumns} rests.

We begin with what Sections~\ref{sec:series} and~\ref{sec:critical} provide. The tilted kernel is nonnegative (Lemma~\ref{lem:kernelnonneg}), and in the free range it is the tilted free kernel,
\[
\hat\kappa_m=\varepsilon_m\,x_k^{m}\qquad(1\le m\le k),
\]
by Lemma~\ref{lem:kernelfree}. By Lemma~\ref{lem:epsdown} we have $\varepsilon_{m+1}\,x_k^{m+1}\le\varepsilon_m\,x_k^{m}$ for $1\le m<k$, so that Lemma~\ref{lem:harmonic} and \eqref{eq:critical-kerneldown} apply with $y=k$ and $x=x_k$. They give
\[
\hat\kappa_{m+1}\le\hat\kappa_m\quad(m\ge 1)\qquad\text{and}\qquad(m+1)\,\hat\kappa_{m+1}\le m\,\hat\kappa_m\quad(m\ge k).
\]
Thus the tilted kernel is nonincreasing, and the second inequality yields, by induction on $L$,
\begin{equation}\label{eq:left-harmonic}
L\,\hat\kappa_L\le k\,\hat\kappa_k\qquad(L\ge k).
\end{equation}
Finally, the recursion of Lemma~\ref{lem:kernelrecursion}, multiplied by $x_k^{N}$, reads
\[
N\,\hat\kappa_N=\sum_{j=1}^{k-1}(k-j)\,\hat a_j\,\hat\kappa_{N-j}\qquad(N\ge k),
\]
and its weights $(k-j)\,\hat a_j=(k-j)\,a_j\,x_k^{j}$ are positive.

Beyond the free range the kernel decays the faster the further out its index lies: for all integers $r\ge 1$ and $m\ge k+(r-1)(k-1)$,
\begin{equation}\label{eq:left-rate}
(m+1)\,\hat\kappa_{m+1}\le(m+1-r)\,\hat\kappa_m .
\end{equation}
We prove this by induction on $r$. For $r=1$ it is the inequality $(m+1)\,\hat\kappa_{m+1}\le m\,\hat\kappa_m$ for $m\ge k$ recorded above. Suppose that the assertion holds for some $r\ge 1$, and let $m\ge k+r\,(k-1)$. For $1\le j\le k-1$ the index $i=m-j$ satisfies $i\ge m-k+1\ge k+(r-1)(k-1)$, so that $(i+1)\,\hat\kappa_{i+1}\le(i+1-r)\,\hat\kappa_i$ by the induction hypothesis. Since moreover $\hat\kappa_{i+1}\le\hat\kappa_i$ and $m-i-1\ge 0$, we get
\begin{align*}
m\,\hat\kappa_{i+1}&=(i+1)\,\hat\kappa_{i+1}+(m-i-1)\,\hat\kappa_{i+1}\\
&\le(i+1-r)\,\hat\kappa_i+(m-i-1)\,\hat\kappa_i=(m-r)\,\hat\kappa_i ,
\end{align*}
i.e.\ $m\,\hat\kappa_{m+1-j}\le(m-r)\,\hat\kappa_{m-j}$. We multiply this by $(k-j)\,\hat a_j$, sum over $1\le j\le k-1$, and use the tilted recursion for $N=m+1$ and for $N=m$. This gives $m\,(m+1)\,\hat\kappa_{m+1}\le(m-r)\,m\,\hat\kappa_m$, which is the assertion for $r+1$.

In the free range the size of the kernel is given by Stirling's formula. Let $2\le L\le k$. Since $\hat\kappa_L=\varepsilon_L\,x_k^{L}$ and $L\,\varepsilon_L=(L-1)^{L-1}/(L-1)!$,
\[
\hat\kappa_L=\frac{(e\,x_k)^{L}}{e\,L}\cdot e^{-(L-1)}\,\frac{(L-1)^{L-1}}{(L-1)!},
\]
and the second factor is the Poisson mass at $L-1\ge 1$. By \eqref{eq:prelim-stirling} it lies between $\frac{11}{12}\,(2\pi(L-1))^{-1/2}$ and $(2\pi(L-1))^{-1/2}$, since $1-\frac{1}{12(L-1)}\ge\frac{11}{12}$. Hence
\begin{equation}\label{eq:left-free}
\frac{11}{12}\cdot\frac{(e\,x_k)^{L}}{e\,L\,\sqrt{2\pi\,(L-1)}}\le\hat\kappa_L\le\frac{(e\,x_k)^{L}}{e\,L\,\sqrt{2\pi\,(L-1)}}\qquad(2\le L\le k).
\end{equation}
Since $e\,x_k\le e^{3/(2k)}$ by Lemma~\ref{lem:tiltupper}, the upper bound gives
\begin{equation}\label{eq:left-freeupper}
\hat\kappa_L\le\frac{\exp\bigl(\frac{3L}{2k}-1\bigr)}{L\,\sqrt{2\pi\,(L-1)}}\qquad(2\le L\le k).
\end{equation}

The next two lemmas order the values $\hat\kappa_{n-j}$ of the tilted kernel, measured against the weight $\Psi(j)=\max(j,2k-1)$ of Section~\ref{sec:layer}, on the two sides of a row $q\le k+27k/100$.

\begin{lemma}\label{lem:orderbefore}
Let $k$, $q$, $n$ and $j$ be nonnegative integers with $k\ge 2$, $100\,(q-k)\le 27\,k$, $n\ge 2k$ and $j<q$. Then
\[
\hat\kappa_{n-j}\,(2k-1)\le\hat\kappa_{n-q+1}\,\Psi(j).
\]
\end{lemma}

\begin{proof}
Since $k\ge 2$, we have $j<q\le 1.27\,k\le 2k-1$, so that $\Psi(j)=2k-1$. Moreover $1\le n-q+1\le n-j$, and since the tilted kernel is nonincreasing, $\hat\kappa_{n-j}\le\hat\kappa_{n-q+1}$, which is the claim.
\end{proof}

\begin{lemma}\label{lem:orderafter}
Let $k$, $q$, $n$ and $j$ be integers with $k\ge 2$, $k+1\le q$, $100\,(q-k)\le 27\,k$, $n\ge 2k$ and $q\le j\le n-1$. Then
\[
\hat\kappa_{n-q+1}\,\Psi(j)\le\hat\kappa_{n-j}\,(2k-1).
\]
\end{lemma}

\begin{proof}
Put $L=n-q+1$ and $l=n-j$, so that the claim is $\hat\kappa_L\,\Psi(j)\le(2k-1)\,\hat\kappa_l$. The hypothesis on $q$ gives $q\le 1.27\,k$; in particular $q\le 2k\le n$, and $1\le l<L$ because $q\le j\le n-1$. If $j\le 2k-1$, then $\Psi(j)=2k-1$, and the claim holds because the tilted kernel is nonincreasing, so that $\hat\kappa_L\le\hat\kappa_l$.

Let now $j\ge 2k$. Then $\Psi(j)=j$, and the claim is
\[
j\,\hat\kappa_L\le(2k-1)\,\hat\kappa_l .
\]
Whether the kernel grows fast enough for this as its index falls from $L$ to $l$ depends on where the two indices lie relative to the free range $1\le m\le k$.

Beyond the free range the kernel falls at least by a fixed factor from each index to the next:
\[
(2k-1)\,\hat\kappa_{m+1}\le(2k-2)\,\hat\kappa_m\qquad(m\ge k).
\]
Indeed, let $m\ge k$, and write $m-k=(k-1)\,d+s$ with integers $d\ge 0$ and $0\le s\le k-2$. Then $m\ge k+d\,(k-1)$, and the inequality \eqref{eq:left-rate} with $r=d+1$ gives $(m+1)\,\hat\kappa_{m+1}\le(m-d)\,\hat\kappa_m$. Since $m+1\le k+(d+1)(k-1)\le(d+1)(2k-1)$, we have
\[
(2k-1)\,(m-d)-(2k-2)\,(m+1)=(m+1)-(d+1)(2k-1)\le 0 ,
\]
and therefore $(2k-1)\,(m+1)\,\hat\kappa_{m+1}\le(2k-1)\,(m-d)\,\hat\kappa_m\le(2k-2)\,(m+1)\,\hat\kappa_m$, because $\hat\kappa_m\ge 0$. Division by $m+1$ gives the asserted inequality.

We deduce by induction on $i$ that, for every integer $i\ge 2k-1$ with $n-i\ge k$,
\begin{equation}\label{eq:left-beyond}
i\,\hat\kappa_{n-2k+1}\le(2k-1)\,\hat\kappa_{n-i}.
\end{equation}
For $i=2k-1$ the two sides are equal. If the inequality holds for some $i\ge 2k-1$ and $n-(i+1)\ge k$, then $(2k-1)\,\hat\kappa_{n-i}\le(2k-2)\,\hat\kappa_{n-i-1}$, by the inequality of the preceding paragraph with $m=n-i-1$. Hence $i\,\hat\kappa_{n-2k+1}\le(2k-2)\,\hat\kappa_{n-i-1}$, and the inequality for $i+1$ follows because $(2k-2)(i+1)\le(2k-1)\,i$.

In the free range we compare the values of the kernel at two indices that are far apart: if $l_1$ and $l_2$ are integers with $1\le l_1<l_2\le k$ and $100\,(l_2-l_1)\ge 73\,k$, then
\begin{equation}\label{eq:left-drop}
1.22\,\hat\kappa_{l_2}\le\hat\kappa_{l_1}.
\end{equation}
Suppose first that $l_1=1$. In the free range the kernel is the tilted free kernel, so that $\hat\kappa_1=\varepsilon_1x_k=x_k$ and $\hat\kappa_2=\varepsilon_2x_k^2=\frac12x_k^2$. Since $l_2\ge 2$ and $0<x_k\le 1$, it follows that $\hat\kappa_{l_2}\le\hat\kappa_2\le\frac12x_k=\frac12\hat\kappa_1$, which implies the asserted inequality. Let now $l_1\ge 2$. Applying the two-sided bound \eqref{eq:left-free} to $l_2$ from above and to $l_1$ from below, we obtain
\[
\frac{\hat\kappa_{l_2}}{\hat\kappa_{l_1}}\le\frac{12}{11}\,(e\,x_k)^{l_2-l_1}\,\frac{l_1}{l_2}\,\sqrt{\frac{l_1-1}{l_2-1}}.
\]
Here $(e\,x_k)^{l_2-l_1}\le e^{3(l_2-l_1)/(2k)}\le e^{3/2}$ by Lemma~\ref{lem:tiltupper}, because $l_2-l_1\le k$. The hypothesis gives $l_2-l_1\ge 0.73\,k\ge 0.73\,l_2$, hence $l_1\le 0.27\,l_2$ and also $l_1-1\le 0.27\,(l_2-1)$. So the right-hand side is at most $\frac{12}{11}\,e^{3/2}\cdot 0.27^{3/2}$, which is smaller than $0.69<1/1.22$, since $e^{3/2}<4.482$ and $0.27^{3/2}<0.1403$.

We now prove the claim $j\,\hat\kappa_L\le(2k-1)\,\hat\kappa_l$, in three cases according to the position of $l$ and $L$ relative to $k$: either $l\ge k$, or $L\le k$, or $l<k<L$. Suppose first that $l\ge k$. Since $j\ge 2k-1$ and $n-j=l\ge k$, the inequality \eqref{eq:left-beyond} with $i=j$ gives $j\,\hat\kappa_{n-2k+1}\le(2k-1)\,\hat\kappa_l$. Since $q\le 2k$, we have $L\ge n-2k+1\ge 1$, hence $\hat\kappa_L\le\hat\kappa_{n-2k+1}$, and the claim follows.

Suppose next that $L\le k$. Then $1\le l<L\le k$, and $L-l=j-q+1\ge 2k-1.27\,k+1>0.73\,k$. So the comparison \eqref{eq:left-drop} gives $1.22\,\hat\kappa_L\le\hat\kappa_l$. Moreover $n=L+q-1\le k+q-1$, so that $j\le n-1\le 2.27\,k-2\le 1.22\,(2k-1)$. Hence $j\,\hat\kappa_L\le 1.22\,(2k-1)\,\hat\kappa_L\le(2k-1)\,\hat\kappa_l$, which is the claim.

It remains to treat the case $l<k<L$. By the bound \eqref{eq:left-harmonic} we have $L\,\hat\kappa_L\le k\,\hat\kappa_k$, and we compare $\hat\kappa_k$ with $\hat\kappa_l$. Note that $j=n-l=L+q-1-l$ and that $L\ge k+1$. If $100\,l\le 27\,k$, then $100\,(k-l)\ge 73\,k$, and \eqref{eq:left-drop} with $l_1=l$ and $l_2=k$ gives $1.22\,\hat\kappa_k\le\hat\kappa_l$. Since $l\ge 1$, we have $j\le L+1.27\,k-2$, and
\[
k\,j\le k\,(L+1.27\,k-2)\le 1.22\,(2k-1)\,L ,
\]
where the second inequality holds because $L\ge k+1$, as a direct computation shows. Therefore
\[
L\cdot j\,\hat\kappa_L\le j\cdot k\,\hat\kappa_k\le 1.22\,(2k-1)\,L\,\hat\kappa_k\le(2k-1)\,L\,\hat\kappa_l ,
\]
and division by $L$ gives the claim. If, on the other hand, $100\,l>27\,k$, then $j<L+1.27\,k-1-0.27\,k=L+k-1$, i.e.\ $j\le L+k-2$, and
\[
k\,j\le k\,(L+k-2)\le(2k-1)\,L ,
\]
again because $L\ge k+1$. Since $\hat\kappa_k\le\hat\kappa_l$, the kernel being nonincreasing, we obtain $L\cdot j\,\hat\kappa_L\le j\cdot k\,\hat\kappa_k\le(2k-1)\,L\,\hat\kappa_k\le(2k-1)\,L\,\hat\kappa_l$, and the claim follows as before.
\end{proof}

Finally, from index $2k-1$ on the kernel decays at least like $1/L^2$.

\begin{lemma}\label{lem:quadraticdecay}
Let $k\ge 2$ be an integer. Then, for every integer $L\ge 2k-1$,
\[
(L-1)\,L\,\hat\kappa_L\le 2k\,(k-1)\,\hat\kappa_k .
\]
\end{lemma}

\begin{proof}
We use induction on $L$. For $L=2k-1$, the bound \eqref{eq:left-harmonic} gives $(2k-1)\,\hat\kappa_{2k-1}\le k\,\hat\kappa_k$, and multiplying by $2k-2$ gives the claim. Suppose that the claim holds for some $L\ge 2k-1$. Since $L\ge k+(k-1)$, the inequality \eqref{eq:left-rate} with $r=2$ and $m=L$ gives $(L+1)\,\hat\kappa_{L+1}\le(L-1)\,\hat\kappa_L$. Multiplying by $L$ and using the induction hypothesis, we obtain
\[
L\,(L+1)\,\hat\kappa_{L+1}\le(L-1)\,L\,\hat\kappa_L\le 2k\,(k-1)\,\hat\kappa_k .\qedhere
\]
\end{proof}

\subsection{Columns $k\ge 82$}\label{subsec:left-induction}

We now treat the columns $k\ge 82$, for which we prove the strict inequality $b_{n,k-1}<b_{n,k}$ in every row $n\ge 2k$. By Proposition~\ref{prop:onecrossing}, such a column is nonpositive in rows $k\le j<q$, before its crossing row $q$, and positive in rows $q\le j\le 2k-1$, and we show by induction on $n\ge 2k$ that it is positive in row $n$. As in the chain range, the difficulty is that the rows before $q$ are fed into row $n$, here by the column inequality of Proposition~\ref{prop:columninequality}, and that we have no bound for a single step in these rows. What we have is the bound of the weighted debt (Proposition~\ref{prop:debt}) for the sum of the steps with the weights $\Psi(j)$, and Lemmas~\ref{lem:orderbefore} and~\ref{lem:orderafter} allow us to pass from the kernel weights of the column inequality to these weights.

\begin{proposition}\label{prop:largecolumns}
Let $k\ge 82$ and $n\ge 2k$ be integers. Then $b_{n,k-1}<b_{n,k}$.
\end{proposition}

\begin{proof}
Fix $k\ge 82$. As in Section~\ref{sec:critical}, a hat denotes the tilted form of a sequence at the critical tilt $x_k$. Since $\hat\delta_n=d_{n,k}\,x_k^{n}/n!$ and $x_k>0$, we have to show that $\hat\delta_n>0$ for every $n\ge 2k$.

By Proposition~\ref{prop:onecrossing} there is an integer $q$ with $k+1\le q$ and $100\,(q-k)\le 27\,k$ such that
\[
\hat\delta_j\le 0\quad\text{for } k\le j<q\qquad\text{and}\qquad \hat\delta_j>0\quad\text{for } q\le j\le 2k-1 .
\]
In particular $q\le 1.27\,k<2k-1$. We show that $\hat\delta_n>0$ for every $n\ge q$, by induction on $n$. For $q\le n\le 2k-1$ this holds by the choice of $q$. Let $n\ge 2k$, and suppose that $\hat\delta_j>0$ for $q\le j\le n-1$. We put
\[
L=n-q+1\qquad\text{and}\qquad\theta=\frac{\hat\kappa_L}{2k-1},
\]
so that $L\ge 1$ and $\theta\ge 0$, and we show that $\hat\delta_n>0$.

We start from the column inequality. Multiplying the inequality of Proposition~\ref{prop:columninequality} for row $n$ by $x_k^{n}$ and indexing the sum by $j=n-m$, we obtain
\begin{equation}\label{eq:left-column}
n\,\hat\delta_n\ge(k+1)\,x_k\,\hat\delta_{n-1}+2\,x_k\,\hat\beta^{(k-1)}_{n-1}+n\sum_{j=k}^{n-1}\hat\kappa_{n-j}\,\hat\delta_j .
\end{equation}

We first bound the sum in this inequality from below. Recall from Section~\ref{sec:layer} that the weighted debt of the column at a row $J\ge k$ is
\[
\mathfrak D_J=\Psi(J)\,\hat\beta^{(k-1)}_J+\Psi(J-1)\,\hat\beta^{(k-1)}_{J-1}+\bigl(\Psi(J)-\Psi(J-1)\bigr)\,\hat z^{(k)}_J .
\]
For $k\le j<q$, Lemma~\ref{lem:orderbefore} gives $\hat\kappa_{n-j}\le\theta\,\Psi(j)$, and $\hat\delta_j\le 0$ by the choice of $q$. For $q\le j\le n-1$, Lemma~\ref{lem:orderafter} gives $\hat\kappa_{n-j}\ge\theta\,\Psi(j)$, and $\hat\delta_j>0$ by the choice of $q$ and the induction hypothesis. In both cases
\[
\hat\kappa_{n-j}\,\hat\delta_j\ge\theta\,\Psi(j)\,\hat\delta_j .
\]
We sum over $k\le j\le n-1$ and apply Proposition~\ref{prop:debt}, which holds for $k\ge 7$, with $J=n-1\ge k$. As $\theta\ge 0$, this gives
\[
\sum_{j=k}^{n-1}\hat\kappa_{n-j}\,\hat\delta_j\ge\theta\sum_{j=k}^{n-1}\Psi(j)\,\hat\delta_j\ge-\theta\,\mathfrak D_{n-1}.
\]

Next we bound $\mathfrak D_{n-1}$ from above by a combination of $\hat\beta^{(k-1)}_{n-1}$ and $\hat\delta_{n-1}$, the two quantities in the first two terms on the right of the column inequality \eqref{eq:left-column}. Since $n-1\ge 2k-1$, we have $\Psi(n-1)=n-1$ and $n-2\le\Psi(n-2)\le n-1$. As the tilted counts are nonnegative, the definition of $\mathfrak D_{n-1}$ gives
\begin{equation}\label{eq:left-debtfirst}
\mathfrak D_{n-1}\le(n-1)\,\bigl(\hat\beta^{(k-1)}_{n-1}+\hat\beta^{(k-1)}_{n-2}\bigr)+\hat z^{(k)}_{n-1}.
\end{equation}
For the count of row $n-2$ we use Corollary~\ref{cor:diagonalcritical} with $n-1$ in place of $n$, which states that $k\,\hat\beta^{(k-1)}_{n-2}\le(k+1)\,\hat\beta^{(k)}_{n-1}$. Since $\hat\beta^{(k)}_{n-1}=\hat\beta^{(k-1)}_{n-1}+\hat\delta_{n-1}$, this is
\[
\hat\beta^{(k-1)}_{n-2}\le\Bigl(1+\frac1k\Bigr)\,\bigl(\hat\beta^{(k-1)}_{n-1}+\hat\delta_{n-1}\bigr).
\]
For $\hat z^{(k)}_{n-1}$ we show that
\[
x_k\,\Bigl(\sum_{j=0}^{8}\hat c_j\Bigr)\,\hat z^{(k)}_{n-1}\le(n-1)\,\hat\beta^{(k-1)}_{n-1}.
\]
The sum on the left has nine terms, the numbers $\hat c_j=c_j\,x_k^{j}$ with $0\le j\le 8$, whatever $k$ is; nine terms are enough for the numerical verification at the end of the proof. To prove the inequality, we subtract the tilted last-car recursions of thresholds $k$ and $k-1$ at length $n-1$. By \eqref{eq:recursion-tilted} with $N=n-2$ and by $\hat\beta^{(k-1)}_{n-1}=\hat z^{(k)}_{n-1}-\hat z^{(k-1)}_{n-1}$,
\[
(n-1)\,\hat\beta^{(k-1)}_{n-1}=x_k\sum_{i=0}^{n-2}H_i ,
\]
where
\[
H_i=\min(i+1,k)\;\hat z^{(k)}_i\,\hat z^{(k)}_{n-2-i}-\min(i+1,k-1)\;\hat z^{(k-1)}_i\,\hat z^{(k-1)}_{n-2-i}.
\]
Every $H_i$ is nonnegative, because $\min(i+1,k-1)\le\min(i+1,k)$ and $0\le\hat z^{(k-1)}_N\le\hat z^{(k)}_N$ for every $N$ (Section~\ref{sec:recursion}). Consider the nine indices $i=n-2-j$ with $0\le j\le 8$. For these we have $i+1\ge n-9\ge k$, because $n\ge 2k$ and $k\ge 9$, so that the two minima are $k$ and $k-1$, and
\[
H_{n-2-j}=k\,\hat z^{(k)}_{n-2-j}\,\hat z^{(k)}_j-(k-1)\,\hat z^{(k-1)}_{n-2-j}\,\hat z^{(k-1)}_j\ge\hat z^{(k)}_{n-2-j}\,\hat z^{(k)}_j .
\]
Here $\hat z^{(k)}_j=\hat c_j$ by Lemma~\ref{lem:freerange}, since $j\le k$, and $\hat z^{(k)}_{n-2-j}\ge\hat z^{(k)}_{n-1}$ by Proposition~\ref{prop:criticalmonotone}. Keeping only these nine terms of the sum, we obtain the asserted bound for $\hat z^{(k)}_{n-1}$.

We insert the bounds for $\hat\beta^{(k-1)}_{n-2}$ and for $\hat z^{(k)}_{n-1}$ into the bound \eqref{eq:left-debtfirst}. With
\[
B=2+\frac1k+\Bigl(x_k\sum_{j=0}^{8}\hat c_j\Bigr)^{-1}
\]
the result is
\[
\mathfrak D_{n-1}\le(n-1)\,\Bigl(B\,\hat\beta^{(k-1)}_{n-1}+\Bigl(1+\frac1k\Bigr)\,\hat\delta_{n-1}\Bigr).
\]

We now return to the column inequality \eqref{eq:left-column}. Its sum is at least $-\theta\,\mathfrak D_{n-1}$, as we have seen, and with the bound for $\mathfrak D_{n-1}$ just obtained and the definition of $\theta$ we find
\begin{multline*}
n\,\hat\delta_n\ge\hat\beta^{(k-1)}_{n-1}\,\Bigl(2\,x_k-\frac{n(n-1)\,\hat\kappa_L}{2k-1}\,B\Bigr)\\
+\hat\delta_{n-1}\,(k+1)\,\Bigl(x_k-\frac{n(n-1)\,\hat\kappa_L}{(2k-1)\,k}\Bigr).
\end{multline*}
Here $\hat\beta^{(k-1)}_{n-1}\ge 0$, and $\hat\delta_{n-1}>0$ because $q\le n-1$. So $\hat\delta_n>0$ as soon as
\begin{equation}\label{eq:left-conditions}
n(n-1)\,\hat\kappa_L\cdot B\le 2\,x_k\,(2k-1)\qquad\text{and}\qquad n(n-1)\,\hat\kappa_L<x_k\,k\,(2k-1).
\end{equation}

It remains to verify these two inequalities. To this end we show that
\[
n(n-1)\,\hat\kappa_L\le 3.762\,\sqrt k .
\]
For this we need the size of the kernel at two places:
\[
\begin{gathered}
\hat\kappa_k\le\frac{0.662}{k\sqrt k}\\
\text{and}\qquad\hat\kappa_M\le\frac{0.73}{k\sqrt k}\quad\text{for every integer $M$ with } 100\,(M-1)\ge 73\,k .
\end{gathered}
\]
For the first bound, the upper bound \eqref{eq:left-freeupper} at the index $k$ gives
\[
\hat\kappa_k\le\frac{e^{1/2}}{k\,\sqrt{2\pi\,(k-1)}},
\]
and $e^{1/2}\sqrt k\le 0.662\,\sqrt{2\pi\,(k-1)}$, because $e\,k\le 0.662^2\cdot 2\pi\,(k-1)$ for $k\ge 79$; here we use that $e<2.7183$ and $0.662^2\cdot 2\pi>2.7534$.

For the second bound let $M$ be an integer with $100\,(M-1)\ge 73\,k$, and let $t$ be the least integer with $100\,t\ge 73\,k$. Then $t\le M-1$ and $0.73\,k\le t<0.73\,k+1$, so that $2\le t+1\le M$ and $t+1<0.73\,k+2\le k$. Hence $\hat\kappa_M\le\hat\kappa_{t+1}$, because the tilted kernel is nonincreasing, and since $t+1$ lies in the free range, the upper bound \eqref{eq:left-freeupper} at the index $t+1$ gives
\[
\hat\kappa_{t+1}\le\frac{1}{(t+1)\,\sqrt{2\pi\,t}}\,\exp\Bigl(\frac{3\,(t+1)}{2k}-1\Bigr).
\]
Put $\rho=(t+1)/k$, so that the exponent is $\frac32\rho-1$. We have $0.73\le\rho\le 0.755$, because $0.73\,k+1\le t+1<0.73\,k+2$ and $2/k<0.025$; hence $0<\frac32\rho-1<1$. We use the elementary inequality $e^{s}\le 1+s+s^{2}$ for real numbers $s$ with $0\le s\le 1$, which holds because $e^{s}-1-s=s^{2}\sum_{j\ge 0}s^{j}/(j+2)!$ and the last sum is at most $e-2<1$ on this interval. For $s=\frac32\rho-1$ we have $1+s+s^{2}=1-\frac32\rho+\frac94\rho^{2}$, and since $t+1=\rho\,k$, it follows that
\[
\frac{1}{t+1}\,\exp\Bigl(\frac32\rho-1\Bigr)\le\frac{1}{k}\,\Bigl(\frac1\rho-\frac32+\frac94\,\rho\Bigr)\le\frac{1.5233}{k}.
\]
The second inequality holds because the bracket is a convex function of $\rho>0$ whose values at $\rho=0.73$ and at $\rho=0.755$ are smaller than $1.5233$. Finally $\sqrt{2\pi\,t}\ge\sqrt{2\pi\cdot 0.73\,k}\ge 2.1416\,\sqrt k$, because $2\pi\cdot 0.73>4.5865>2.1416^2$. Together,
\[
\hat\kappa_M\le\hat\kappa_{t+1}\le\frac{1.5233}{2.1416}\cdot\frac{1}{k\sqrt k}\le\frac{0.73}{k\sqrt k}.
\]

We now prove that $n(n-1)\,\hat\kappa_L\le 3.762\,\sqrt k$. Recall that $n=L+q-1$ and $q\le 1.27\,k$. Suppose first that $L\le k$. Then $L-1=n-q\ge 2k-1.27\,k=0.73\,k$, so that the bound for $\hat\kappa_M$ applies to $M=L$, and $n<2.27\,k$. Hence
\[
n(n-1)\,\hat\kappa_L\le n^2\,\hat\kappa_L\le 2.27^2\cdot 0.73\,\sqrt k\le 3.762\,\sqrt k .
\]
Suppose next that $k<L\le 2k$. By the bound \eqref{eq:left-harmonic} and the bound for $\hat\kappa_k$, we have $L\,\hat\kappa_L\le k\,\hat\kappa_k\le 0.662/\sqrt k$. Moreover $n\le L+1.27\,k$, and it is easy to check that $(L+1.27\,k)^2\le 5.35\,k\,L$ for $k\le L\le 2k$: the difference of the two sides is convex in $L$, and the inequality holds at $L=k$ and at $L=2k$. Therefore
\[
n(n-1)\,\hat\kappa_L\le\frac{(L+1.27\,k)^2}{L}\cdot L\,\hat\kappa_L\le 5.35\,k\cdot\frac{0.662}{\sqrt k}<3.542\,\sqrt k .
\]
Suppose finally that $L>2k$. By Lemma~\ref{lem:quadraticdecay} and the bound for $\hat\kappa_k$,
\[
(L-1)\,L\,\hat\kappa_L\le 2k\,(k-1)\,\hat\kappa_k\le 1.324\,\sqrt k .
\]
Since $1.27\,k<0.635\,L$, we have $n<1.635\,L-1$, and hence $n(n-1)\le(1.635\,L-1)(1.635\,L-2)\le 2.675\,(L-1)\,L$, the last inequality being elementary for $L\ge 1$. Therefore $n(n-1)\,\hat\kappa_L\le 2.675\cdot 1.324\,\sqrt k<3.542\,\sqrt k$. This proves the bound in all three cases.

We can now verify the two conditions \eqref{eq:left-conditions}. The quantities that enter satisfy
\[
x_k>0.3678,\qquad\sum_{j=0}^{8}\hat c_j\ge 2.01,\qquad B<3.3652 .
\]
Indeed, $x_k\bigl(1+\frac1k\bigr)^{k-1}=1$ and $\bigl(1+\frac1k\bigr)^{k-1}\le e^{(k-1)/k}<e$, so that $x_k>e^{-1}>0.3678$. Each term $c_jx^{j}$ of the sum $\sum_{j=0}^{8}c_jx^{j}$ is nondecreasing in $x\ge 0$, so that $\sum_{j=0}^{8}\hat c_j=\sum_{j=0}^{8}c_j\,x_k^{j}$ is at least the value of this sum at $x=0.3678$, which is $2.0162\ldots$\,. Finally,
\[
B\le 2+\frac{1}{82}+\frac{1}{0.3678\cdot 2.01}<3.3652 .
\]
Since $n(n-1)\,\hat\kappa_L\le 3.762\,\sqrt k$, it follows that
\[
n(n-1)\,\hat\kappa_L\cdot B<3.762\cdot 3.3652\,\sqrt k<12.66\,\sqrt k\le 0.7356\,(2k-1)<2\,x_k\,(2k-1).
\]
The third inequality holds for every $k\ge 82$: at $k=82$ its two sides are $114.64\ldots$ and $119.90\ldots$, and their difference increases with $k$ for $k\ge 19$. This is the first of the two conditions. The second follows from the same bound for $n(n-1)\,\hat\kappa_L$, since $3.762\,\sqrt k\le 3.762\,k<0.3678\,k\,(2k-1)$. Hence $\hat\delta_n>0$, which completes the induction.
\end{proof}

\subsection{Columns $2\le k\le 81$}\label{subsec:left-small}

For $2\le k\le 81$ the inequality $b_{n,k-1}\le b_{n,k}$ is obtained from the computation of Lemma~\ref{lem:columntable}. The next proposition turns finitely many inequalities for the numbers $\xi^{(k)}_t$ of column $k$ (Section~\ref{sec:series}), checked at a single row $n_0$, into the sign of the step in every row $n\ge n_0$; recall that $\delta^{(k)}_n\ge 0$ means that $b_{n,k-1}\le b_{n,k}$.

\begin{proposition}\label{prop:propagation}
Let $k\ge 2$ and $n_0\ge k$ be integers, and let $x>0$ be a real number. Suppose that
\begin{enumerate}
\item $z^{(k)}_{N+1}\,x^{N+1}\le z^{(k)}_N\,x^{N}$ for every integer $N\ge 0$, and
\item $\sum_{t=s}^{n_0}\xi^{(k)}_t\,x^{t}\ge 0$ for every integer $s$ with $k\le s\le n_0$.
\end{enumerate}
Then $\delta^{(k)}_n\ge 0$ for every integer $n\ge n_0$.
\end{proposition}

\begin{proof}
We write $\xi_t=\xi^{(k)}_t$ and $\delta_n=\delta^{(k)}_n$.

For an integer $n\ge k$ let $I(n)$ be the statement that
\[
\sum_{t=s}^{n}\xi_t\,x^{t}\ge 0\qquad\text{for every integer $s$ with } k\le s\le n .
\]
Hypothesis~(2) is $I(n_0)$. We show that, for every $n\ge k$, the statement $I(n)$ implies that $\delta_n\ge 0$ and that $I(n+1)$ holds. By induction, $I(n)$ then holds for every $n\ge n_0$, and the proposition follows.

Let $n\ge k$, and suppose that $I(n)$ holds. Put $u_t=\xi_t\,x^{t}$ and $v_t=z^{(k)}_{n-t}\,x^{n-t}$ for $0\le t\le n$. Multiplying the inequality of Proposition~\ref{prop:onetree} by $x^{n}$, we obtain $\sum_{t=0}^{n}u_t\,v_t\le\delta_n\,x^{n}$. The terms with $t<k$ vanish, because $\xi_t=0$ for $t<k$ (Section~\ref{sec:series}). Hence
\[
\sum_{t=k}^{n}u_t\,v_t\le\delta_n\,x^{n}.
\]
The weights $v_t$ are nonnegative, and $v_s\le v_{s+1}$ for $k\le s<n$ by hypothesis~(1) with $N=n-s-1$. We claim that
\[
v_s\sum_{t=s}^{n}u_t\le\sum_{t=s}^{n}u_t\,v_t\qquad\text{for every integer $s$ with } k\le s\le n ,
\]
and we prove this by downward induction on $s$. For $s=n$ both sides are $u_n\,v_n$. If the inequality holds for $s+1$, where $k\le s<n$, then
\[
\sum_{t=s}^{n}u_t\,v_t\ge u_s\,v_s+v_{s+1}\sum_{t=s+1}^{n}u_t\ge u_s\,v_s+v_s\sum_{t=s+1}^{n}u_t=v_s\sum_{t=s}^{n}u_t ,
\]
where the second step uses that $v_{s+1}\ge v_s$ and that $\sum_{t=s+1}^{n}u_t\ge 0$ by $I(n)$. For $s=k$ the inequality, together with $I(n)$ and the one-tree bound in the form displayed above, gives $0\le v_k\sum_{t=k}^{n}u_t\le\sum_{t=k}^{n}u_t\,v_t\le\delta_n\,x^{n}$, and hence $\delta_n\ge 0$.

It remains to derive $I(n+1)$. Lemma~\ref{lem:curvatureidentity} gives $(n+1)\,\xi_{n+1}=(k+1)\,\delta_n+2\,\beta^{(k-1)}_n\ge 0$, because $\delta_n\ge 0$ and $\beta^{(k-1)}_n\ge 0$. Each suffix sum in $I(n+1)$ is the nonnegative term $\xi_{n+1}\,x^{n+1}$, alone or added to a suffix sum in $I(n)$. Hence $I(n+1)$ holds.
\end{proof}

\begin{proposition}\label{prop:lowcolumns}
Let $k$ and $n$ be integers with $2\le k\le 81$, $2k\le n$ and $17\le n$. Then $b_{n,k-1}\le b_{n,k}$.
\end{proposition}

\begin{proof}
Let $n_0=n_0(k)$ be the number of Definition~\ref{def:startrow}. Lemma~\ref{lem:columntable} supplies three facts: we have $2k\le n_0$; the inequality $b_{m,k-1}\le b_{m,k}$ holds for every integer $m$ with $2k\le m$, $17\le m$ and $m<n_0$; and $\sum_{t=s}^{n_0}\xi^{(k)}_t\,x_k^{t}\ge 0$ for every integer $s$ with $k\le s\le n_0$. If $n<n_0$, the claim is the second of these facts. If $n\ge n_0$, we apply Proposition~\ref{prop:propagation} with this $n_0$ and with $x=x_k$: its hypothesis~(1) is Proposition~\ref{prop:criticalmonotone}, its hypothesis~(2) is the third fact, and its conclusion $\delta^{(k)}_n\ge 0$ is the claim, because $\delta^{(k)}_n=(b_{n,k}-b_{n,k-1})/n!$.
\end{proof}

\subsection{The first column and the conclusion}\label{subsec:left-proof}

The preceding arguments involve the three thresholds $k+1$, $k$ and $k-1$ of a column and need $k\ge 2$. For $k=1$ we compare the last-car recursions of $A_{n,0}$ and $A_{n,1}$ term by term.

\begin{lemma}\label{lem:firstcolumn}
Let $n\ge 3$ be an integer. Then $b_{n,0}\le b_{n,1}$.
\end{lemma}

\begin{proof}
Since $b_{n,0}=A_{n,0}$ and $b_{n,1}=A_{n,1}-A_{n,0}$, the claim is that $2A_{n,0}\le A_{n,1}$, and we prove this by induction on $n\ge 3$. In the last-car recursion of Proposition~\ref{prop:lastcar} the factor $\min(j+1,\ell+1)$ is $1$ for $\ell=0$, and for $\ell=1$ it is $1$ for $j=0$ and $2$ for $j\ge 1$. Since $A_{0,0}=A_{0,1}=1$, the recursion gives, for every $n\ge 0$,
\begin{align*}
2A_{n+1,0}&=2A_{n,0}+\sum_{j=1}^{n}2\binom{n}{j}A_{j,0}\,A_{n-j,0},\\
A_{n+1,1}&=A_{n,1}+\sum_{j=1}^{n}2\binom{n}{j}A_{j,1}\,A_{n-j,1}.
\end{align*}
From these we find $A_{3,0}=6$ and $A_{3,1}=13$, so that the claim holds for $n=3$. If $2A_{n,0}\le A_{n,1}$ for some $n\ge 3$, then the right-hand side of the first line is at most that of the second, term by term: for the first terms this is the induction hypothesis, and for the others it holds because $A_{m,0}\le A_{m,1}$ for every $m\ge 0$. Hence $2A_{n+1,0}\le A_{n+1,1}$.
\end{proof}

The three ranges of columns together give Proposition~\ref{prop:leftpart}.

\begin{proof}[Proof of Proposition~\ref{prop:leftpart}]
For $k=1$ the claim is Lemma~\ref{lem:firstcolumn}, for $2\le k\le 81$ it is Proposition~\ref{prop:lowcolumns}, and for $k\ge 82$ it follows from Proposition~\ref{prop:largecolumns}.
\end{proof}

\section{Compositions weighted by Cayley numbers}\label{sec:comps}
The proof of Proposition~\ref{prop:nextcolumn}, in Sections~\ref{sec:chains} and~\ref{sec:next}, is written in terms of sums over the compositions of an integer, in which a part of size $p$ carries the Cayley coefficient $a_p$ as its weight and each cut carries a factor that depends on its position. We introduce these sums and prove three rules that compare the total weight of the compositions with $q+1$ parts with that of the compositions with $q$ parts: an identity when the cuts carry no factor, which shows that the number of parts then follows the law of the birthday problem; an inequality when each cut carries a factor between $0$ and $1$; and an inequality in the opposite direction for the factors that the parking problem produces.

\subsection{Compositions and their weights}\label{subsec:comps-weights}

Let $m\ge 0$ and $q\ge 1$ be integers. A \emph{composition} of $m$ with $q$ \emph{parts} is a sequence $\mathbf p=(p_1,\dots,p_q)$ of positive integers with $p_1+\dots+p_q=m$. Its \emph{cut points}, or \emph{cuts}, are $s_j=p_1+\dots+p_j$ for $1\le j\le q-1$. We also put $s_0=0$ and $s_q=m$, so that
\[
0=s_0<s_1<\dots<s_{q-1}<s_q=m\qquad\text{and}\qquad p_j=s_j-s_{j-1}\quad(1\le j\le q).
\]
Thus a composition of $m$ with $q$ parts is a way of cutting a segment of length $m$ at $q-1$ of its interior integer points $1,\dots,m-1$; the parts are the lengths of the pieces, and $s_j$ lies between $p_j$ and $p_{j+1}$. We write $\operatorname{Comp}_q(m)$ for the set of compositions of $m$ with $q$ parts. This set is empty if $q>m$, and in particular if $m=0$.

A \emph{cut profile} is a real-valued function $f$ on the nonnegative integers. Recall from Section~\ref{sec:prelim} the Cayley coefficients $a_0=0$ and $a_p=p^{p-1}/p!$ for $p\ge 1$. The \emph{weight} of a composition $\mathbf p\in\operatorname{Comp}_q(m)$ for the cut profile $f$ is
\[
\operatorname{wt}_f(\mathbf p)=\prod_{j=1}^{q}a_{p_j}\ \prod_{j=1}^{q-1}f(s_j),
\]
the product of the Cayley coefficients of its parts and of the values of $f$ at its cut points. For instance, the compositions of $3$ are $(3)$, $(1,2)$, $(2,1)$ and $(1,1,1)$, and since $a_1=a_2=1$ and $a_3=\frac32$ their weights are $\frac32$, $f(1)$, $f(2)$ and $f(1)\,f(2)$.

\begin{definition}\label{def:compositions}
Let $f$ be a cut profile. For integers $q\ge 0$ and $m\ge 0$ the real number $V_q(m;f)$ is defined by recursion on $q$:
\begin{gather*}
V_0(m;f)=0,\qquad V_1(m;f)=a_m,\\
V_{q+2}(m;f)=\sum_{0<s<m}V_{q+1}(s;f)\,f(s)\,a_{m-s},
\end{gather*}
where the sum runs over all integers $s$ with $0<s<m$. Unfolding the recursion gives, for $q\ge 1$,
\[
V_q(m;f)=\sum_{0=s_0<s_1<\dots<s_{q-1}<s_q=m}\ \prod_{i=1}^{q}a_{s_i-s_{i-1}}\ \prod_{i=1}^{q-1}f(s_i),
\]
where the sum runs over the sequences of integers $s_1<\dots<s_{q-1}$ that lie strictly between $s_0=0$ and $s_q=m$. Thus $V_q(m;f)$ is the total weight of the compositions of $m$ with exactly $q$ parts for the cut profile $f$, in which a part of size $p$ weighs $a_p$ and a cut at position $s$ weighs $f(s)$.
\end{definition}

In the notation introduced above, the second formula of the definition reads
\begin{equation}\label{eq:comps-total}
V_q(m;f)=\sum_{\mathbf p\in\operatorname{Comp}_q(m)}\operatorname{wt}_f(\mathbf p)\qquad(q\ge 1,\ m\ge 0).
\end{equation}
It follows from the recursion by induction on $q$: the case $q=1$ is clear, since $a_0=0$, and a composition of $m$ with $q+2$ parts consists of a composition of some integer $s$ with $0<s<m$ into $q+1$ parts, of a cut at $s$, and of a last part of size $m-s$.

By this formula we have $V_q(m;f)=0$ whenever $q>m$, and $V_q(m;f)$ depends on the profile only through its values $f(s)$ with $0<s<m$, the possible cut points of a composition of $m$. Moreover, $a_p>0$ for $p\ge 1$, and therefore $V_q(m;f)\ge 0$ if $f(s)\ge 0$ for all $s$ with $0<s<m$.

We denote by $\mathbf 1$ the constant profile with value $1$, for which the cuts carry no factor, and we write
\[
T_{m,q}=V_q(m;\mathbf 1)\qquad(m\ge 0,\ q\ge 0).
\]
For $q\ge 1$ this is the sum of the products $a_{p_1}\cdots a_{p_q}$ over the $q$-tuples $(p_1,\dots,p_q)$ of positive integers with sum $m$, and $T_{m,0}=0$. In particular $T_{m,1}=a_m$, and $T_{m,q}>0$ for $1\le q\le m$, since there is a composition of $m$ with $q$ parts and the weights are positive. More generally, we denote a constant profile by its value $h$, a real number. A composition with $q$ parts has $q-1$ cuts, so that
\begin{equation}\label{eq:comps-constant}
V_q(m;h)=h^{q-1}\,T_{m,q}\qquad(q\ge 1).
\end{equation}

For a cut profile $f$ and an integer $m\ge 0$ we write
\[
\Lambda_m[f]=\sum_{q=1}^{m}V_q(m;f)
\]
for the total weight of all compositions of $m$. The sum stops at $q=m$ because there are no compositions with more parts, and $\Lambda_0[f]=0$.

For real numbers $y$ and $N$ the \emph{threshold profile} $\lambda^{y,N}$ is the cut profile
\[
\lambda^{y,N}(s)=\frac{y}{N-s},
\]
where we put $\lambda^{y,N}(N)=0$ if $N$ is a nonnegative integer; wherever we use the profile, the cut points are smaller than $N$, so that this value plays no role. The total weight of all compositions of $m$ for the threshold profile is written
\[
\Lambda_m(y;N)=\Lambda_m[\lambda^{y,N}]
\]
and called a \emph{chain sum}. Chain sums are what the parking problem produces. In the rows that are considered in the proof of Proposition~\ref{prop:nextcolumn}, the normalised counts of a threshold $y$ satisfy a linear equation, and solving it by iteration, as we do in Section~\ref{sec:chains}, leads to sums over decreasing chains of lengths $N>t_1>t_2>\cdots$ in which every step from $t'$ down to $t$ contributes $a_{t'-t}$ and every visited length $t$ contributes a factor $y/t$. Such a chain is a composition of the distance that it covers: the steps are the parts, and a visited length $t$ is a cut at position $s=N-t$, with factor $y/(N-s)$. Since $y$ enters the factor of every cut once, and a composition with $q$ parts has $q-1$ cuts, we have
\begin{equation}\label{eq:comps-scaling}
V_q(m;\lambda^{y,N})=y^{q-1}\,V_q(m;\lambda^{1,N})\qquad(q\ge 1).
\end{equation}

\subsection{Erasing a cut}\label{subsec:comps-erasure}

We compare $V_{q+1}(m;f)$ with $V_q(m;f)$ by erasing a cut: erasing a cut of a composition with $q+1$ parts merges the two parts next to it into one part and leaves a composition with $q$ parts. When the cuts carry no factor, the comparison is an identity.

\begin{lemma}\label{lem:erasure}
For all integers $m\ge 0$ and $q\ge 1$,
\[
q\,m\,T_{m,q+1}=(q+1)(m-q)\,T_{m,q}.
\]
\end{lemma}

\begin{proof}
We prove an identity for an arbitrary cut profile $f$, the identity \eqref{eq:comps-defect} below, and obtain the lemma as its case $f=\mathbf 1$; the two inequalities of this section come from the same identity.

Let $m\ge 0$ and $q\ge 1$ be integers and $f$ a cut profile. For a composition $\mathbf p=(p_1,\dots,p_{q+1})$ of $m$ with $q+1$ parts, with cut points $s_1<\dots<s_q$, and for $1\le i\le q$, we put
\[
\operatorname{mult}_i(\mathbf p)=(q+1-i)\,p_i+i\,p_{i+1},
\]
which is a positive integer, and we write
\[
\operatorname{wt}_f^{(i)}(\mathbf p)=\prod_{j=1}^{q+1}a_{p_j}\ \prod_{\substack{1\le l\le q\\ l\ne i}}f(s_l)
\]
for the weight of $\mathbf p$ from which the factor of the cut $s_i$ has been removed, so that
\[
\operatorname{wt}_f(\mathbf p)=f(s_i)\,\operatorname{wt}_f^{(i)}(\mathbf p)\qquad(1\le i\le q).
\]
We count in two ways the pairs $(\mathbf p,i)$ formed by a composition $\mathbf p$ of $m$ with $q+1$ parts and one of its cuts $s_i$: with the weight of $\mathbf p$, and with the weight from which the factor of the chosen cut has been removed. In both counts the pair $(\mathbf p,i)$ is counted with the multiplicity $\operatorname{mult}_i(\mathbf p)$, whose two coefficients are chosen so that each count gives a factor that does not depend on the composition.

For the first count we note that, for every composition $\mathbf p$ of $m$ with $q+1$ parts,
\begin{equation}\label{eq:comps-first}
\sum_{i=1}^{q}\operatorname{mult}_i(\mathbf p)=q\,m.
\end{equation}
Indeed, in this sum $p_j$ occurs with coefficient $q+1-j$ as the left neighbour of $s_j$ and with coefficient $j-1$ as the right neighbour of $s_{j-1}$; for the first and for the last part, which have only one neighbouring cut, the missing coefficient is $0$. So every part occurs with coefficient $q$, and the parts add up to $m$. We multiply this identity by $\operatorname{wt}_f(\mathbf p)$, sum over $\mathbf p$, and obtain by \eqref{eq:comps-total}
\begin{equation}\label{eq:comps-full}
\sum_{\mathbf p\in\operatorname{Comp}_{q+1}(m)}\ \sum_{i=1}^{q}\operatorname{mult}_i(\mathbf p)\,\operatorname{wt}_f(\mathbf p)=q\,m\,V_{q+1}(m;f).
\end{equation}

In the second count the factor of the chosen cut is removed from the weight. The same pairs then weigh
\begin{equation}\label{eq:comps-erased}
\sum_{\mathbf p\in\operatorname{Comp}_{q+1}(m)}\ \sum_{i=1}^{q}\operatorname{mult}_i(\mathbf p)\,\operatorname{wt}_f^{(i)}(\mathbf p)=(q+1)(m-q)\,V_q(m;f).
\end{equation}
To see this, we sort the pairs by the composition that is left when the chosen cut is erased. For a pair $(\mathbf p,i)$ this is the composition $\mathbf u=(u_1,\dots,u_q)$ of $m$ with $q$ parts in which $u_i=p_i+p_{i+1}$, while $u_j=p_j$ for $j<i$ and $u_j=p_{j+1}$ for $j>i$. Its cut points are the $s_l$ with $l\ne i$. The pair is determined by $\mathbf u$, by $i$, and by $v=p_i$, the size of the left one of the two pieces into which $s_i$ divides $u_i$; and conversely, every composition $\mathbf u\in\operatorname{Comp}_q(m)$, every $i$ with $1\le i\le q$ and every integer $v$ with $1\le v\le u_i-1$ belong to exactly one pair. In these terms $\operatorname{mult}_i(\mathbf p)=(q+1-i)\,v+i\,(u_i-v)$, and
\[
\operatorname{wt}_f^{(i)}(\mathbf p)=\frac{a_v\,a_{u_i-v}}{a_{u_i}}\,\operatorname{wt}_f(\mathbf u),
\]
because $\mathbf p$ and $\mathbf u$ have the same parts apart from $u_i$, which is replaced in $\mathbf p$ by its two pieces, and because the cuts of $\mathbf u$ are the cuts of $\mathbf p$ other than $s_i$. We now use the convolution identity of the Cayley coefficients: for every integer $u\ge 1$,
\[
\sum_{v=1}^{u-1}v\,a_v\,a_{u-v}=(u-1)\,a_u=\sum_{v=1}^{u-1}(u-v)\,a_v\,a_{u-v}.
\]
The first equality is Lemma~\ref{lem:cayleyconvolution}, and the second follows from the first on replacing $v$ by $u-v$. For a fixed composition $\mathbf u$ and a fixed $i$, the pairs that belong to them therefore contribute to the second count the number
\[
\frac{\operatorname{wt}_f(\mathbf u)}{a_{u_i}}\sum_{v=1}^{u_i-1}\bigl((q+1-i)\,v+i\,(u_i-v)\bigr)\,a_v\,a_{u_i-v}=(q+1)\,(u_i-1)\,\operatorname{wt}_f(\mathbf u),
\]
by the two equalities of the convolution identity for $u=u_i$. Since $\sum_{i=1}^{q}(u_i-1)=m-q$, the pairs that belong to $\mathbf u$ contribute $(q+1)(m-q)\,\operatorname{wt}_f(\mathbf u)$ in all, and summing over $\mathbf u\in\operatorname{Comp}_q(m)$ gives the value $(q+1)(m-q)\,V_q(m;f)$ of the second count, by \eqref{eq:comps-total}.

We now compare the two counts. Since $\operatorname{wt}_f(\mathbf p)=f(s_i)\,\operatorname{wt}_f^{(i)}(\mathbf p)$, the left-hand sides of \eqref{eq:comps-full} and of \eqref{eq:comps-erased} differ only in the factor $f(s_i)$ of the chosen cut. Subtracting the first of these identities from the second, we obtain
\begin{multline}\label{eq:comps-defect}
(q+1)(m-q)\,V_q(m;f)-q\,m\,V_{q+1}(m;f)\\
=\sum_{\mathbf p\in\operatorname{Comp}_{q+1}(m)}\ \sum_{i=1}^{q}\operatorname{mult}_i(\mathbf p)\,\bigl(1-f(s_i)\bigr)\,\operatorname{wt}_f^{(i)}(\mathbf p).
\end{multline}
This identity holds for all integers $m\ge 0$ and $q\ge 1$ and for every cut profile $f$. For the constant profile $f=\mathbf 1$ its right-hand side vanishes, and since $V_q(m;\mathbf 1)=T_{m,q}$ by definition, this is the identity of the lemma.
\end{proof}

\subsection{The birthday law}\label{subsec:comps-birthday}

For integers $m\ge 1$ and $j\ge 0$ let
\[
P_{m,j}=\prod_{i=0}^{j-1}\Bigl(1-\frac{i}{m}\Bigr)=\frac{m\,(m-1)\cdots(m-j+1)}{m^{j}} .
\]
We give $P_{m,j}$ no value for $m=0$. We have $P_{m,0}=P_{m,1}=1$, the definition gives the recursion
\[
P_{m,j+1}=P_{m,j}\,\Bigl(1-\frac{j}{m}\Bigr)\qquad(j\ge 0),
\]
and $P_{m,j}=0$ for $j\ge m+1$, because the product then contains the factor with $i=m$. The number $P_{m,j}$ is the probability that $j$ independent draws from a set of $m$ values, each draw uniformly distributed, are pairwise distinct. The \emph{birthday law} of $m$ is the law of the number $Q$ of draws that are made before the first draw that repeats an earlier value. Since $Q\ge j$ holds exactly when the first $j$ draws are pairwise distinct, we have $\Pr(Q\ge j)=P_{m,j}$ for every $j\ge 1$. Thus $Q$ takes the values $1,\dots,m$, and by the recursion for $P_{m,j}$
\[
\Pr(Q=q)=P_{m,q}-P_{m,q+1}=\frac{q}{m}\,P_{m,q}\qquad(q\ge 1).
\]

Since $T_{m,1}=a_m$, Lemma~\ref{lem:erasure} determines $T_{m,q}$ for all $q\ge 1$ when $m\ge 1$, and the result is expressed by the numbers $P_{m,q}$.

\begin{proposition}\label{prop:birthday}
Let $m\ge 0$ and $q\ge 1$ be integers. If $m\ge 1$, then
\[
T_{m,q}=a_m\,q\,P_{m,q}.
\]
If $m=0$, then $T_{0,q}=0$.
\end{proposition}

\begin{proof}
There is no composition of $0$, and hence $T_{0,q}=0$ by \eqref{eq:comps-total}. Let $m\ge 1$. We prove the formula by induction on $q$. For $q=1$ we have $T_{m,1}=a_m$ and $P_{m,1}=1$. Suppose that the formula holds for some $q\ge 1$. By Lemma~\ref{lem:erasure}, by the induction hypothesis, and by the recursion for $P_{m,q}$ in the form $(m-q)\,P_{m,q}=m\,P_{m,q+1}$,
\begin{align*}
q\,m\,T_{m,q+1}&=(q+1)(m-q)\,T_{m,q}=(q+1)(m-q)\,a_m\,q\,P_{m,q}\\
&=q\,m\cdot a_m\,(q+1)\,P_{m,q+1},
\end{align*}
and dividing by $q\,m$ gives the formula for $q+1$.
\end{proof}

Because $\Pr(Q=q)=\frac qm\,P_{m,q}$, the proposition says that, for $m\ge 1$, the number $T_{m,q}$ is $m\,a_m$ times the probability that $Q=q$. Since these probabilities add up to $1$, the compositions of $m$ have total weight $\Lambda_m[\mathbf 1]=m\,a_m$, and if one of them is drawn with probability proportional to its weight for the profile $\mathbf 1$, its number of parts follows the birthday law of $m$. For $1\le q\le m$, where $a_m=m^{m-1}/m!$ and $P_{m,q}=m!/((m-q)!\,m^{q})$, the formula reads $T_{m,q}=q\,m^{\,m-q-1}/(m-q)!$. This is the value that the Lagrange inversion formula gives for the coefficient of $X^m$ in the $q$-th power of the tree function $XC=\sum_{j\ge 1}a_jX^j$ of Section~\ref{sec:prelim}. We derive it from Lemma~\ref{lem:erasure} because the identity behind that lemma is needed in any case, for the profiles that are not constant.

\subsection{Cut factors that depend on the position}\label{subsec:comps-profiles}

We return to a general cut profile. If every cut carries a factor between $0$ and $1$, the identity of Lemma~\ref{lem:erasure} becomes an inequality: such factors, whatever their dependence on the position, can only lower the weight of the compositions with $q+1$ parts relative to the weight of those with $q$ parts, in comparison with the profile $\mathbf 1$.

\begin{lemma}\label{lem:erasureupper}
Let $m\ge 0$ be an integer, and let $f$ be a cut profile such that $0\le f(s)\le 1$ for every integer $s$ with $0<s<m$. Then, for every integer $q\ge 1$,
\[
q\,m\,V_{q+1}(m;f)\le(q+1)(m-q)\,V_q(m;f).
\]
\end{lemma}

\begin{proof}
We use the notation of the proof of Lemma~\ref{lem:erasure}. By the identity \eqref{eq:comps-defect} it suffices to show that every term on its right-hand side is nonnegative. The numbers $\operatorname{mult}_i(\mathbf p)$ and the Cayley coefficients of the parts are positive. The cut points of a composition of $m$ lie strictly between $0$ and $m$, so that the hypothesis applies to each of them: it gives $f(s_l)\ge 0$ for the cuts with $l\ne i$, whose factors occur in $\operatorname{wt}_f^{(i)}(\mathbf p)$, and $1-f(s_i)\ge 0$ for the erased cut.
\end{proof}

For a constant profile $h>0$, the formula \eqref{eq:comps-constant} and Lemma~\ref{lem:erasure} give
\[
(q+1)(m-q)\,V_q(m;h)=q\,m\,\frac{1}{h}\,V_{q+1}(m;h)\qquad(q\ge 1).
\]
This shows that the upper bound in Lemma~\ref{lem:erasureupper} cannot be dropped: for $1\le q<m$, where $V_{q+1}(m;h)=h^{q}\,T_{m,q+1}>0$, the inequality of that lemma holds for the constant profile $h$ if and only if $h\le 1$.

The second inequality concerns the threshold profile $\lambda^{1,N}$ and goes in the opposite direction: if $m<N$, the weight of the compositions of $m$ with $q+1$ parts, relative to the weight of those with $q$ parts, is at least what the identity displayed above gives for the constant profile $h=1/(N-\frac m2)$, the value of $1/(N-s)$ at the midpoint $s=\frac m2$ of the segment that is cut. By the scaling relation \eqref{eq:comps-scaling} the same then holds for $\lambda^{y,N}$ with any $y>0$ and $h=y/(N-\frac m2)$.

\begin{lemma}\label{lem:midpoint}
Let $m\ge 0$ be an integer and $N$ a real number with $m<N$. Then the threshold profile $\lambda^{1,N}$ satisfies, for every integer $q\ge 1$,
\[
(q+1)(m-q)\,V_q(m;\lambda^{1,N})\le q\,m\,\Bigl(N-\frac{m}{2}\Bigr)\,V_{q+1}(m;\lambda^{1,N}).
\]
\end{lemma}

\begin{proof}
We use the notation of the proof of Lemma~\ref{lem:erasure}, and we write $f=\lambda^{1,N}$ and $V_q=V_q(m;f)$. Both sides vanish if $q\ge m$, and we assume from now on that $q<m$.

Every cut point $s$ of a composition of $m$ satisfies $s<m<N$. At the cut points we therefore have
\[
f(s)>0,\qquad (N-s)\,f(s)=1,\qquad f(s)\le f(s')\ \text{ for cut points } s\le s'.
\]

We first express the difference of the two sides of the lemma through the parts. Let $\mathbf p$ be a composition of $m$ with $q+1$ parts and let $1\le i\le q$. Then
\[
\operatorname{wt}_f^{(i)}(\mathbf p)=(N-s_i)\,f(s_i)\,\operatorname{wt}_f^{(i)}(\mathbf p)=(N-s_i)\,\operatorname{wt}_f(\mathbf p),
\]
by the relation $(N-s_i)\,f(s_i)=1$ and because $\operatorname{wt}_f^{(i)}(\mathbf p)$ is the weight of $\mathbf p$ without the factor $f(s_i)$. We insert this into \eqref{eq:comps-erased}, and we subtract the result from \eqref{eq:comps-full} multiplied by $N-\frac m2$. This gives
\begin{multline*}
q\,m\,\Bigl(N-\frac{m}{2}\Bigr)\,V_{q+1}-(q+1)(m-q)\,V_q\\
=\sum_{\mathbf p\in\operatorname{Comp}_{q+1}(m)}\operatorname{wt}_f(\mathbf p)\,\sum_{i=1}^{q}\operatorname{mult}_i(\mathbf p)\,\Bigl(s_i-\frac{m}{2}\Bigr).
\end{multline*}
The inner sum can be written in terms of the parts alone: for every composition $\mathbf p$ of $m$ with $q+1$ parts,
\[
\sum_{i=1}^{q}\operatorname{mult}_i(\mathbf p)\,\Bigl(s_i-\frac{m}{2}\Bigr)=\frac12\sum_{j=1}^{q+1}(q+2-2j)\,p_j^{2}.
\]
To prove this, we argue as for the identity \eqref{eq:comps-first}. In the sum $\sum_{i=1}^{q}\operatorname{mult}_i(\mathbf p)\,s_i$ the part $p_j$ occurs with coefficient $(q+1-j)\,s_j$ as the left neighbour of $s_j$ and with coefficient $(j-1)\,s_{j-1}$ as the right neighbour of $s_{j-1}$, where $s_0=0$ and $s_{q+1}=m$. Since
\[
(q+1-j)\,s_j+(j-1)\,s_{j-1}=\frac{q}{2}\,(s_j+s_{j-1})+\frac{q+2-2j}{2}\,(s_j-s_{j-1})
\]
and $p_j=s_j-s_{j-1}$, so that $p_j\,(s_j+s_{j-1})=s_j^{2}-s_{j-1}^{2}$, we get
\[
\sum_{i=1}^{q}\operatorname{mult}_i(\mathbf p)\,s_i=\frac{q}{2}\,m^{2}+\frac12\sum_{j=1}^{q+1}(q+2-2j)\,p_j^{2},
\]
the sum of the differences $s_j^{2}-s_{j-1}^{2}$ being telescopic. By \eqref{eq:comps-first} the term $\frac q2\,m^{2}$ equals $\frac m2\sum_{i=1}^{q}\operatorname{mult}_i(\mathbf p)$, and the formula for the inner sum follows. For $1\le j\le q+1$ let
\[
E_j=\sum_{\mathbf p\in\operatorname{Comp}_{q+1}(m)}p_j^{2}\,\operatorname{wt}_f(\mathbf p)
\]
be the total weight of the compositions of $m$ with $q+1$ parts when each of them is counted with the square of its $j$-th part. Then the formula for the difference of the two sides of the lemma, with its inner sum written in terms of the parts, becomes
\[
q\,m\,\Bigl(N-\frac{m}{2}\Bigr)\,V_{q+1}-(q+1)(m-q)\,V_q=\frac12\sum_{j=1}^{q+1}(q+2-2j)\,E_j .
\]
It remains to show that the sum on the right is nonnegative. Its coefficients $q+2-2j$ are positive for the early parts and negative for the late ones, and they add up to $0$.

We therefore sum by parts. For arbitrary real numbers $E_1,\dots,E_{q+1}$,
\[
\sum_{j=1}^{q+1}(q+2-2j)\,E_j=\sum_{j=1}^{q}j\,(q+1-j)\,\bigl(E_j-E_{j+1}\bigr),
\]
since on the right-hand side $E_j$ has coefficient $j\,(q+\nobreak 1-\nobreak j)-(j-\nobreak 1)\,(q+\nobreak 2-\nobreak j)=q+2-2j$. The coefficients $j\,(q+1-j)$ with $1\le j\le q$ are positive. It is therefore enough to prove that
\[
E_{j+1}\le E_j\qquad(1\le j\le q).
\]

For the proof, let $1\le j\le q$. By definition,
\[
E_j-E_{j+1}=\sum_{\mathbf p\in\operatorname{Comp}_{q+1}(m)}\bigl(p_j^{2}-p_{j+1}^{2}\bigr)\,\operatorname{wt}_f(\mathbf p).
\]
We sort the compositions by their cut points other than $s_j$, and we show that the compositions with given points $s_l$, $l\ne j$, contribute a nonnegative number to this sum. Among the given points are $s_{j-1}$ and $s_{j+1}$, where again $s_0=0$ and $s_{q+1}=m$. Let $t=s_{j+1}-s_{j-1}$ be the length of the interval that $p_j$ and $p_{j+1}$ fill together. The compositions in question are obtained by placing the remaining cut at $s_j=s_{j-1}+v$ with $1\le v\le t-1$, and then $p_j=v$ and $p_{j+1}=t-v$. The weight of such a composition is $f(s_{j-1}+v)\,a_v\,a_{t-v}$ times the product of the Cayley coefficients of the other parts and of the factors of the other cuts. This product is the same for all the compositions in question, and it is positive, because $f$ is positive at the cut points. So it suffices to show that
\[
\sum_{v=1}^{t-1}f(s_{j-1}+v)\,a_v\,a_{t-v}\,\bigl(v^{2}-(t-v)^{2}\bigr)\ge 0 .
\]
Replacing $v$ by $t-v$ reflects the cut about the midpoint of the interval and exchanges the sizes of the two parts; it leaves $a_v\,a_{t-v}$ unchanged and changes the sign of $v^{2}-(t-v)^{2}$. Hence the sum on the left is also equal to the negative of the same sum with $f(s_{j-1}+t-v)$ in the place of $f(s_{j-1}+v)$, and taking the mean of the two expressions we find that it equals
\[
\frac12\sum_{v=1}^{t-1}\bigl(f(s_{j-1}+v)-f(s_{j-1}+t-v)\bigr)\,a_v\,a_{t-v}\,\bigl(v^{2}-(t-v)^{2}\bigr).
\]
Every term of this sum is nonnegative. Indeed, $s_{j-1}+v$ and $s_{j-1}+t-v$ are possible positions of the cut $s_j$, so that the monotonicity of the profile applies to them. If $v\ge t-v$, then $v^{2}-(t-v)^{2}\ge 0$ and $f(s_{j-1}+v)\ge f(s_{j-1}+t-v)$. If $v<t-v$, both of these factors are nonpositive. This proves that $E_{j+1}\le E_j$, and with it the lemma.
\end{proof}

\section{The birthday law and the homogeneous law}\label{sec:birthday}
By Proposition~\ref{prop:birthday} the number of parts of a composition of $m$, drawn with a probability proportional to its weight, follows the birthday law of $m$. For this law we bound the mean $s_m$ and the ratio $D_m$ of the third to the second factorial moment. We then give every cut the same factor $h\ge 0$; for the law that results, the homogeneous law, we bound the mean number of cuts from below and compare the total weights of the compositions of $r+1$ and of $r$. These are the estimates on the number of parts that enter the proof of Proposition~\ref{prop:nextcolumn} in Sections~\ref{sec:chains} and~\ref{sec:next}.

\subsection{A correlation inequality}\label{subsec:birthday-correlation}

We begin with the inequality by which two laws of the number of parts are compared, here and in Section~\ref{sec:chains}. Let $u$ and $v$ be positive weights on a finite set $S$ of integers such that the quotient $u/v$ does not increase on $S$; one says that $u$ is smaller than $v$ in the \emph{likelihood-ratio order}. In this case the average of a nondecreasing function $\varphi$ on $S$ with respect to $u$ is at most its average with respect to $v$. This is the classical inequality of Chebyshev for the two oppositely ordered sequences $u/v$ and $\varphi$ with weights $v$, and it is of the kind that Fortuin, Kasteleyn and Ginibre prove for a log-supermodular measure on a finite distributive lattice, under which two nondecreasing functions are nonnegatively correlated \cite[Proposition~1]{fortuin1971}. Since the weights to which we apply it vanish at some points of $S$, we state it with the hypothesis and the conclusion cross-multiplied; in this form no condition is imposed on the signs of $u$, $v$ and $\varphi$, and the usual proof by symmetrisation applies.

\begin{lemma}\label{lem:chebyshev}
Let $S$ be a finite set of nonnegative integers and let $u$, $v$ and $\varphi$ be real-valued functions on the nonnegative integers. Suppose that
\[
u(q')\,v(q)\le u(q)\,v(q')\qquad\text{and}\qquad\varphi(q)\le\varphi(q')
\]
for all $q,q'\in S$ with $q\le q'$. Then
\[
\Bigl(\sum_{q\in S}u(q)\,\varphi(q)\Bigr)\Bigl(\sum_{q\in S}v(q)\Bigr)\le\Bigl(\sum_{q\in S}v(q)\,\varphi(q)\Bigr)\Bigl(\sum_{q\in S}u(q)\Bigr).
\]
\end{lemma}

\begin{proof}
Let $Y$ be the left-hand side minus the right-hand side. Expanding the products of the sums and symmetrising in the two indices of summation, we obtain
\[
2Y=\sum_{q\in S}\sum_{q'\in S}\bigl(u(q)\,v(q')-u(q')\,v(q)\bigr)\bigl(\varphi(q)-\varphi(q')\bigr).
\]
By the two hypotheses, applied to the pair $(q,q')$ if $q\le q'$ and to the pair $(q',q)$ otherwise, one of the two factors of each term is nonnegative and the other is nonpositive. Hence $Y\le 0$.
\end{proof}

\subsection{The moments of the birthday law}\label{subsec:birthday-moments}

Let $m\ge 1$ be an integer. Recall from Section~\ref{sec:comps} the number
\[
P_{m,j}=\prod_{i=0}^{j-1}\Bigl(1-\frac{i}{m}\Bigr)=\frac{m(m-1)\cdots(m-j+1)}{m^{j}}\qquad(j\ge 0),
\]
and the birthday law of $m$, the law of the number $Q$ defined there, for which $\Pr(Q\ge j)=P_{m,j}$ for $j\ge 1$. Clearly
\begin{equation}\label{eq:birthday-tail}
\begin{gathered}
P_{m,0}=P_{m,1}=1,\qquad P_{m,j+1}=\Bigl(1-\frac{j}{m}\Bigr)P_{m,j}\quad(j\ge 0),\\
P_{m,j}=0\quad(j\ge m+1),\qquad 0\le P_{m,j}\le 1\quad(j\ge 0),
\end{gathered}
\end{equation}
so that $Q$ takes the value $q\ge 1$ with probability $P_{m,q}-P_{m,q+1}=\frac qm\,P_{m,q}$.

The sums of this section are evaluated by a summation by parts: for every real-valued function $\psi$ on the nonnegative integers,
\begin{equation}\label{eq:birthday-abel}
\sum_{q=1}^{m}\psi(q)\,q\,P_{m,q}=m\,\Bigl(\psi(0)+\sum_{q=1}^{m}\bigl(\psi(q)-\psi(q-1)\bigr)\,P_{m,q}\Bigr).
\end{equation}
Indeed, $q\,P_{m,q}=m\,(P_{m,q}-P_{m,q+1})$, as we have just seen, and it remains to shift the index in the terms that contain $P_{m,q+1}$, using $P_{m,1}=1$ and $P_{m,m+1}=0$. Divided by $m$, the identity expresses the expectation $\mathbb{E}[\psi(Q)]$ through the tail probabilities $\Pr(Q\ge j)$.

The mean of the birthday law of $m$ is the sum of the tail probabilities,
\[
s_m=\sum_{j=1}^{m}P_{m,j}\qquad(m\ge 1),
\]
and we put $s_0=0$; note that $s_1=1$. Each of the $m$ terms of $s_m$ is at most $1$, and for $m\ge 2$ the term $P_{m,2}=1-\frac1m$ is smaller; hence
\[
s_m<m\qquad(m\ge 2).
\]
Applying the summation by parts to $\psi(j)=1$, $j$, $j(j-1)$ and $j(j-1)(j-2)$ in turn, we obtain
\begin{equation}\label{eq:birthday-moments}
\begin{aligned}
\sum_{q=1}^{m}q\,P_{m,q}&=m,\\
\sum_{q=1}^{m}q^{2}\,P_{m,q}&=m\,s_m,\\
\sum_{q=1}^{m}q^{2}(q-1)\,P_{m,q}&=2m\,(m-s_m),\\
\sum_{q=1}^{m}q^{2}(q-1)(q-2)\,P_{m,q}&=3m\,\bigl((m+2)\,s_m-3m\bigr).
\end{aligned}
\end{equation}
For the last two functions the right-hand side of the summation by parts is evaluated with the first two identities and the definition of $s_m$.

By Proposition~\ref{prop:birthday} we have $\sum_{q=1}^{m}\psi(q)\,T_{m,q}=a_m\sum_{q=1}^{m}\psi(q)\,q\,P_{m,q}$ for every function $\psi$, so that these four identities describe the number of parts of the compositions of $m$. By the first identity the total weight of all compositions of $m$ is $m\,a_m$, and those with $q$ parts carry the fraction $\frac qm\,P_{m,q}=\Pr(Q=q)$ of it. By the second identity the mean number of parts is $s_m$. The last two identities give the second and the third factorial moment of the number of parts, $\mathbb{E}[Q(Q-1)]=2\,(m-s_m)$ and $\mathbb{E}[Q(Q-1)(Q-2)]=3\,((m+2)\,s_m-3m)$, or, in terms of the weights,
\begin{equation}\label{eq:birthday-factorial}
\begin{aligned}
\sum_{q=1}^{m}q(q-1)\,T_{m,q}&=2m\,a_m\,(m-s_m),\\
\sum_{q=1}^{m}q(q-1)(q-2)\,T_{m,q}&=3m\,a_m\,\bigl((m+2)\,s_m-3m\bigr).
\end{aligned}
\end{equation}

The quotient of these two moments enters the comparison of two second differences in Section~\ref{sec:chains}. For an integer $m\ge 2$ let
\[
D_m=\frac{3\,\bigl((m+2)\,s_m-3m\bigr)}{2\,(m-s_m)},
\]
where the denominator is positive because $s_m<m$. For $m=0$ and $m=1$ the formula reads $0/0$, since $s_0=0$ and $s_1=1$, and we set $D_0=D_1=0$. By the two identities \eqref{eq:birthday-factorial}, the definition of $D_m$, multiplied by $2m\,a_m\,(m-s_m)$, says that
\begin{equation}\label{eq:birthday-ratio}
\sum_{q=1}^{m}q(q-1)(q-2)\,T_{m,q}=D_m\sum_{q=1}^{m}q(q-1)\,T_{m,q}\qquad(m\ge 2).
\end{equation}
The following two lemmas give what else we use about $D_m$: a bound by the mean, and its monotonicity in $m$, which the formula does not show and which lets the one constant $D_r$ serve for all $m\le r$.

\begin{lemma}\label{lem:ratiobound}
For every integer $m\ge 2$,
\[
D_m\le\tfrac32\,(s_m-1).
\]
\end{lemma}

\begin{proof}
The sum $\sum_{q=1}^{m}(q-s_m)^2\,q\,P_{m,q}$, which is $m$ times the variance of the birthday law of $m$, is nonnegative, and a direct computation with the first three identities of \eqref{eq:birthday-moments} shows that it equals $m\,(2m-s_m-s_m^2)$. Hence $s_m^2+s_m\le 2m$. Now $m-s_m>0$ because $s_m<m$ for $m\ge 2$, so that the assertion of the lemma is equivalent to $(m+2)\,s_m-3m\le(s_m-1)(m-s_m)$, and the right-hand side of this inequality minus its left-hand side is $2m-s_m-s_m^2\ge 0$.
\end{proof}

\begin{lemma}\label{lem:ratiomonotone}
For all integers $m$ and $r$ with $2\le m\le r$,
\[
D_m\le D_r.
\]
\end{lemma}

\begin{proof}
By induction on $r$ it suffices to prove that $D_m\le D_{m+1}$ for every integer $m\ge 2$. We fix such an $m$.

We claim that for all integers $q$ and $q'$ with $1\le q\le q'\le m+1$,
\[
P_{m,q'}\,P_{m+1,q}\le P_{m,q}\,P_{m+1,q'}.
\]
For $q'=q$ the two sides are equal, and by the recursion \eqref{eq:birthday-tail} the passage from $q'=j$ to $q'=j+1$, where $q\le j\le m$, multiplies the left-hand side by $1-\frac jm$ and the right-hand side by $1-\frac{j}{m+1}$. Since $0\le 1-\frac jm\le 1-\frac{j}{m+1}$ and the right-hand side is nonnegative, the claim follows by induction on $q'$.

By Proposition~\ref{prop:birthday}, multiplying the inequality just proved by $a_m\,a_{m+1}\,q\,q'\ge 0$ gives $T_{m,q'}\,T_{m+1,q}\le T_{m,q}\,T_{m+1,q'}$ for $1\le q\le q'\le m+1$. The hypotheses of Lemma~\ref{lem:chebyshev} therefore hold for $S=\{1,\dots,m+1\}$ and
\[
u(q)=q(q-1)\,T_{m,q},\qquad v(q)=q(q-1)\,T_{m+1,q},\qquad\varphi(q)=q-2,
\]
and the lemma gives
\begin{multline*}
\Bigl(\sum_{q=1}^{m+1}q(q-1)(q-2)\,T_{m,q}\Bigr)\Bigl(\sum_{q=1}^{m+1}q(q-1)\,T_{m+1,q}\Bigr)\\
\le\Bigl(\sum_{q=1}^{m+1}q(q-1)(q-2)\,T_{m+1,q}\Bigr)\Bigl(\sum_{q=1}^{m+1}q(q-1)\,T_{m,q}\Bigr),
\end{multline*}
in which the two sums that contain $T_{m,q}$ end at $q=m$, since a composition of $m$ has at most $m$ parts.

Put $F_m=\sum_{q=1}^{m}q(q-1)\,T_{m,q}$. By \eqref{eq:birthday-ratio}, applied to $m$ and to $m+1$, the last inequality reads $D_m\,F_m\,F_{m+1}\le D_{m+1}\,F_{m+1}\,F_m$. Since $F_m=2m\,a_m\,(m-s_m)>0$ by the first identity of \eqref{eq:birthday-factorial} and $s_m<m$, and $F_{m+1}>0$ in the same way, we conclude that $D_m\le D_{m+1}$.
\end{proof}

\subsection{The mean of the birthday law}\label{subsec:birthday-mean}

Since $D_m$ is an explicit function of $m$ and $s_m$, the numerical information about the birthday law that we need is an upper bound for its mean. The sum $s_m=1+\frac{m-1}{m}+\frac{(m-1)(m-2)}{m^2}+\cdots$ is known as Ramanujan's $Q$-function, and our bound for it rests on a known inequality for the Poisson law. For an integer $m\ge 0$ and a real number $x\ge 0$ let
\[
\varpi_m(x)=e^{-x}\,\frac{x^m}{m!},\qquad H_m(x)=e^{-x}\sum_{i=0}^{m}\frac{x^i}{i!},
\]
the probability that a Poisson variable with mean $x$ takes the value $m$, resp.\ a value at most $m$; thus $\varpi_m(m)=e^{-m}m^m/m!$ is the Poisson mass at $m$ of Section~\ref{sec:prelim}. The inequality in question is
\begin{equation}\label{eq:birthday-third}
H_m(m)\le\frac12+\frac23\,\varpi_m(m)\qquad(m\ge 1).
\end{equation}

In the classical notation, in which a number $\theta_m$ is defined by the equation $\frac12\,e^{m}=\sum_{i<m}\frac{m^i}{i!}+\theta_m\,\frac{m^m}{m!}$, the inequality \eqref{eq:birthday-third} is the statement $\theta_m\ge\frac13$, since this equation, multiplied by $e^{-m}$, reads $\frac12=H_m(m)-(1-\theta_m)\,\varpi_m(m)$. That $\theta_m$ lies between $\frac13$ and $\frac12$ was asked by Ramanujan, and Flajolet, Grabner, Kirschenhofer and Prodinger, who recount the history of the question, prove more: $\theta_n=\frac13+\frac{4}{135\,(n+k_n)}$ with $\frac{2}{21}\le k_n\le\frac{8}{45}$ for all integers $n\ge 0$ \cite[Theorem~7]{flajolet1995}. In particular $\theta_m\ge\frac13$, which is \eqref{eq:birthday-third}. The constant $\frac23$ in \eqref{eq:birthday-third} is the source of the term $-\frac13$ in the following bound, and this term is needed in the proof of Proposition~\ref{prop:criterionpositive}.

\begin{lemma}\label{lem:birthdaymean}
For every integer $r\ge 1$,
\[
s_r\le\sqrt{\frac{\pi r}{2}}\;e^{1/(12r)}-\frac13.
\]
\end{lemma}

\begin{proof}
For $0\le j\le r$ we have $P_{r,j}=r!/((r-j)!\,r^{j})$, i.e.\ $r^{\,r-j}/(r-j)!=P_{r,j}\,r^r/r!$. Summing over $j$ and multiplying by $e^{-r}$, we obtain $H_r(r)=\varpi_r(r)\,(1+s_r)$, because $P_{r,0}=1$. It follows from \eqref{eq:birthday-third} that
\[
s_r=\frac{H_r(r)}{\varpi_r(r)}-1\le\frac{\frac12+\frac23\,\varpi_r(r)}{\varpi_r(r)}-1=\frac{1}{2\,\varpi_r(r)}-\frac13.
\]
Finally, $\varpi_r(r)=(r/e)^r/r!$, and the upper bound in \eqref{eq:prelim-robbins} gives
\[
\frac{1}{2\,\varpi_r(r)}=\frac{r!}{2\,(r/e)^r}\le\frac{\sqrt{2\pi r}}{2}\;e^{1/(12r)}=\sqrt{\frac{\pi r}{2}}\;e^{1/(12r)}.\qedhere
\]
\end{proof}

\subsection{The homogeneous law}\label{subsec:birthday-homogeneous}

We now give every cut of a composition the same factor $h$. The chain sums are total weights for the threshold profile, which is not constant; in Sections~\ref{sec:chains} and~\ref{sec:next} they are compared with the total weights for constant profiles, and the two lemmas of this subsection are what these comparisons use.

For an integer $m\ge 0$ and a real number $h$ let
\[
\tau_m(h)=\Lambda_m[h]
\]
be the total weight $\Lambda_m[f]=\sum_{q=1}^{m}V_q(m;f)$ of all compositions of $m$ (Section~\ref{sec:comps}) for the constant cut profile $f=h$. A composition with $q$ parts has $q-1$ cuts, so that its weight for the constant profile $h$ is $h^{q-1}$ times its weight for the profile $\mathbf 1$. Therefore
\[
\tau_m(h)=\sum_{q=1}^{m}T_{m,q}\,h^{q-1}.
\]
For $m\ge 1$ and $h\ge 0$ this number is positive, since every $T_{m,q}$ is nonnegative by Proposition~\ref{prop:birthday} and the term with $q=1$ is $T_{m,1}=a_m>0$.

Let $r\ge 1$ be an integer and let $h\ge 0$. The \emph{homogeneous law with parameter $h$} on the compositions of $r$ gives each composition a probability proportional to its weight for the constant profile $h$. Under it the number of parts $Q$ takes the value $q$ with probability $T_{r,q}\,h^{q-1}/\tau_r(h)$; for $h=1$ this is the birthday law of $r$. The mean number of cuts under the homogeneous law is therefore
\[
\mu_r(h)=\frac{\sum_{q=1}^{r}(q-1)\,T_{r,q}\,h^{q-1}}{\tau_r(h)}\qquad(r\ge 1,\ h\ge 0).
\]
It is nonnegative, and $\mu_r(1)=s_r-1$ by the first two identities of \eqref{eq:birthday-moments}. Inserting Proposition~\ref{prop:birthday} into the formulas for $\tau_r(h)$ and $\mu_r(h)$ and writing $q=j+1$, we express both quantities through the numbers $P_{r,j}$: for $r\ge 1$ and $h\ge 0$,
\begin{equation}\label{eq:birthday-cuts}
\begin{aligned}
\tau_r(h)&=a_r\sum_{j=0}^{r-1}(j+1)\,h^{j}\,P_{r,j+1},\\
\mu_r(h)\,\tau_r(h)&=a_r\sum_{j=0}^{r-1}j\,(j+1)\,h^{j}\,P_{r,j+1}.
\end{aligned}
\end{equation}

For the lower bound of the mean number of cuts we define, for an integer $r\ge 1$ and a real number $X\ge 0$,
\[
\mu^0_r(X)=\frac13\,\Bigl(\sqrt{(4+X)^2+12\,(r-1)}-(4+X)\Bigr),
\]
which is the nonnegative solution $\mu$ of the quadratic equation $\frac32\,\mu^2+(4+X)\,\mu=2\,(r-1)$. Unlike the mean itself, this bound is explicit in $r$ and $h$.

\begin{lemma}\label{lem:meancuts}
For every integer $r\ge 1$ and every real number $h$ with $0<h\le 1$,
\[
\mu^0_r\bigl(r\,(1/h-1)\bigr)\le\mu_r(h).
\]
\end{lemma}

\begin{proof}
We fix $r$ and $h$ and write $\mu=\mu_r(h)$ and $X=r\,(1/h-1)\ge 0$. We shall prove the quadratic inequality
\[
2\,(r-1)\le\tfrac32\,\mu^2+(4+X)\,\mu,
\]
of which the lemma is the solution for $\mu$.

For an integer $j\ge 0$ put $\sigma_j=h^{j}P_{r,j+1}$. By the properties \eqref{eq:birthday-tail} we have $\sigma_j\ge 0$ for every $j$, $\sigma_0=1$, $\sigma_j=0$ for $j\ge r$, and
\[
\sigma_{j+1}=\rho_j\,\sigma_j\qquad\text{with}\qquad\rho_j=h\,\Bigl(1-\frac{j+1}{r}\Bigr)\qquad(j\ge 0).
\]
The $\rho_j$ do not increase as $j$ grows, and $\rho_j\ge 0$ for $j\le r-1$. We work with the three sums
\[
M_0=\sum_{j=0}^{r-1}(j+1)\,\sigma_j,\qquad M_1=\sum_{j=0}^{r-1}j\,(j+1)\,\sigma_j,\qquad M_2=\sum_{j=0}^{r-1}j^2\,(j+1)\,\sigma_j.
\]
By the expressions \eqref{eq:birthday-cuts} we have $\tau_r(h)=a_r\,M_0$ and $\mu\,\tau_r(h)=a_r\,M_1$, and since $a_r>0$ and $M_0\ge\sigma_0=1$,
\[
M_0>0\qquad\text{and}\qquad M_1=\mu\,M_0.
\]
Thus $M_1/M_0$ is the mean of the number of cuts under the homogeneous law, and $M_2/M_0$ is the mean of its square.

We first prove the identity
\[
r\,M_1=h\,\bigl(2\,(r-1)\,M_0+(r-3)\,M_1-M_2\bigr).
\]
Since $\sigma_r=0$, a shift of the index gives $M_1=\sum_{j=0}^{r-1}(j+1)(j+2)\,\sigma_{j+1}$. By the recursion $\sigma_{j+1}=\rho_j\,\sigma_j$, for every $j\ge 0$,
\begin{align*}
r\,(j+1)(j+2)\,\sigma_{j+1}&=h\,(j+1)(j+2)(r-1-j)\,\sigma_j\\
&=h\,(j+1)\,\sigma_j\,\bigl(2\,(r-1)+(r-3)\,j-j^2\bigr),
\end{align*}
and summing over $0\le j\le r-1$ gives the identity.

Next we prove the inequality
\[
\tfrac16\,(M_2-M_1)\,M_0\le\tfrac14\,M_1^2.
\]
The factor $j+1$ in the three sums is the number of pairs $(i,l)$ of nonnegative integers with $i+l=j$. The sums are therefore sums over such pairs with weight $\sigma_{i+l}$, and the inequality says that under this weight the two coordinates $i$ and $l$ are negatively correlated. This holds because $(\sigma_j)$ is log-concave, and we derive it from two applications of Lemma~\ref{lem:chebyshev}, one in each coordinate.

We use the log-concavity in the following form: for all integers $i$, $j$ and $d$ with $0\le i\le j$ and $d\ge 0$,
\[
\sigma_i\,\sigma_{j+d}\le\sigma_{i+d}\,\sigma_j.
\]
We argue by induction on $d$; for $d=0$ the two sides are equal. Suppose that the inequality holds for some $d$. If $j+d+1\ge r$, then $\sigma_{j+d+1}=0$, and the inequality for $d+1$ holds because its right-hand side is nonnegative. Otherwise $j+d\le r-2$ and $i\le j$, so that $0\le\rho_{j+d}\le\rho_{i+d}$, and by the recursion for $\sigma_j$, by the induction hypothesis and by $\sigma_{i+d}\,\sigma_j\ge 0$,
\[
\sigma_i\,\sigma_{j+d+1}=\rho_{j+d}\,\sigma_i\,\sigma_{j+d}\le\rho_{j+d}\,\sigma_{i+d}\,\sigma_j\le\rho_{i+d}\,\sigma_{i+d}\,\sigma_j=\sigma_{i+d+1}\,\sigma_j.
\]

We write the three sums as sums over pairs in the following notation. For a real-valued function $g$ of two nonnegative integers put
\[
\langle g\rangle=\sum_{i=0}^{r-1}\sum_{l=0}^{r-1}\sigma_{i+l}\,g(i,l).
\]
Since $\sigma_j=0$ for $j\ge r$, only pairs with $i+l\le r-1$ contribute, and collecting the pairs with $i+l=j$ we get $\langle g\rangle=\sum_{j=0}^{r-1}\sigma_j\sum_{i=0}^{j}g(i,j-i)$. The inner sum is $j+1$ for $g=1$, it is $\frac12\,j\,(j+1)$ for $g(i,l)=i$ and for $g(i,l)=l$, and it is $\frac16\,(j-1)\,j\,(j+1)=\frac16\,\bigl(j^2(j+1)-j\,(j+1)\bigr)$ for $g(i,l)=i\,l$. This gives
\[
\langle 1\rangle=M_0,\qquad\langle i\rangle=\langle l\rangle=\tfrac12\,M_1,\qquad\langle i\,l\rangle=\tfrac16\,(M_2-M_1),
\]
where $\langle i\rangle$, $\langle l\rangle$ and $\langle i\,l\rangle$ stand for $\langle g\rangle$ with $g(i,l)=i$, with $g(i,l)=l$ and with $g(i,l)=i\,l$.

It remains to apply Lemma~\ref{lem:chebyshev}, first in the coordinate $l$ and then in the coordinate $i$. For an integer $i\ge 0$ put
\[
U_0(i)=\sum_{l=0}^{r-1}\sigma_{i+l},\qquad U_1(i)=\sum_{l=0}^{r-1}l\,\sigma_{i+l}.
\]
Let $i$ and $i'$ be integers with $0\le i\le i'$. For integers $l$ and $l'$ with $0\le l\le l'$, the log-concavity in the form proved above, applied with $i+l\le i'+l$ in the place of $i\le j$ and with $d=l'-l$, reads $\sigma_{i+l}\,\sigma_{i'+l'}\le\sigma_{i+l'}\,\sigma_{i'+l}$. This is the first hypothesis of Lemma~\ref{lem:chebyshev} for $S=\{0,\dots,r-1\}$, $u(l)=\sigma_{i'+l}$ and $v(l)=\sigma_{i+l}$, and with $\varphi(l)=l$ the lemma gives
\[
U_1(i')\,U_0(i)\le U_1(i)\,U_0(i')\qquad(0\le i\le i').
\]
This in turn is the first hypothesis of Lemma~\ref{lem:chebyshev} for the same set $S$, $u=U_1$ and $v=U_0$, and with $\varphi(i)=i$ the lemma gives
\[
\Bigl(\sum_{i=0}^{r-1}i\,U_1(i)\Bigr)\Bigl(\sum_{i=0}^{r-1}U_0(i)\Bigr)\le\Bigl(\sum_{i=0}^{r-1}i\,U_0(i)\Bigr)\Bigl(\sum_{i=0}^{r-1}U_1(i)\Bigr),
\]
i.e.\ $\langle i\,l\rangle\,\langle 1\rangle\le\langle i\rangle\,\langle l\rangle$. With the values of these four sums found above, this is the inequality $\frac16\,(M_2-M_1)\,M_0\le\frac14\,M_1^2$ that we had to prove.

We now deduce the quadratic inequality from the identity and the inequality just proved. Multiplied by $6h$, the inequality reads $h\,(M_2-M_1)\,M_0\le\frac32\,h\,M_1^2$, and the identity gives $h\,M_2=2\,(r-1)\,h\,M_0+(r-3)\,h\,M_1-r\,M_1$. Eliminating $M_2$, we obtain
\[
2\,(r-1)\,h\,M_0^2\le\tfrac32\,h\,M_1^2+(4h+r-rh)\,M_1\,M_0.
\]
Dividing by $h\,M_0^2>0$ and using $M_1=\mu\,M_0$ and $(4h+r-rh)/h=4+X$, we arrive at the quadratic inequality $2\,(r-1)\le\frac32\,\mu^2+(4+X)\,\mu$.

Finally we solve the quadratic inequality for $\mu$. Completing the square, we write it as $(4+X)^2+12\,(r-1)\le(3\mu+4+X)^2$. Since $\mu\ge 0$ and $X\ge 0$, we have $3\mu+4+X>0$, and taking square roots we obtain $\mu^0_r(X)\le\mu$.
\end{proof}

For $h=1$ the total weights satisfy $\tau_{r+1}(1)=\bigl(1+\frac1r\bigr)^{r}\,\tau_r(1)$, since $\tau_m(1)=m\,a_m$ by the first identity of \eqref{eq:birthday-moments} and $(r+1)\,a_{r+1}=\bigl(1+\frac1r\bigr)^{r}\,r\,a_r$ by \eqref{eq:prelim-ratio}. For $h<1$ the comparison of $\tau_{r+1}(h)$ with $\tau_r(h)$ gains a factor that contains the mean number of cuts.

\begin{lemma}\label{lem:homogeneous}
For every integer $r\ge 1$ and every real number $h$ with $0<h\le 1$,
\[
\tau_{r+1}(h)\le\Bigl(1+\frac1r\Bigr)^{r}\,\Bigl(1-\frac{1-h}{2h\,(r+1)}\,\mu_r(h)\Bigr)\,\tau_r(h).
\]
\end{lemma}

\begin{proof}
We fix $r$ and $h$, and for an integer $m\ge 1$ we put
\[
B_m=\sum_{j=1}^{m}h^{j-1}\,P_{m,j+1}.
\]

We first express the total weights: for every integer $m\ge 1$,
\[
\tau_m(h)=m\,a_m\,\bigl(1-(1-h)\,B_m\bigr).
\]
Indeed, since $\tau_m(h)=\sum_{q=1}^{m}T_{m,q}\,h^{q-1}$, Proposition~\ref{prop:birthday} gives
\[
\tau_m(h)=a_m\sum_{q=1}^{m}h^{q-1}\,q\,P_{m,q},
\]
and the summation by parts \eqref{eq:birthday-abel}, applied to the function $\psi$ with $\psi(0)=1$ and $\psi(q)=h^{q-1}$ for $q\ge 1$, shows that this sum equals $m\,\bigl(1-(1-h)\sum_{q=2}^{m}h^{q-2}\,P_{m,q}\bigr)$. The last sum is $B_m$, because $P_{m,m+1}=0$.

To compare the tails $P_{r,j+1}$ and $P_{r+1,j+1}$ of the birthday laws of $r$ and of $r+1$, write $w=2r(r+1)$. We claim that for every integer $j$ with $0\le j\le r$,
\[
\bigl(w+j\,(j+1)\bigr)\,P_{r,j+1}\le w\,P_{r+1,j+1}.
\]
We argue by induction on $j$; for $j=0$ the two sides are equal. Let $0\le j<r$ and suppose that the inequality holds for $j$. By the recursion \eqref{eq:birthday-tail} we have $P_{r,j+2}=\alpha\,P_{r,j+1}$ and $P_{r+1,j+2}=\beta\,P_{r+1,j+1}$ with
\[
\alpha=1-\frac{j+1}{r},\qquad\beta=1-\frac{j+1}{r+1}>0.
\]
A direct computation gives
\[
\Bigl(\bigl(w+j\,(j+1)\bigr)\,\beta-\bigl(w+(j+1)(j+2)\bigr)\,\alpha\Bigr)\,r(r+1)=(j+1)^2\,(2r+j+2)\ge 0,
\]
so that $(w+(j+1)(j+2))\,\alpha\le(w+j\,(j+1))\,\beta$. Multiplying this inequality by $P_{r,j+1}\ge 0$ and the induction hypothesis by $\beta$, we obtain
\begin{align*}
\bigl(w+(j+1)(j+2)\bigr)\,P_{r,j+2}&\le\bigl(w+j\,(j+1)\bigr)\,\beta\,P_{r,j+1}\\
&\le w\,\beta\,P_{r+1,j+1}=w\,P_{r+1,j+2},
\end{align*}
which is the claim for $j+1$.

A lower bound for $B_{r+1}-B_r$ follows by summation. We write the inequality just proved as $j\,(j+1)\,P_{r,j+1}\le w\,(P_{r+1,j+1}-P_{r,j+1})$, multiply by $h^{j}$ and sum over $1\le j\le r$. Since $P_{r,r+1}=P_{r+1,r+2}=0$, the left-hand side becomes $\mu_r(h)\,\tau_r(h)/a_r$ by \eqref{eq:birthday-cuts}, and the right-hand side becomes $h\,w\,(B_{r+1}-B_r)$. Therefore
\[
\mu_r(h)\,\tau_r(h)\le 2h\,r(r+1)\,a_r\,(B_{r+1}-B_r).
\]
On the other hand, the formula $\tau_m(h)=m\,a_m\,(1-(1-h)\,B_m)$ for $m=r+1$ and for $m=r$ and the relation $(r+\nobreak 1)\,a_{r+1}=\bigl(1+\frac1r\bigr)^{r}\,r\,a_r$, which is \eqref{eq:prelim-ratio} multiplied by $r+1$, give
\begin{align*}
\tau_{r+1}(h)&=\Bigl(1+\frac1r\Bigr)^{r}\,r\,a_r\,\bigl(1-(1-h)\,B_{r+1}\bigr)\\
&=\Bigl(1+\frac1r\Bigr)^{r}\,\Bigl(\tau_r(h)-(1-h)\,r\,a_r\,(B_{r+1}-B_r)\Bigr).
\end{align*}
Since $1-h\ge 0$ and $h>0$, the lower bound for $B_{r+1}-B_r$ gives
\[
(1-h)\,r\,a_r\,(B_{r+1}-B_r)\ge\frac{1-h}{2h\,(r+1)}\,\mu_r(h)\,\tau_r(h),
\]
and the assertion follows.
\end{proof}

\section{Chain sums and the frozen polynomial}\label{sec:chains}
Let $k\ge 2$, and let $n=k+r$ with $1\le r\le k-1$ be a row of the chain range of column $k$. We show that each of the normalised steps $\delta^{(k)}_n$ and $\delta^{(k+1)}_n$ is a second difference of one polynomial in a real threshold minus a multiple of a chain sum (Proposition~\ref{prop:curvatureboundary}); this is the form in which Section~\ref{sec:next} compares the two steps. We also compare the two second differences with each other, and the two chain sums with a third.

\subsection{Compositions and their first part}\label{subsec:chains-first}

We use the notation of Sections~\ref{sec:comps} and~\ref{sec:birthday}. Recall that the weight of a composition for a cut profile $f$ is the product of the Cayley coefficients $a_p$ of its parts and of the values of $f$ at its cut points, that $V_q(m;f)$ is the total weight of the compositions of $m$ with $q$ parts, and that $\Lambda_m[f]=\sum_{q=1}^{m}V_q(m;f)$. For the constant profiles $\mathbf 1$ and $h$ we have $T_{m,q}=V_q(m;\mathbf 1)$ and $\tau_m(h)=\Lambda_m[h]=\sum_{q=1}^{m}T_{m,q}\,h^{q-1}$, and for the threshold profile $\lambda^{y,N}(s)=y/(N-s)$ the total weight is the chain sum $\Lambda_m(y;N)$. In every chain sum $\Lambda_m(y;N)$ of this section $N$ is an integer with $N>m$, so that $N-s\ge 2$ at each cut point $s$ of a composition of $m$.

We shall use the following properties, which are clear from this description. The cut points of a composition of $m$ lie in $\{1,\dots,m-1\}$, so that $V_q(m;f)$ and $\Lambda_m[f]$ depend only on $f(1),\dots,f(m-1)$, and $V_q(m;f)=0$ for $q>m$. If $f$ is nonnegative at these points, then $V_q(m;f)\ge 0$ for every $q$, the total weight $\Lambda_m[f]$ does not decrease when $f$ is raised at any of them, and
\[
\Lambda_m[f]\ge V_1(m;f)=a_m>0\qquad(m\ge 1).
\]
If $f$ is positive at these points, then $V_q(m;f)>0$ for $1\le q\le m$; in particular $T_{m,q}>0$ for $1\le q\le m$. Finally, a composition with $q$ parts has $q-1$ cuts, so that for the threshold profile
\begin{equation}\label{eq:chains-homogeneous}
\begin{aligned}
V_q(m;\lambda^{y,N})&=y^{q-1}\,V_q(m;\lambda^{1,N})\qquad(q\ge 1),\\
\Lambda_m(y;N)&=\sum_{q=1}^{m}y^{q-1}\,V_q(m;\lambda^{1,N}),
\end{aligned}
\end{equation}
and a chain sum is a polynomial in $y$ of degree at most $m-1$.

Definition~\ref{def:compositions} comes with a recursion that removes the last part of a composition. Removing the first part instead shifts the profile.

\begin{lemma}\label{lem:firstpart}
Let $f$ be a cut profile and let $m\ge 0$ be an integer. Then
\[
\Lambda_m[f]=a_m+\sum_{j=1}^{m-1}a_j\,f(j)\,\Lambda_{m-j}[f(j+\cdot)],
\]
where $f(j+\cdot)$ denotes the profile $s\mapsto f(j+s)$.
\end{lemma}

\begin{proof}
We may assume that $m\ge 1$. The composition of $m$ with one part has weight $a_m$. A composition of $m$ with at least two parts consists of its first part $j$, where $1\le j\le m-1$, of the cut at $j$, and of a composition of $m-j$ whose cut points are the remaining cut points shifted by $j$; its weight for $f$ is therefore $a_j\,f(j)$ times the weight of this composition of $m-j$ for $f(j+\cdot)$. Since every composition of $m-j$ arises in this way exactly once, the total weight of the compositions of $m$ with first part $j$ equals
\[
a_j\,f(j)\,\Lambda_{m-j}[f(j+\cdot)]\qquad(1\le j\le m-1),
\]
and the sum of these numbers and of $a_m$ is $\Lambda_m[f]$.
\end{proof}

For the threshold profile we have $\lambda^{y,N}(j+s)=y/(N-j-s)=\lambda^{y,N-j}(s)$, so that Lemma~\ref{lem:firstpart} reads
\begin{equation}\label{eq:chains-firstchain}
\Lambda_m(y;N)=a_m+\sum_{j=1}^{m-1}a_j\,\frac{y}{N-j}\,\Lambda_{m-j}(y;N-j).
\end{equation}

\subsection{Shortening a chain sum and raising its threshold}\label{subsec:chains-boundary}

The two chain sums that occur in Proposition~\ref{prop:curvatureboundary}, $\Lambda_{r+1}(k-1;N)$ and $\Lambda_r(k+2;N)$ with $N=k+r+1$, differ in their length and in their threshold. We compare each of them with the chain sum $\Lambda_r(k;N)$, by means of two homogeneous laws on the compositions of $r$. Recall from Section~\ref{sec:birthday} that $\mu_r(h)$ denotes the mean number of cuts of a composition of $r$ under the homogeneous law with parameter $h$. For integers $k\ge 2$ and $r\ge 1$ we put
\[
h_-=\frac{k-1}{k+r}\qquad\text{and}\qquad\mu_-=\mu_r(h_-),
\]
and for integers $k\ge 1$ and $r\ge 1$ we put
\[
h_{\mathrm{mid}}=\frac{k}{k+1+\frac r2}\qquad\text{and}\qquad\mu_{\mathrm{mid}}=\mu_r(h_{\mathrm{mid}}).
\]
Thus $h_-$ is the value of the threshold profile $\lambda^{k-1,N}$ at the first cut point $s=1$, and $h_{\mathrm{mid}}$ is the value of the formula $k/(N-s)$ of the threshold profile $\lambda^{k,N}$ at the midpoint $s=\frac r2$ of the segment that is cut.

The first comparison rests on the correlation inequality of Fortuin, Kasteleyn and Ginibre \cite[Proposition~1]{fortuin1971}: under a log-supermodular measure on a finite distributive lattice, two nondecreasing functions are nonnegatively correlated. We use it for the lattice of all subsets of a finite set $U$, where it reads as follows. Let $\mu$, $F$ and $G$ be real functions on the set of all subsets of $U$. Suppose that $\mu$ takes positive values and is log-supermodular, that is,
\[
\mu(S)\,\mu(T)\le\mu(S\cap T)\,\mu(S\cup T)
\]
for all subsets $S$ and $T$ of $U$, and that $F$ and $G$ are nondecreasing, that is, $F(S)\le F(T)$ and $G(S)\le G(T)$ whenever $S\subseteq T\subseteq U$. Then $F$ and $G$ are nonnegatively correlated under the probability measure proportional to $\mu$, which means that
\[
\Bigl(\sum_{S}\mu(S)\,F(S)\Bigr)\Bigl(\sum_{S}\mu(S)\,G(S)\Bigr)\le\Bigl(\sum_{S}\mu(S)\Bigr)\Bigl(\sum_{S}\mu(S)\,F(S)\,G(S)\Bigr),
\]
where the four sums are over all subsets $S$ of $U$. We shall apply this with $U=\{1,\dots,r-1\}$, whose subsets are the sets of cut points of the compositions of $r$.

The chain sum $\Lambda_{r+1}(k-1;N)$ can be shortened by one, and its threshold raised from $k-1$ to $k$ at the same time, at the price of the ratio $\tau_{r+1}(h_-)/\tau_r(h_-)$ by which the same shortening changes the total weight for the constant profile $h_-$. The proof compares the threshold profile $\lambda^{k-1,N}$ with the constant profile $h_-$ through the first part of a composition: a composition of $r+1$ arises from a composition of $r$ either by putting a part $1$ in front or by raising the first part $p$ by one, which multiplies the weight by $a_{p+1}/a_p$. This reduces the lemma to a correlation inequality for the compositions of $r$, which the proof deduces from the inequality of Fortuin, Kasteleyn and Ginibre and from the log-convexity of the Cayley coefficients.

\begin{lemma}\label{lem:firstcell}
Let $k\ge 2$ and $r\ge 1$ be integers, and put $N=k+r+1$. Then
\[
\Lambda_{r+1}(k-1;N)\;\tau_r(h_-)\le\tau_{r+1}(h_-)\;\Lambda_r(k;N).
\]
\end{lemma}

\begin{proof}
We write $\psi(j)=(1+\frac1j)^{j-1}$ for integers $j\ge 1$. Thus $a_{j+1}=\psi(j)\,a_j$ by \eqref{eq:prelim-ratio}, and $\psi$ is positive and nondecreasing by Lemma~\ref{lem:powermono}.

For a real function $\varphi$ on the positive integers and a cut profile $f$ let $\Lambda^{\varphi}_r[f]$ be the sum, over all compositions of $r$, of the weight of the composition for $f$ times the value of $\varphi$ at its first part. Grouping the compositions by their first part, as in the proof of Lemma~\ref{lem:firstpart}, we find that
\[
\Lambda^{\varphi}_r[f]=\varphi(r)\,a_r+\sum_{j=1}^{r-1}\varphi(j)\,a_j\,f(j)\,\Lambda_{r-j}[f(j+\cdot)],
\]
and for the constant function $\varphi=1$ this is $\Lambda_r[f]$. We write $f^{+}$ for the shifted profile $f(1+\cdot)$; a constant profile $h$ is its own shift, and for it $\Lambda_r[h]=\tau_r(h)$.

We first express $\Lambda_{r+1}[f]$ through sums over the compositions of $r$: for every cut profile $f$,
\begin{equation}\label{eq:chains-firstunit}
\Lambda_{r+1}[f]=f(1)\,\Lambda_r[f^{+}]+\Lambda^{\psi}_r[f^{+}].
\end{equation}
This is Lemma~\ref{lem:firstpart} with $m=r+1$. Since $a_1=1$, its term with first part $1$ is $f(1)\,\Lambda_r[f^{+}]$. Since $a_{j+1}=\psi(j)\,a_j$, its term with first part $j+1$ is, for $1\le j\le r-1$, the term with first part $j$ in the formula for $\Lambda^{\psi}_r[f^{+}]$, and the remaining term $a_{r+1}=\psi(r)\,a_r$ is the term with first part $r$ in that formula.

Next we prove a correlation inequality. Let $h>0$ be a real number, let $f$ be a cut profile with $f(s)\ge h$ for $1\le s\le r-1$, and let $\varphi$ be a real function on the positive integers that is nondecreasing on $\{1,\dots,r\}$. We claim that
\begin{equation}\label{eq:chains-correlation}
\Lambda^{\varphi}_r[f]\;\tau_r(h)\le\Lambda^{\varphi}_r[h]\;\Lambda_r[f].
\end{equation}
We describe a composition of $r$ by its set $S$ of cut points, which is a subset of $U=\{1,\dots,r-1\}$; every subset of $U$ is the set of cut points of exactly one composition. Let $w(S)>0$ be the product of the $a_p$ over the parts $p$ of this composition, and let $\operatorname{fp}(S)$ be its first part, which is the least element of $S$, or $r$ if $S$ is empty. The weight of the composition for the constant profile $h$ is $\mu(S)=w(S)\,h^{|S|}$, which is positive, and its weight for $f$ is $\mu(S)\,F(S)$, where
\[
F(S)=\prod_{s\in S}\frac{f(s)}{h}.
\]
We put $G(S)=-\varphi(\operatorname{fp}(S))$. The four quantities in the claim are then sums over the subsets $S$ of $U$:
\begin{align*}
\tau_r(h)&=\sum_{S}\mu(S), & \Lambda_r[f]&=\sum_{S}\mu(S)\,F(S),\\
\Lambda^{\varphi}_r[h]&=-\sum_{S}\mu(S)\,G(S), & \Lambda^{\varphi}_r[f]&=-\sum_{S}\mu(S)\,F(S)\,G(S).
\end{align*}
With these expressions, the conclusion of the inequality of Fortuin, Kasteleyn and Ginibre, in the form given before the lemma and for these $\mu$, $F$ and $G$, reads
\[
-\Lambda_r[f]\;\Lambda^{\varphi}_r[h]\le-\tau_r(h)\;\Lambda^{\varphi}_r[f],
\]
which is the claim. In words, when every set of cut points is counted with its weight for the constant profile $h$, the factor $F$ by which the profile $f$ changes this weight and the value of $\varphi$ at the first part are nonpositively correlated: the first does not decrease when cut points are added, and the second does not increase.

So it suffices to verify the hypotheses of that inequality. Let $S\subseteq T\subseteq U$. Then $F(T)$ is the product of $F(S)\ge 0$ and of the factors $f(s)/h\ge 1$ with $s\in T\setminus S$, so that $F(S)\le F(T)$. Moreover $\operatorname{fp}(T)\le\operatorname{fp}(S)$, because the least element of a set does not increase when the set is enlarged and no first part exceeds $r$; since both first parts lie in $\{1,\dots,r\}$, where $\varphi$ is nondecreasing, $G(S)\le G(T)$. Thus $F$ and $G$ are nondecreasing. It remains to show that $\mu$ is log-supermodular. Since $|S|+|T|=|S\cap T|+|S\cup T|$ for all subsets $S$ and $T$ of $U$, the powers of $h$ on the two sides of this condition agree, and what has to be shown is that
\[
w(S)\,w(T)\le w(S\cap T)\,w(S\cup T)
\]
for all subsets $S$ and $T$ of $U$. This is where the log-convexity of the Cayley coefficients enters.

We first consider one additional cut. Let $A\subseteq\{1,\dots,r-1\}$, and let $x\in\{1,\dots,r-1\}$ be a point outside $A$. Let $i<x<j$ be the neighbours of $x$ in $A\cup\{0,r\}$. A cut at $x$ replaces a part of size $j-i$ by two parts of sizes $x-i$ and $j-x$, so that
\[
\frac{w(A\cup\{x\})}{w(A)}=\frac{a_{x-i}\,a_{j-x}}{a_{j-i}}.
\]
This quotient does not decrease when $i$ is raised or $j$ is lowered: for integers $i\le i'<x<j'\le j$,
\[
\frac{a_{x-i}\,a_{j-x}}{a_{j-i}}\le\frac{a_{x-i'}\,a_{j-x}}{a_{j-i'}}\le\frac{a_{x-i'}\,a_{j'-x}}{a_{j'-i'}}.
\]
Both inequalities are cases of the log-convexity of the Cayley coefficients, $a_{u+s}\,a_v\le a_u\,a_{v+s}$ for $1\le u\le v$ and $s\ge 0$ (Lemma~\ref{lem:cayleylogconvex}): the first with $(u,v,s)=(x-i',\,j-i',\,i'-i)$ and the second with $(u,v,s)=(j'-x,\,j'-i',\,j-j')$. Now let $B$ be a set with $A\subseteq B\subseteq\{1,\dots,r-1\}$ that does not contain $x$. The neighbours $i'<x<j'$ of $x$ in $B\cup\{0,r\}$ satisfy $i\le i'$ and $j'\le j$, and therefore
\[
\frac{w(A\cup\{x\})}{w(A)}\le\frac{w(B\cup\{x\})}{w(B)}.
\]
Finally, let $S$ and $T$ be arbitrary subsets of $\{1,\dots,r-1\}$, let $x_1,\dots,x_t$ be the elements of $S\setminus T$, and put $A_l=(S\cap T)\cup\{x_1,\dots,x_l\}$ and $B_l=T\cup\{x_1,\dots,x_l\}$ for $0\le l\le t$. Then $A_l\subseteq B_l$ and $x_{l+1}\notin B_l$ for $l<t$, while $A_0=S\cap T$, $A_t=S$, $B_0=T$ and $B_t=S\cup T$. Applying the last display with $A=A_l$, $B=B_l$ and $x=x_{l+1}$ for each $l<t$, we obtain
\[
\frac{w(S)}{w(S\cap T)}=\prod_{l=0}^{t-1}\frac{w(A_{l+1})}{w(A_l)}\le\prod_{l=0}^{t-1}\frac{w(B_{l+1})}{w(B_l)}=\frac{w(S\cup T)}{w(T)},
\]
i.e.\ $w(S)\,w(T)\le w(S\cap T)\,w(S\cup T)$. Thus $\mu$ is log-supermodular, and the correlation inequality \eqref{eq:chains-correlation} is proved.

We now combine the identity \eqref{eq:chains-firstunit} and the correlation inequality \eqref{eq:chains-correlation}. Let $h>0$ be a real number, and let $f$ be a cut profile with $f(1)=h$ and $f(s)\ge h$ for $2\le s\le r$. We claim that
\[
\Lambda_{r+1}[f]\;\tau_r(h)\le\tau_{r+1}(h)\;\Lambda_r[f^{+}].
\]
By the identity \eqref{eq:chains-firstunit}, applied to $f$ and to the constant profile $h$,
\begin{align*}
\Lambda_{r+1}[f]&=h\,\Lambda_r[f^{+}]+\Lambda^{\psi}_r[f^{+}],\\
\tau_{r+1}(h)&=h\,\tau_r(h)+\Lambda^{\psi}_r[h].
\end{align*}
Multiplying the first equation by $\tau_r(h)$ and the second by $\Lambda_r[f^{+}]$, we see that the claim is equivalent to $\Lambda^{\psi}_r[f^{+}]\,\tau_r(h)\le\Lambda^{\psi}_r[h]\,\Lambda_r[f^{+}]$. This is the correlation inequality \eqref{eq:chains-correlation} for $\varphi=\psi$ and for the profile $f^{+}$, which is at least $h$ at the points $1,\dots,r-1$.

Finally, let $f$ be the threshold profile $\lambda^{k-1,N}$. For $1\le s\le r$ we have $0<N-s\le k+r$, hence $f(s)=(k-1)/(N-s)\ge h_-$, with equality for $s=1$. Since $k\ge 2$, we have $h_->0$, so that the inequality just proved applies with $h=h_-$ and gives
\[
\Lambda_{r+1}(k-1;N)\;\tau_r(h_-)\le\tau_{r+1}(h_-)\;\Lambda_r[f^{+}].
\]
For $1\le s\le r-1$,
\[
0\le f^{+}(s)=\frac{k-1}{N-1-s}\le\frac{k}{N-s}=\lambda^{k,N}(s),
\]
where the second inequality amounts to $k\le N-s$. Since the total weight does not decrease when a nonnegative profile is raised, $\Lambda_r[f^{+}]\le\Lambda_r(k;N)$, and the assertion follows.
\end{proof}

\begin{lemma}\label{lem:shorten}
Let $k\ge 2$ and $r\ge 1$ be integers, and put $N=k+r+1$. Then
\[
\Lambda_{r+1}(k-1;N)\le\Bigl(1+\frac1r\Bigr)^{r}\Bigl(1-\frac{\mu_-}{2(k-1)}\Bigr)\,\Lambda_r(k;N).
\]
\end{lemma}

\begin{proof}
Put $h=h_-$, so that $0<h\le 1$ and $(1-h)/(2h\,(r+1))=1/(2(k-1))$. Lemma~\ref{lem:homogeneous} then reads
\[
\tau_{r+1}(h)\le\Bigl(1+\frac1r\Bigr)^{r}\Bigl(1-\frac{\mu_-}{2(k-1)}\Bigr)\,\tau_r(h).
\]
We insert this into the inequality of Lemma~\ref{lem:firstcell} and cancel $\tau_r(h)$; both $\Lambda_r(k;N)$ and $\tau_r(h)$ are positive, because the total weight of the compositions of $r$ for a nonnegative profile is at least $a_r>0$.
\end{proof}

The second comparison raises the threshold from $k$ to $k+2$ at a fixed length.

\begin{lemma}\label{lem:raise}
Let $k\ge 1$ and $r\ge 1$ be integers, and put $N=k+r+1$. Then
\[
\Bigl(1+\frac2k\Bigr)^{\mu_{\mathrm{mid}}}\,\Lambda_r(k;N)\le\Lambda_r(k+2;N).
\]
\end{lemma}

\begin{proof}
Put $\sigma=1+\frac2k\ge 1$ and $h=h_{\mathrm{mid}}>0$, write $V_q=V_q(r;\lambda^{1,N})$, and for $1\le q\le r$ let
\[
u_q=T_{r,q}\,h^{q-1}\qquad\text{and}\qquad v_q=k^{q-1}\,V_q.
\]
These numbers are positive, because the profiles $\mathbf 1$ and $\lambda^{1,N}$ are positive at the cut points of a composition of $r$. We have $\sum_{q=1}^{r}u_q=\tau_r(h)$, and by the expansion \eqref{eq:chains-homogeneous} of a chain sum also $\sum_{q=1}^{r}v_q=\Lambda_r(k;N)$ and, since $k\sigma=k+2$, $\sum_{q=1}^{r}v_q\,\sigma^{q-1}=\Lambda_r(k+2;N)$. Thus the quotient of the two chain sums is a mean of the $\sigma^{q-1}$ with weights $v_q$.

For $1\le q\le r-1$ we have
\[
u_{q+1}\,v_q\le u_q\,v_{q+1}.
\]
Indeed, Lemma~\ref{lem:erasure} gives $q\,r\,T_{r,q+1}=(q+1)(r-q)\,T_{r,q}$, and Lemma~\ref{lem:midpoint}, which applies because $r<N$, gives $(q+\nobreak 1)(r-\nobreak q)\,V_q\le q\,r\,(N-\frac r2)\,V_{q+1}$. Therefore
\begin{align*}
q\,r\,u_{q+1}\,v_q&=(q+1)(r-q)\,T_{r,q}\,h^{q}\,k^{q-1}\,V_q\\
&\le q\,r\,T_{r,q}\,h^{q-1}\,k^{q-1}\cdot h\,\Bigl(N-\frac r2\Bigr)\,V_{q+1}=q\,r\,u_q\,v_{q+1},
\end{align*}
where the last step uses that $h\,(N-\frac r2)=k$ by the choice of $h$. Since the $v_q$ are positive, the inequalities $u_{q+1}\,v_q\le u_q\,v_{q+1}$ say that $u_q/v_q$ does not increase in $q$, and hence
\[
u_{q'}\,v_q\le u_q\,v_{q'}\qquad(1\le q\le q'\le r).
\]
This is the first hypothesis of Lemma~\ref{lem:chebyshev} on $\{1,\dots,r\}$ for $u$ and $v$, and its second hypothesis holds for $\varphi(q)=\sigma^{q-1}$, because $\sigma\ge 1$. The lemma gives
\[
\Bigl(\sum_{q=1}^{r}u_q\,\sigma^{q-1}\Bigr)\,\Lambda_r(k;N)\le\Lambda_r(k+2;N)\,\tau_r(h).
\]

Now let $\mu=\mu_{\mathrm{mid}}=\mu_r(h)$, so that $\mu\,\tau_r(h)=\sum_{q=1}^{r}(q-1)\,u_q$ by the definition of the mean number of cuts. By Jensen's inequality for the convex function $t\mapsto\sigma^{t}$ and the positive weights $u_q$,
\[
\sum_{q=1}^{r}u_q\,\sigma^{q-1}\ge\sigma^{\mu}\,\tau_r(h).
\]
Since $\Lambda_r(k;N)>0$, the last two displays give the assertion multiplied by $\tau_r(h)$.
\end{proof}

\subsection{The frozen polynomial}\label{subsec:chains-frozen}

We now turn to the steps. Let $k\ge 2$, and let $n=k+r$ be a row of the chain range of column $k$. The normalised step of column $k$ is a second difference in the threshold (Section~\ref{sec:recursion}),
\[
\delta^{(k)}_n=z^{(k+1)}_n-2\,z^{(k)}_n+z^{(k-1)}_n,
\]
and the step of column $k+1$ is the same expression with the three thresholds raised by one. If the four counts $z^{(k-1)}_n,\dots,z^{(k+2)}_n$ were the values of one function of a real variable at four neighbouring points, the two steps would be second differences of that function at neighbouring points, and they could be compared by the methods of calculus. But the counts are defined for integer thresholds only, and a threshold $y$ enters the last-car recursion through the weight $\min(t,y)$, that is, both through the index at which the weight stops growing and through the value that it then keeps. We separate the two roles. With the index fixed at $k$, the value becomes a real variable, and the counts are values of one polynomial in it: exactly for the thresholds $k$ and $k+1$, and up to an explicit multiple of a chain sum for the thresholds $k-1$ and $k+2$ (Proposition~\ref{prop:curvatureboundary}).

We begin with the linear equation. Let $k\ge 0$ be an integer, and let $j\ge 1$ be one of the two integers $k$ and $k+1$. The sequence $z^{(j)}$ is a solution of the weighted last-car recursion with weight $\min(t,j)$ (Corollary~\ref{cor:normalised}), and this weight equals $t$ for $1\le t\le k$ and $j$ for $t\ge k+1$. The single-chain formula (Proposition~\ref{prop:forward}) therefore gives, for $0\le n\le 2k+1$,
\[
c_n-z^{(j)}_n=\sum_{t=k+1}^{n}c_{n-t}\Bigl(c_{t-1}-\frac{j}{t}\,z^{(j)}_{t-1}\Bigr).
\]
We write this with the Cayley coefficients, using $c_i=a_{i+1}$, and with $M=n+1$ in place of $n$, and we collect the terms that do not contain $z^{(j)}$. For integers $k\ge 0$ and $t\ge 0$ let
\begin{equation}\label{eq:chains-nu}
\nu_k(t)=a_t-\sum_{k<i<t}a_i\,a_{t-i}.
\end{equation}
Then the formula becomes
\begin{equation}\label{eq:chains-linear}
z^{(j)}_{M-1}=\nu_k(M)+\sum_{k<t<M}a_{M-t}\,\frac{j}{t}\,z^{(j)}_{t-1}\qquad(1\le M\le 2k+2).
\end{equation}
Here $\nu_k(M)$ and the coefficients $a_{M-t}/t$ depend on $k$ alone. Solving this linear equation by iteration expresses $z^{(j)}_{M-1}$ through the $\nu_k(i)$ with $k<i\le M$ and through chain sums with threshold $j$, as we show in the proof of Proposition~\ref{prop:curvatureboundary}, and the resulting expression is defined for every real number in place of $j$. We take it as a definition.

\begin{definition}\label{def:frozen}
Let $k\ge 0$ and $n\ge 0$ be integers, and put $N=n+1$. The \emph{frozen polynomial} of position $(n,k)$ is the function $\mathcal K_{n,k}$ of a real variable $y$ given by
\[
\mathcal K_{n,k}(y)=\nu_k(N)+\sum_{m=1}^{n-k}\frac{\nu_k(N-m)}{N-m}\;y\;\Lambda_m(y;N),
\]
with $\nu_k$ as in \eqref{eq:chains-nu}.
\end{definition}

In each term of the sum we have $N-m\ge k+1$, so that no denominator vanishes. Since a chain sum $\Lambda_m(y;N)$ is a polynomial in $y$ of degree at most $m-1$, the function $\mathcal K_{n,k}$ is a polynomial in $y$, and for $n\ge k$ its degree is at most $n-k$. For $n\le k$ the sum in the definition is empty, and $\mathcal K_{n,k}$ is the constant $\nu_k(n+1)=a_{n+1}=c_n$, because the sum that defines $\nu_k(t)$ is empty for $t\le k+1$.

Now let $k\ge 2$ and $r$ be integers with $1\le r\le k-1$, and put $n=k+r$, $N=n+1$ and $\mathcal K=\mathcal K_{n,k}$. We write
\begin{align*}
\mathcal C_-&=\mathcal K(k+1)-2\,\mathcal K(k)+\mathcal K(k-1),\\
\mathcal C_+&=\mathcal K(k+2)-2\,\mathcal K(k+1)+\mathcal K(k)
\end{align*}
for the two second differences of $\mathcal K$, and
\[
\mathcal B_-=\frac{a_k}{k}\,\Lambda_{r+1}(k-1;N)\qquad\text{and}\qquad\mathcal B_+=\frac{a_{k+1}}{k+1}\,\Lambda_{r}(k+2;N)
\]
for the two multiples of chain sums that are subtracted from them. In this notation the following proposition says that $\delta^{(k)}_n=\mathcal C_--\mathcal B_-$ and $\delta^{(k+1)}_n=\mathcal C_+-\mathcal B_+$.

\begin{proposition}\label{prop:curvatureboundary}
Let $k\ge 2$ and $r$ be integers with $1\le r\le k-1$, and put $n=k+r$, $N=n+1$ and $\mathcal K=\mathcal K_{n,k}$. Then
\begin{align*}
\delta^{(k)}_n&=\bigl(\mathcal K(k+1)-2\,\mathcal K(k)+\mathcal K(k-1)\bigr)-\frac{a_k}{k}\,\Lambda_{r+1}(k-1;N),\\
\delta^{(k+1)}_n&=\bigl(\mathcal K(k+2)-2\,\mathcal K(k+1)+\mathcal K(k)\bigr)-\frac{a_{k+1}}{k+1}\,\Lambda_{r}(k+2;N).
\end{align*}
\end{proposition}

\begin{proof}
Let $k\ge 0$ be an arbitrary integer, and let $j\ge 1$ be one of the integers $k$ and $k+1$. We first show that
\begin{equation}\label{eq:chains-values}
z^{(j)}_n=\mathcal K_{n,k}(j)\qquad(0\le n\le 2k+1).
\end{equation}
With $M=n+1$, and with $i=M-m$ as the index of summation in Definition~\ref{def:frozen}, this is the assertion that, for every integer $M$ with $1\le M\le 2k+2$,
\[
z^{(j)}_{M-1}=\nu_k(M)+\sum_{k<i<M}\nu_k(i)\,\frac{j}{i}\,\Lambda_{M-i}(j;M).
\]
We prove this by induction on $M$. For $M\le k+1$ the sum in the assertion and the sum in the linear equation \eqref{eq:chains-linear} are empty, and the two equations coincide. Let $M\ge k+2$, and suppose that the assertion holds for all smaller values of $M$. We insert it for $z^{(j)}_{t-1}$ with $k<t<M$ into the linear equation and collect, for each $i$ with $k<i<M$, the terms that contain $\nu_k(i)$. This gives
\[
z^{(j)}_{M-1}=\nu_k(M)+\sum_{k<i<M}\nu_k(i)\,\frac{j}{i}\,\Bigl(a_{M-i}+\sum_{i<t<M}a_{M-t}\,\frac{j}{t}\,\Lambda_{t-i}(j;t)\Bigr).
\]
The bracket is the right-hand side of the formula \eqref{eq:chains-firstchain}, with first part $M-t$, for the chain sum $\Lambda_{M-i}(j;M)$. So it equals $\Lambda_{M-i}(j;M)$, and the assertion holds for~$M$.

We turn to two consecutive columns. Let $k$ and $n$ be integers with $1\le k\le n$, put $N=n+1$, and let $y$ be a real number. We show that
\begin{equation}\label{eq:chains-knot}
\mathcal K_{n,k-1}(y)-\mathcal K_{n,k}(y)=(y-k)\,\frac{a_k}{k}\,\Lambda_{N-k}(y;N).
\end{equation}
By the definition of $\nu_k$, lowering $k$ by one adds the term $i=k$ to the sum:
\[
\nu_{k-1}(t)=\nu_k(t)-a_k\,a_{t-k}\quad(t>k),\qquad\nu_{k-1}(k)=a_k.
\]
The sum that defines $\mathcal K_{n,k-1}(y)$ has one term more than the sum that defines $\mathcal K_{n,k}(y)$, the term $m=N-k$, which is $\frac{a_k}{k}\,y\,\Lambda_{N-k}(y;N)$. In the constant term and in the terms with $1\le m\le n-k$ the argument of $\nu$ is larger than $k$, so that these terms contribute to the difference of the two polynomials
\[
-a_k\Bigl(a_{N-k}+\sum_{m=1}^{n-k}\Lambda_m(y;N)\,\frac{y}{N-m}\,a_{N-k-m}\Bigr).
\]
The bracket equals $\Lambda_{N-k}(y;N)$. Indeed, the recursion of Definition~\ref{def:compositions}, which removes the last part of a composition, gives when summed over the number of parts
\[
\Lambda_{m'}[f]=a_{m'}+\sum_{s=1}^{m'-1}\Lambda_s[f]\,f(s)\,a_{m'-s}
\]
for every cut profile $f$ and every integer $m'\ge 1$, and for $m'=N-k$ and $f=\lambda^{y,N}$ the right-hand side is the bracket. The two contributions add up to $(y-k)\,\frac{a_k}{k}\,\Lambda_{N-k}(y;N)$, as claimed.

We can now prove the two formulas. Let $k$, $r$, $n$ and $N$ be as in the proposition. By \eqref{eq:chains-values}, applied with $k$, with $k-1$ and with $k+1$ in place of $k$,
\begin{align*}
z^{(k)}_n&=\mathcal K_{n,k}(k), & z^{(k+1)}_n&=\mathcal K_{n,k}(k+1),\\
z^{(k-1)}_n&=\mathcal K_{n,k-1}(k-1), & z^{(k+2)}_n&=\mathcal K_{n,k+1}(k+2);
\end{align*}
for the third of these we use that $n\le 2(k-1)+1$, which holds because $r\le k-1$. Next we apply \eqref{eq:chains-knot}. With $k$ itself and $y=k-1$, where $N-k=r+1$, it gives
\[
z^{(k-1)}_n=\mathcal K_{n,k}(k-1)-\frac{a_k}{k}\,\Lambda_{r+1}(k-1;N).
\]
With $k+1$ in place of $k$, which is admissible because $r\ge 1$, and with $y=k+2$, where $N-(k+1)=r$, it gives
\[
z^{(k+2)}_n=\mathcal K_{n,k}(k+2)-\frac{a_{k+1}}{k+1}\,\Lambda_{r}(k+2;N).
\]
By \eqref{eq:recursion-second}, divided by $n!$, we have $\delta^{(k)}_n=z^{(k+1)}_n-2\,z^{(k)}_n+z^{(k-1)}_n$ and $\delta^{(k+1)}_n=z^{(k+2)}_n-2\,z^{(k+1)}_n+z^{(k)}_n$, and inserting the four values gives the two formulas of the proposition.
\end{proof}

\subsection{The ratio of the two second differences}\label{subsec:chains-curvature}

By Proposition~\ref{prop:curvatureboundary}, column $k$ is nonpositive in row $n=k+r$ if and only if $\mathcal C_-\le\mathcal B_-$, and column $k+1$ is negative in this row if and only if $\mathcal C_+<\mathcal B_+$. The chain sums in $\mathcal B_-$ and $\mathcal B_+$ were compared with $\Lambda_r(k;N)$ in Lemmas~\ref{lem:shorten} and~\ref{lem:raise}. We now compare the two second differences: for $r\ge 1$ the following proposition says that $\mathcal C_+\le e^{D_r/(k-1)}\,\mathcal C_-$. Recall from Section~\ref{sec:birthday} that, for $m\ge 2$, $D_m$ is the ratio of the third to the second factorial moment of the birthday law of $m$, and that $D_0=D_1=0$.

\begin{proposition}\label{prop:curvature}
Let $k\ge 2$ and $r$ be integers with $0\le r\le k-1$, and let $\mathcal K=\mathcal K_{k+r,k}$ be the frozen polynomial of position $(k+r,k)$. Then
\[
\mathcal K(k+2)-2\,\mathcal K(k+1)+\mathcal K(k)\le e^{D_r/(k-1)}\,\bigl(\mathcal K(k+1)-2\,\mathcal K(k)+\mathcal K(k-1)\bigr).
\]
\end{proposition}

\begin{proof}
Put $n=k+r$ and $N=n+1$.

We begin with the coefficients of $\mathcal K$. Let $\alpha_m=\nu_k(N-m)/(N-m)$ for $1\le m\le r$. We insert the expansion \eqref{eq:chains-homogeneous} of the chain sums into Definition~\ref{def:frozen} and exchange the two summations:
\begin{align*}
\mathcal K(y)&=\nu_k(N)+\sum_{m=1}^{r}\alpha_m\sum_{q=1}^{m}y^{q}\,V_q(m;\lambda^{1,N})\\
&=\nu_k(N)+\sum_{q=1}^{r}p_q\,y^q,
\end{align*}
where $p_q=\sum_{m=q}^{r}\alpha_m\,V_q(m;\lambda^{1,N})$. The numbers $V_q(m;\lambda^{1,N})$ are nonnegative, because the profile $s\mapsto 1/(N-s)$ is positive at the cut points. To see that the $\alpha_m$ are positive, we show that
\[
\nu_k(t)>0\qquad(1\le t\le 2k+1).
\]
This applies to $t=N-m$, because $k+1\le N-m\le k+r\le 2k-1$ for $1\le m\le r$, since $r\le k-1$. For the proof, let $1\le t\le 2k+1$. The substitution of $t-i$ for $i$ turns the sum of the terms $a_i\,a_{t-i}$ over $k<i<t$, which occurs in the definition of $\nu_k(t)$, into the sum of the same terms over $0<i<t-k$. Since $t-k\le k+1$, the two sums run over disjoint sets of indices in $\{0,\dots,t\}$, and since all terms are nonnegative,
\[
2\sum_{k<i<t}a_i\,a_{t-i}\le\sum_{i=0}^{t}a_i\,a_{t-i}.
\]
By Lemma~\ref{lem:cayleyconvolution}, to whose sum the vanishing term $i=t$ may be added, and by the same substitution, each of the sums $\sum_{i=0}^{t}i\,a_i\,a_{t-i}$ and $\sum_{i=0}^{t}(t-i)\,a_i\,a_{t-i}$ equals $(t-1)\,a_t$, so that $t\sum_{i=0}^{t}a_i\,a_{t-i}=2\,(t-1)\,a_t$. Therefore $t\sum_{k<i<t}a_i\,a_{t-i}\le(t-1)\,a_t$, and $t\,\nu_k(t)\ge a_t>0$. Consequently $\alpha_m>0$ for $1\le m\le r$ and $p_q\ge 0$ for $1\le q\le r$.

Next we compare two factorial moments. We claim that, for every real number $y$ with $0<y\le k+2$,
\begin{equation}\label{eq:chains-moments}
\sum_{q=1}^{r}q(q-1)(q-2)\,p_q\,y^q\le D_r\sum_{q=1}^{r}q(q-1)\,p_q\,y^q.
\end{equation}
By the definition of the coefficients $p_q$, and since $y^{q-1}\,V_q(m;\lambda^{1,N})=V_q(m;\lambda^{y,N})$ by \eqref{eq:chains-homogeneous}, the left-hand side is
\[
\sum_{m=1}^{r}\alpha_m\,y\sum_{q=1}^{m}q(q-1)(q-2)\,V_q(m;\lambda^{y,N}),
\]
and the sum on the right-hand side is the same expression with $q(q-1)$ in place of $q(q-1)(q-2)$. Since $\alpha_m\,y\ge 0$, it suffices to prove, for each $m$ with $1\le m\le r$ and with $f=\lambda^{y,N}$,
\[
\sum_{q=1}^{m}q(q-1)(q-2)\,V_q(m;f)\le D_r\sum_{q=1}^{m}q(q-1)\,V_q(m;f).
\]
We may assume that $m\ge 2$, so that $r\ge 2$. A cut point $s$ of a composition of $m$ satisfies $s\le m-1\le r-1$, hence $N-s\ge k+2\ge y$, and therefore $0<f(s)\le 1$. For such a profile Lemma~\ref{lem:erasureupper} gives $q\,m\,V_{q+1}(m;f)\le(q+1)(m-q)\,V_q(m;f)$ for $q\ge 1$, while Lemma~\ref{lem:erasure} gives $q\,m\,T_{m,q+1}=(q+1)(m-q)\,T_{m,q}$. Multiplying the inequality by $T_{m,q}\ge 0$ and using the equality, we obtain
\[
V_{q+1}(m;f)\,T_{m,q}\le V_q(m;f)\,T_{m,q+1}\qquad(q\ge 1).
\]
Since $T_{m,q}>0$ for $1\le q\le m$, this says that $V_q(m;f)/T_{m,q}$ does not increase in $q$ on $\{1,\dots,m\}$. Hence
\[
V_{q'}(m;f)\,T_{m,q}\le V_q(m;f)\,T_{m,q'}\qquad(1\le q\le q'\le m).
\]
Multiplied by the nonnegative number $q(q-1)\,q'(q'-1)$, this is the first hypothesis of Lemma~\ref{lem:chebyshev} on $\{1,\dots,m\}$ for $u(q)=q(q-1)\,V_q(m;f)$ and $v(q)=q(q-1)\,T_{m,q}$, and its second hypothesis holds for $\varphi(q)=q-2$. The lemma gives
\begin{multline*}
\Bigl(\sum_{q=1}^{m}q(q-1)(q-2)\,V_q(m;f)\Bigr)\Bigl(\sum_{q=1}^{m}q(q-1)\,T_{m,q}\Bigr)\\
\le\Bigl(\sum_{q=1}^{m}q(q-1)(q-2)\,T_{m,q}\Bigr)\Bigl(\sum_{q=1}^{m}q(q-1)\,V_q(m;f)\Bigr).
\end{multline*}
By \eqref{eq:birthday-ratio} the first sum on the right equals $D_m\sum_{q=1}^{m}q(q-1)\,T_{m,q}$. Cancelling this last sum, which is positive because $T_{m,2}>0$, we obtain the required inequality with $D_m$ in place of $D_r$. Since $D_m\le D_r$ by Lemma~\ref{lem:ratiomonotone} and $\sum_{q=1}^{m}q(q-1)\,V_q(m;f)\ge 0$, the inequality with $D_r$ follows, and with it the claim.

It remains to deduce the inequality between the second differences. We regard $\mathcal K$ as a function on the real line and write $\mathcal K''$ and $\mathcal K'''$ for its second and third derivative. Since $\mathcal K(y)=\nu_k(N)+\sum_{q=1}^{r}p_q\,y^q$,
\begin{align*}
y^2\,\mathcal K''(y)&=\sum_{q=1}^{r}q(q-1)\,p_q\,y^q,\\
y^3\,\mathcal K'''(y)&=\sum_{q=1}^{r}q(q-1)(q-2)\,p_q\,y^q,
\end{align*}
so that the inequality \eqref{eq:chains-moments} between the factorial moments, divided by $y^2>0$, reads $y\,\mathcal K'''(y)\le D_r\,\mathcal K''(y)$ for $0<y\le k+2$. Moreover $\mathcal K'''(y)\ge 0$ for $y\ge 0$, because $p_q\ge 0$. Put $\theta=D_r/(k-1)$. For $k-1\le y\le k+2$ we have
\[
(k-1)\,\mathcal K'''(y)\le y\,\mathcal K'''(y)\le D_r\,\mathcal K''(y),
\]
i.e.\ $\mathcal K'''(y)\le\theta\,\mathcal K''(y)$. Hence $t\mapsto e^{-\theta t}\,\mathcal K''(t)$ does not increase on $[k-1,k+2]$, so that
\[
\mathcal K''(y+1)\le e^{\theta}\,\mathcal K''(y)\qquad(k-1\le y\le k+1).
\]
Consequently $H(t)=e^{\theta}\,\mathcal K(t)-\mathcal K(t+1)$ has a nonnegative second derivative on $[k-1,k+1]$ and is convex there. In particular $2\,H(k)\le H(k-1)+H(k+1)$, and rearranging we obtain the inequality of the proposition.
\end{proof}

\section{The next column}\label{sec:next}
In this section we prove Proposition~\ref{prop:nextcolumn}: for $k\ge 9$ and $n\ge k+2$, if column $k$ is nonpositive in rows $k+1,\dots,n$, then column $k+1$ is nonpositive in row $n$. In the proof of Theorem~\ref{thm:main} this proposition carries a nonpositive step from each column to the next along a row. For $k\ge 265$ we deduce it, with a strict conclusion, from the comparisons of Section~\ref{sec:chains} and from the bound of Proposition~\ref{prop:strip} for the crossing row; for $9\le k\le 264$ it follows from the computed signs of Lemma~\ref{lem:signpairs} and from Proposition~\ref{prop:chainrange}, by which a positive column stays positive in its chain range.

\subsection{A criterion for the next column}\label{subsec:next-criterion}

By Proposition~\ref{prop:curvatureboundary} the normalised steps of columns $k$ and $k+1$ in a row $n=k+r$ of the chain range of column $k$ are each of the form $\mathcal C-\mathcal B$, a second difference of a polynomial minus a multiple of a chain sum, and Section~\ref{sec:chains} bounds the growth of $\mathcal C$ from above, and that of $\mathcal B$ from below, when one passes from column $k$ to column $k+1$. The function that we now define is $k$ times an explicit lower bound for the logarithm of the quotient of the second bound by the first.

\begin{definition}\label{def:criterionfunction}
For integers $r\ge 1$ and $k\ge 2$, the \emph{criterion function} is
\begin{multline*}
\mathsf G(r,k)=k\Bigl(\frac{1}{2r}-\frac{1}{3r^2}\Bigr)-\frac52+\frac12\,\mu^0_r\Bigl(\frac{r(r+1)}{k-1}\Bigr)\\
+2\Bigl(1-\frac1k\Bigr)\mu^0_r\Bigl(\frac{r(r+2)}{2k}\Bigr)-\frac{k}{k-1}\,D_r .
\end{multline*}
Here, for a real number $X\ge 0$,
\[
\mu^0_r(X)=\frac13\Bigl(\sqrt{(4+X)^2+12\,(r-1)}-(4+X)\Bigr),
\]
and $D_r$ is the number defined in Section~\ref{sec:birthday}: for $r\ge 2$,
\[
D_r=\frac{3\,\bigl((r+2)\,s_r-3r\bigr)}{2\,(r-s_r)},\qquad\text{where}\quad s_r=\sum_{j=1}^{r}\,\prod_{i=0}^{j-1}\Bigl(1-\frac ir\Bigr)
\]
is the mean of the birthday law of $r$, and $D_1=0$.
\end{definition}

We write
\[
X_-=\frac{r(r+1)}{k-1},\qquad X_{\mathrm{mid}}=\frac{r(r+2)}{2k}
\]
for the two arguments of $\mu^0_r$ in $\mathsf G(r,k)$; both are nonnegative. Recall that $\delta^{(j)}_n=d_{n,j}/n!$ denotes the normalised step of column $j$ in row $n$, which has the sign of $b_{n,j}-b_{n,j-1}$. In terms of the counts, the following proposition therefore says that $b_{n,k}\le b_{n,k-1}$ implies $b_{n,k+1}<b_{n,k}$ whenever $\mathsf G(r,k)\ge 0$.

\begin{proposition}\label{prop:nextcriterion}
Let $k\ge 3$ and $r$ be integers with $2\le r\le k-1$, and put $n=k+r$. If $\mathsf G(r,k)\ge 0$ and $\delta^{(k)}_n\le 0$, then $\delta^{(k+1)}_n<0$.
\end{proposition}

\begin{proof}
Put $N=n+1$ and $\Lambda=\Lambda_r(k;N)$. Recall from Section~\ref{sec:chains} the notation for row $n=k+r$: with the frozen polynomial $\mathcal K=\mathcal K_{n,k}$ of Definition~\ref{def:frozen},
\begin{align*}
\mathcal C_-&=\mathcal K(k+1)-2\,\mathcal K(k)+\mathcal K(k-1), & \mathcal B_-&=\frac{a_k}{k}\,\Lambda_{r+1}(k-1;N),\\
\mathcal C_+&=\mathcal K(k+2)-2\,\mathcal K(k+1)+\mathcal K(k), & \mathcal B_+&=\frac{a_{k+1}}{k+1}\,\Lambda_r(k+2;N).
\end{align*}
Since $k\ge 2$ and $1\le r\le k-1$, Proposition~\ref{prop:curvatureboundary} gives $\delta^{(k)}_n=\mathcal C_--\mathcal B_-$ and $\delta^{(k+1)}_n=\mathcal C_+-\mathcal B_+$. Thus $\delta^{(k)}_n\le 0$ means that $\mathcal C_-\le\mathcal B_-$, and we have to show that $\mathcal C_+<\mathcal B_+$.

Let
\[
h_-=\frac{k-1}{k+r},\qquad h_{\mathrm{mid}}=\frac{k}{k+1+r/2},
\]
and let $\mu_-=\mu_r(h_-)$ and $\mu_{\mathrm{mid}}=\mu_r(h_{\mathrm{mid}})$ be the mean numbers of cuts of a composition of $r$ under the homogeneous laws with these parameters (Section~\ref{sec:birthday}). Proposition~\ref{prop:curvature} and Lemmas~\ref{lem:shorten} and~\ref{lem:raise}, whose hypotheses on $k$ and $r$ hold as well, give the three bounds
\begin{align*}
\mathcal C_+&\le e^{D_r/(k-1)}\,\mathcal C_-,\\
\Lambda_{r+1}(k-1;N)&\le\Bigl(1+\frac1r\Bigr)^{r}\Bigl(1-\frac{\mu_-}{2(k-1)}\Bigr)\,\Lambda,\\
\Bigl(1+\frac2k\Bigr)^{\mu_{\mathrm{mid}}}\Lambda&\le\Lambda_r(k+2;N).
\end{align*}
By \eqref{eq:prelim-ratio}, $a_{k+1}=(1+1/k)^{k-1}a_k$, and therefore
\[
\frac{a_{k+1}}{k+1}=\Bigl(1+\frac1k\Bigr)^{k-2}\,\frac{a_k}{k}.
\]
Moreover $a_k>0$ and $\Lambda>0$, since the Cayley coefficients $a_p$ with $p\ge 1$ and the factors $k/(N-s)$ of the cut points $s<r<N$ are positive.

Suppose for the moment that
\[
e^{D_r/(k-1)}\Bigl(1+\frac1r\Bigr)^{r}e^{-\mu_-/(2(k-1))}<\Bigl(1+\frac1k\Bigr)^{k-2}\Bigl(1+\frac2k\Bigr)^{\mu_{\mathrm{mid}}}.
\]
Then the first of the three bounds and the assumption $\mathcal C_-\le\mathcal B_-$ give
\[
\mathcal C_+\le e^{D_r/(k-1)}\,\mathcal C_-\le e^{D_r/(k-1)}\,\frac{a_k}{k}\,\Lambda_{r+1}(k-1;N),
\]
and by the second bound and the inequality $1-t\le e^{-t}$ at $t=\mu_-/(2(k-1))$ the right-hand side is at most
\begin{multline*}
e^{D_r/(k-1)}\Bigl(1+\frac1r\Bigr)^{r}e^{-\mu_-/(2(k-1))}\,\frac{a_k}{k}\,\Lambda\\
<\Bigl(1+\frac1k\Bigr)^{k-2}\Bigl(1+\frac2k\Bigr)^{\mu_{\mathrm{mid}}}\,\frac{a_k}{k}\,\Lambda ,
\end{multline*}
where the strict inequality is the supposed one, multiplied by $\frac{a_k}{k}\,\Lambda>0$. By the formula for $a_{k+1}/(k+1)$ and the third bound, the last expression is at most $\frac{a_{k+1}}{k+1}\,\Lambda_r(k+2;N)=\mathcal B_+$. Thus $\mathcal C_+<\mathcal B_+$, i.e.\ $\delta^{(k+1)}_n<0$.

It remains to prove the supposed inequality. Both of its sides are positive, and the logarithm of the right-hand side minus the logarithm of the left-hand side is
\[
(k-2)\log\Bigl(1+\frac1k\Bigr)+\mu_{\mathrm{mid}}\log\Bigl(1+\frac2k\Bigr)-r\log\Bigl(1+\frac1r\Bigr)+\frac{\mu_-}{2(k-1)}-\frac{D_r}{k-1}.
\]
We show that this difference is at least $\mathsf G(r,k)/k+1/k^2$, which is positive because $\mathsf G(r,k)\ge 0$.

We first bound the two means from below. Both $h_-$ and $h_{\mathrm{mid}}$ lie in $(0,1]$, and $r\,(1/h_--1)=X_-$ and $r\,(1/h_{\mathrm{mid}}-1)=X_{\mathrm{mid}}$. Hence Lemma~\ref{lem:meancuts} and the definition of $\mu^0_r$, which is nonnegative, give
\[
0\le\mu^0_r(X_-)\le\mu_-,\qquad 0\le\mu^0_r(X_{\mathrm{mid}})\le\mu_{\mathrm{mid}}.
\]

Next we bound the logarithms. Put
\[
u=\frac1r,\qquad v=\frac1k,\qquad w=\frac{1}{k-1},
\]
so that $0<v<w$ and, since $k\ge 3$, $v\le\frac13$. For $t\ge 0$ the Taylor polynomial of $\log(1+t)$ of even (resp.\ odd) degree bounds it from below (resp.\ above). This gives
\begin{align*}
r\log\Bigl(1+\frac1r\Bigr)&\le 1-\frac u2+\frac{u^2}{3},\\
(k-2)\log\Bigl(1+\frac1k\Bigr)&\ge 1-\frac52\,v+v^2,\\
\log\Bigl(1+\frac2k\Bigr)&\ge 2v\,(1-v).
\end{align*}
Indeed, the first is the bound of degree $3$ at $t=u$, multiplied by $r=1/u$, and the third is the bound of degree $2$ at $t=2v$. For the second we multiply the bound of degree $4$ at $t=v$ by $k-2=(1-2v)/v\ge 0$; a direct computation shows that the result exceeds $1-\frac52\,v+v^2$ by $\frac{1}{12}\,v^2\,(4-3v)(1-2v)$, which is nonnegative because $0<v\le\frac13$.

We insert the three bounds just obtained into the difference of the logarithms, using that $\mu_{\mathrm{mid}}\ge\mu^0_r(X_{\mathrm{mid}})\ge 0$, and find that it is at least
\[
\Bigl(1-\frac52\,v+v^2\Bigr)+2v\,(1-v)\,\mu_{\mathrm{mid}}-\Bigl(1-\frac u2+\frac{u^2}{3}\Bigr)+\frac w2\,\mu_--w\,D_r .
\]
Since $0\le\mu^0_r(X_-)\le\mu_-$, $\mu^0_r(X_{\mathrm{mid}})\le\mu_{\mathrm{mid}}$ and $0<v<w$, we have $w\,\mu_-\ge v\,\mu^0_r(X_-)$ and $2v\,(1-v)\,\mu_{\mathrm{mid}}\ge 2v\,(1-v)\,\mu^0_r(X_{\mathrm{mid}})$. Hence the difference of the logarithms is at least
\[
v^2+\Bigl(\frac u2-\frac{u^2}{3}\Bigr)-\frac52\,v+\frac v2\,\mu^0_r(X_-)+2v\,(1-v)\,\mu^0_r(X_{\mathrm{mid}})-w\,D_r .
\]
Since $u=1/r$, $v=1/k$ and $w=v\cdot\frac{k}{k-1}$, this expression is $v^2+v\,\mathsf G(r,k)$, i.e.\ $\mathsf G(r,k)/k+1/k^2$. This proves the supposed inequality.
\end{proof}

\subsection{Positivity of the criterion function in the strip}\label{subsec:next-strip}

The criterion is of use where $\mathsf G(r,k)\ge 0$. We prove this for a column $k\ge 265$ in the rows of its strip, i.e.\ in rows $k+r$ with $25\,r^3\le 121\,k^2$; by Proposition~\ref{prop:strip}, every row up to which such a column has not turned positive lies in the strip.

The function has three positive terms, namely $k\bigl(\frac{1}{2r}-\frac{1}{3r^2}\bigr)$ and the two terms with $\mu^0_r$, and two negative ones, the constant $\frac52$ and the term with $D_r$. For small $r$ the first term outweighs the negative terms by itself. For large $r$ it does not, and the term with $D_r$ can only be bounded by a multiple of $\sqrt r$, through the mean $s_r$ of the birthday law (Lemma~\ref{lem:ratiobound}). It has to be balanced by the two terms with $\mu^0_r$. These are of order $\sqrt r$ as well, but only as long as their arguments $X_-$ and $X_{\mathrm{mid}}$, which are of size $r^2/k$, are at most multiples of $\sqrt r$, because $\mu^0_r(X)$ decreases as $X$ grows; the condition $25\,r^3\le 121\,k^2$, i.e.\ $5\,r^{3/2}\le 11\,k$, bounds the arguments in this way. The proof first replaces the coefficients $2\,(1-\frac1k)$ and $\frac{k}{k-1}$ in $\mathsf G(r,k)$ by constants and records two bounds for the mean $s_r$, and then treats $2\le r\le 13$ and $r\ge 14$ separately.

\begin{proposition}\label{prop:criterionpositive}
Let $k\ge 265$ and $r\ge 2$ be integers with $25\,r^3\le 121\,k^2$. Then the criterion function of Definition~\ref{def:criterionfunction} satisfies $\mathsf G(r,k)>0$.
\end{proposition}

\begin{proof}
As before, let $X_-=\frac{r(r+1)}{k-1}$ and $X_{\mathrm{mid}}=\frac{r(r+2)}{2k}$ be the two arguments of $\mu^0_r$ in $\mathsf G(r,k)$, and note that $\frac{1}{2r}-\frac{1}{3r^2}=\frac{3r-2}{6r^2}$. The numerical comparisons that occur are checked by direct evaluation.

We begin by replacing the coefficients $2\,(1-\frac1k)$ and $\frac{k}{k-1}$ in $\mathsf G(r,k)$ by constants. The numbers $\mu^0_r(X_-)$ and $\mu^0_r(X_{\mathrm{mid}})$ are nonnegative, and so is $D_r$: by \eqref{eq:birthday-ratio},
\[
\sum_{q=1}^{r}q\,(q-1)(q-2)\,T_{r,q}=D_r\sum_{q=1}^{r}q\,(q-1)\,T_{r,q},
\]
where the total weights $T_{r,q}$ of the compositions of $r$ with $q$ parts (Section~\ref{sec:comps}) are nonnegative and $T_{r,2}=\sum_{p=1}^{r-1}a_p\,a_{r-p}>0$ because $r\ge 2$. Since $k\ge 265$, we have $2\,(1-\frac1k)\ge\frac{528}{265}$ and $\frac{k}{k-1}\le\frac{265}{264}$, and therefore
\begin{equation}\label{eq:next-assemble}
\mathsf G(r,k)\ge k\,\frac{3r-2}{6r^2}-\frac52+\frac12\,\mu^0_r(X_-)+\frac{528}{265}\,\mu^0_r(X_{\mathrm{mid}})-\frac{265}{264}\,D_r .
\end{equation}

Next we record two bounds for the mean of the birthday law. We have
\[
1\le s_r<r .
\]
The upper bound is shown in Section~\ref{sec:birthday}, where $D_r$ is defined, and the lower bound holds because the first term of the sum $s_r$ is $1$ and the others are nonnegative. Further, $s_m\le\sqrt{\pi m/2}\;e^{1/(12m)}-\frac13$ for every integer $m\ge 1$, by Lemma~\ref{lem:birthdaymean}. Here $e^{1/(12m)}\le\frac{12m}{12m-1}$, because $e^{-1/(12m)}\ge 1-\frac{1}{12m}$, and $\frac{12m}{12m-1}$ decreases in $m$; its values at $m=13$ and at $m=14$ are $\frac{156}{155}$ and $\frac{168}{167}$. Since
\[
\sqrt{\pi/2}\le 1.25332,\qquad 1.25332\cdot\frac{156}{155}<1.2615,\qquad 1.25332\cdot\frac{168}{167}<1.261,
\]
we obtain
\[
s_{13}\le 1.2615\,\sqrt{13}-\frac13,\qquad s_m\le 1.261\,\sqrt m-\frac13\quad(m\ge 14).
\]

Suppose first that $2\le r\le 13$. One checks that $\frac{3r-2}{6r^2}\ge\frac{37}{1014}$ for these $r$, with equality at $r=13$. Since $k\ge 265$, the first term of \eqref{eq:next-assemble} is therefore at least $265\cdot\frac{37}{1014}>9.669$. We bound the two terms with $\mu^0_r$ below by $0$. For the last term, $D_r\le D_{13}\le\frac32\,(s_{13}-1)$ by Lemmas~\ref{lem:ratiomonotone} and~\ref{lem:ratiobound}, and since $s_{13}\le 1.2615\,\sqrt{13}-\frac13<4.216$, we get $D_r\le\frac32\cdot 3.216=4.824$. Hence
\[
\mathsf G(r,k)>9.669-\frac52-\frac{265}{264}\cdot 4.824>9.669-2.5-4.843>2.32 .
\]

Suppose now that $r\ge 14$, and put $t=\sqrt r$, so that $t\ge\sqrt{14}>3.74$. The condition $25\,r^3\le 121\,k^2$ then reads
\[
5\,t^3\le 11\,k .
\]
We prove the four bounds
\begin{gather*}
k\,\frac{3r-2}{6r^2}\ge 1.88,\qquad\mu^0_r(X_{\mathrm{mid}})\ge 0.837\,t-0.92,\\
\mu^0_r(X_-)\ge 0.628\,t-0.62,\qquad D_r\le 1.8915\,t-2.6 .
\end{gather*}
The first bound, and the bounds for $X_-$ and $X_{\mathrm{mid}}$ on which the second and the third rest, need a lower bound for $k$ in terms of $r$: for $r\le 69$ we use $k\ge 265$, and for $r\ge 70$ we use $5\,t^3\le 11\,k$.

For the first bound, let $14\le r\le 69$. Then $\frac{3r-2}{6r^2}\ge\frac{205}{28566}$, with equality at $r=69$, as one checks in the same way, so that $k\,\frac{3r-2}{6r^2}\ge 265\cdot\frac{205}{28566}>1.9$. If $r\ge 70$, then $t\ge\sqrt{70}>8.36$, and since $k\ge\frac{5}{11}\,t^3$,
\[
k\,\frac{3r-2}{6r^2}\ge\frac{5}{11}\,t^3\cdot\frac{3t^2-2}{6t^4}=\frac{5}{11}\Bigl(\frac t2-\frac{1}{3t}\Bigr),
\]
which increases with $t$ and is larger than $1.88$ at $t=8.36$.

Next we show that
\[
X_{\mathrm{mid}}\le 1.1315\,t,\qquad X_-\le 2.24\,t .
\]
Since $r=t^2$, we have $X_{\mathrm{mid}}=t\cdot\frac{t\,(t^2+2)}{2k}$ and $X_-=t\cdot\frac{t\,(t^2+1)}{k-1}$. If $r\le 69$, then $t\le\sqrt{69}<8.31$, and $k\ge 265$ gives
\[
\frac{t\,(t^2+2)}{2k}\le\frac{8.31\cdot 71}{530}<1.12,\qquad\frac{t\,(t^2+1)}{k-1}\le\frac{8.31\cdot 70}{264}<2.21 .
\]
If $r\ge 70$, then $\frac1k\le\frac{11}{5\,t^3}$, again because $5\,t^3\le 11\,k$, and $\frac{1}{k-1}\le\frac{265}{264}\cdot\frac1k$ since $k\ge 265$. With $t^2\ge 70$ we get
\begin{gather*}
\frac{t\,(t^2+2)}{2k}\le\frac{11}{10}\Bigl(1+\frac{2}{t^2}\Bigr)\le\frac{11}{10}\cdot\frac{72}{70}<1.1315,\\
\frac{t\,(t^2+1)}{k-1}\le\frac{265}{264}\cdot\frac{11}{5}\cdot\frac{71}{70}<2.24 .
\end{gather*}

For the second and the third bound we use that, for real numbers $X\ge 0$ and $m$,
\[
\frac32\,m^2+(4+X)\,m\le 2\,(r-1)\qquad\text{implies}\qquad m\le\mu^0_r(X).
\]
Indeed, the first inequality says that $(3m+4+X)^2\le(4+X)^2+12\,(r-1)$, and taking square roots gives the second. We apply this implication with $X=X_{\mathrm{mid}}$ and $m=0.837\,t-0.92$, and with $X=X_-$ and $m=0.628\,t-0.62$. In both cases $m>0$, because $t\ge 3.74$, so that the bounds for $X_{\mathrm{mid}}$ and $X_-$ just shown give $(4+X_{\mathrm{mid}})\,m\le(4+1.1315\,t)\,m$ in the first case and $(4+X_-)\,m\le(4+2.24\,t)\,m$ in the second. Since $2\,(r-1)=2\,(t^2-1)$, it therefore suffices to note that, for these two values of $m$, respectively,
\begin{align*}
2\,(t^2-1)-\frac32\,m^2-(4+1.1315\,t)\,m&=0.002081\,t^2+0.0031\,t+0.4104,\\
2\,(t^2-1)-\frac32\,m^2-(4+2.24\,t)\,m&=0.001704\,t^2+0.04488\,t-0.0966,
\end{align*}
as one finds by expanding, and that both right-hand sides are positive for $t\ge 3.74$, the second because $0.04488\cdot 3.74>0.0966$. This gives the second and the third of the four bounds.

For the fourth bound, put $\bar s=1.261\,t-\frac13$, so that $1\le s_r\le\bar s$ by the bounds for the mean of the birthday law recorded above. Since $r-s_r>0$, an elementary calculation shows that the inequality $D_r\le\frac32\,(s_r-\frac75)$ is equivalent to
\[
s_r^2+\frac35\,s_r\le\frac85\,r .
\]
This holds, because $0<s_r\le\bar s$ and
\[
\bar s^{\,2}+\frac35\,\bar s=1.590121\,t^2-\frac{1.261}{15}\,t-\frac{4}{45}\le\frac85\,t^2 .
\]
Consequently $D_r\le\frac32\,(s_r-\frac75)\le\frac32\,(\bar s-\frac75)=1.8915\,t-2.6$.

Finally, we insert the four bounds into \eqref{eq:next-assemble} and obtain
\begin{align*}
\mathsf G(r,k)&\ge 1.88-\frac52+\frac12\,(0.628\,t-0.62)\\
&\quad+\frac{528}{265}\,(0.837\,t-0.92)-\frac{265}{264}\,(1.8915\,t-2.6)\\
&=0.083018\ldots\cdot t-0.153208\ldots\ge 0.083\,t-0.1533 .
\end{align*}
Since $t\ge 3.74$ and $0.083\cdot 3.74-0.1533>0.157$, we have $\mathsf G(r,k)>0$ in this range as well.
\end{proof}

\subsection{From a column to the next}\label{subsec:next-column}

For $k\ge 265$ we can now prove Proposition~\ref{prop:nextcolumn}, with a strict conclusion: by Proposition~\ref{prop:strip} the hypothesis confines row $n$ to the strip of column $k$, where the criterion applies. In terms of the steps $d_{n,k}=b_{n,k}-b_{n,k-1}$, the hypothesis of the following proposition is that $d_{j,k}\le 0$ for $k+1\le j\le n$, and its conclusion is that $d_{n,k+1}<0$.

\begin{proposition}\label{prop:nextlarge}
Let $k\ge 265$ and $n\ge k+2$ be integers. If $b_{j,k}\le b_{j,k-1}$ for every integer $j$ with $k+1\le j\le n$, then $b_{n,k+1}<b_{n,k}$.
\end{proposition}

\begin{proof}
By Proposition~\ref{prop:strip} there is an integer $j>k$ with $25\,(j-k-1)^3\le 121\,k^2$ and $d_{j,k}>0$, and the hypothesis shows that $j>n$. Write $n=k+r$. Then $2\le r\le j-k-1$, so that
\[
25\,r^3\le 25\,(j-k-1)^3\le 121\,k^2 ;
\]
that is, row $n$ lies in the strip of column $k$. In particular $r\le k-1$, because $25\,k^3>121\,k^2$. Hence $\mathsf G(r,k)>0$ by Proposition~\ref{prop:criterionpositive}, and since $\delta^{(k)}_n\le 0$ by the hypothesis at row $n$, Proposition~\ref{prop:nextcriterion} gives $\delta^{(k+1)}_n<0$, i.e.\ $b_{n,k+1}<b_{n,k}$.
\end{proof}

For $9\le k\le 264$ exact computation takes the place of Propositions~\ref{prop:strip} and~\ref{prop:criterionpositive}: the conclusion follows, without strictness, from the two computed signs of Lemma~\ref{lem:signpairs} and from Proposition~\ref{prop:chainrange}, applied in the chain range of column $k+1$.

\begin{proof}[Proof of Proposition~\ref{prop:nextcolumn}]
Let $k\ge 9$ and $n\ge k+2$ be integers such that $b_{j,k}\le b_{j,k-1}$ for every integer $j$ with $k+1\le j\le n$. We have to show that $b_{n,k+1}\le b_{n,k}$. For $k\ge 265$ this is contained in Proposition~\ref{prop:nextlarge}.

Let $9\le k\le 264$, and let $\zeta=\zeta(k)$ be the number of Definition~\ref{def:zeta}. By Lemma~\ref{lem:signpairs} we have $\zeta+2\le k$ and
\[
b_{k+\zeta+1,\,k-1}<b_{k+\zeta+1,\,k},\qquad b_{k+\zeta,\,k+1}<b_{k+\zeta,\,k};
\]
that is, column $k$ is positive in row $k+\zeta+1$, and column $k+1$ is negative in row $k+\zeta$. By the first inequality and the hypothesis, $n\le k+\zeta$.

Suppose now that $b_{n,k}<b_{n,k+1}$. Rows $n,\dots,k+\zeta$ lie in the chain range of column $k+1$, which consists of rows $k+2,\dots,2k+1$, because $n\ge k+2$ and $\zeta\le k-2$. By Proposition~\ref{prop:chainrange}, applied to column $k+1$ in the rows $n,\dots,k+\zeta-1$ one after another, this column is then positive in row $k+\zeta$ as well, which contradicts the second inequality. Hence $b_{n,k+1}\le b_{n,k}$.
\end{proof}

\section{Proof of Theorem~\ref{thm:main}}\label{sec:proof}
In this section we derive Theorem~\ref{thm:main} from Propositions~\ref{prop:leftpart}, \ref{prop:chainrange} and~\ref{prop:nextcolumn} and from the computed rows $n\le 16$ (Lemma~\ref{lem:sixteenrows}). The section closes with a bound for the position of the maximum of a row $n\ge 335$ (Proposition~\ref{prop:rowpeak}), which follows from statements proved in the earlier sections and is not used in the proof of the theorem.

The following lemma shows that, to the right of the middle of a row, a nonpositive step is followed only by nonpositive steps. Its conclusion is stated for all rows up to $n$, because the induction in the proof requires them.

\begin{lemma}\label{lem:rightpart}
Let $n\ge 0$ and $i\ge 9$ be integers with $n\le 2i-1$, and suppose that $b_{n,i}\le b_{n,i-1}$. Then
\[
b_{m,j}\le b_{m,j-1}
\]
for all integers $j$ and $m$ with $j\ge i$ and $j+1\le m\le n$.
\end{lemma}

\begin{proof}
We prove by induction on $j\ge i$ that $b_{m,j}\le b_{m,j-1}$ for every $m$ with $j+1\le m\le n$.

Let $j=i$, and let $m$ be an integer with $i+1\le m\le n$. If we had $b_{m,i-1}<b_{m,i}$, then Proposition~\ref{prop:chainrange}, applied with $k=i$ in the rows $m,\dots,n-1$ one after another, which lie between $i+1$ and $2i-2$, would give $b_{n,i-1}<b_{n,i}$, contrary to the assumption of the lemma. Hence $b_{m,i}\le b_{m,i-1}$.

Suppose now that the claim holds for some $j\ge i$, and let $m$ be an integer with $j+2\le m\le n$. By the induction hypothesis $b_{m',j}\le b_{m',j-1}$ for every $m'$ with $j+1\le m'\le m$. Hence Proposition~\ref{prop:nextcolumn}, applied with $k=j\ge 9$ in row $m$, gives $b_{m,j+1}\le b_{m,j}$. This is the claim for $j+1$.
\end{proof}

Together with Proposition~\ref{prop:leftpart} and Lemma~\ref{lem:sixteenrows}, Lemma~\ref{lem:rightpart} excludes, in every row, a strict descent that is followed by a strict ascent.

\begin{proposition}\label{prop:descentascent}
Let $n$, $i$ and $j$ be integers with $n\ge 0$ and $1\le i<j\le n-1$. Then the inequalities $b_{n,i}<b_{n,i-1}$ and $b_{n,j-1}<b_{n,j}$ do not both hold.
\end{proposition}

\begin{proof}
Suppose that $b_{n,i}<b_{n,i-1}$ and $b_{n,j-1}<b_{n,j}$. By Lemma~\ref{lem:sixteenrows} this is impossible for $n\le 16$, so let $n\ge 17$. If $2i\le n$, then Proposition~\ref{prop:leftpart} with $k=i$ gives $b_{n,i-1}\le b_{n,i}$, contrary to the first inequality. If $n\le 2i-1$, then $i\ge 9$, and Lemma~\ref{lem:rightpart} applies to $n$ and $i$. Since $j\ge i$ and $j+1\le n$, its conclusion for column $j$ and row $m=n$ reads $b_{n,j}\le b_{n,j-1}$, contrary to the second inequality.
\end{proof}

A finite sequence in which no strict descent is followed by a strict ascent is unimodal, and so the theorem follows.

\begin{proof}[Proof of Theorem~\ref{thm:main}]
Let $n\ge 1$ be an integer, and let $p$ be the largest integer with $0\le p\le n-1$ such that $b_{n,\ell-1}\le b_{n,\ell}$ for $1\le\ell\le p$. If $p<n-1$, then $b_{n,p+1}<b_{n,p}$, and Proposition~\ref{prop:descentascent} with $i=p+1$ shows that $b_{n,j}\le b_{n,j-1}$ for $p+2\le j\le n-1$. Hence $b_{n,\ell+1}\le b_{n,\ell}$ for every $\ell$ with $p\le\ell\le n-2$, which is the assertion of the theorem with $\ell^*=p$.
\end{proof}

\subsection{The position of the maximum}\label{subsec:proof-peak}

Theorem~\ref{thm:main} does not say where the maximum of a row stands. In the rows $n\ge 335$ the statements of the earlier sections bound its position (Proposition~\ref{prop:rowpeak}). We first note two facts about a column $k\ge 265$ and its strip, which consists of the rows $k+r$ with $r\ge 1$ and $25\,r^3\le 121\,k^2$ (Section~\ref{sec:strip}).

The first fact is that the column is positive in every row beyond its strip: if $k\ge 265$ and $r$ are integers with $121\,k^2<25\,r^3$, then $b_{k+r,k-1}<b_{k+r,k}$. Indeed, by Proposition~\ref{prop:strip} there is an integer $j>k$ with $25\,(j-k-1)^3\le 121\,k^2$ and $b_{j,k-1}<b_{j,k}$, and $j\le k+r$, since $j>k+r$ would imply $25\,r^3\le 25\,(j-k-1)^3\le 121\,k^2$, against the hypothesis. Put $n=k+r$. If $n\le 2k-1$, then $k+1\le j\le n\le 2k-1$, and Proposition~\ref{prop:chainrange}, applied in the rows $j,\dots,n-1$ one after another, gives $b_{n,k-1}<b_{n,k}$; if $n\ge 2k$, this inequality is Proposition~\ref{prop:largecolumns}.

The second fact is that, in a row of the strip other than the first, a nonpositive step of the column is followed by a negative step of the next column: for integers $k\ge 265$ and $r\ge 2$ with $25\,r^3\le 121\,k^2$ and for $n=k+r$, the inequality $b_{n,k}\le b_{n,k-1}$ implies $b_{n,k+1}<b_{n,k}$. Indeed, $r\le k-1$ because $25\,r^3\le 121\,k^2<25\,k^3$, so that $\mathsf G(r,k)>0$ by Proposition~\ref{prop:criterionpositive}; and since the normalised step $\delta^{(j)}_n=d_{n,j}/n!$ of column $j$ in row $n$ has the sign of $b_{n,j}-b_{n,j-1}$, the inequality $b_{n,k}\le b_{n,k-1}$ says that $\delta^{(k)}_n\le 0$, and Proposition~\ref{prop:nextcriterion} gives $\delta^{(k+1)}_n<0$, i.e.\ $b_{n,k+1}<b_{n,k}$.

Applied along a row $n\ge 335$, the two facts give the shape of the row.

\begin{proposition}\label{prop:rowpeak}
Let $n\ge 335$ be an integer. Then there is an integer $p$ with $0\le p\le n-1$ and $25\,(n-1-p)^3\le 121\,(p+1)^2$ such that
\begin{align*}
b_{n,\ell-1}&\le b_{n,\ell}\quad(1\le\ell\le p), & b_{n,\ell-1}&<b_{n,\ell}\quad(82\le\ell\le p),\\
b_{n,\ell+1}&\le b_{n,\ell}\quad(p\le\ell\le n-2), & b_{n,\ell+1}&<b_{n,\ell}\quad(p+1\le\ell\le n-2).
\end{align*}
\end{proposition}

\begin{proof}
Let first $k$ be an integer with $1\le k\le 264$; we show that $b_{n,k-1}\le b_{n,k}$, and that $b_{n,k-1}<b_{n,k}$ if $k\ge 82$. If $2k\le n$, this is Proposition~\ref{prop:leftpart} and, for $k\ge 82$, Proposition~\ref{prop:largecolumns}. If $n\le 2k-1$, then $k\ge 168$, and Corollary~\ref{cor:earlypositive} gives an integer $r\ge 1$ with $100\,r\le 27\,k\le 7128$, hence $r\le 71$, and $b_{k+r,k-1}<b_{k+r,k}$. Since $k+1\le k+r\le 335\le n\le 2k-1$, Proposition~\ref{prop:chainrange}, applied in the rows $k+r,\dots,n-1$ one after another, gives $b_{n,k-1}<b_{n,k}$.

Let $p$ be the least integer $p\ge 81$ such that $p=n-1$ or $b_{n,p+1}\le b_{n,p}$; since $n-1\ge 81$, it exists and $p\le n-1$. For $82\le\ell\le p$ the integer $\ell-1$ is at least $81$ and less than $p$, so that $b_{n,\ell-1}<b_{n,\ell}$ by the minimality of $p$; and $b_{n,\ell-1}\le b_{n,\ell}$ for $1\le\ell\le 81$ by the first paragraph of the proof. If $p=n-1$, then $n-1-p=0$ and no $\ell$ satisfies $p\le\ell\le n-2$, so that nothing remains to be shown.

Let $p\le n-2$, and put $q=p+1$, so that $82\le q\le n-1$ and $b_{n,q}\le b_{n,q-1}$. By the first paragraph of the proof, $q\ge 265$. If we had $121\,q^2<25\,(n-q)^3$, the first fact with $k=q$ and $r=n-q$ would give $b_{n,q-1}<b_{n,q}$; hence $25\,(n-q)^3\le 121\,q^2$, which is the inequality for $p$. For every integer $j$ with $q\le j\le n-2$ we have $j\ge 265$, $n-j\ge 2$ and $25\,(n-j)^3\le 25\,(n-q)^3\le 121\,q^2\le 121\,j^2$, so that, by the second fact with $k=j$ and $r=n-j$, the inequality $b_{n,j}\le b_{n,j-1}$ implies $b_{n,j+1}<b_{n,j}$. Since $b_{n,q}\le b_{n,q-1}$, induction on $j$ gives $b_{n,j+1}<b_{n,j}$ for $q\le j\le n-2$. Together with $b_{n,p+1}\le b_{n,p}$, these are the inequalities for $\ell\ge p$.
\end{proof}

The inequality for $p$ forces $p\ge 265$, since $p\le 264$ would give $n-1-p\ge 70$ and $p+1\le 265$, whereas $25\cdot 70^3>121\cdot 265^2$. Hence the maximum of the row is $b_{n,p}$, and it is attained at column $p$ and at no other column except possibly $p+1$. Since $p+1\le n$, the inequality for $p$ also gives
\[
n-1-p\le\Bigl(\frac{121}{25}\Bigr)^{1/3}(p+1)^{2/3}\le\Bigl(\frac{121}{25}\Bigr)^{1/3}n^{2/3}<1.692\,n^{2/3}.
\]
This bounds the distance of the maximum from the last column of the row from above. It does not determine the position of the maximum, and we have no bound for the distance from below. The constant $\frac{121}{25}$ is that of the strip, the range of rows which the estimates behind Propositions~\ref{prop:strip} and~\ref{prop:criterionpositive} reach, and it has no meaning beyond these estimates; the bound of Proposition~\ref{prop:strip} is not sharp, as was remarked after it. The other constants arise in the same way: $265$ is the first column to which these two propositions apply, $82$ is the first to which Proposition~\ref{prop:largecolumns} applies, and $335=264+71$ comes from the bound of Corollary~\ref{cor:earlypositive}, as the first paragraph of the proof shows.

\bibliographystyle{amsplain}
\bibliography{references}
\end{document}
